\documentclass[a4paper,10pt]{article}
\usepackage[top=24mm,bottom=24mm,left=26mm,right=26mm,headheight=14pt]{geometry}
\usepackage{fontspec,amsmath,amsthm,unicode-math}
\usepackage[spanish,es-nodecimaldot,es-noquoting]{babel}
\usepackage{microtype,xcolor,graphicx,tikz,booktabs,array,needspace,fancyhdr}
\usepackage{longtable,tabularx,enumitem,float,import,xurl,multicol}
\usetikzlibrary{arrows.meta,calc,positioning,decorations.pathreplacing,matrix}
\usepackage[unicode,hidelinks]{hyperref}
\usepackage[nameinlink,noabbrev]{cleveref}
\crefname{theorem}{teorema}{teoremas}\Crefname{theorem}{Teorema}{Teoremas}
\crefname{proposition}{proposición}{proposiciones}\Crefname{proposition}{Proposición}{Proposiciones}
\crefname{lemma}{lema}{lemas}\Crefname{lemma}{Lema}{Lemas}
\crefname{section}{sección}{secciones}\Crefname{section}{Sección}{Secciones}
\crefname{equation}{ecuación}{ecuaciones}\Crefname{equation}{Ecuación}{Ecuaciones}
\setmainfont{STIXTwoText-Regular.otf}[BoldFont=STIXTwoText-Bold.otf,ItalicFont=STIXTwoText-Italic.otf,BoldItalicFont=STIXTwoText-BoldItalic.otf]
\setmathfont{STIXTwoMath-Regular.otf}
\newtheorem{theorem}{Teorema}[section]
\newtheorem{proposition}[theorem]{Proposición}
\newtheorem{lemma}[theorem]{Lema}
\newtheorem{corollary}[theorem]{Corolario}
\newtheorem{falsifier}[theorem]{Criterio de refutación}
\newtheorem{ablation}[theorem]{Prueba de necesidad estructural}
\newtheorem{criterion}[theorem]{Criterio}
\theoremstyle{definition}\newtheorem{definition}[theorem]{Definición}
\newtheorem{axiom}[theorem]{Axioma}
\newtheorem{principle}[theorem]{Principio}
\theoremstyle{remark}\newtheorem{remark}[theorem]{Observación}
\newcommand{\RR}{\mathbb R}\newcommand{\QQ}{\mathbb Q}
\newcommand{\ZZ}{\mathbb Z}\newcommand{\CC}{\mathbb C}
\newcommand{\Gr}{\operatorname{Gr}}\newcommand{\Res}{\operatorname{Res}}
\newcommand{\APP}{\mathrm{APP}}\newcommand{\TPK}{\mathrm{TPK}}
\newcommand{\TRIT}{\mathrm{TRIT}}\newcommand{\id}{\operatorname{id}}
\newcommand{\diag}{\operatorname{diag}}\newcommand{\conv}{\operatorname{conv}}
\newcommand{\FF}{\mathbb F}\newcommand{\Id}{\operatorname{id}}
\newcommand{\tr}{\operatorname{tr}}\newcommand{\dd}{\mathrm d}
\newcommand{\Span}{\operatorname{span}}\newcommand{\Cyl}{\operatorname{Cyl}}
\providecommand{\llbracket}{\lBrack}\providecommand{\rrbracket}{\rBrack}
\providecommand{\square}{\mdlgwhtsquare}
\providecommand{\dashrightarrow}{\mathrel{\rightdasharrow}}
\providecommand{\index}[1]{}
\newcommand{\HMT}{\mathrm{HMT}}\newcommand{\Tr}{\operatorname{Tr}}\newcommand{\Spec}{\operatorname{Spec}}
\newcommand{\Cstar}{C^{*}}\newcommand{\R}{\mathbb R}\newcommand{\Z}{\mathbb Z}\newcommand{\I}{\mathrm I}
\setcounter{tocdepth}{2}
\setlist{nosep,topsep=4pt}
\numberwithin{equation}{section}
\setlength{\parskip}{4pt plus 1pt}\setlength{\parindent}{1em}
\setlength{\emergencystretch}{2em}
\renewcommand{\normalsize}{\fontsize{10.5}{13.5}\selectfont}\normalsize
\clubpenalty=10000\widowpenalty=10000\displaywidowpenalty=10000
\AddToHook{cmd/section/before}{\Needspace{10\baselineskip}}
\AddToHook{cmd/subsection/before}{\Needspace{6\baselineskip}}
\AddToHook{env/theorem/before}{\Needspace{6\baselineskip}}
\AddToHook{env/proposition/before}{\Needspace{6\baselineskip}}
\pagestyle{fancy}\fancyhf{}
\fancyhead[L]{\footnotesize Grassmannianos positivos y amplituedros: formas canónicas y refinamiento}
\fancyfoot[C]{\thepage}\renewcommand{\headrulewidth}{.2pt}
\hypersetup{pdftitle={Grassmannianos positivos y amplituedros: formas canónicas y refinamiento},pdfsubject={Dualidad reticular y salto de masa de Yang–Mills en Holografía Modular Triádica},pdfauthor={Oumar Haidara Fall; Rubén Ramos Balsa}}
\makeatletter
\newcommand{\compactcontents}{%
  \begingroup
  \fontsize{9.2}{11.1}\selectfont
  \setlength{\parskip}{0pt}\setlength{\columnsep}{7mm}
  \renewcommand{\l@section}{\@dottedtocline{1}{0pt}{1.8em}}
  \renewcommand{\l@subsection}{\@dottedtocline{2}{1.5em}{2.7em}}
  \renewcommand{\l@part}[2]{\addpenalty{-\@highpenalty}\addvspace{6pt}%
    {\bfseries\noindent ##1\nobreak\hfill\nobreak ##2\par}\nobreak}
  \section*{Índice}
  \@starttoc{toc}
  \endgroup}
\makeatother
% BEGIN NUCLEO_ADITIVO_20260921_PREAMBULO
\makeatletter
\@ifundefined{example}{\theoremstyle{definition}\newtheorem{example}[theorem]{Ejemplo}}{}
\@ifundefined{notatabla}{\newcommand{\notatabla}[1]{\par\smallskip{\footnotesize #1\par}}}{}
\makeatother
\providecommand{\FF}{\mathbb F}

% END NUCLEO_ADITIVO_20260921_PREAMBULO
% Apertura independiente del índice: corrección quirúrgica 20260928.
\let\HMTIndiceOriginal\tableofcontents
\renewcommand{\tableofcontents}{\clearpage\begingroup\setlength{\parskip}{0pt}\fontsize{9.2}{10.5}\selectfont\HMTIndiceOriginal\endgroup\clearpage}
\begin{document}
\thispagestyle{plain}
\begin{center}
{\small Manuscrito de recopilación para transferencia académica}\par\vspace{6mm}
{\LARGE\bfseries Grassmannianos positivos y amplituedros:\\formas canónicas y refinamiento}\par\vspace{3mm}
{\large Dualidad reticular y salto de masa de Yang--Mills\\en Holografía Modular Triádica}\par\vspace{5mm}
Oumar Haidara Fall \textperiodcentered\ Rubén Ramos Balsa\par
{\fontsize{8}{10}\selectfont Coautoría sin prelación de contribución\par
Integración y recopilación documental generadas por ChatGPT - Astra.}
\end{center}
\begin{abstract}
La Holografía Modular Triádica construye la estructura discreta conjunta
del continuo mediante APP, TRIT y el transporte con memoria del TPK.
La reducción compatible de cilindros produce la realización aritmética,
mientras la proyección conjugada conserva la incidencia excepcional.
Sobre esta genealogía se demuestra una colección de realizaciones positivas
recuperables. La carga uniforme y dos memorias nonádicas determinan el
registro de doce coordenadas; sus separaciones ordenadas producen una
medición planar de frontera y una representación grassmanniana con
inversa marcada. Los lectores que factorizan por el registro se transportan
a esa representación y se conservan bajo cambios de coordenadas del mismo
punto. El refinamiento nonádico construye banderas positivas y puntos
amplituhedrales en todo rango y dimensión externa finitos admisibles.
En rango uno se obtiene cobertura simplicial completa y cancelación de
residuos interiores en toda dimensión finita. Las formas con valores en
la recta de acción satisfacen una condición exacta de pegado en las
fronteras comunes. Los promedios celulares conservan una medida
cronológica y descomponen el incremento de los momentos en detalles
ortogonales; su límite determina un operador de Jacobi cuyas compresiones
finitas tienen espectro simple, residuos positivos y resolventes en
fracción continua.
La reducción de Schur conserva la energía efectiva de la partición,
con reconstrucción del detalle y realización límite de Sobolev.
El flujo del registro induce además una deformación reticular de
veinticuatro dimensiones que conserva volumen y simetría ternaria,
y transforma la inversión del parámetro en dualidad reticular y
reciprocidad theta. Se reúnen los Hamiltonianos elíptico, hiperbólico
y parabólico, sus coordenatizaciones simplécticas, términos de conexión y
transformaciones de Legendre, junto con los lectores materiales
correspondientes. El desarrollo de Yang--Mills incorpora la medida
común, la escala areal--volumétrica, la coercividad uniforme, el
testigo espectral persistente y la reconstrucción del Hamiltoniano.
La articulación de estos resultados distingue y relaciona las
positividades de Plücker, de los momentos, de las constitutivas
y del espectro mediante operaciones explícitas sobre el mismo
estado enriquecido.
La representación incidencial conserva las doce
coordenadas mediante una isometría ponderada con inversa explícita.
La transición héxada--óctada permite distinguir el área total de la
composición sectorial y reconstruir las ocupaciones; su purificación
determina una entropía de entrelazamiento con ley modular finita.
La composición longitudinal del retorno marcado determina la escala
areal y conserva su naturalidad bajo los refinamientos heredados;
las formas canónicas con coeficiente operatorial de área transmiten
la aditividad y los residuos de frontera.
\par\smallskip\noindent\textbf{Palabras clave:} Grassmannianos positivos; redes planares; coordenadas de Plücker; geometría de incidencia; Holografía Modular Triádica.
\end{abstract}
\clearpage
\compactcontents
\clearpage

\section{Introducción: genealogía de la forma y conservación del refinamiento}
\label{sec:genealogy}
La cuestión rectora de este trabajo es la conservación de una estructura
generada cuando cambia su realización matemática. La Holografía Modular
Triádica produce, mediante operaciones aritméticas levantadas, estados que
retienen simultáneamente orientación, incidencia, frontera y memoria. Su
estructura discreta del continuo reúne las prolongaciones compatibles de
esos estados. La realización arquimediana expresa coordenadas mediante el
refinamiento de cilindros; la realización conjugada conserva relaciones
excepcionales de incidencia. La genealogía precede, por tanto, a los valores
y a las formas en que se los representa. Ésta es la tesis de partida que el
desarrollo formal siguiente expone con sus operaciones y dependencias.

El problema geométrico propio del artículo comienza cuando una salida
positiva de aquella construcción se convierte en separación entre nodos,
peso de una red, coordenada de Plücker, momento de una medida, coeficiente
de energía o dirección de deformación reticular. Para justificar ese paso
se necesita determinar qué información recupera cada representación, cómo
se compara con las restantes y qué conserva su refinamiento. El interés
de una realización positiva aumenta de forma decisiva cuando admite una
inversa marcada: la geometría deja entonces de ofrecer sólo una evaluación
del estado y pasa a conservar la información que utilizan sus lectores.
La marca comprende el orden de las posiciones, la escala total y los datos
enriquecidos que el lector particular requiere.

La aportación central se organiza alrededor de tres operaciones. La primera
es la reconstrucción: la carga y la memoria recuperan las coordenadas
del registro, y los cocientes de menores recuperan sus separaciones.
La segunda es el refinamiento compatible: una partición se subdivide
conservando sus extremos, sus marcas y un archivo exacto del detalle.
La tercera es el transporte de estructuras: una forma, un operador o un
lector se lleva a otra representación mediante un mapa explícito. Estas
operaciones permiten demostrar conservación de residuos, energía efectiva,
momentos y espectros en los dominios concretos donde cada identidad tiene
sentido. El calificativo positivo adquiere así varias realizaciones
matemáticas precisas, articuladas por la reconstrucción del mismo dato.

La geometría positiva posee antecedentes propios que permiten identificar
con precisión el nivel de esta contribución. Las mediciones de frontera
de redes dirigidas parametrizan celdas del Grassmanniano totalmente no
negativo en el trabajo de Postnikov~\cite{postnikov}; Scott describe la
estructura de álgebra de conglomerados del anillo de coordenadas
grassmanniano~\cite{scott}. Arkani-Hamed y Trnka introducen el amplituedro
en la descripción geométrica de amplitudes de la teoría supersimétrica
\(\mathcal N=4\) de Yang--Mills en el límite planar~\cite{amplituhedron}.
La formulación de geometrías positivas caracteriza después sus formas
canónicas por polos logarítmicos y residuos recursivos de
frontera~\cite{positive}. El presente desarrollo compone esos lenguajes
de realización con datos previamente generados por HMT. Su resultado
diferencial es una colección marcada, recuperable y refinable, junto con
los mapas que enlazan su forma positiva y sus lecturas materiales.

El refinamiento amplía además el alcance dimensional. El número de nodos
crece mediante subdivisiones nonádicas y permite construir matrices
positivas de todo rango y dimensión externa finitos compatibles con ese
número. La prueba se apoya en determinantes de Vandermonde y formas de
momentos de rango completo. En rango uno se obtiene la geometría convexa
completa, con cobertura por símplices y residuos en toda dimensión finita.
Para rangos superiores se construye una colección explícita de imágenes
amplituhedrales y se preservan sus banderas marcadas. La cobertura global
de un amplituedro de rango superior constituye un enunciado de alcance
distinto del teorema de existencia y compatibilidad de esta colección; las
pruebas posteriores conservan esa distinción en su formulación.

La dimensión física se introduce mediante los operadores correspondientes.
El calendario trítico distingue rotación elíptica, apertura hiperbólica y
transporte parabólico de memoria. Sus realizaciones hamiltonianas conservan
la forma simpléctica y permiten calcular la transformación de Legendre y
los términos de conexión originados por una coordenatización variable. En el sector
gauge, la escala areal--volumétrica, la medida común y las formas de
reflexión anteceden a la reconstrucción del Hamiltoniano de Yang--Mills.
El argumento espectral se presenta junto a la coercividad uniforme, el
testigo que alcanza el primer nivel positivo y los entrelazadores que lo
preservan. La relación con las formas canónicas se establece mediante sus
lectores y sus residuos, mientras la brecha pertenece al Hamiltoniano y a
su espacio físico. Esta jerarquía sitúa cada resultado en el lugar donde
su demostración produce efectivamente la conclusión.

La exposición incorpora primero la construcción formal necesaria para
obtener el estado y su registro. Desarrolla después las realizaciones
positivas, la extensión dimensional, la energía y la dualidad reticular;
reúne finalmente las realizaciones hamiltonianas y el desarrollo gauge.
Los programas de comprobación acompañan a identidades determinadas y a
controles negativos explícitos. La demostración simbólica conserva el
alcance general; la ejecución exacta comprueba las especializaciones
declaradas y su correspondencia con las fórmulas impresas.

\subsection{Genealogía aritmética y sistema de realizaciones geométricas}
La construcción empieza por las dos hojas aritméticas de APP. Si
\(p,q\in\{1,\ldots,9\}\), cada operación conserva una división completa,
\[
 p+q=R_\Sigma+9q_\Sigma,\qquad
 pq=R_\Pi+9q_\Pi.
\]
TRIT añade el régimen local y la orientación. El selector trítico de fase
operativa elige la operación compatible con ese estado, mientras ambas hojas
permanecen en el soporte. Sobre el espacio de estados enriquecidos, el
Topological Phase Kernel compone selección, transporte y actualización:
\[
 U_t=\operatorname{Upd}_t\circ\operatorname{Tra}_t\circ\operatorname{Sel}_t.
\]
El registro conserva los residuos, cocientes, acarreos, hojas, rutas,
fronteras, memoria y supervivencia que utiliza cada prolongación. La
holonomía nonádica \(\Gamma_9=U_{t+8}\circ\cdots\circ U_t\) pertenece
a esta dinámica enriquecida. Su representación reducida de memoria se efectúa
después. El retorno de fase conserva el avance del registro y por ello
distingue estados situados a profundidades diferentes.

La prolongación
\[
 w_6\longrightarrow w_{12}\longrightarrow w_{18}\longrightarrow
 w_{24}\longrightarrow w_{30}\longrightarrow R_{36}\longrightarrow G_9
\]
organiza los refinamientos compatibles del mismo estado. Las cinco
construcciones del continuo emergen conjuntamente, con sus aspectos
espectral, solenoidal, gauge, cohomológico y determinantal. Los lectores
de ese estado expresan correlativamente \(\pi,\varphi,e,\alpha\).
Su origen común admite un orden interno de evaluación: las coordenadas
regionales ya generadas actúan junto con el registro dodecafásico en los
lectores de clausura de \(\alpha\). En este artículo, Grassmannianos,
polígonos y formas diferenciales son realizaciones posteriores de datos
generados por esta cadena. Sus parámetros quedan fijados antes de toda
comparación geométrica o metrológica.

El problema concreto consiste en describir lo que conserva una realización
positiva y lo que puede variar al cambiar su forma. Una misma lista positiva
admite varios usos matemáticos: pesos de aristas, longitudes de intervalos,
coeficientes de una forma cuadrática o coordenadas de un lector material.
Cada uso posee un dominio, una operación y unos invariantes propios. La
expresión \emph{multimorfología} designa aquí el sistema de estas realizaciones
con sus mapas explícitos de compatibilidad. Una identidad de sus números
requiere una comprobación menor que una identidad de sus operadores; el
desarrollo siguiente formula y demuestra cada nivel por separado.

La base de procedencia es el artículo de generación coinductiva de las
constantes y estructura algebraica del electrón, junto con el desarrollo
de extrema y media razón holográfica. Sus construcciones anteriores fijan
el registro utilizado abajo. Los resultados geométricos de este texto comienzan
con esa salida finita y prueban sus consecuencias mediante datos explícitos. Las
proposiciones de transferencia material declaran adicionalmente los
lectores que emplean. Esta organización preserva el alcance específico de
cada prueba.

\clearpage
\part{Construcción aritmética, estado enriquecido y continuo conjunto}
\begingroup
\sloppy
\renewcommand{\TRIT}{\{-1,0,+1\}}
\import{foundation/}{sections/nucleo.tex}
% BEGIN NUCLEO_ADITIVO_20260921_CENSO
\clearpage
\section{Núcleo topológico de fase y generación del catálogo de semillas}
\label{act21:tpk:capitulo}

Las dos evaluaciones de APP adquieren eficacia generativa cuando se especifica
qué evaluación actúa, en qué posición se efectúa, cómo se desplaza esa posición
y qué información conserva cada transición. El \emph{Topological Phase Kernel},
abreviado TPK y denominado en español \emph{núcleo topológico de fase}, reúne
estas operaciones en una dinámica sobre estados enriquecidos. Su definición
comprende selección, transporte, actualización aritmética y conservación del
registro. La construcción finita desarrollada en este capítulo proporciona
las semillas sobre las que actuarán las elevaciones regionales y la prolongación
coinductiva.

El resultado principal de esta etapa es un catálogo generado y clasificado:
las condiciones iniciales de los dos cursores producen emisiones enteras;
su lectura ternaria determina palabras orientadas; la acción diedral organiza
esas palabras en órbitas. Cada paso conserva sus preimágenes y las relaciones
que permiten regresar a la descripción precedente. La cadena de cardinales
\[
 104\,976\longrightarrow468\longrightarrow243\subset729,
 \qquad 243\longrightarrow43,
\]
abrevia objetos diferentes. A continuación se construye cada objeto y se
justifica cada recuento, antes de utilizar cualquiera de ellos como dato de
la selección regional.

\subsection{Selección de fase, transporte y estado enriquecido}
\label{act21:tpk:operaciones}

\subsubsection{Selector trítico de fase operativa}

La fase nonádica toma sus valores en \(D_9=\{1,\ldots,9\}\). El selector
trítico es la función
\begin{equation}
 M_{\mathrm{ph}}:D_9\longrightarrow\{-1,0,+1\},\qquad
 M_{\mathrm{ph}}(g)=
 \begin{cases}
 +1,&g\in\{1,4,7\},\\
 -1,&g\in\{2,5,8\},\\
 0,&g\in\{3,6,9\}.
 \end{cases}
 \label{act21:tpk:selector}
\end{equation}
En la realización emisora, el valor \(+1\) activa la evaluación aditiva,
\(-1\) activa la evaluación multiplicativa y \(0\) registra el evento
neutral. La orientación cardinal del cursor es otra coordenada del estado.
Esta separación permite fijar simultáneamente qué operación se lee y en qué
dirección se transporta su punto de evaluación.

Para las transiciones numeradas por \(t\geq1\), la fase previa a la lectura es
\begin{equation}
 g_t=1+[(t-1)]_9,\qquad \epsilon_t=M_{\mathrm{ph}}(g_t),
 \label{act21:tpk:fase}
\end{equation}
donde \([n]_b\) designa el representante de \(n\) módulo \(b\) en
\(\{0,\ldots,b-1\}\). Una ventana de nueve transiciones contiene la sucesión
\((+1,-1,0,+1,-1,0,+1,-1,0)\).

\subsubsection{Cursores orientados y transporte cardinal}

Se utilizan índices \(I_9=\{0,\ldots,8\}\) para las posiciones y etiquetas
positivas \((a,b)=(i+1,j+1)\) para las evaluaciones APP. Un cursor finito es
\[
 c=(i,j;d)\in\mathcal S:=I_9^2\times\mathcal D,
 \qquad \mathcal D=\{N,E,S,O\}.
\]
Las cuatro direcciones actúan mediante
\[
 v_N=(-1,0),\quad v_E=(0,1),\quad
 v_S=(1,0),\quad v_O=(0,-1).
\]
La aplicación cardinal correspondiente a \(d\) es
\[
 T_d(i,j)=\bigl([i+(v_d)_1]_9,[j+(v_d)_2]_9\bigr).
\]
Así, \(M_{\mathrm{ph}}\) y \(T_d\) tienen dominios, codominios y acciones
diferentes. La activación trítica selecciona uno de los dos cursores;
el transporte cardinal modifica su posición en la dirección que ese cursor
ya conserva.

El levantamiento entero del cursor es
\(\widetilde c=(x,y;d)\in\mathbb Z^2\times\mathcal D\), con posición finita
\(([x]_9,[y]_9)\). Sus números de vuelta se recuperan en
\begin{equation}
 x=9\lfloor x/9\rfloor+[x]_9,\qquad
 y=9\lfloor y/9\rfloor+[y]_9.
 \label{act21:tpk:vueltas}
\end{equation}
Estas dos divisiones conciernen al transporte espacial. La división de una
evaluación aditiva o multiplicativa constituye un registro aritmético
distinto, que se conservará junto a ellas.

\subsubsection{Composición de la transición}

El estado de la realización finita incluye los dos cursores levantados,
sus direcciones, la fase, los acumuladores de la ventana activa, la lista
ordenada de bloques ya emitidos y el registro de transiciones. Cada registro
local contiene las posiciones previas, las evaluaciones enteras de ambas
hojas, sus residuos y cocientes, el régimen activo y las orientaciones.
Los campos de prefijo y frontera que aparezcan en una prolongación se
incorporarán al mismo estado con sus reglas de compatibilidad.

La transición se ordena como
\begin{equation}
 \mathcal U_t=\mathsf{Mem}_t\circ\mathsf{Tra}_t\circ\mathsf{Sel}_t.
 \label{act21:tpk:composicion}
\end{equation}
La selección determina el régimen y conserva la muestra previa a todo
desplazamiento. El transporte aplica el movimiento activo y la inversión
cardinal prescrita por el calendario. La actualización incorpora esa muestra
a los acumuladores y al registro; al terminar una ventana, emite su bloque.
La composición utiliza por tanto una muestra anterior al transporte aunque
su archivo definitivo se efectúe al final de la transición.

\subsection{Condiciones iniciales y enumeración reversible}
\label{act21:tpk:semillas}

Las hojas aditiva y multiplicativa disponen de cursores independientes. El
dominio de condiciones iniciales es el producto ordenado
\[
 \Omega_{\mathrm{seed}}=\mathcal S_+\times\mathcal S_\times,
 \qquad |\mathcal S|=9^2\cdot4=324.
\]
Por consiguiente,
\begin{equation}
 |\Omega_{\mathrm{seed}}|=324^2=104\,976.
 \label{act21:tpk:censo-inicial}
\end{equation}
El número \(81\) cuenta posiciones; \(324\) cuenta posiciones orientadas
de un cursor; \(104\,976\) cuenta parejas ordenadas de cursores. La elección
de una posición en APP conserva únicamente una parte de una condición
inicial del emisor.

Se fija la numeración
\[
 \operatorname{ind}(N)=0,\quad\operatorname{ind}(E)=1,\quad
 \operatorname{ind}(S)=2,\quad\operatorname{ind}(O)=3,
\]
y se define
\begin{equation}
 \nu(i,j;d)=4(9i+j)+\operatorname{ind}(d).
 \label{act21:tpk:indice-semilla}
\end{equation}

\begin{proposition}[Enumeración exacta de las condiciones iniciales]
\label{act21:tpk:enumeracion}
La aplicación \(\nu\) es una biyección de \(\mathcal S\) sobre
\(\{0,\ldots,323\}\). La pareja de índices se codifica de manera reversible
por \(324\nu_++\nu_\times\in\{0,\ldots,104\,975\}\).
\end{proposition}
\begin{proof}
La división por cuatro recupera
\[
 \operatorname{ind}(d)=[\nu]_4,\qquad
 j=[\lfloor\nu/4\rfloor]_9,\qquad i=\lfloor\nu/36\rfloor.
\]
Las cotas del dominio garantizan que estas tres cantidades pertenecen a
sus conjuntos de índices y reconstruyen \(\nu\) mediante
\eqref{act21:tpk:indice-semilla}. Para la pareja, la división euclídea por
\(324\) recupera cociente \(\nu_+\) y residuo \(\nu_\times\).
\end{proof}

La enumeración recorre el dominio completo antes de cualquier agrupación
regional. Los índices se utilizan para localizar las preimágenes de una
emisión, manteniendo explícita la condición que produjo cada trayectoria.

\subsection{Evaluación de las hojas y registro aritmético}
\label{act21:tpk:lecturas}

En una posición positiva \((a,b)\), las evaluaciones enteras son
\(A_+(a,b)=a+b\) y \(A_\times(a,b)=ab\). Para \(n>0\), se mantienen
\begin{equation}
 \rho_9(n)=1+[n-1]_9,\qquad
 q_9(n)=\lfloor(n-1)/9\rfloor,\qquad
 n=\rho_9(n)+9q_9(n).
 \label{act21:tpk:division-evaluacion}
\end{equation}
La activación aditiva incorpora \(\rho_9(A_+)\) al acumulador
correspondiente. La activación multiplicativa utiliza un observable del
residuo \(\rho_9(A_\times)\), definido a continuación.

\subsubsection{Exponente discreto en las unidades módulo nueve}

Las potencias de \(2\) módulo \(9\) recorren
\[
 2^0,2^1,\ldots,2^5\equiv1,2,4,8,7,5\pmod9,
 \qquad 2^6\equiv1\pmod9.
\]
El logaritmo discreto en este grupo de seis unidades se representa por
\begin{equation}
 \ell_9(1)=0,\quad\ell_9(2)=1,\quad\ell_9(4)=2,
 \quad\ell_9(8)=3,\quad\ell_9(7)=4,\quad\ell_9(5)=5.
 \label{act21:tpk:logaritmo}
\end{equation}
La ley emisora extiende el observable mediante
\(\ell_9(3)=\ell_9(6)=\ell_9(9)=0\). Para conservar su procedencia, la
muestra registrada incluye el residuo original y el indicador de pertenencia
al grupo de unidades. De ese modo, un valor observado cero puede proceder
de la unidad \(1\) o de alguno de los tres residuos no invertibles, y su
origen permanece recuperable.

\subsubsection{Orden de lectura y movimiento}

Las evaluaciones se efectúan en las posiciones anteriores a la transición.
Si \(\epsilon_t=+1\), se acumula la lectura aditiva y se desplaza su cursor.
Si \(\epsilon_t=-1\), se acumula el exponente discreto multiplicativo y se
desplaza el segundo cursor. Si \(\epsilon_t=0\), se incrementa el contador
neutral y se conservan ambas posiciones.

Al finalizar los pasos \(27\) y \(54\), las dos direcciones se sustituyen
por sus opuestas. Las ventanas se encadenan conservando posiciones y
orientaciones. Al inicio de cada ventana sólo se ponen a cero sus tres
acumuladores locales; el registro anterior permanece disponible.

\begin{proposition}[Inversión del transporte levantado]
\label{act21:tpk:inversion}
Sean \(a,b\in\{0,1\}\) los indicadores de activación e inversión de un
cursor. Si \(\operatorname{opp}\) intercambia \(N\leftrightarrow S\) y
\(E\leftrightarrow O\), la transformación
\[
 F_{a,b}(z,d)=\bigl(z+a v_d,\operatorname{opp}^{\,b}(d)\bigr)
\]
tiene inversa
\[
 F_{a,b}^{-1}(z',d')=
 \bigl(z'-a v_{\operatorname{opp}^{\,b}(d')},
       \operatorname{opp}^{\,b}(d')\bigr).
\]
La reducción espacial módulo nueve entrelaza esta transformación con el
transporte finito.
\end{proposition}
\begin{proof}
Aplicar dos veces la oposición recupera la dirección inicial. Una vez
recuperada esa dirección, la resta de \(a v_d\) deshace el desplazamiento.
Para la proyección espacial basta la igualdad
\([z+a v_d]_9=[[z]_9+a v_d]_9\), componente a componente. La inversión de
dirección conserva esta igualdad.
\end{proof}

La fase determina los indicadores utilizados por la inversa. El registro
espacial y el aritmético pueden así conservarse conjuntamente, incluso al
atravesar el borde del representante \(I_9^2\).

\subsection{Seis ventanas y emisión decimal por bloques}
\label{act21:tpk:emision}

\subsubsection{Acumulación en una ventana nonádica}

La ventana \(m\), con \(1\leq m\leq6\), comprende los pasos
\(9m-8,\ldots,9m\). Sean \(S_m\) la suma de sus residuos positivos
aditivos activos, \(E_m\) la suma de sus exponentes discretos activos y
\(Z_m\) su número de eventos neutrales.

\begin{lemma}[Multiplicidades de activación]
\label{act21:tpk:tres-eventos}
Cada ventana contiene tres activaciones aditivas, tres multiplicativas y
tres neutrales. En particular, \(Z_m=3\).
\end{lemma}
\begin{proof}
Las fases de una ventana recorren exactamente \(1,\ldots,9\). Cada uno
de los tres conjuntos de \eqref{act21:tpk:selector} tiene cardinal tres.
\end{proof}

Las seis ventanas forman \(54\) transiciones y producen seis bloques.
Esta longitud pertenece al emisor inicial aquí definido. La continuación
en profundidad actúa después sobre las palabras y sus registros compatibles.

\subsubsection{Regla de codificación decimal}

La ley emisora especifica las cifras
\begin{align}
 d_{m,1}&=[S_m]_{10},\notag\\
 d_{m,2}&=[S_m+E_m]_{10},\notag\\
 d_{m,3}&=[3S_m+5E_m+7Z_m]_{10},\notag\\
 U_m&=100d_{m,1}+10d_{m,2}+d_{m,3}.
 \label{act21:tpk:bloque-decimal}
\end{align}
Los pesos \(100,10,1\) son las posiciones centena, decena y unidad.
El subíndice \(10\) indica el residuo decimal, y cada bloque satisface
\(0\leq U_m\leq999\). La palabra emitida es la sucesión ordenada
\(\mathbf U=(U_1,\ldots,U_6)\), con ancho fijo de tres cifras por bloque.

Los coeficientes \(3,5,7\) constituyen los pesos explícitos de esta ley
emisora. Una vez fijada la regla, cada cifra se calcula a partir de los
acumuladores, y los resultados de este capítulo se deducen de esa regla.
El número \(1\) que aparece en su forma reducida tiene una procedencia
adicional concreta: \(7Z_m=21\equiv1\pmod{10}\).

Si \(s_m=[S_m]_{10}\) y \(e_m=[E_m]_{10}\), se obtiene
\begin{equation}
 G(s_m,e_m)=100s_m+10[s_m+e_m]_{10}
                   +[3s_m+5e_m+1]_{10}=U_m.
 \label{act21:tpk:acoplamiento-firmas}
\end{equation}
La letra \(e_m\) denota aquí una componente de la firma de exponentes
discretos. Su definición pertenece al emisor finito.

\begin{proposition}[Recuperación de las dos firmas]
\label{act21:tpk:recuperacion-firmas}
La aplicación \(G:\{0,\ldots,9\}^2\to\{0,\ldots,999\}\) es inyectiva.
Para un bloque \(U=100a+10b+c\) de su imagen, la inversa es
\[
 s=a,\qquad e=[b-a]_{10},\qquad
 c=[3s+5e+1]_{10}.
\]
En consecuencia, la palabra \(\mathbf U\) determina las dos firmas de
seis componentes \(\mathbf s\) y \(\mathbf e\).
\end{proposition}
\begin{proof}
Los dos últimos sumandos de \eqref{act21:tpk:acoplamiento-firmas} pertenecen a
\(\{0,\ldots,99\}\), de modo que la centena recupera \(s\). La decena
es \([s+e]_{10}\); al restar \(s\) y reducir módulo diez se recupera
\(e\). La unidad verifica la tercera congruencia. La reconstrucción se
aplica independientemente a las seis posiciones de la palabra.
\end{proof}

Los totales enteros \(S_m,E_m\) se recuperan de las lecturas registradas;
las centenas y decenas recuperan sus residuos. Por ejemplo, un total
\(S_m=15\) produce \(s_m=5\). Esta diferencia identifica exactamente
qué información conserva cada nivel de codificación.

\subsubsection{Codificación posicional de la palabra de seis bloques}

Los seis bloques admiten una codificación entera de ancho fijo,
\begin{equation}
 \mathscr N(\mathbf U)=\sum_{m=1}^{6}U_m1000^{6-m},
 \qquad0\leq\mathscr N(\mathbf U)<1000^6.
 \label{act21:tpk:entero-base1000}
\end{equation}
El primer bloque ocupa la posición de mayor peso y el sexto, la de unidades
de base mil. Un bloque \(058\) ocupa una posición completa con valor
entero \(58\); su ancho de tres cifras forma parte de la presentación
decimal de esa posición.

\begin{proposition}[Recuperación posicional de las emisiones]
\label{act21:tpk:inversa-base1000}
La codificación \eqref{act21:tpk:entero-base1000} determina la palabra completa
mediante
\[
 U_m=\left[\left\lfloor
 \frac{\mathscr N(\mathbf U)}{1000^{6-m}}\right\rfloor\right]_{1000},
 \qquad1\leq m\leq6.
\]
\end{proposition}
\begin{proof}
Al dividir por \(1000^{6-m}\), los bloques anteriores producen múltiplos
enteros de \(1000\), el bloque \(m\) produce \(U_m\), y la contribución
de los posteriores pertenece a \([0,1)\). La parte entera elimina esta
última contribución y el residuo módulo \(1000\) elimina las primeras.
\end{proof}

Esta codificación conserva las seis emisiones. Su dominio posee posiciones
ordenadas de base mil, mientras el dominio APP posee posiciones espaciales
módulo nueve. La composición mantiene ambos índices: \(m\) identifica una
ventana de emisión y \((i,j)\) identifica la celda visitada durante una
transición de esa ventana.

\subsection{Cálculo completo de una emisión que comienza por 501}
\label{act21:tpk:ejemplo501}

Se fijan las condiciones iniciales
\[
 c_{+,0}=(0,4;N),\qquad c_{\times,0}=(2,3;N),
 \qquad \nu_+=16,\quad\nu_\times=84.
\]
Las posiciones positivas iniciales son \((1,5)\) y \((3,4)\). El índice
cronológico previo es \(t=0\), la primera fase leída es \(g_1=1\) y el
registro está vacío. Estas coordenadas son las entradas
del cálculo; los bloques se obtienen ejecutando la recurrencia.

\begin{table}[htbp]
\caption{Primera ventana: lecturas previas y acumulación activa.}
\label{act21:tpk:tabla-501}
\small
\begin{tabular}{@{}rcccrrr@{}}
\toprule
\(t\)&\(\epsilon_t\)&posición \(+\)&posición \(\times\)&
\(\Delta S\)&\(\Delta E\)&\((S,E,Z)\)\\
\midrule
1&\(+1\)&\((1,5)\)&\((3,4)\)&6&0&\((6,0,0)\)\\
2&\(-1\)&\((9,5)\)&\((3,4)\)&0&0&\((6,0,0)\)\\
3&0&\((9,5)\)&\((2,4)\)&0&0&\((6,0,1)\)\\
4&\(+1\)&\((9,5)\)&\((2,4)\)&5&0&\((11,0,1)\)\\
5&\(-1\)&\((8,5)\)&\((2,4)\)&0&3&\((11,3,1)\)\\
6&0&\((8,5)\)&\((1,4)\)&0&0&\((11,3,2)\)\\
7&\(+1\)&\((8,5)\)&\((1,4)\)&4&0&\((15,3,2)\)\\
8&\(-1\)&\((7,5)\)&\((1,4)\)&0&2&\((15,5,2)\)\\
9&0&\((7,5)\)&\((9,4)\)&0&0&\((15,5,3)\)\\
\bottomrule
\end{tabular}
\notatabla{Las posiciones son etiquetas positivas de las celdas anteriores
al movimiento. Los eventos neutrales incrementan únicamente \(Z\).}
\end{table}

Las tres evaluaciones aditivas activas son
\[
 1+5=6,\qquad9+5=14=5+9,\qquad8+5=13=4+9,
\]
de modo que \(S_1=6+5+4=15\). Las tres evaluaciones multiplicativas activas
son \(3\cdot4=12\), \(2\cdot4=8\) y \(1\cdot4=4\). Sus residuos
positivos son \(3,8,4\), con observables \(0,3,2\); por tanto,
\(E_1=5\). Finalmente \(Z_1=3\). La regla decimal da
\[
 d_{1,1}=[15]_{10}=5,\quad
 d_{1,2}=[15+5]_{10}=0,\quad
 d_{1,3}=[45+25+21]_{10}=1,
\]
y se obtiene \(U_1=501\).

La segunda lectura aditiva anterior tiene cociente aritmético \(1\),
mientras su cursor levantado está en \((-1,4;N)\), con número de vuelta
vertical \(-1\). Son dos datos diferentes del mismo registro. Ambos
intervienen en la recuperación íntegra de la evaluación y de su posición.

\begin{table}[htbp]
\caption{Las seis ventanas de la condición inicial \((16,84)\).}
\label{act21:tpk:tabla-seis501}
\begin{tabular}{@{}rrrrrrrr@{}}
\toprule
\(m\)&pasos&\(S_m\)&\(E_m\)&\(Z_m\)&\(s_m\)&\(U_m\)&\([-U_m]_3\)\\
\midrule
1&1--9&15&5&3&5&501&0\\
2&10--18&6&5&3&6&614&1\\
3&19--27&24&5&3&4&498&0\\
4&28--36&21&5&3&1&169&2\\
5&37--45&12&5&3&2&272&1\\
6&46--54&12&5&3&2&272&1\\
\bottomrule
\end{tabular}
\end{table}

La palabra resultante y sus firmas son
\begin{align*}
 \mathbf U&=(501,614,498,169,272,272),\\
 \mathbf s&=(5,6,4,1,2,2),&
 \mathbf e&=(5,5,5,5,5,5).
\end{align*}
La inversión cardinal tras el paso \(27\) modifica los recorridos de las
tres ventanas finales. El bloque repetido \(272\) corresponde a dos
ventanas sucesivas con posiciones y registro cronológico propios.

Como segundo caso, la condición inicial mínima \((\nu_+,\nu_\times)=(0,0)\)
produce
\[
 \mathbf U=(252,109,252,903,850,850),\quad
 \mathbf s=(2,1,2,9,8,8),\quad
 \mathbf e=(3,9,3,1,7,7).
\]
En este caso, ambos cursores regresan a su posición y orientación iniciales
tras \(54\) transiciones. La fase nonádica completa seis vueltas, el
contador cronológico avanza \(54\) unidades y el registro contiene
\(54\) muestras. La lectura posicional del retorno y
la historia registrada describen propiedades distintas de la ejecución.

\subsection{Factorización del censo finito}
\label{act21:tpk:censo}

Las firmas \(f_+(\nu)=\mathbf s\) y \(f_\times(\nu)=\mathbf e\) se
calculan ejecutando el cursor correspondiente durante las seis ventanas.
La independencia de posiciones y direcciones entre los dos cursores da
\begin{equation}
 T_{\mathrm{em}}(\nu_+,\nu_\times)
   =G^{(6)}\bigl(f_+(\nu_+),f_\times(\nu_\times)\bigr),
 \label{act21:tpk:factorizacion}
\end{equation}
donde \(G^{(6)}\) aplica \eqref{act21:tpk:acoplamiento-firmas} en cada posición.
La factorización permite enumerar primero \(324+324\) trayectorias y
combinar después las firmas distintas, conservando la multiplicidad de
sus preimágenes.

\par\noindent\begin{minipage}{\linewidth}
\subsubsection{Imágenes de las dos firmas}

Las tablas siguientes incluyen todas las firmas. Su multiplicidad es el
número de condiciones iniciales que las producen; el índice mínimo localiza
una de esas condiciones. La preimagen completa se define por
\[
 f_\sigma^{-1}(a)=\{\nu\in\{0,\ldots,323\}:f_\sigma(\nu)=a\}.
\]
La recurrencia de las secciones anteriores determina cada elemento de este
conjunto mediante un cálculo entero finito.

\begin{table}[H]
\centering\fontsize{9.5}{11.4}\selectfont
\setlength{\abovecaptionskip}{8pt}\setlength{\belowcaptionskip}{6pt}
\renewcommand{\arraystretch}{1.12}
\caption{Las dieciocho firmas aditivas.}
\label{act21:tpk:firmas-aditivas}
\begin{tabular}{@{}lrr@{\qquad}lrr@{}}
\toprule
firma&multiplicidad&\(\nu_{\min}\)&firma&multiplicidad&\(\nu_{\min}\)\\
\midrule
\texttt{122564}&18&17&\texttt{564122}&18&16\\
\texttt{122898}&18&24&\texttt{645212}&18&4\\
\texttt{212645}&18&5&\texttt{654889}&18&33\\
\texttt{212988}&18&0&\texttt{889221}&18&13\\
\texttt{221456}&18&29&\texttt{889654}&18&32\\
\texttt{221889}&18&12&\texttt{898122}&18&25\\
\texttt{456221}&18&28&\texttt{898465}&18&20\\
\texttt{465898}&18&21&\texttt{988212}&18&1\\
\texttt{546988}&18&9&\texttt{988546}&18&8\\
\bottomrule
\end{tabular}

\par\vspace{12pt}

\caption{Las veintiséis firmas multiplicativas.}
\label{act21:tpk:firmas-multiplicativas}
\begin{tabular}{@{}lrr@{\qquad}lrr@{}}
\toprule
firma&multiplicidad&\(\nu_{\min}\)&firma&multiplicidad&\(\nu_{\min}\)\\
\midrule
\texttt{000000}&108&8&\texttt{555555}&24&84\\
\texttt{177393}&8&1&\texttt{555717}&8&14\\
\texttt{177555}&8&54&\texttt{555771}&8&18\\
\texttt{177771}&8&11&\texttt{555933}&8&24\\
\texttt{339555}&8&6&\texttt{717555}&8&12\\
\texttt{339771}&8&15&\texttt{717717}&8&21\\
\texttt{339933}&8&59&\texttt{717933}&8&19\\
\texttt{393177}&8&0&\texttt{771177}&8&9\\
\texttt{393393}&8&45&\texttt{771339}&8&13\\
\texttt{393555}&8&40&\texttt{771555}&8&16\\
\texttt{555177}&8&52&\texttt{933339}&8&57\\
\texttt{555339}&8&4&\texttt{933555}&8&26\\
\texttt{555393}&8&41&\texttt{933717}&8&17\\
\bottomrule
\end{tabular}
\end{table}
\end{minipage}\par


\begin{proposition}[Imagen y multiplicidades del emisor]
\label{act21:tpk:468}
La imagen \(\mathcal U_6=\operatorname{im}T_{\mathrm{em}}\) contiene
\(468\) palabras distintas. Sus preimágenes se distribuyen en \(432\)
fibras de cardinal \(144\), \(18\) de cardinal \(432\) y \(18\) de
cardinal \(1944\).
\end{proposition}
\begin{proof}
La evaluación exhaustiva del algoritmo de la sección~\ref{act21:tpk:algoritmo}
sobre los índices \(0,\ldots,323\) da las dos tablas anteriores.
La primera partición tiene \(18\) clases de cardinal \(18\), y la segunda
\(24\) clases de cardinal \(8\), una de cardinal \(24\) y una de
cardinal \(108\). Las sumas \(18\cdot18=324\) y
\(24\cdot8+24+108=324\) comprueban la masa total de ambas particiones.
La proposición~\ref{act21:tpk:recuperacion-firmas} hace inyectivo el acoplamiento
de sus imágenes; por tanto, hay \(18\cdot26=468\) emisiones.
Por \eqref{act21:tpk:factorizacion}, la preimagen de una pareja de firmas es
el producto de sus dos preimágenes. Se obtienen \(18\cdot24=432\)
productos de cardinal \(18\cdot8=144\), \(18\) de cardinal
\(18\cdot24=432\) y \(18\) de cardinal \(18\cdot108=1944\).
Finalmente,
\[
 432\cdot144+18\cdot432+18\cdot1944=104\,976.
\]
La parte enumerativa es una comprobación finita exacta de la recurrencia;
la factorización explica por qué ese recuento cubre todas las parejas.
\end{proof}

La emisión del ejemplo \ref{act21:tpk:ejemplo501} tiene una fibra de cardinal
\(18\cdot24=432\). El ejemplo de índices \((0,0)\) tiene una fibra de
cardinal \(18\cdot8=144\). En ambos casos la palabra emitida recupera
las firmas, mientras la identificación de una condición inicial particular
utiliza su índice o su registro de trayectoria.

\subsection{Lectura ternaria orientada y conservación de las fibras}
\label{act21:tpk:reduccion-ternaria}

La reducción orientada actúa separadamente sobre cada bloque:
\begin{equation}
 \Psi:\mathcal U_6\longrightarrow\mathbb F_3^6,\qquad
 \Psi(U_1,\ldots,U_6)=([-U_1]_3,\ldots,[-U_6]_3).
 \label{act21:tpk:psi}
\end{equation}
Su signo fija una orientación del catálogo. La identificación balanceada
\(0\mapsto0\), \(1\mapsto+1\), \(2\mapsto-1\) se aplica después,
componente a componente.

El módulo nueve organiza las evaluaciones y posiciones de APP; los residuos
módulo diez forman cada cifra; la base mil reúne tres posiciones decimales;
el módulo tres determina la firma ternaria de cada bloque. Estos cuatro usos
tienen objetos y funciones definidos. En particular, \(\Psi\) lee seis
bloques y produce seis trits; una conversión posicional del entero formado
por dieciocho cifras sería otra operación.

Para las dos condiciones calculadas se obtiene
\begin{align*}
 \Psi(501,614,498,169,272,272)&=(0,1,0,2,1,1),\\
 \Psi(252,109,252,903,850,850)&=(0,2,0,0,2,2).
\end{align*}
El emisor decimal está definido antes de esta reducción. La función de
\(\Psi\) es producir el objeto ternario sobre el que actúan la simetría
orbital y los selectores regionales. Esa función explica su posición en la
genealogía: enlaza una emisión ya calculada con una clasificación estructural.

La imagen \(\mathcal W_6:=\Psi(\mathcal U_6)\) contiene exactamente
\(243\) palabras dentro del conjunto ambiente de \(3^6=729\) palabras.
La enumeración del algoritmo final da el histograma
\begin{equation}
\begin{array}{c|rrrrrr}
 |\Psi^{-1}(w)|&1&2&3&4&5&6\\\hline
 \text{número de palabras }w&132&66&18&3&6&18.
\end{array}
\label{act21:tpk:fibras-ternarias}
\end{equation}
Las sumas de ambas lecturas son
\[
132+66+18+3+6+18=243,
\qquad132+2\cdot66+3\cdot18+4\cdot3+5\cdot6+6\cdot18=468.
\]
La primera cuenta palabras ternarias; la segunda recupera el número de
emisiones distribuidas entre sus fibras.

\begin{example}[Dos emisiones con la misma firma ternaria]
\label{act21:tpk:colision}
Las emisiones
\[
 (498,501,614,252,252,109),\qquad
 (870,810,923,252,252,109)
\]
son distintas y ambas tienen imagen \((0,0,1,0,0,2)\) por \(\Psi\).
Cada una tiene \(144\) condiciones iniciales. La palabra ternaria conserva
una clase de lectura; las emisiones y sus preimágenes distinguen los datos
que confluyen en esa clase.
\end{example}

La conservación de fibras se expresa concretamente mediante
\[
 \mathcal F_w=\{\mathbf U\in\mathcal U_6:\Psi(\mathbf U)=w\},
 \qquad
 \Omega_{\mathbf U}=T_{\mathrm{em}}^{-1}(\mathbf U).
\]
Por tanto, la preimagen completa de \(w\) en el dominio inicial es la
unión disjunta de los conjuntos \(\Omega_{\mathbf U}\), con
\(\mathbf U\in\mathcal F_w\). El registro enriquecido puede conservar la
emisión, sus firmas y la condición inicial junto a su lectura ternaria.
La inversa exacta se refiere entonces al dato registrado con su tipo;
la palabra de seis trits aislada mantiene el alcance preciso de
\eqref{act21:tpk:psi}.

\subsection{Acción diedral y las cuarenta y tres órbitas}
\label{act21:tpk:orbitas}

Sobre seis coordenadas se definen las permutaciones
\begin{align}
 r(u_1,u_2,u_3,u_4,u_5,u_6)
   &=(u_3,u_1,u_2,u_5,u_6,u_4),\\
 s(u_1,u_2,u_3,u_4,u_5,u_6)
   &=(u_4,u_5,u_6,u_1,u_2,u_3).
 \label{act21:tpk:generadores-diedrales}
\end{align}
La primera rota los dos grupos de tres coordenadas con sentidos conjugados;
la segunda intercambia ambos grupos. Su composición satisface
\(r^3=s^2=1\) y \(srs=r^{-1}\). Generan así una acción del grupo diedral
\(D_3\), de seis elementos.

\begin{proposition}[Equivariancia y estabilidad del catálogo]
\label{act21:tpk:equivarianza}
La reducción \(\Psi\) conmuta con la acción de \(D_3\). Los conjuntos
\(\mathcal U_6\) y \(\mathcal W_6\) son estables bajo sus generadores.
\end{proposition}
\begin{proof}
La igualdad \(\Psi(g\mathbf U)=g\Psi(\mathbf U)\) se verifica en cada
coordenada porque \(g\) permuta posiciones y \(\Psi\) aplica la misma
reducción a cada posición. La estabilidad del catálogo constituye otra
comprobación: al generar sus \(468\) elementos mediante
\eqref{act21:tpk:factorizacion}, se verifica que \(r\mathbf U\) y
\(s\mathbf U\) pertenecen de nuevo a ese conjunto. El algoritmo de la
sección~\ref{act21:tpk:algoritmo} efectúa estas comprobaciones de pertenencia.
La equivariancia transmite después la estabilidad a su imagen ternaria.
\end{proof}

\begin{proposition}[Descomposición orbital completa]
\label{act21:tpk:43}
El conjunto \(\mathcal W_6\) se descompone en \(43\) órbitas: \(38\)
de cardinal seis y \(5\) de cardinal tres.
\end{proposition}
\begin{proof}
Para cada \(w\in\mathcal W_6\) se forma
\[
 \mathcal O(w)=\{w,rw,r^2w,sw,srw,sr^2w\}.
\]
Se escoge como representante la palabra lexicográficamente mínima de ese
conjunto. Dos palabras tienen el mismo representante exactamente cuando
pertenecen a la misma órbita. El algoritmo genera la lista de representantes
de la tabla~\ref{act21:tpk:representantes}; sus cardinales suman
\(38\cdot6+5\cdot3=243\). Como cada palabra del catálogo participa en
esta operación y el catálogo es estable, la lista cubre toda la imagen.
\end{proof}

\begin{table}[htbp]
\caption{Representantes mínimos y cardinales de las cuarenta y tres órbitas.}
\label{act21:tpk:representantes}
\small
\begin{tabular}{@{}lr@{\qquad}lr@{\qquad}lr@{}}
\toprule
representante&cardinal&representante&cardinal&representante&cardinal\\
\midrule
\texttt{001002}&6&\texttt{002112}&6&\texttt{012222}&6\\
\texttt{001011}&6&\texttt{002211}&6&\texttt{021022}&6\\
\texttt{001012}&6&\texttt{002220}&6&\texttt{021121}&6\\
\texttt{001021}&6&\texttt{011021}&6&\texttt{021202}&6\\
\texttt{001022}&6&\texttt{011112}&6&\texttt{022022}&3\\
\texttt{001100}&3&\texttt{011120}&6&\texttt{022112}&6\\
\texttt{001101}&6&\texttt{011201}&6&\texttt{022121}&6\\
\texttt{001110}&6&\texttt{011211}&6&\texttt{022202}&3\\
\texttt{001112}&6&\texttt{011220}&6&\texttt{022211}&6\\
\texttt{001202}&6&\texttt{011222}&6&\texttt{022222}&6\\
\texttt{001210}&6&\texttt{012112}&6&\texttt{112112}&3\\
\texttt{001211}&6&\texttt{012121}&6&\texttt{112211}&3\\
\texttt{001220}&6&\texttt{012202}&6&\texttt{112222}&6\\
\texttt{001222}&6&\texttt{012210}&6&&\\
\texttt{002012}&6&\texttt{012220}&6&&\\
\bottomrule
\end{tabular}
\end{table}

Las órbitas heredan invariantes enteros de sus fibras. Para cada palabra,
\(n(w)=|\mathcal F_w|\) cuenta emisiones, mientras
\[
 M(w)=\sum_{\mathbf U\in\mathcal F_w}|\Omega_{\mathbf U}|
\]
cuenta condiciones iniciales. El censo de símbolos
\(c(w)=(n_0,n_1,n_2)\) registra cuántas veces aparece cada clase ternaria.
Estos datos, junto con la orientación y el soporte de los recorridos,
proporcionan los predicados de la selección regional posterior. La
clasificación conserva así más información que la forma de una palabra
representante.

\subsection{Algoritmo autocontenido de enumeración y control}
\label{act21:tpk:algoritmo}

La enumeración utiliza enteros, listas finitas y las operaciones ya definidas.
El siguiente procedimiento precisa su ejecución; \(\sigma\) toma los valores
\(+\) y \(\times\).

\begin{enumerate}
\item Para cada \(\nu=0,\ldots,323\), recuperar \((i,j;d)\) mediante
la inversa de \eqref{act21:tpk:indice-semilla}. Iniciar una lista vacía de firmas.
\item Para cada ventana \(m=1,\ldots,6\), iniciar el acumulador \(A=0\).
Recorrer \(t=9m-8,\ldots,9m\), calculando \(\epsilon_t\) por
\eqref{act21:tpk:fase}.
\item Si \(\sigma=+\) y \(\epsilon_t=+1\), añadir
\(\rho_9(i+j+2)\) a \(A\) y actualizar
\((i,j)\leftarrow T_d(i,j)\). Si \(\sigma=\times\) y
\(\epsilon_t=-1\), añadir
\(\ell_9(\rho_9((i+1)(j+1)))\) y efectuar el mismo transporte.
\item Al terminar \(t=27\) o \(t=54\), sustituir \(d\) por su opuesta.
Al concluir la ventana, añadir \([A]_{10}\) a la lista de firmas.
Conservar la posición y la dirección para la siguiente ventana.
\item Agrupar los \(324\) índices por igualdad de sus listas de seis
componentes. Estas dos agrupaciones dan \(f_+\), \(f_\times\) y todas
sus preimágenes.
\item Para cada pareja de firmas distintas, aplicar \(G^{(6)}\).
Asignar a la emisión la multiplicidad producto de las dos preimágenes.
Comprobar la recuperación de firmas en cada bloque y la suma de
multiplicidades \(104\,976\).
\item Aplicar \(\Psi\) a cada emisión y agrupar por igualdad de palabras.
Contar los \(243\) valores y las multiplicidades de
\eqref{act21:tpk:fibras-ternarias}.
\item Aplicar \(r,s\) a cada emisión y verificar pertenencia al catálogo.
Para cada palabra ternaria, formar sus seis imágenes diedrales, eliminar
repeticiones y registrar su mínimo lexicográfico. Contar las \(43\)
órbitas y los dos tamaños de la tabla~\ref{act21:tpk:representantes}.
\end{enumerate}

La ejecución termina: el primer tramo evalúa \(648\) trayectorias de
\(54\) transiciones; el segundo combina \(468\) parejas de firmas; el
tercero agrupa un conjunto de \(243\) palabras. Una comprobación directa
adicional puede ejecutar las \(104\,976\) parejas y contrastar la
factorización \eqref{act21:tpk:factorizacion}; su alcance coincide con el mismo
dominio finito.

\subsection{Conservación del registro y función de las semillas}
\label{act21:tpk:conclusion}

Sea \(\mathsf{Sample}(q_t)\) la muestra previa a la transición y
\(\mathcal M_t\) el registro acumulado. La actualización de memoria es
\begin{equation}
 \mathcal M_{t+1}=\mathcal M_t\mathbin{\Vert}\mathsf{Sample}(q_t),
 \label{act21:tpk:archivo}
\end{equation}
donde \(\Vert\) denota concatenación ordenada. La inducción da
\[
 \mathcal M_n=\mathcal M_0\mathbin{\Vert}
 (\mathsf{Sample}(q_0),\ldots,\mathsf{Sample}(q_{n-1})).
\]
Los acumuladores se reconstruyen sumando las contribuciones activas de
cada segmento de nueve registros. Sus cocientes, las vueltas de los
cursores y las orientaciones permanecen asociados a las mismas transiciones.
La ejecución determina ese registro mediante la recurrencia; el archivo
preserva la procedencia concreta de las lecturas.

Al finalizar esta etapa están construidos el dominio completo de condiciones
iniciales, la emisión decimal, sus dos firmas recuperables, la imagen ternaria,
sus fibras y su clasificación orbital. Los censos \(104\,976\), \(468\),
\(243\), \(729\) y \(43\) tienen ahora dominio y mecanismo de obtención.
Las selecciones regionales recibirán estas estructuras con sus invariantes,
y las elevaciones actuarán sobre palabras orientadas con procedencia
conservada. Esa continuidad convierte el catálogo finito en antecedente
constructivo del desarrollo coinductivo.

% END NUCLEO_ADITIVO_20260921_CENSO
\import{foundation/}{sections/extension.tex}
\import{foundation/}{sections/continuo_conjunto.tex}
\import{foundation/}{sections/generacion.tex}
% BEGIN NUCLEO_ADITIVO_20260921_SELECCION
\clearpage
% Ampliación demostrativa DF003--DF007. No modifica la fuente integral.
% Propietario común de los archivos Lean citados en estos comentarios:
% output/PAQUETE_ARTICULO_I_INTEGRACION_ESTADO_CAMPOS_20260919/.
% Residencia causal de la primera sección: después de los lectores regionales
% coinductivos y antes de emplear sus registros como entrada de la selección K.
% Residencia causal de la segunda sección: después de la construcción marcada
% Paley--Witt--Golay--Leech y antes de la reconstrucción de la estructura VOA.
% Los dos cortes expositivos conservan el mismo origen APP--TRIT--TPK.

\section[Generación regional y refinamiento]{Compatibilidad integral de la generación regional, el refinamiento cilíndrico y la incidencia}
\label{act21:sec:df-lectores-cilindros}

La prolongación regional produce sucesiones ordenadas de bloques, junto con
sus profundidades de lectura. La proyección arquimediana y la incidencia
ternaria se calculan sobre esas mismas sucesiones. El vínculo entre ambas
realizaciones puede formularse enteramente mediante operaciones enteras:
selección por intervalos racionales, división euclídea, elevación de seis
trits, lectura de acarreo y transformación de Witt. Esta formulación permite
seguir cada coeficiente desde su bloque generador hasta su empleo posterior.

La construcción utiliza los tres lectores regionales ya obtenidos del estado
APP--TRIT--TPK. Los denotaremos por los índices
\(c\in\{\mathrm P,\mathrm E,\mathrm F\}\), correspondientes, respectivamente,
a las regiones de clausura, propagación y autoescala. Su realización real se
escribirá \(v_c\). Las identificaciones posteriores
\(v_{\mathrm P}=\pi_{\mathrm{HMT}}\),
\(v_{\mathrm E}=e_{\mathrm{HMT}}\) y
\(v_{\mathrm F}=\varphi_{\mathrm{HMT}}\) conservan ese orden causal.
La operación que emite un bloque inspecciona los intervalos racionales del
lector regional; el valor real interviene en la demostración de corrección
de dicha operación.

\subsection{Emisión posicional mediante separación racional de celdas}
\label{act21:subsec:df-emision-racional}

% Fuente: lean/biblioteca/RegionalPublicationComposition.lean,
% produced_brackets, eventual_publication, joint_eventual_publication,
% search_terminates, stopped_cell_correct, publications_compatible.

Para cada región se dispone de dos sucesiones racionales
\(\ell_c(j),u_c(j)\), construidas por su lector, que cumplen
\[
 \ell_c(j)\leq v_c\leq u_c(j).
\]
La propiedad de separación de celdas demostrada para esos lectores afirma
que, para cada entero \(s>0\), existe \(J\) tal que, para todo \(j\geq J\),
\begin{equation}
 \lfloor s\ell_c(j)\rfloor
 =\lceil su_c(j)\rceil-1
 =\lfloor sv_c\rfloor.
 \label{act21:eq:df-separacion-celdas}
\end{equation}
La evitación de las fronteras racionales por los valores regionales es la
hipótesis de los teoremas previos que garantiza esta separación. Aquí se
emplea su conclusión operativa, con los lectores y los intervalos que la
producen.

\begin{definition}[Índice de separación y prefijo emitido]
Para \(s>0\), sea
\[
 j_c(s)=\min\{j\in\mathbb N:
 \lfloor s\ell_c(j)\rfloor=\lceil su_c(j)\rceil-1\}.
\]
Para una base entera \(b\geq2\) y una profundidad \(n\geq0\), se define
\begin{equation}
 P_c(b,n)=\lfloor b^n\ell_c(j_c(b^n))\rfloor.
 \label{act21:eq:df-publicacion-prefijo}
\end{equation}
El prefijo contiene la parte entera y las primeras \(n\) posiciones en base
\(b\); por consiguiente, \(P_c(b,0)\) es la parte entera regional.
\end{definition}

\begin{theorem}[Terminación y corrección de la emisión]
\label{act21:thm:df-emision-correcta}
El índice \(j_c(s)\) existe para todo \(s>0\). Los prefijos emitidos cumplen
\begin{equation}
 P_c(b,n)=\lfloor b^n v_c\rfloor,
 \qquad
 P_c(b,n)\leq b^n v_c<P_c(b,n)+1.
 \label{act21:eq:df-prefijo-correcto}
\end{equation}
Para cada escala \(s\) puede elegirse además un único índice de suficiencia
\(J(s)\) válido simultáneamente para las tres regiones.
\end{theorem}

\begin{proof}
La igualdad eventual de \eqref{act21:eq:df-separacion-celdas} hace no vacío el
conjunto que define \(j_c(s)\). Su mínimo existe por buen orden de
\(\mathbb N\). En cualquier índice donde se produzca la igualdad, la cota
inferior racional proporciona
\(\lfloor s\ell_c(j)\rfloor\leq\lfloor sv_c\rfloor\).
La cota superior, junto con la evitación de una frontera de celda, da
\(\lfloor sv_c\rfloor\leq\lceil su_c(j)\rceil-1\).
Al coincidir ambos extremos enteros, el entero intermedio coincide con
ellos. Esto demuestra \eqref{act21:eq:df-prefijo-correcto}. Para la afirmación
conjunta basta tomar el máximo de los tres índices de suficiencia obtenidos
para la misma escala. El índice conjunto controla simultáneamente las tres
emisiones, mientras que cada región conserva sus propios intervalos.
\end{proof}

\begin{proposition}[Compatibilidad exacta de los prefijos]
\label{act21:prop:df-prefijos-compatibles}
Para \(n,m\geq0\),
\begin{equation}
 \left\lfloor\frac{P_c(b,n+m)}{b^m}\right\rfloor=P_c(b,n).
 \label{act21:eq:df-truncamiento}
\end{equation}
En particular, el bloque siguiente
\(d_c(b,n)=P_c(b,n+1)\bmod b\) satisface
\begin{equation}
 P_c(b,n+1)=bP_c(b,n)+d_c(b,n),\qquad 0\leq d_c(b,n)<b.
 \label{act21:eq:df-paso-posicional}
\end{equation}
\end{proposition}

\begin{proof}
Escribamos \(b^n v_c=N+\theta\), con
\(N=P_c(b,n)\) y \(0\leq\theta<1\). Entonces
\(P_c(b,n+m)=b^mN+\lfloor b^m\theta\rfloor\), y el segundo sumando
pertenece a \(\{0,\ldots,b^m-1\}\). La división euclídea por \(b^m\)
recupera \(N\). La especialización \(m=1\) produce
\eqref{act21:eq:df-paso-posicional} y su cota residual.
\end{proof}

\subsection{Bloques regionales de seis trits a profundidad arbitraria}
\label{act21:subsec:df-bloques-arbitrarios}

% Fuente: deltas/integracion/JointRegionalFrontier.lean,
% publish_base729_eq_base3, generated_block_from_longer_prefix,
% generated_block_eq_finite, generated_signature_eq_finite.

La igualdad entera \(729=3^6\) determina la agrupación de seis posiciones
ternarias en un bloque. Para \(t\geq0\), definimos
\begin{equation}
 u_c(t)=P_c(729,t+1)\bmod729,
 \qquad
 B_{c,j}(t)=
 \left\lfloor\frac{u_c(t)}{3^{5-j}}\right\rfloor\bmod3,
 \quad 0\leq j<6.
 \label{act21:eq:df-bloque-y-banda}
\end{equation}
De este modo, \(B_c(t)\) es una palabra ordenada de seis trits y
\begin{equation}
 u_c(t)=\sum_{j=0}^{5}B_{c,j}(t)3^{5-j}.
 \label{act21:eq:df-elevacion-banda}
\end{equation}
La palabra y su elevación entera son dos representaciones mutuamente
recuperables del mismo bloque emitido.

\begin{theorem}[Estabilidad de extracción a cualquier profundidad]
\label{act21:thm:df-bloque-estable}
Se cumplen, para todo \(n\),
\[
 P_c(729,n)=P_c(3,6n).
\]
Si \(t<n\), entonces
\begin{equation}
 u_c(t)=
 \left\lfloor\frac{P_c(729,n)}{729^{n-t-1}}\right\rfloor\bmod729.
 \label{act21:eq:df-bloque-prefijo-largo}
\end{equation}
En consecuencia, cualquier ampliación del prefijo preserva todos los
bloques previamente emitidos y todas las firmas que se calculan sobre ellos.
\end{theorem}

\begin{proof}
La primera igualdad resulta de que ambos prefijos se separan a la misma
escala, \(729^n=3^{6n}\), y del
teorema~\ref{act21:thm:df-emision-correcta}. Para la segunda aplicamos
\eqref{act21:eq:df-truncamiento} a las profundidades \(n\) y \(t+1\), y después
tomamos el residuo módulo \(729\). La extracción ternaria de
\eqref{act21:eq:df-bloque-y-banda} es determinista. Cualquier función explícita
de esas dieciocho coordenadas conserva, por tanto, el mismo resultado al
ampliar el prefijo.
\end{proof}

La construcción está cuantificada sobre \(t\in\mathbb N\). Una tabla que
materialice cien bloques constituye una evaluación finita de esta
construcción; la fórmula \eqref{act21:eq:df-bloque-prefijo-largo} determina
exactamente su compatibilidad con cualquier extensión posterior.

\begin{example}[Primer bloque conjunto]
La evaluación de los lectores regionales seleccionados da las palabras
\[
 B_{\mathrm P}(0)=010211,\qquad
 B_{\mathrm E}(0)=201101,\qquad
 B_{\mathrm F}(0)=121200.
\]
Su elevación por \eqref{act21:eq:df-elevacion-banda} produce
\[
 u_{\mathrm P}(0)=103,\qquad
 u_{\mathrm E}(0)=523,\qquad
 u_{\mathrm F}(0)=450.
\]
Por ejemplo,
\(201101_3=2\cdot243+27+9+1=523\).
Las seis posiciones, incluidos los ceros iniciales, se conservan en la
representación de banda. El cero inicial de \(010211\) tiene una posición
definida aunque la escritura decimal de \(103\) carezca de él.
\end{example}

\subsection{Registros de transición y reconstrucción de las elevaciones lineales}
\label{act21:subsec:df-registros-transicion}

% Fuentes: deltas/integracion/RegionalTransitionRecords.lean y
% GeneratedTransitionRecords.lean; lean/biblioteca/TPKLifts.lean.
% Los enlaces reconstruidos son R(0)->R(1) y R(2)->R(3).

Sobre \(\mathbb F_3\), reunimos dos emisiones consecutivas de cada región
en la matriz
\begin{equation}
 R(t)=
 \begin{pmatrix}
 B_{\mathrm P}(t)\\ B_{\mathrm P}(t+1)\\
 B_{\mathrm E}(t)\\ B_{\mathrm E}(t+1)\\
 B_{\mathrm F}(t)\\ B_{\mathrm F}(t+1)
 \end{pmatrix}.
 \label{act21:eq:df-registro-matricial}
\end{equation}
Esta definición fija tanto el orden regional como el orden temporal. Los
cuatro primeros registros son
\[
 X_0=R(0),\quad Y_0=R(1),\quad X_1=R(2),\quad Y_1=R(3).
\]
Sus valores, obtenidos por extracción de los prefijos regionales, son
\begin{align*}
 X_0&=\begin{pmatrix}
0&1&0&2&1&1\\0&1&2&2&2&2\\2&0&1&1&0&1\\
1&2&1&2&2&1\\1&2&1&2&0&0\\1&1&2&2&0&2
\end{pmatrix},&
 Y_0&=\begin{pmatrix}
0&1&2&2&2&2\\0&1&0&2&1&1\\1&2&1&2&2&1\\
1&0&2&0&1&1\\1&1&2&2&0&2\\1&2&1&0&2&0
\end{pmatrix},\\[1ex]
 X_1&=\begin{pmatrix}
0&1&0&2&1&1\\0&0&2&1&1&1\\1&0&2&0&1&1\\
0&1&2&2&2&2\\1&2&1&0&2&0\\0&1&0&2&1&0
\end{pmatrix},&
 Y_1&=\begin{pmatrix}
0&0&2&1&1&1\\1&1&0&2&2&1\\0&1&2&2&2&2\\
1&0&2&0&1&1\\0&1&0&2&1&0\\0&1&0&2&0&0
\end{pmatrix}.
\end{align*}

\begin{theorem}[Reconstrucción única de los dos enlaces iniciales]
\label{act21:thm:df-elevaciones-reconstruidas}
Las ecuaciones \(X_0L_0=Y_0\) y \(X_1L_1=Y_1\) tienen, cada una, una
única solución en \(M_6(\mathbb F_3)\). Estas soluciones son
\begin{align*}
 L_0&=\begin{pmatrix}
2&2&2&1&2&1\\2&1&2&2&1&1\\1&0&0&1&0&2\\
0&0&1&1&1&0\\2&2&2&0&2&1\\2&1&2&1&0&0
\end{pmatrix},&
 L_1&=\begin{pmatrix}
2&1&2&0&2&2\\1&2&1&1&2&0\\1&2&1&1&1&2\\
0&1&0&0&2&1\\2&0&2&1&0&1\\0&2&2&2&1&1
\end{pmatrix}.
\end{align*}
Ambas matrices son invertibles.
\end{theorem}

\begin{proof}
Las matrices
\begin{align*}
 A_0&=\begin{pmatrix}
1&1&2&0&1&2\\0&0&1&0&0&1\\2&1&2&1&0&1\\
0&2&0&1&0&2\\2&1&1&0&2&2\\2&1&1&1&1&2
\end{pmatrix},&
 A_1&=\begin{pmatrix}
1&0&1&2&0&0\\0&0&1&1&2&2\\0&1&2&0&1&1\\
1&1&0&2&0&0\\1&1&2&1&1&2\\1&0&0&0&0&2
\end{pmatrix}
\end{align*}
satisfacen \(A_iX_i=X_iA_i=I\), por multiplicación en \(\mathbb F_3\).
Por tanto, las soluciones están forzadas por \(L_i=A_iY_i\); la
multiplicación produce las dos matrices mostradas. Sus inversas son
\begin{align*}
 L_0^{-1}&=\begin{pmatrix}
1&2&0&2&0&2\\2&0&2&2&0&0\\2&1&2&0&2&0\\
1&0&0&0&2&0\\0&2&1&1&2&0\\2&2&2&2&2&2
\end{pmatrix},&
 L_1^{-1}&=\begin{pmatrix}
2&0&2&1&2&1\\2&1&0&0&2&0\\2&0&2&2&0&2\\
2&0&0&1&1&0\\1&1&1&1&1&0\\2&0&1&2&2&0
\end{pmatrix}.
\end{align*}
La verificación bilateral \(L_iL_i^{-1}=L_i^{-1}L_i=I\) demuestra la
recuperación exacta de cada palabra transformada. Todos los productos
utilizan residuos \(0,1,2\), con su interpretación en \(\mathbb F_3\).
\end{proof}

La lectura causal de este resultado es precisa: las emisiones regionales
producen \(R(t)\); los registros \(R(0),\ldots,R(3)\) determinan las dos
elevaciones anteriores; sus inversas recuperan las palabras transformadas.
La continuidad para todo \(t\) pertenece a los lectores coinductivos y al
teorema~\ref{act21:thm:df-bloque-estable}. La reconstrucción matricial aquí
demostrada tiene por dominio los dos enlaces especificados en su enunciado.

\subsection{Refinamiento simultáneo en las bases ternaria y decimal}
\label{act21:subsec:df-cilindros-mixtos}

% Fuente: deltas/integracion/RegionalCylinderTransition.lean.

Una emisión ternaria de seis trits y una emisión decimal de tres cifras
constituyen dos lecturas de distinta base. Para conservar su relación se
registra el estado
\[
 s=(n,k,N,D),
\]
donde \(N\) es un prefijo a profundidad \(n\) en base \(729\), y \(D\) es
un prefijo a profundidad \(k\) en base \(1000\). Sus cilindros
arquimedianos son
\begin{equation}
 I_{729}(s)=\left[\frac{N}{729^n},\frac{N+1}{729^n}\right),\qquad
 I_{1000}(s)=\left[\frac{D}{1000^k},\frac{D+1}{1000^k}\right).
 \label{act21:eq:df-cilindros}
\end{equation}
El extremo superior semiabierto asigna cada valor a una única celda de una
base y una profundidad dadas.

\begin{definition}[Defectos enteros de compatibilidad]
Definimos
\begin{align}
 A(s)&=N1000^k-D729^n,\label{act21:eq:df-defecto-inferior}\\
 B(s)&=(N+1)1000^k-D729^n=A(s)+1000^k.
 \label{act21:eq:df-defecto-superior}
\end{align}
El estado es compatible cuando
\begin{equation}
 A(s)<729^n,\qquad B(s)>0.
 \label{act21:eq:df-compatibilidad}
\end{equation}
\end{definition}

\begin{lemma}[Criterio entero de intersección]
Los cilindros de \eqref{act21:eq:df-cilindros} tienen intersección no vacía si y
sólo si se cumple \eqref{act21:eq:df-compatibilidad}.
\end{lemma}

\begin{proof}
Dos intervalos semiabiertos \([a,b)\) y \([c,d)\), con \(a<b\) y
\(c<d\), se cortan exactamente cuando \(a<d\) y \(c<b\). Aplicando esta
condición a \eqref{act21:eq:df-cilindros} y multiplicando por los denominadores
positivos, obtenemos
\(N1000^k<(D+1)729^n\) y
\(D729^n<(N+1)1000^k\). Son las dos desigualdades de
\eqref{act21:eq:df-compatibilidad}.
\end{proof}

\begin{definition}[Actualización de profundidades y prefijos]
Para \(0\leq u<729\) y \(0\leq d<1000\), definimos
\begin{align*}
 T_u(n,k,N,D)&=(n+1,k,729N+u,D),\\
 T_{u,d}(n,k,N,D)&=(n+1,k+1,729N+u,1000D+d).
\end{align*}
La primera operación refina únicamente el cilindro ternario. La segunda
refina ambas lecturas y conserva las dos profundidades en el estado.
\end{definition}

\begin{proposition}[Recurrencias exactas de frontera]
\label{act21:prop:df-rec-frontera}
Las actualizaciones anteriores satisfacen
\begin{align}
 A(T_u s)&=729A(s)+u1000^k,\label{act21:eq:df-rec-ternaria}\\
 A(T_{u,d}s)&=729000A(s)+u1000^{k+1}-d729^{n+1}.
 \label{act21:eq:df-rec-conjunta}
\end{align}
\end{proposition}

\begin{proof}
Para la primera identidad,
\[
 (729N+u)1000^k-D729^{n+1}
 =729(N1000^k-D729^n)+u1000^k.
\]
Para la segunda,
\begin{align*}
 &(729N+u)1000^{k+1}-(1000D+d)729^{n+1}\\
 &\qquad=729000(N1000^k-D729^n)
      +u1000^{k+1}-d729^{n+1}.
\end{align*}
Ambas son identidades enteras anteriores a cualquier aproximación real.
\end{proof}

\begin{theorem}[Refinamiento de los cilindros regionales generados]
\label{act21:thm:df-refinamiento-generado}
Sea
\[
 s_c(n,k)=(n,k,P_c(729,n),P_c(1000,k)).
\]
Entonces \(s_c(n,k)\) es compatible para cualesquiera \(n,k\). Además,
\begin{align*}
 T_{u_c(n)}s_c(n,k)&=s_c(n+1,k),\\
 T_{u_c(n),d_c(1000,k)}s_c(n,k)&=s_c(n+1,k+1).
\end{align*}
La división euclídea de los prefijos refinados recupera exactamente los
prefijos anteriores. Para todo \(a,b\geq0\),
\[
 \frac{P_c(729,n+a)-P_c(729,n+a)\bmod729^a}{729^a}
 =P_c(729,n),
\]
y se cumple la identidad análoga en base \(1000\).
\end{theorem}

\begin{proof}
El teorema~\ref{act21:thm:df-emision-correcta} sitúa \(v_c\) simultáneamente en
los dos intervalos \eqref{act21:eq:df-cilindros}; el criterio de intersección da
la compatibilidad. Las identidades de actualización son
\eqref{act21:eq:df-paso-posicional} en cada base. La recuperación se obtiene de
\eqref{act21:eq:df-truncamiento}. Por tanto, la recurrencia del defecto y la
recuperación del prefijo son consecuencias de una misma actualización.
\end{proof}

\begin{example}[Primer par de cilindros regionales]
Para \(n=k=1\), la evaluación de las tres regiones da
\[
\begin{array}{c|rr|rr}
 c&N&D&A&B\\\hline
 \mathrm P&2290&3141&211&1211\\
 \mathrm E&1981&2718&-422&578\\
 \mathrm F&1179&1618&-522&478
\end{array}
\]
En las tres filas \(A<729\) y \(B>0\). La primera columna de prefijos
se descompone como \(2290=3\cdot729+103\),
\(1981=2\cdot729+523\) y \(1179=729+450\): reaparecen exactamente los
bloques de seis trits ya emitidos. Los signos diferentes de \(A\) expresan
la posición relativa de los dos extremos inferiores, sin alterar la
pertenencia común del valor regional.
\end{example}

\subsection{Selección del bloque siguiente mediante el lector regional}

La compatibilidad de dos intervalos expresa una relación entre lecturas.
La determinación del siguiente bloque utiliza además el intervalo racional
producido por la región. Esta operación puede formularse sin eliminar la
historia de profundidad.

\begin{definition}[Admisibilidad de la prolongación regional]
Fijados \(c,n\), sea \(s=729^{n+1}\) y \(j=j_c(s)\). Un residuo
\(u\in\{0,\ldots,728\}\) es admisible cuando
\begin{equation}
 \lfloor s\ell_c(j)\rfloor
 \leq729P_c(729,n)+u
 \leq\lceil su_c(j)\rceil-1.
 \label{act21:eq:df-admisibilidad}
\end{equation}
\end{definition}

\begin{theorem}[Unicidad de la prolongación compatible con el lector]
\label{act21:thm:df-prolongacion-unica}
La condición \eqref{act21:eq:df-admisibilidad} equivale a \(u=u_c(n)\).
Por tanto, existe exactamente una prolongación admisible de cada región
a cada profundidad.
\end{theorem}

\begin{proof}
En el índice de separación, los dos extremos enteros de
\eqref{act21:eq:df-admisibilidad} coinciden con \(P_c(729,n+1)\).
La condición equivale así a
\(729P_c(729,n)+u=P_c(729,n+1)\).
Por \eqref{act21:eq:df-paso-posicional}, el lado derecho es
\(729P_c(729,n)+u_c(n)\). La cancelación del término común da la
igualdad de residuos. Recíprocamente, esa igualdad satisface la condición.
\end{proof}

La dependencia efectiva de esta selección queda completamente determinada:
región seleccionada, lector racional, profundidad, índice de separación,
prefijo anterior y residuo siguiente. El defecto \(A\) resume una relación
de frontera; el estado \((n,k,N,D)\) y los intervalos regionales conservan
los datos empleados por la actualización.

\subsection{Firma ternaria, acarreo, incidencia de Witt y reloj de lectura}
\label{act21:subsec:df-firma-conjunta}

% Fuentes: lean/regional/N69RegionalSignature.lean,
% deltas/integracion/JointRegionalFrontier.lean y JointCylinderSignature.lean.

Fijada una profundidad \(t\), escribamos \(B_{c,j}=B_{c,j}(t)\). Para
cada columna \(j\) se calculan los siguientes enteros y residuos:
\begin{align}
 S_j&=B_{\mathrm P,j}+B_{\mathrm E,j}+B_{\mathrm F,j},
 &q_j&=S_j\bmod3,
 &c_j&=\left\lfloor S_j/3\right\rfloor,
 \label{act21:eq:df-firma-qc}\\
 a_j&=(B_{\mathrm P,j}+B_{\mathrm E,j}-B_{\mathrm F,j})\bmod3,
 &w_j&=\sum_c\mathbf1_{B_{c,j}\ne0},
 &\bar w_j&=w_j\bmod3.
 \label{act21:eq:df-firma-axial}
\end{align}
El residuo de la expresión axial se toma en los enteros. Una implementación
con naturales lo calcula como
\((B_{\mathrm P,j}+B_{\mathrm E,j}+3-B_{\mathrm F,j})\bmod3\), que conserva
exactamente ese residuo.

La matriz de incidencia ternaria utilizada por el lector es
\begin{equation}
 W=\begin{pmatrix}
0&1&1&1&1&1\\
1&0&1&1&2&2\\
1&1&0&2&1&2\\
1&1&2&0&2&1\\
1&2&1&2&0&1\\
1&2&2&1&1&0
\end{pmatrix}.
\label{act21:eq:df-matriz-witt}
\end{equation}
Para cada región se conservan dos cargas residuales:
\begin{equation}
 v_c^{\mathrm{res}}=\sum_j B_{c,j}\bmod3,
 \qquad
 v_c^{\mathrm W}=\sum_j(B_cW)_j\bmod3.
 \label{act21:eq:df-cargas-residuales}
\end{equation}
La firma completa del lector es
\(\mathcal S(B)=(q,a,c,w,\bar w,v^{\mathrm{res}},v^{\mathrm W})\).
Su estructura conserva por separado el residuo de columna, el acarreo, el
contraste axial, el peso de soporte y las dos lecturas de carga.

\begin{theorem}[Recuperación exacta y cotas de la firma]
\label{act21:thm:df-recuperacion-firma}
Para cada columna y cada región se cumplen
\begin{align}
 S_j&=q_j+3c_j,&0\leq q_j,a_j&<3,&0\leq c_j&\leq2,
 \label{act21:eq:df-recuperacion-total}\\
 B_{\mathrm F,j}&=2(q_j+3-a_j)\bmod3,&0\leq w_j&\leq3,
 &0\leq\bar w_j&<3.
 \label{act21:eq:df-recuperacion-eje}
\end{align}
Además,
\begin{equation}
 v_c^{\mathrm W}
 =\left(\sum_j\bigl((B_cW)_j\bmod3\bigr)\right)\bmod3.
 \label{act21:eq:df-witt-reduccion}
\end{equation}
Las recuperaciones enunciadas determinan el total de cada columna y su
coordenada de autoescala; el registro regional conserva adicionalmente
las tres bandas ordenadas.
\end{theorem}

\begin{proof}
La identidad \eqref{act21:eq:df-recuperacion-total} es la división euclídea de
\(S_j\) por \(3\). Cada sumando de \(S_j\) está entre \(0\) y \(2\),
de modo que \(0\leq S_j\leq6\) y \(c_j\leq2\). En \(\mathbb F_3\),
la resta de las dos congruencias que definen \(q_j\) y \(a_j\) da
\(q_j-a_j=2B_{\mathrm F,j}\). Como el inverso de \(2\) es \(2\), se
obtiene \eqref{act21:eq:df-recuperacion-eje}; el sumando \(3\) permite
escribirla con naturales. Las cotas de \(w_j\) resultan de sumar tres
indicadores. Finalmente, reducir una suma entera módulo \(3\) equivale
a reducir cada sumando y volver a reducir la suma, lo que prueba
\eqref{act21:eq:df-witt-reduccion}.
\end{proof}

\begin{example}[Firma del primer bloque conjunto]
Las tres bandas iniciales producen
\begin{align*}
 q&=(0,0,2,2,1,2),&a&=(1,2,0,1,1,2),\\
 c&=(1,1,0,1,0,0),&w&=(2,2,2,3,1,2),\\
 v^{\mathrm{res}}&=(2,2,0),&v^{\mathrm W}&=(2,1,1).
\end{align*}
La cuarta columna tiene total \(2+1+2=5\); su firma registra
\(q_3=2\) y \(c_3=1\), por lo que \(q_3+3c_3=5\).
La recuperación axial da
\(2(2+3-1)\bmod3=2\), precisamente la coordenada de autoescala.
Las imágenes de Witt, reducidas componente a componente, son
\[
 (2,0,2,1,1,2),\qquad
 (0,0,0,2,0,2),\qquad
 (2,1,1,2,1,0).
\]
Su suma residual por fila produce las tres cargas \((2,1,1)\).
\end{example}

\begin{definition}[Reloj entero de lectura decimal]
El índice de transición \(t\) y la profundidad decimal de representación se
relacionan mediante
\begin{equation}
 \kappa(t)=\max\{k\in\mathbb N:1000^k\leq3^{6(t+1)}\}.
 \label{act21:eq:df-reloj}
\end{equation}
\end{definition}

\begin{proposition}[Caracterización y monotonía del reloj]
\label{act21:prop:df-reloj}
El entero \(\kappa(t)\) queda determinado unívocamente por
\begin{equation}
 1000^{\kappa(t)}\leq729^{t+1}<1000^{\kappa(t)+1}.
 \label{act21:eq:df-reloj-cotas}
\end{equation}
La función \(\kappa\) es monótona y satisface
\(\kappa(20)=\kappa(21)=20\).
\end{proposition}

\begin{proof}
Las potencias de \(1000\) son estrictamente crecientes y no acotadas;
existe, por tanto, un único intervalo consecutivo de potencias que
contiene \(729^{t+1}\). La monotonía de esta última expresión demuestra
la de \(\kappa\). Las dos evaluaciones finales son las comparaciones
enteras
\[
 1000^{20}\leq729^{21}<1000^{21},\qquad
 1000^{20}\leq729^{22}<1000^{21}.
\]
Estas comparaciones se efectúan con potencias enteras exactas.
\end{proof}

La igualdad de profundidad decimal en dos transiciones distintas explica
la necesidad de conservar ambos índices: el reloj de lectura puede
permanecer constante mientras avanza la banda ternaria y se actualiza su
firma. La información cronológica se encuentra en el registro ordenado
de las transiciones.

\subsection{Teorema de composición del bloque, el cilindro y la firma}

\begin{theorem}[Compatibilidad conjunta de las lecturas regionales]
\label{act21:thm:df-composicion-conjunta}
Para cada profundidad \(n\), existe una única terna de bloques
\((u_{\mathrm P},u_{\mathrm E},u_{\mathrm F})\) que satisface la
admisibilidad regional \eqref{act21:eq:df-admisibilidad} en las tres regiones.
Esta terna es \((u_{\mathrm P}(n),u_{\mathrm E}(n),u_{\mathrm F}(n))\).
Simultáneamente:
\begin{enumerate}
\item cada bloque actualiza su cilindro por las recurrencias
\eqref{act21:eq:df-rec-ternaria} y \eqref{act21:eq:df-rec-conjunta};
\item sus seis trits producen la firma
\(\mathcal S(B_{\mathrm P}(n),B_{\mathrm E}(n),B_{\mathrm F}(n))\);
\item los totales de columna se recuperan como \(q+3c\), y las cargas
de Witt se obtienen de las mismas bandas mediante
\eqref{act21:eq:df-witt-reduccion};
\item todo prefijo de mayor profundidad devuelve exactamente estos
bloques, cilindros truncados y firmas.
\end{enumerate}
\end{theorem}

\begin{proof}
Aplicamos el teorema~\ref{act21:thm:df-prolongacion-unica} a cada uno de los tres
índices regionales. La unicidad coordenada a coordenada da la unicidad
conjunta. Sustituimos esos mismos residuos en las operaciones
\(T_u,T_{u,d}\), lo que permite aplicar el
teorema~\ref{act21:thm:df-refinamiento-generado}. La aplicación determinista de
\eqref{act21:eq:df-bloque-y-banda} produce las bandas utilizadas por
\eqref{act21:eq:df-firma-qc}--\eqref{act21:eq:df-cargas-residuales}. Las identidades
de recuperación son las del teorema~\ref{act21:thm:df-recuperacion-firma}.
Finalmente, el teorema~\ref{act21:thm:df-bloque-estable} y el truncamiento de
cada base demuestran la estabilidad a mayor profundidad. En cada paso se
ha conservado la identidad del bloque emitido; ninguna de las dos
lecturas requiere seleccionar otro bloque.
\end{proof}

La conclusión es una compatibilidad operacional entre la contracción
cilíndrica y la conservación de incidencia. Ambas utilizan el mismo
registro generado, mantienen las profundidades y admiten recuperación
entera. Este resultado es la forma explícita del enlace
\[
 \text{bloque regional emitido}
 \longrightarrow
 \begin{cases}
 \text{cilindro y frontera entera},\\
 \text{banda, firma, acarreo y carga de Witt}.
 \end{cases}
\]
El desarrollo posterior de la incidencia y de los campos emplea la
segunda lectura, conservando la procedencia común con la primera.

\begingroup\fontsize{10}{11.8}\selectfont\setlength{\abovedisplayskip}{5pt}\setlength{\belowdisplayskip}{5pt}
\subsection{Censo de calendarios y compatibilidad simultánea de las biografías regionales}
Las filas anteriores se reúnen, sin aplicar todavía ningún lector decimal,
en las tres biografías de bloques producidas por las transiciones:
\begin{align}
 B^{\rm cl}
  &=010211|012222|010211|002111|110221,\notag\\
 B^{\rm pr}
  &=201101|121221|102011|012222|102011,
 \label{act21:eq:biografias-tpk-tres-caracteres}\\
 B^{\rm au}
  &=121200|112202|121020|010210|010200.\notag
\end{align}
En efecto, los pares consecutivos de los dos primeros tramos son las filas
de \(X_0,Y_0\), y los de los dos últimos son las filas de \(X_1,Y_1\).
Así, las biografías son otra escritura del registro de doce transiciones, no
una colección de prefijos convencionales.

Para comprobar la palabra operativa del calendario, sea
\(c=c_1c_2c_3c_4\in\{0,1\}^4\) y, para cada régimen
\(\chi\in\{{\rm cl},{\rm pr},{\rm au}\}\), defínase
\[
 u_0^\chi=b_0^\chi,\qquad
 u_j^\chi=u_{j-1}^\chi L_{c_j}\quad(1\le j\le4).
\]
Para condensar el censo, denotemos por
\[
 N_{\mathrm{ar}}(c):=
 \#\{(j,\chi):u_j^\chi=b_j^\chi,\ 1\le j\le4\},\qquad
 N_{\mathrm{bio}}(c):=
 \#\{\chi:(u_0^\chi,\ldots,u_4^\chi)=B^\chi\}.
\]
La multiplicación exacta en \(\FF_3\) de los dieciséis calendarios da:
\[
\begin{array}{c|c|c}
 c&N_{\mathrm{ar}}(c)&N_{\mathrm{bio}}(c)\\ \hline
0000,0001&6&0\\
0010&9&0\\
0011&12&3\\
0100,0101,0110,0111&3&0\\
1000,1001,1010,1011&0&0\\
1100,1101,1110,1111&0&0
\end{array}
\tag{D}
\label{act21:eq:censo-dieciseis-calendarios}
\]
La tabla agota \(2^4=16\) posibilidades y cada una de sus entradas se obtiene
por cuatro multiplicaciones de una fila por una de las dos matrices impresas.
Por consiguiente, \(0011\) es el único calendario que reconstruye
simultáneamente las tres biografías completas. En notación operatoria,
\begin{equation}
 (L_0,L_0,L_1,L_1)
 \label{act21:eq:calendario-0011-interno}
\end{equation}
no es una prescripción independiente: es la única lectura binaria compatible
con las doce aristas ya producidas por \(U_t\).

El calendario \eqref{act21:eq:calendario-0011-interno} produce cuatro nuevos
bloques de seis componentes. Sus concatenaciones forman
\[
 w_6\longrightarrow w_{12}\longrightarrow w_{18}
 \longrightarrow w_{24}\longrightarrow w_{30}.
\]
Los subíndices de \(w\) indican longitud acumulada. La frontera \(R_{36}\) conserva la estructura terminal y abre la relación de prolongación; no se renombra como un último bloque periódico que pudiera repetirse indefinidamente.

La identidad \(L=X^{-1}Y\) demuestra unicidad una vez fijadas las
transiciones; no sustituye su producción por \(U_t\). Recíprocamente,
\eqref{act21:eq:censo-dieciseis-calendarios} prueba dentro del artículo la unicidad
del orden de los cuatro transportes. La cadena causal queda así cerrada:
el censo APP produce las regiones iniciales, el TRIT fija régimen y
orientación, el TPK produce el registro de transiciones, y ese registro
determina \(L_0,L_1\) y \(0011\) antes de toda realización arquimediana.


\par\endgroup

\clearpage
% Desarrollo reunido K31-07--K31-15. Inserción anterior a registro_k y alpha.
% Mantiene explícita la interfaz del repertorio terminal de ocho bloques.
% Propietarios y localizadores: metadata/K_ALPHA_PROPIETARIOS.md.
\section{Determinación regional del descriptor dodecafásico y selección del registro de acoplamiento}
\label{act21:chap:despliegue-selector-k}

La determinación del registro de acoplamiento articula dos construcciones
que conviene desplegar en su orden efectivo. La primera produce los
antecedentes del selector a partir de las representaciones regionales del TPK:
bandas de seis trits, firmas de columna, calendario de conversión,
descriptor orientado y panel funcional. La segunda opera sobre esos
antecedentes y sobre el repertorio terminal de ocho bloques, enumera sus
ordenaciones admisibles y selecciona un único registro mediante dos
condiciones simultáneas: incidencia excepcional y heterotipia de lectura.
El valor del registro aparece al término de esta composición.

El presente desarrollo presupone las representaciones regionales obtenidas
en los capítulos precedentes desde APP--TRIT--TPK. Su propósito es hacer
explícita cada operación intermedia que conduce desde tales representaciones
al selector finito. Después se incorporan la realización del mismo
registro por incidencias firmadas, la inversión de Hadamard y la
representación de la constante de estructura fina. Las tres construcciones
conservan sus respectivos dominios: selección del registro, recuperación
desde sus canales y evaluación aritmética de sus representaciones.

\subsection{Bandas regionales, codificación y firmas enteras}

\subsubsection{Dominios y codificación de las bandas regionales}

Se fija el orden de canales
\[
 \chi\in\{\pi,e,\varphi\},
 \qquad B_\chi(t)\in\mathbb F_3^6,
 \qquad t\in\mathbb N.
\]
Los símbolos de canal identifican las regiones ya generadas. Cada
$B_\chi(t)$ conserva el lugar que ocupa en su sucesión de representaciones;
el índice $t$ cuenta bloques ternarios. La evaluación natural de una
banda $b=(b_0,\ldots,b_5)$ es
\begin{equation}
 \operatorname{enc}_6(b)=
 243b_0+81b_1+27b_2+9b_3+3b_4+b_5,
 \qquad 0\leq\operatorname{enc}_6(b)<729.
 \label{act21:eq:kdes-enc6}
\end{equation}
Su inversa, con anchura fijada, es
\[
 \operatorname{dig}_j(n)=
 \left\lfloor\frac{n}{3^{5-j}}\right\rfloor\bmod3,
 \qquad 0\leq j<6.
\]
Así se retienen los ceros iniciales: las palabras $002001$ y $2001$
podrían compartir una escritura entera abreviada, pero el tipo de la
primera contiene seis posiciones. Todas las concatenaciones que siguen
utilizan esa anchura.

En la ventana de cien bloques, la representación de seiscientos trits del
canal $\chi$ determina un entero $P_\chi^{(600)}$. La extracción se efectúa
mediante
\begin{equation}
 b_\chi(t)=
 \left\lfloor\frac{P_\chi^{(600)}}{729^{99-t}}\right\rfloor\bmod729,
 \qquad 0\leq t<100,
 \qquad B_\chi(t)=\operatorname{enc}_6^{-1}(b_\chi(t)).
 \label{act21:eq:kdes-extraccion}
\end{equation}
El entero de seiscientos trits es una representación del productor regional;
la tabla de cien filas es la evaluación de
\eqref{act21:eq:kdes-extraccion}. Esta separación permite sustituir una tabla
histórica por su productor sin modificar el selector que la consume.

\subsubsection{Residuos, cocientes, ocupación y orientación}

Para cada columna $j$ y tiempo $t$, se introducen
\begin{align}
 T_j(t)&=B_\pi(t)_j+B_e(t)_j+B_\varphi(t)_j,\\
 q_j(t)&=T_j(t)\bmod3,
 &c_j(t)&=\left\lfloor T_j(t)/3\right\rfloor,\\
 a_j(t)&=\bigl(B_\pi(t)_j+B_e(t)_j-B_\varphi(t)_j\bigr)\bmod3,\\
 w_j(t)&=\mathbf1_{B_\pi(t)_j\ne0}
       +\mathbf1_{B_e(t)_j\ne0}
       +\mathbf1_{B_\varphi(t)_j\ne0},
 &\bar w_j(t)&=w_j(t)\bmod3.
 \label{act21:eq:kdes-firmas}
\end{align}
La resta de la tercera expresión se realiza en los enteros y se reduce
después. En una implementación sobre naturales, la expresión equivalente
es $(B_\pi+B_e+3-B_\varphi)\bmod3$; la resta truncada alteraría el lector.

\begin{proposition}[Reconstrucción de la suma de columna]
Para todas las bandas admisibles,
\[
 T_j(t)=q_j(t)+3c_j(t),\qquad
 0\leq q_j(t)<3,\quad 0\leq c_j(t)\leq2,
 \quad0\leq w_j(t)\leq3.
\]
\end{proposition}
\begin{proof}
Cada columna suma tres representantes de $\{0,1,2\}$, de modo que
$0\leq T_j(t)\leq6$. La división euclídea por tres proporciona el residuo
y el cociente indicados. La ocupación cuenta tres condiciones binarias;
su cota se deduce directamente de esa definición.
\end{proof}

Esta identidad conserva el cociente de la suma local. La reconstrucción
de las tres bandas individuales requiere las bandas mismas o un lector
adicional que las distinga: una suma de columna tiene menos coordenadas
que el triple que la produce. Por ello, la fila temporal utilizada más
adelante conserva conjuntamente las tres bandas y los cuatro campos
$q,a,c,\bar w$.

\subsection{Reloj de conversión y primer avance nulo}

\subsubsection{Capacidad ternaria y profundidad decimal}

Una representación de $t+1$ bandas contiene $6(t+1)$ trits. El calendario
de conversión asociado se define por el entero
\begin{equation}
 d_t=\max\{k\in\mathbb N:1000^k\leq3^{6(t+1)}\}
     =\max\{k\in\mathbb N:1000^k\leq729^{t+1}\}.
 \label{act21:eq:kdes-reloj}
\end{equation}
La definición compara capacidades enteras. El contador $t$ y la
profundidad $d_t$ registran magnitudes distintas: una transición añade
una banda de seis trits; el reloj cuenta potencias completas de la base
decimal mil contenidas en la capacidad ternaria acumulada.

\begin{lemma}[Caracterización entera del reloj]
Para $t,k\in\mathbb N$ se tiene
\[
 d_t=k\quad\Longleftrightarrow\quad
 1000^k\leq729^{t+1}<1000^{k+1}.
\]
Además, $d_t\leq d_{t+1}$.
\end{lemma}
\begin{proof}
Las potencias de mil forman una sucesión estrictamente creciente y no
acotada. La potencia $729^{t+1}$ queda entre dos términos consecutivos;
el exponente inferior es exactamente el máximo de
\eqref{act21:eq:kdes-reloj}. La desigualdad
$729^{t+1}<729^{t+2}$ implica la monotonía del máximo.
\end{proof}

\begin{proposition}[Primer avance nulo]
El menor tiempo que satisface $d_{t+1}=d_t$ es
\[
 t_*=20.
\]
En particular,
\[
 d_0=0,\quad d_{19}=19,\quad d_{20}=20,
 \quad d_{21}=20,\quad d_{22}=21.
\]
\end{proposition}
\begin{proof}
Las comparaciones enteras exactas
\[
 1000^{20}\leq729^{21},\qquad
 729^{22}<1000^{21},\qquad
 1000^{21}\leq729^{23}<1000^{22}
\]
dan los valores en los tiempos veinte, veintiuno y veintidós. Para
$0\leq t\leq20$, la expresión
$729^{t+1}/1000^t=729(729/1000)^t$ es decreciente y su valor en veinte
es mayor o igual que uno. Por tanto,
$1000^t\leq729^{t+1}<1000^{t+1}$ y $d_t=t$. Ningún tiempo menor que
veinte puede satisfacer $d_{t+1}=d_t$; en veinte sí se cumple.
Todas las comparaciones utilizadas son de potencias enteras.
\end{proof}

El número veinte es, por tanto, una salida del criterio mínimo. La regla
que selecciona el instante es ``primer avance nulo del reloj''; el
algoritmo y su prueba pueden emplear esta regla sin recibir veinte como
parámetro externo.

\subsubsection{Retorno nonádico y antecedente temporal}

La longitud del retorno de fase determina
\begin{equation}
 r=t_*-9=11.
 \label{act21:eq:kdes-retorno}
\end{equation}
Las consultas del descriptor se efectúan en $t_*+1$, $r+1$ y $r$.
Esta distribución conserva la relación entre el instante de pausa, la
traslación nonádica y la lectura de fase inmediatamente asociada. Los
tiempos once, doce y veintiuno designan posiciones del calendario
regional; su diferencia temporal permanece registrada aun cuando dos
tiempos produzcan la misma profundidad decimal.

\subsection{Construcción orientada del descriptor de veinticuatro trits}

\subsubsection{Defectos opuestos de semibanda}

Se toman los defectos orientados de la coordenatización regional
\[
 \delta_+=(0,0,0,1,1,1),\qquad
 \delta_-=(0,0,0,2,2,2)\in\mathbb F_3^6.
\]
Su oposición se verifica componente a componente:
\[
 \delta_++\delta_-=0\quad\text{en }\mathbb F_3^6.
\]
El representante $2$ realiza $-1$ en la coordenatización ternaria. La concatenación
de las colas de tres posiciones, en el orden negativo--positivo, produce
la palabra $222111$. La suma de defectos y la concatenación de sus
colas son operaciones distintas: la primera tiene codominio
$\mathbb F_3^6$, mientras que la segunda conserva el orden de dos
segmentos de longitud tres.

\subsubsection{Carga dual y selección de fase}

En la coordenatización regional, la transformación de seis coordenadas utiliza
\[
 A_{\rm reg}=
 \begin{pmatrix}
 0&1&1&1&1&1\\
 1&0&1&1&2&2\\
 1&1&0&2&1&2\\
 1&1&2&0&2&1\\
 1&2&1&2&0&1\\
 1&2&2&1&1&0
 \end{pmatrix}.
\]
Con vectores fila, la carga dual de $b\in\mathbb F_3^6$ es la suma
de las coordenadas de $bA_{\rm reg}$. Las sumas enteras de las filas de
$A_{\rm reg}$ son $(5,7,7,7,7,7)$; de ahí resulta la expresión local
\begin{equation}
 \eta^\vee(b)=2b_0+b_1+b_2+b_3+b_4+b_5\pmod3.
 \label{act21:eq:kdes-cargadual}
\end{equation}
La fase de lectura en el instante $r$ es
$\theta_r=(r-1)\bmod3$. Para cada canal se define
\[
 \varepsilon_\chi(r)=
 \begin{cases}
 1,&\eta^\vee(B_\chi(r))=\theta_r,\\
 2,&\eta^\vee(B_\chi(r))\ne\theta_r.
 \end{cases}
\]
La representación fase--carga es
\begin{equation}
 D(r)=\varepsilon_\pi(r)B_\pi(r)
      +\varepsilon_e(r)B_e(r)
      +\varepsilon_\varphi(r)B_\varphi(r)
      \quad\text{en }\mathbb F_3^6.
 \label{act21:eq:kdes-fasecarga}
\end{equation}
El coeficiente se determina por una igualdad de cargas, antes de evaluar
la palabra que resultará de la combinación.

\subsubsection{Composición del descriptor}

Se denota por $\Vert$ la concatenación de palabras de anchura fijada.
La construcción completa es
\begin{equation}
 \begin{split}
 \Omega_{24}={}&
 (B_e(t_*+1)+\delta_-)
 \Vert (\delta_-|_{3,4,5}\Vert\delta_+|_{3,4,5})\\
 &\Vert(B_\pi(r+1)+B_e(r+1))\Vert D(r).
 \end{split}
 \label{act21:eq:kdes-omega}
\end{equation}
Cada sumando de longitud seis se evalúa en $\mathbb F_3^6$; la
concatenación final tiene longitud $6+6+6+6=24$. Esta construcción
separa el descriptor $\Omega_{24}$ de las etapas regionales
$w_{24}^{\chi}$: ambos objetos comparten longitud, pero sus productores
y sus lugares en la composición son diferentes.

\begin{proposition}[Evaluación regional del descriptor]
Las representaciones regionales y el criterio del primer avance nulo
determinan
\[
 \Omega_{24}=020022\,222111\,211201\,021101,
 \qquad [\Omega_{24}]_3=66241910521.
\]
\end{proposition}
\begin{proof}
Las bandas necesarias, conservando el orden $(\pi,e,\varphi)$, son
\[
 \begin{array}{c|ccc}
 t&B_\pi(t)&B_e(t)&B_\varphi(t)\\\hline
 11&022212&110001&100021\\
 12&012012&202222&000120\\
 21&221022&020100&211021
 \end{array}
\]
y proceden de la extracción \eqref{act21:eq:kdes-extraccion}. En el tiempo
veintiuno,
$020100+000222=020022$ en $\mathbb F_3^6$. Las colas de los defectos
aportan $222111$. En el tiempo doce,
$012012+202222=211201$.

Para el tiempo once, las cargas duales son, respectivamente, $0,1,2$;
la fase vale $(11-1)\bmod3=1$. Los coeficientes resultan $(2,1,2)$ y
\[
 2(022212)+(110001)+2(100021)=021101
 \quad\text{en }\mathbb F_3^6.
\]
La concatenación de estos cuatro resultados da la palabra anunciada.
Su evaluación entera se obtiene por Horner en base tres, o por tres
concatenaciones sucesivas en base $729$.
\end{proof}

La tabla anterior desempeña una función comprobatoria: cada una de sus
entradas se obtiene del productor regional. La regla de construcción
\eqref{act21:eq:kdes-omega} es la que permite reproducir el descriptor sin
introducir como dato la cadena final de veinticuatro trits.

\subsection{Panel funcional y repertorio terminal}

\subsubsection{Nueve posiciones regionales}

El panel funcional conserva las tres primeras bandas de cada uno de los
tres canales, en el orden regional declarado:
\begin{equation}
 \mathcal P_9=
 \bigl(b_\pi(0),b_\pi(1),b_\pi(2),
       b_e(0),b_e(1),b_e(2),
       b_\varphi(0),b_\varphi(1),b_\varphi(2)\bigr).
 \label{act21:eq:kdes-panel}
\end{equation}
Su evaluación es
\[
 \mathcal P_9=(103,161,103,523,457,301,450,398,438).
\]
Las nueve posiciones conservan procedencia regional, aunque el valor
103 aparece dos veces. El panel tiene nueve incidencias ordenadas y
ocho valores distintos. La enumeración ulterior emplea las incidencias;
la deduplicación de candidatos se efectúa en una etapa posterior.

\subsubsection{Repertorio no ordenado de ocho bloques}

El repertorio terminal que interviene en esta construcción es
\begin{equation}
 \mathcal S_{\rm term}=
 \{002001,020111,200110,121012,
   010122,001001,221110,222110\}\subset\mathbb F_3^6.
 \label{act21:eq:kdes-sterm}
\end{equation}
Sus ocho elementos son distintos. La codificación entera, ordenada por
valor únicamente para fijar una representación informática, es
\[
 \{28,55,98,175,437,498,687,714\}.
\]
La procedencia documental de este repertorio es la segmentación terminal
de la palabra de setenta y dos trits y la recuperación por posiciones
mediante incidencias de los lectores fase--carga y afín. El selector que
se desarrolla a continuación recibe el repertorio sin su ordenación
histórica y somete a prueba sus $8!$ ordenaciones.

La formulación exacta de la implementación reunida mantiene aquí una
interfaz explícita: los productores regionales sustituyen materialmente
el descriptor, el panel y las filas temporales; el repertorio
\eqref{act21:eq:kdes-sterm} conserva su productor documental anterior. Esta
distinción identifica dónde entra cada antecedente. La pertenencia de
ocho palabras a un repertorio temporal de lectores, la selección de esas
ocho palabras y la determinación de su orden terminal son tres
enunciados diferentes. La prueba de selección terminal demuestra el
tercero para el repertorio declarado.

La recuperación documental por posiciones puede expresarse de forma
exacta. Sea $\mathcal I$ el registro de incidencias fase--carga y sea
$\mathcal J$ el registro de la extensión afín. Cada incidencia conserva
la posición $\ell$ de la palabra de setenta y dos trits y la palabra
emitida $u$. Para $5\leq\ell\leq12$ se forma
\[
 \mathcal S_\ell=
 \{u:(\ell,u)\in\mathcal I\cup\mathcal J\}.
\]
El recuperador comprueba que las ocho posiciones aparecen y que cada
$\mathcal S_\ell$ es un singleton. Si $s_\ell$ es su elemento único,
la salida entregada al selector es
\[
 \mathcal S_{\rm term}=\{s_5,s_6,\ldots,s_{12}\}.
\]
La prueba de singleton por posición conserva la coherencia de los
testigos históricos; la exportación posterior conserva las ocho
palabras y suprime únicamente la ordenación histórica. El selector
reconstruye el orden mediante la enumeración completa y las condiciones
terminales, en lugar de recibirlo del registro de procedencia.

% RESULTADO_RECUPERADO: K31-10/P10 del inventario REV05.
% Fuente de las incidencias:
% 16_CIERRE_GLOBAL_HMT_MD_2026-07-22/02_LECTURA_DODECAFASICA/continuacion_fases/
% verificar_obstruccion_y_dato_minimo_w72.py, lineas 45--60 y 141--233.
% RESULTADO_OBSTRUCCION_Y_DATO_MINIMO_W72.json, phase_charge_reader.
% Fuente de las bandas: N69_signatures_t0_t99.csv, tiempos impresos abajo.
% Consumidor: PAQUETE_CONTINUIDAD_K_UNIDAD_20260918_BASE_COMPARTIDA/python/
% selector_orbital_original.py, verify_historical_inputs, lineas 159--193.
% No se transfiere la restriccion del lector comparativo al generador completo.
\subsubsection{Incidencias temporales del repertorio terminal y recuperación por posición}
\label{act21:sec:repertorio-incidencias-completas}

La representación del repertorio terminal utiliza los mismos bloques regionales,
los mismos acarreos y el mismo calendario que producen el descriptor inicial
de veinticuatro trits. Conviene explicitar por separado dos operaciones de
esta composición. La primera recupera ocho palabras a partir de incidencias
temporales registradas. La segunda presenta esas ocho palabras al selector
dodecafásico como un conjunto no ordenado. El selector vuelve a establecer
el orden mediante sus condiciones de incidencia, de cilindro y de cierre.
La distinción permite conservar los testigos de procedencia sin convertir
el orden de una segmentación histórica en una entrada del selector.

Sea $B_t=(b_{t,0},b_{t,1},b_{t,2})\in(\mathbb F_3^6)^3$ la frontera
regional ya generada en el tiempo $t$, con $0\leq t<100$. En la coordenatización
utilizada por el registro histórico, la matriz de incidencia es
\[
 A_{\mathrm{hist}}=
 \begin{pmatrix}
 0&1&1&1&1&1\\
 1&0&1&1&2&2\\
 1&1&0&2&1&2\\
 1&1&2&0&2&1\\
 1&2&1&2&0&1\\
 1&2&2&1&1&0
 \end{pmatrix}.
\]
Las palabras se consideran vectores fila. Sus dos cargas y la fase
cronológica son
\[
 r_{t,i}=\sum_{j=1}^{6}b_{t,i,j},\qquad
 d_{t,i}=\sum_{j=1}^{6}(b_{t,i}A_{\mathrm{hist}})_j,
 \qquad \vartheta_t=t-1\pmod 3.
\]
Todas las operaciones de estas tres expresiones tienen lugar en
$\mathbb F_3$. El subíndice histórico identifica una coordenatización concreta de
la incidencia de Witt; mantiene explícita la transformación de coordenadas
cuando se utiliza otra presentación de esa misma incidencia.

Para $c\in\{r_t,d_t\}$, $\delta\in\mathbb F_3$ y
$\epsilon\in\{1,-1\}$, definimos
\[
 \eta_i(c,\delta,\epsilon)=
 \begin{cases}
  \epsilon,&c_i=\vartheta_t+\delta,\\
  -\epsilon,&c_i\ne\vartheta_t+\delta,
 \end{cases}
 \qquad
 L_{t,c,\delta,\epsilon}
   =\sum_{i=0}^{2}\eta_i(c,\delta,\epsilon)b_{t,i}.
\]
La orientación $-1$ se escribe como $2$ al imprimir los coeficientes
ternarios. Junto a estos doce lectores comparativos por frontera, el
registro dispone del lector afín
\[
 L^{\mathrm{af}}_t
    =\sum_{i=0}^{2}(d_{t,i}+\vartheta_t)b_{t,i}.
\]
La adición de la fase en el lector afín y la comparación de cargas en
el lector anterior son operaciones distintas, ambas sobre datos ya
producidos por la frontera TPK.

\begin{table}[htbp]
\centering\small
\begin{tabular}{rccc cc}
\toprule
$t$&$b_{t,0}$&$b_{t,1}$&$b_{t,2}$&$r_t$&$d_t$\\
\midrule
8 &200121&212212&110110&011&202\\
11&022212&110001&100021&001&012\\
30&210112&110202&211110&100&012\\
34&121022&000011&112110&220&021\\
37&120121&011202&112221&100&201\\
49&110212&200212&100122&110&201\\
58&111121&022121&122201&122&220\\
67&010012&001220&011120&122&122\\
75&220000&020111&100002&120&021\\
81&122201&020202&201122&202&001\\
90&022112&022110&121010&202&200\\
\bottomrule
\end{tabular}
\caption{Fronteras temporales que realizan las incidencias impresas del
repertorio terminal. Cada entrada de seis cifras conserva el orden de
las seis coordenadas de la frontera regional.}
\label{act21:tab:terminal-bandas-testigo}
\end{table}

La tabla siguiente despliega todas las incidencias comparativas del
registro sobre los doce bloques de la segmentación considerada. Las
letras $r$ y $d$ designan respectivamente la carga visible y la carga
dual. La columna $j$ conserva el índice de posición del testigo histórico;
la operación de entrega al selector olvida este índice después de formar
el conjunto de palabras.

\begin{table}[htbp]
\centering\small
\begin{tabular}{rrcrrc c}
\toprule
$j$&$t$&carga&$\delta$&$\epsilon$&$(\eta_0\eta_1\eta_2)$&$L_{t,c,\delta,\epsilon}$\\
\midrule
4&11&$d$&0& 1&212&021101\\
5&67&$d$&1&$-1$&211&002001\\
5&67&$d$&2& 1&211&002001\\
5&67&$r$&1&$-1$&211&002001\\
5&67&$r$&2& 1&211&002001\\
6&81&$d$&0& 1&222&020111\\
6&81&$r$&2& 1&222&020111\\
7&34&$r$&1&$-1$&111&200110\\
7&75&$d$&1&$-1$&211&200110\\
7&75&$r$&2&$-1$&211&200110\\
8&90&$r$&0& 1&121&121012\\
8&90&$r$&1&$-1$&121&121012\\
9&49&$d$&0&$-1$&121&010122\\
11&11&$d$&2&$-1$&211&221110\\
11&37&$d$&0&$-1$&121&221110\\
12&8&$d$&0&$-1$&111&222110\\
12&8&$r$&1&$-1$&111&222110\\
12&58&$d$&1&$-1$&111&222110\\
12&58&$r$&0&$-1$&111&222110\\
\bottomrule
\end{tabular}
\caption{Incidencias fase--carga completas sobre la segmentación terminal
en las cien fronteras del registro. Las coincidencias de palabras entre
filas conservan sus tiempos, orientaciones y tipos de carga.}
\label{act21:tab:terminal-incidencias-completas}
\end{table}

El décimo bloque dispone de un testigo afín. En $t=30$ se tiene
$d_{30}=(0,1,2)$ y $\vartheta_{30}=2$, de donde
\[
 (d_{30,0}+\vartheta_{30},d_{30,1}+\vartheta_{30},
 d_{30,2}+\vartheta_{30})=(2,0,1)
\]
y, usando las bandas de la tabla~\ref{act21:tab:terminal-bandas-testigo},
\[
 L^{\mathrm{af}}_{30}
 =2(210112)+0(110202)+(211110)
 =(001001)\quad\text{en }\mathbb F_3^6.
\]
La igualdad es coordenada a coordenada: antes de reducir módulo tres,
el vector de la suma es $(6,3,1,3,3,4)$.

\begin{proposition}[Recuperación íntegra del repertorio por incidencias]
\label{act21:prop:repertorio-terminal-ocho-incidencias}
Considérese el registro de incidencias comparativas de la
tabla~\ref{act21:tab:terminal-incidencias-completas}, aumentado con el testigo
afín $(j,t,L)=(10,30,001001)$. Para cada posición $5\leq j\leq12$, el
conjunto de palabras de salida de sus incidencias es unitario. La
extracción de esas ocho palabras y el olvido posterior de su orden
producen exactamente
\[
 \mathcal S_{\mathrm{term}}=
 \{002001,020111,200110,121012,010122,001001,221110,222110\}.
\]
Su codificación como enteros de base tres es
\[
 \{28,55,98,175,437,498,687,714\}.
\]
\end{proposition}
\begin{proof}
La multiplicidad de incidencias por posición $j=5,6,7,8,9,10,11,12$
es, respectivamente,
\[
 (4,2,3,2,1,1,2,4).
\]
En cada posición, todas las filas imprimen la misma palabra. Cada fila
se verifica sustituyendo sus tres bandas y los tres coeficientes en la
fórmula de $L$; el décimo caso se ha calculado explícitamente mediante
$L^{\mathrm{af}}_{30}$. Por ejemplo, en $t=67$ y con coeficientes
$(2,1,1)$ se obtiene
\[
 2(010012)+(001220)+(011120)=(002001)\pmod3.
\]
Los ocho resultados son distintos como palabras de longitud seis. Su
codificación $\sum_{k=1}^{6}w_k3^{6-k}$ da, en el orden de las
posiciones, $(55,175,498,437,98,28,714,687)$. Al olvidar el orden se
obtiene el conjunto de enteros enunciado. Así, la extracción preserva
cada palabra completa, incluidas sus posiciones nulas, y entrega ocho
elementos distintos al selector.
\end{proof}

La prueba anterior es una recuperación exacta desde incidencias
identificadas. Su dato de entrada comprende el registro de incidencias
con sus posiciones y testigos. La selección orbital posterior recibe
únicamente $\mathcal S_{\mathrm{term}}$ y el descriptor regional inicial;
recorre las $8!$ ordenaciones y determina su compatibilidad con las
fronteras funcionales y las incidencias orientadas de la rueda
dodecafásica. Ambas operaciones quedan así contiguas, con dominios
explícitos: procedencia temporal de las ocho palabras y determinación
constructiva de su orden.

El cuadro comparativo contiene las diecinueve incidencias que su fórmula
produce sobre la segmentación de doce bloques en los cien tiempos
declarados. Esa exhaustividad finita se comprueba recorriendo
$100\cdot2\cdot3\cdot2=1200$ evaluaciones. El lector afín amplía la
cobertura de las posiciones terminales de siete a ocho. Este recuento
caracteriza el repertorio local de lectores que acaba de definirse;
las demás operaciones del TPK conservan los dominios y funciones
constructivas establecidos en sus respectivas derivaciones.

\subsection{Lectores comparativos, lector afín e incidencia temporal}

\subsubsection{Transformación de Witt del selector}

El selector utiliza la coordenatización de seis coordenadas
\[
 A_{\rm ter}=
 \begin{pmatrix}
 0&1&1&1&1&1\\
 1&0&1&2&2&1\\
 1&1&0&1&2&2\\
 1&2&1&0&1&2\\
 1&2&2&1&0&1\\
 1&1&2&2&1&0
 \end{pmatrix}.
\]
Se define $W(b)=bA_{\rm ter}$ en $\mathbb F_3^6$, así como las cargas
\[
 \eta(b)=\sum_{j=0}^5 b_j\pmod3,
 \qquad\eta_{\rm ter}^\vee(b)=\eta(W(b)).
\]
La coordenatización del descriptor y la del selector se relacionan mediante la
permutación $\sigma=(0,1,2,4,5,3)$, escrita como lista de imágenes:
\begin{equation}
 (A_{\rm ter})_{ij}=(A_{\rm reg})_{\sigma(i),\sigma(j)}.
 \label{act21:eq:kdes-cambio-carta}
\end{equation}

\begin{lemma}[Compatibilidad del funcional de carga]
Para toda banda $b$, se cumple
\[
 \eta_{\rm ter}^\vee(b)=\eta^\vee(b).
\]
\end{lemma}
\begin{proof}
La suma entera de cada fila correspondiente de ambas matrices coincide:
en la primera fila es cinco y en las otras cinco es siete. Por
distributividad,
\[
 \sum_j\sum_i b_i(A_{\rm ter})_{ij}
 =\sum_i b_i\sum_j(A_{\rm ter})_{ij}
 =\sum_i b_i\sum_j(A_{\rm reg})_{ij}.
\]
La reducción módulo tres da la igualdad de cargas. La identidad de
matrices \eqref{act21:eq:kdes-cambio-carta} se verifica por sustitución de sus
treinta y seis entradas y conserva, además, el cambio de coordenadas.
\end{proof}

La igualdad del funcional total de carga permite comparar sus
representaciones. Las transformaciones vectoriales siguen identificadas
por sus matrices y por la permutación explícita; la prueba anterior
conserva esa información de coordenadas.

\subsubsection{Doce lectores comparativos por instante}

Se define la fase temporal
\[
 \theta(t)=(t+2)\bmod3.
\]
Esta escritura realiza $(t-1)\bmod3$ también para $t=0$, donde la fase
vale dos. Para cada tipo de carga
$\nu\in\{\eta,\eta_{\rm ter}^\vee\}$, desplazamiento
$o\in\{0,1,2\}$ y orientación $s\in\{1,2\}$, se toman
\[
 \epsilon_\chi^{\nu,o,s}(t)=s
 \begin{cases}
 1,&\nu(B_\chi(t))=\theta(t)+o\pmod3,\\
 2,&\nu(B_\chi(t))\ne\theta(t)+o\pmod3.
 \end{cases}
\]
La palabra comparativa correspondiente es
\begin{equation}
 C_{\nu,o,s}(t)=
 \sum_{\chi\in\{\pi,e,\varphi\}}
 \epsilon_\chi^{\nu,o,s}(t)B_\chi(t)
 \quad\text{en }\mathbb F_3^6.
 \label{act21:eq:kdes-comparativo}
\end{equation}
Los parámetros generan doce representaciones etiquetadas por fila. Dos
etiquetas pueden producir la misma palabra; los testigos temporales y
las etiquetas permanecen disponibles en el calendario completo.

\subsubsection{Lector afín de fase y carga}

El lector afín emplea coeficientes
\[
 a_\chi^{\rm af}(t)=\eta_{\rm ter}^\vee(B_\chi(t))+\theta(t)
 \pmod3,
\]
y produce
\begin{equation}
 F(t)=\sum_{\chi\in\{\pi,e,\varphi\}}
 a_\chi^{\rm af}(t)B_\chi(t)
 \quad\text{en }\mathbb F_3^6.
 \label{act21:eq:kdes-afin}
\end{equation}
La regla comparativa distingue coincidencia y diferencia de cargas;
la afín suma carga y fase antes de combinar las bandas. Esta diferencia
operativa es la que utiliza la selección heterotípica.

\subsubsection{Incidencia palabra--residuo}

Sea $R$ la rotación cíclica de las seis posiciones de una banda. En la
fila de tiempo $t$ se forma el repertorio residual
\[
 \mathcal R(t)=\{R^jv:
      v\in\{q(t),a(t),c(t),\bar w(t)\},\ 0\leq j<6\}.
\]
Las relaciones de incidencia temporal son
\begin{align}
 \mathcal A_{\rm cmp}
 &=\{(C_{\nu,o,s}(t),u):
      0\leq t<100,\ u\in\mathcal R(t),\ \nu,o,s\text{ admisibles}\},\\
 \mathcal A_{\rm af}
 &=\{(F(t),u):0\leq t<100,\ u\in\mathcal R(t)\}.
 \label{act21:eq:kdes-atlas}
\end{align}
Cada par relaciona una palabra leída con un residuo de la misma fila.
El tiempo se cuantifica de manera común en ambas coordenadas del par.
En el calendario empleado, los tiempos $0,\ldots,99$ son distintos;
identificar una misma fila e identificar un mismo tiempo son, por ello,
operaciones equivalentes.

\begin{proposition}[Compresión existencial del repertorio temporal]
La función
\[
 M(v,u)=\mathbf1_{(v,u)\in\mathcal A_{\rm cmp}}
        +2\mathbf1_{(v,u)\in\mathcal A_{\rm af}}
        \in\{0,1,2,3\}
\]
conserva exactamente la existencia de cada tipo de lectura para el par
$(v,u)$.
\end{proposition}
\begin{proof}
El primer bit vale uno exactamente cuando existe un tiempo, un campo
residual, una rotación y una etiqueta comparativa que producen el par.
El segundo bit expresa la condición correspondiente al lector afín.
La unión por disyunción de bits conserva cada condición de existencia
al recorrer todas las filas. Un par puede admitir ambos tipos, y en ese
caso la máscara vale tres.
\end{proof}

La función $M$ es un lector de existencia. Los tiempos testigos, sus
multiplicidades y las etiquetas individuales permanecen en
\eqref{act21:eq:kdes-atlas} y en el calendario que las produce. La
implementación finita emplea $M$ para decidir heterotipia, mientras que
el archivo genealógico conserva los datos de procedencia más ricos.

\subsection{Enumeración de cilindros y candidatos decimales}

\subsubsection{Concatenación de setenta y dos y setenta y ocho trits}

Para una ordenación $\tau=(s_1,\ldots,s_8)$ de
$\mathcal S_{\rm term}$ se define
\[
 P_{72}(\tau)=[\Omega_{24}\Vert s_1\Vert\cdots\Vert s_8]_3.
\]
En forma recursiva, se parte de $p_0=[\Omega_{24}]_3$ y se calcula
$p_i=729p_{i-1}+[s_i]_3$ para $1\leq i\leq8$. Esta recurrencia
conserva los ceros iniciales de cada bloque por la multiplicación
explícita por $729$. Para cada posición $b$ del panel funcional,
\[
 P_{78}(\tau,b)=729P_{72}(\tau)+b.
\]
Las longitudes son, respectivamente, $24+8\cdot6=72$ y $72+6=78$.
El cilindro ternario correspondiente es
\begin{equation}
 I_{\tau,b}=
 \left[\frac{P_{78}(\tau,b)}{3^{78}},
       \frac{P_{78}(\tau,b)+1}{3^{78}}\right).
 \label{act21:eq:kdes-cilindro}
\end{equation}

\subsubsection{Intersección exacta con la rejilla de doce tríadas}

La lectura decimal utiliza doce tríadas, es decir, la rejilla
$N/10^{36}$. Se abrevia $D=10^{36}$ y $Q=3^{78}$.

\begin{lemma}[Extremos enteros del cilindro]
Para un prefijo entero $p$ de longitud setenta y ocho, los numeradores
de los puntos de la rejilla que pertenecen a su cilindro son
\begin{equation}
 \left\{N\in\mathbb Z:
 \left\lceil\frac{pD}{Q}\right\rceil
 \leq N<
 \left\lceil\frac{(p+1)D}{Q}\right\rceil\right\}.
 \label{act21:eq:kdes-rejilla}
\end{equation}
Cada cilindro contiene a lo sumo un punto de esta rejilla.
\end{lemma}
\begin{proof}
La condición $N/D\in[p/Q,(p+1)/Q)$ equivale a
$pD/Q\leq N<(p+1)D/Q$. Para $N$ entero, la primera desigualdad
equivale a $\lceil pD/Q\rceil\leq N$; la segunda equivale a
$N<\lceil(p+1)D/Q\rceil$. La longitud del intervalo en coordenada
$N$ es $D/Q<1$, pues $10^{36}<3^{78}$. Por ello puede contener,
como máximo, un entero.
\end{proof}

Los techos se evalúan por división entera:
$\lceil a/Q\rceil=\lfloor(a+Q-1)/Q\rfloor$ para $a\geq0$.
Toda la enumeración utiliza enteros exactos; la elección de puntos de
frontera queda fijada por el extremo derecho semiabierto.

Se conserva primero la lista indexada de incidencias
\[
 \mathcal C_{\rm ind}=
 \{(\tau,j,N):N/D\in I_{\tau,\mathcal P_9(j)}\},
 \qquad1\leq j\leq9,
\]
y se obtiene después el conjunto de numeradores
\[
 \mathcal C=\{N:\exists\tau,j\ (\tau,j,N)\in\mathcal C_{\rm ind}\}.
\]
El paso a $\mathcal C$ elimina repeticiones de un numerador, conservadas
en $\mathcal C_{\rm ind}$ con su procedencia.

\begin{lemma}[Recuperación de la ordenación desde el punto del cilindro]
Si $(\tau,j,N)\in\mathcal C_{\rm ind}$, el bloque número $i$ de la
palabra de setenta y ocho trits se recupera mediante
\[
 \left\lfloor\frac{N\,729^i}{D}\right\rfloor\bmod729,
 \qquad1\leq i\leq13.
\]
En particular, los bloques quinto a duodécimo recuperan $\tau$.
\end{lemma}
\begin{proof}
La pertenencia al cilindro significa
$p\leq QN/D<p+1$, donde $p=P_{78}(\tau,\mathcal P_9(j))$.
Por tanto, $\lfloor QN/D\rfloor=p$. La extracción de los bloques
anteriores coincide con la división sucesiva del prefijo $p$ por
potencias de $729$: truncar una expansión posicional antes de extraer
un bloque conserva dicho bloque. La fórmula escrita realiza
exactamente esa extracción. Los cuatro primeros bloques forman
$\Omega_{24}$ y los ocho siguientes son la ordenación $\tau$.
\end{proof}

\begin{proposition}[Censo exacto de candidatos]
Para $\Omega_{24}$, $\mathcal S_{\rm term}$ y $\mathcal P_9$
definidos anteriormente,
\[
 \#\operatorname{Ord}(\mathcal S_{\rm term})=40320,
 \qquad\#\mathcal C_{\rm ind}=21902,
 \qquad\#\mathcal C=19446.
\]
\end{proposition}
\begin{proof}
Los ocho bloques son distintos, por lo que sus ordenaciones son las
$8!=40320$ permutaciones. Para cada permutación y cada una de las
nueve posiciones del panel se calcula $p=P_{78}(\tau,b)$, se aplican
los extremos \eqref{act21:eq:kdes-rejilla} y se añade su único entero cuando
el intervalo entero es no vacío. La suma de los cardinales de esos
intervalos es $21902$. La unión de sus enteros contiene $19446$
elementos distintos. La evaluación exhaustiva de estas operaciones
está certificada por los teoremas
\texttt{indexed\_cylinder\_count} y \texttt{candidate\_count} del
selector finito. El algoritmo realiza exactamente la enumeración
descrita: permutaciones, nueve incidencias del panel, límites por
división entera y eliminación de duplicados.
\end{proof}

Los números $21902$ y $19446$ cuentan objetos diferentes. El primero
cuenta incidencias indexadas que contienen un punto decimal; el segundo
cuenta numeradores distintos. El número total de consultas de cilindro
es $40320\cdot9=362880$, incluidos los cilindros cuya intersección con
la rejilla es vacía.

\subsection{Incidencia excepcional y condición de heterotipia}

\subsubsection{Cara excepcional producida por el panel}

Se eleva una banda a doce coordenadas mediante
\[
 \Gamma(b)=(b,W(b))\in\mathbb F_3^{12}.
\]
El soporte de una palabra es el conjunto de posiciones de sus
coordenadas no nulas, con índices humanos $1,\ldots,12$. Las primeras
bandas regionales del panel producen
\begin{equation}
 \mathcal B=\operatorname{supp}\Gamma(B_\pi(0))
          \cap\operatorname{supp}\Gamma(B_e(0))
          =\{4,6,9,10\}.
 \label{act21:eq:kdes-cara}
\end{equation}
La operación utiliza las palabras regionales $103$ y $523$ y la
matriz de Witt del selector. La cara se calcula antes de consultar las
coordenadas de los candidatos.

Para $N\in\mathcal C$, sus doce tríadas son
\begin{equation}
 K_i(N)=\left\lfloor\frac{N}{1000^{12-i}}\right\rfloor\bmod1000,
 \qquad1\leq i\leq12.
 \label{act21:eq:kdes-triadas}
\end{equation}
El soporte de corona es
\[
 \mathcal H(N)=\{i:K_i(N)\geq729\},
 \qquad
 \mathcal C_{\mathcal B}=\{N\in\mathcal C:\mathcal H(N)=\mathcal B\}.
\]
La frontera $729=3^6$ relaciona la banda de seis trits con la tríada
decimal. La evaluación exhaustiva de la igualdad de soportes da
\begin{equation}
 \#\mathcal C_{\mathcal B}=79.
 \label{act21:eq:kdes-79}
\end{equation}
El soporte conserva posiciones, mientras que el registro completo
conserva también las doce magnitudes. La etapa siguiente utiliza esa
información posicional junto con relaciones entre magnitudes.

\subsubsection{Palabras ternarias y diferencias orientadas del candidato}

La $i$-ésima banda ternaria del punto $N/D$ se obtiene por
\begin{equation}
 T_i(N)=\left\lfloor\frac{N\,729^i}{D}\right\rfloor\bmod729,
 \qquad1\leq i\leq13.
 \label{act21:eq:kdes-tercandidato}
\end{equation}
La arista orientada $a\to b$ tiene residuo
\begin{equation}
 R_{a,b}(N)=(K_a(N)-K_b(N))\bmod729.
 \label{act21:eq:kdes-resarista}
\end{equation}
La resta se realiza en $\mathbb Z$ antes de la reducción. Para esta
arista se consulta en las relaciones \eqref{act21:eq:kdes-atlas} el par
\[
 \bigl(\operatorname{enc}_6^{-1}(T_b(N)),
  \operatorname{enc}_6^{-1}(R_{a,b}(N))\bigr).
\]
Es, por tanto, una consulta
conjunta a la palabra de destino y a la diferencia orientada de las
tríadas.

Se escribe $\mathrm{Cmp}(a,b;N)$ cuando el par pertenece a
$\mathcal A_{\rm cmp}$ y $\mathrm{Af}(a,b;N)$ cuando pertenece a
$\mathcal A_{\rm af}$.

\subsubsection{Determinación de la órbita transversal}

El descriptor de veinticuatro trits determina las tres primeras
tríadas decimales antes de las filtraciones terminales. En efecto, su
cilindro satisface la inclusión exacta
\[
 \left[\frac{66241910521}{3^{24}},
       \frac{66241910522}{3^{24}}\right)
 \subseteq
 \left[\frac{234543140}{10^9},
       \frac{234543141}{10^9}\right).
\]
Las dos desigualdades de extremos se verifican por productos cruzados
enteros. Por tanto, todo candidato comparte las tres posiciones de
referencia $1,2,3$ y sus representaciones $(234,543,140)$.

La traslación de paso cuatro sobre $\mathbb Z/12\mathbb Z$ posee las
cuatro órbitas
\[
 (1,5,9),\qquad(2,6,10),\qquad(3,7,11),\qquad(4,8,12).
\]
Las tres primeras contienen, respectivamente, una de las posiciones de
referencia; la cuarta es la única disjunta de $\{1,2,3\}$. El bosque
orientado de los pasos tres y cuatro se fija mediante las aristas
\[
 \begin{array}{lll}
 1\to4,&1\to5,&2\to6,\\
 3\to7,&4\to8,&5\to9,\\
 6\to10,&7\to11,&8\to12.
 \end{array}
\]
Su primera arista tiene paso tres y las restantes paso cuatro. Las dos
aristas internas de la órbita distinguida son precisamente
$4\to8$ y $8\to12$. Así queda fijado el lugar de la condición de
lectura por la geometría de las posiciones y por el prefijo inicial,
antes de examinar los candidatos de la cara excepcional.

\subsubsection{Condición terminal de heterotipia}

La condición terminal es
\begin{equation}
 \begin{split}
 \mathcal T(N)={}&
 [\mathrm{Cmp}(4,8;N)\wedge\mathrm{Af}(8,12;N)]\\
 &\vee[\mathrm{Af}(4,8;N)\wedge\mathrm{Cmp}(8,12;N)].
 \end{split}
 \label{act21:eq:kdes-heterotipia}
\end{equation}
Se admiten las dos asignaciones de tipos a las dos aristas. La
selección exige que exista una realización con tipos distintos; las
etiquetas individuales y sus tiempos quedan disponibles en los
testigos de incidencia.

\begin{theorem}[Selección terminal del registro de acoplamiento]
Sobre el dominio de candidatos producido por
$\Omega_{24}$, $\mathcal S_{\rm term}$, $\mathcal P_9$ y el calendario
regional de cien filas, existe un único numerador que satisface
simultáneamente la incidencia excepcional y la heterotipia:
\[
 \{N\in\mathcal C_{\mathcal B}:\mathcal T(N)\}
 =\{234543140729659824621058914794146601\}.
\]
Su registro de doce tríadas es
\begin{equation}
 K=(234,543,140,729,659,824,621,058,914,794,146,601).
 \label{act21:eq:kdes-kfinal}
\end{equation}
\end{theorem}
\begin{proof}
El censo precedente produce los $79$ elementos de
$\mathcal C_{\mathcal B}$. Para cada uno se extraen sus doce
coordenadas por \eqref{act21:eq:kdes-triadas}; se calculan las bandas de
destino $T_8,T_{12}$ y los residuos $R_{4,8},R_{8,12}$; se consultan
las dos máscaras del repertorio temporal y se aplica la disyunción
\eqref{act21:eq:kdes-heterotipia}. La lista resultante contiene exactamente
el entero indicado, según la evaluación exhaustiva certificada por
\texttt{terminal\_selection}. La igualdad de lista con un singleton
implica existencia y unicidad en el dominio declarado. La extracción
base mil de ese entero da \eqref{act21:eq:kdes-kfinal}.

El procedimiento de selección recibe los productores y las reglas
enumeradas. El entero final interviene en el enunciado de la igualdad
comprobada, después de ejecutar la selección. Esta posición lógica
permite cotejar el resultado sin utilizarlo como criterio de elección.
\end{proof}

\subsubsection{Testigo aritmético de pertenencia a un cilindro}

Una de las incidencias del numerador seleccionado utiliza la ordenación
\[
 (002001,020111,200110,121012,010122,001001,221110,222110)
\]
y la frontera $438=121020_3$. Para ella,
\[
 P_{72}=5283881584882673136514102926090957,
\]
\[
 P_{78}=3851949675379468716518781033120308091.
\]
La fórmula de extremos enteros produce
\begin{align*}
 \left\lceil\frac{P_{78}10^{36}}{3^{78}}\right\rceil
 &=234543140729659824621058914794146601,\\
 \left\lceil\frac{(P_{78}+1)10^{36}}{3^{78}}\right\rceil
 &=234543140729659824621058914794146602.
\end{align*}
El único entero del intervalo es, por tanto, $N_K$. Las trece bandas
ternarias de su punto decimal son
\[
 \begin{array}{rrrr}
 020022&222111&211201&021101\\
 002001&020111&200110&121012\\
 010122&001001&221110&222110\\
 121020&&&
 \end{array}
\]
Este cálculo exhibe un testigo concreto de pertenencia al dominio
enumerado. La unicidad del registro procede del examen exhaustivo de
todos los candidatos, además de este testigo particular.

\subsection{Composición regional y continuidad de las realizaciones}

\subsubsection{Composición de los generadores regionales}

Sea $\operatorname{Sel}(\Omega,\mathcal S,\mathcal P,\mathcal N)$ el
selector definido por la enumeración y los dos filtros anteriores.
Los productores regionales satisfacen las igualdades
\[
 \Omega_{\rm reg}=\Omega_{24},\qquad
 \mathcal P_{\rm reg}=\mathcal P_9,\qquad
 \mathcal N_{\rm reg}=\mathcal N_{100},
\]
donde $\mathcal N_{\rm reg}$ se calcula por extracción de bandas y
lectura de firmas. Por congruencia de funciones,
\begin{equation}
 \begin{split}
 &\operatorname{Sel}(\Omega_{\rm reg},\mathcal S_{\rm term},
                    \mathcal P_{\rm reg},\mathcal N_{\rm reg})\\
 &\hspace{10mm}=
 \operatorname{Sel}(\Omega_{24},\mathcal S_{\rm term},
                    \mathcal P_9,\mathcal N_{100})=\{N_K\}.
 \end{split}
 \label{act21:eq:kdes-composicion}
\end{equation}
La igualdad reemplaza tres interfaces documentales por su construcción
regional efectiva. El repertorio terminal mantiene la procedencia
especificada en \eqref{act21:eq:kdes-sterm}.

El registro se forma entonces por las doce divisiones euclídeas de
\eqref{act21:eq:kdes-triadas}. Cada coordenada pertenece a
$\{0,\ldots,999\}$ y la longitud es doce por construcción. La
implementación que toma el primer elemento de la lista seleccionada
está justificada por la igualdad con el singleton: el valor por defecto
de una operación informática sobre listas vacías nunca interviene en
este resultado.

\subsubsection{Relación con las incidencias firmadas y la inversión}

La construcción presente proporciona una selección prospectiva en un
dominio finito explícito. El desarrollo ulterior del registro
dodecafásico produce sus canales orientados a partir del estado
terminal enriquecido y demuestra una inversión integral mediante
Hadamard. Aquella inversión recupera un registro desde canales
compatibles; aquí se selecciona un registro a partir de sus
antecedentes regionales e incidenciales. La igualdad de sus resultados
permite componer las dos realizaciones manteniendo el origen de cada
operación.

La unicidad de selección y la unicidad de inversión son propiedades
distintas. La primera cuantifica los candidatos construidos en este
capítulo. La segunda cuantifica los registros que comparten una firma
compatible. Su conjunción justifica que el registro seleccionado pueda
transportarse a la realización firmada y recuperarse desde ella, sin
convertir una recuperación retrospectiva en el productor del registro.

\subsubsection{Relación entre los censos terminales}

La secuencia
\[
 40320\ \longrightarrow\ 21902\ \longrightarrow\ 19446
 \ \longrightarrow\ 79\ \longrightarrow\ 1
\]
registra, sucesivamente, ordenaciones, incidencias indexadas con la
rejilla, numeradores distintos, candidatos de la cara excepcional y
supervivientes heterotípicos. La notación de flechas indica el orden de
las operaciones; la segunda cifra cuenta una clase de objetos
distinta de la primera.

El censo de catálogos regionales
$196\cdot197\cdot196\longrightarrow331\longrightarrow2
\longrightarrow1$, desarrollado en su lugar correspondiente, parte de
otra población: bloques regionales de tres catálogos y condiciones
terminales de Gram, ocupación, distancia y orientación. Ambos censos
conservan sus dominios. Su convergencia en una misma representación se
expresa mediante la identificación de sus lectores y de sus salidas;
las cifras intermedias describen operaciones diferentes y permanecen
documentadas por separado.

\subsubsection{Composición regional de la constante de estructura fina}

La salida $K$ interviene a continuación en la composición regional
ordenada y en el acarreo de base mil. Su empleo posterior conserva las
doce posiciones y la frontera de lectura. El acoplamiento con la raíz
analítica y las representaciones a profundidad arbitraria utiliza esa
misma salida, junto con las coordenadas regionales ya generadas y las
reglas de la coordenatización analítica. Esta dependencia prepara el capítulo
dedicado a las dos realizaciones correlacionadas de la constante de
estructura fina.

El alcance de este capítulo es preciso: producción del descriptor y
de sus antecedentes regionales, determinación del dominio finito,
selección única sobre dicho dominio y composición material con los
productores. La recurrencia analítica de todos los grados y la
compatibilidad de los acarreos infinitos se desarrollan después con sus
demostraciones propias.

\subsection{Certificación ejecutable y trazabilidad de las operaciones}

El desarrollo formal conserva la correspondencia entre las
definiciones matemáticas y sus realizaciones ejecutables. Las pruebas
de firmas y de reloj se encuentran en
\texttt{N69RegionalSignature.lean}; la extracción de bandas y la
reproducción de las cien filas, en
\texttt{GeneratedN69Rows.lean}; el criterio mínimo y la construcción
de $\Omega_{24}$, en \texttt{RegionalW24.lean}. La compatibilidad de
las dos matrices de Witt está reunida en
\texttt{WittReaderCompatibility.lean}.

Los lectores temporales, las rotaciones, los residuos orientados y la
condición heterotípica se formalizan en
\texttt{TerminalReaders.lean}. La enumeración y los censos se
formalizan en \texttt{Terminal}\allowbreak\texttt{Orbital}\allowbreak\texttt{Selection.lean}. El panel
regional se enlaza mediante \texttt{Regional}\allowbreak\texttt{Panel}\allowbreak\texttt{Bridge.lean}, y la
composición \eqref{act21:eq:kdes-composicion} queda expresada en
\texttt{Selected}\allowbreak\texttt{From}\allowbreak%
\texttt{Regional}\allowbreak\texttt{Inputs.lean}.

Las comprobaciones finitas de enumeración emplean la evaluación nativa
de Lean; sus informes registran \texttt{Lean.ofReduceBool}, además
de los axiomas lógicos habituales de la biblioteca. La reconstrucción
regional del descriptor utiliza pruebas del núcleo y reducción
ordinaria. Esta distinción describe los mecanismos de comprobación
realmente empleados. Los recibos de compilación acreditan sus módulos
y dependencias declarados; las pruebas matemáticas conservan los
dominios, productores e hipótesis que se han explicitado en el cuerpo
del capítulo.

\clearpage
% END NUCLEO_ADITIVO_20260921_SELECCION
\import{foundation/}{sections/registro_k.tex}
\import{foundation/}{sections/alpha.tex}
\import{foundation/}{sections/excepcional.tex}
\import{foundation/}{sections/01_hilbert_nonadico.tex}
\endgroup
\clearpage
\part{Geometría positiva, refinamiento y dualidad reticular}
\section{Orientación trítica, prolongación compatible y doble proyección}
\label{sec:trit-projection}

El carácter generado de las coordenadas reales se expresa en la relación
entre profundidad aritmética y evaluación. El dato primitivo es una historia
admisible de operaciones con sus residuos, cocientes y orientaciones. Una
lectura arquimediana asocia a sus prefijos intervalos cada vez más estrechos;
una lectura incidencial conserva relaciones de frontera del mismo estado.
La doble proyección del corpus \cite{hmtI,HMT-IX} reúne esas dos operaciones
sin convertir una evaluación escalar en la totalidad de la historia. El
artículo utiliza ese antecedente para estudiar las realizaciones posteriores
del registro dodecafásico.

\subsection{Cilindros posicionales y determinación de la coordenada real}

Para un bloque ternario \(b=(b_1,\ldots,b_6)\), con representantes
\(b_j\in\{0,1,2\}\) de \(\mathbb F_3\), la evaluación entera
\(\nu(b)=\sum_{j=1}^6b_j3^{6-j}\) pertenece a
\(\{0,\ldots,728\}\). Si \(b_{n+1}\) es el bloque admisible siguiente,
la prolongación posicional satisface
\[
 N_{n+1}=729N_n+\nu(b_{n+1}),\qquad
 a_n=\frac{N_n}{729^n},\qquad b_n^+=\frac{N_n+1}{729^n}.
\]
Aquí \(N_0\) es la parte entera que determina el lector regional. La letra
\(b_n^+\) designa un extremo real y se distingue del bloque ternario \(b\).
El cilindro de representación es \([a_n,b_n^+)\); su clausura se utiliza
en el argumento de existencia del límite. El acarreo conserva el paso de
un prefijo al anterior mediante división euclídea.

\begin{proposition}[Límite de los centros cilíndricos]
Para una prolongación compatible, las clausuras de los cilindros están
encajadas y determinan una única coordenada \(x\). Si
\(c_n=(a_n+b_n^+)/2\), entonces
\[
 x=\lim_{n\to\infty}a_n=\lim_{n\to\infty}b_n^+
 =\lim_{n\to\infty}c_n,\qquad
 |x-c_n|\leq\frac1{2\,729^n}.
\]
\end{proposition}
\begin{proof}
La desigualdad \(0\leq\nu(b_{n+1})<729\) sitúa el cilindro sucesor en
la clausura del predecesor. Sus extremos son racionales, su anchura es
\(729^{-n}\) y permanecen dentro del primer intervalo cerrado acotado.
El encajamiento y la anchura decreciente producen una única clase de
Cauchy de extremos. En la realización arquimediana, esa clase tiene el
representante \(x\); la distancia al punto medio es a lo sumo la mitad
de la anchura. La convención de frontera asigna una escritura cuando dos
expansiones posicionales representan el mismo límite.
\end{proof}

La media aritmética interviene, por tanto, en una evaluación precisa: el
centro de cada intervalo que acota la coordenada. El límite conserva la
exactitud del valor aunque la anchura desaparezca. La representación
enriquecida conserva adicionalmente la sucesión de prefijos y su incidencia.
La proyección numérica y la proyección excepcional son dos aplicaciones
con codominios diferentes; una conserva el valor límite y la otra transporta
las relaciones que habilitan la construcción reticular. El paso a
\(\RR\) se sitúa así en la realización del lector, después de la
generación aritmética de los datos que evalúa.

Para hacer explícita la compatibilidad, sean \(p_{n+1,n}\) las restricciones
del estado enriquecido y \(x_n=p_{n+1,n}(x_{n+1})\). Si \(I_n(x_n)\)
es el cilindro numérico y \(\mathcal B_n(x_n)\) la incidencia de frontera,
los dos transportes conservan simultáneamente
\[
 \overline{I_{n+1}(x_{n+1})}\subseteq\overline{I_n(x_n)},\qquad
 r_{n+1,n}\mathcal B_{n+1}(x_{n+1})=\mathcal B_n(x_n).
\]
La segunda igualdad utiliza los mapas incidenciales del corpus; la
proposición~\ref{prop:lf-incidence-naturality} explicita su realización
mediante marcos y códigos transportados. Sus cinco
construcciones consustanciales ---espectral, solenoidal, gauge, cohomológica
y determinantal--- pertenecen a ese mismo sistema compatible. El estudio
geométrico posterior conserva esta unidad de origen al especificar cuál
de sus observaciones utiliza en cada teorema.

\subsection{Trayectorias orientadas y temporalidad de la actualización}

Sea \(x\xrightarrow{\rm adm}y\) la relación de admisibilidad del TPK en
el espacio enriquecido \(\widehat X\). Una trayectoria parametrizada por
un conjunto discretamente ordenado \(\mathcal T\) es
\[
 \gamma:\mathcal T\longrightarrow\widehat X,
 \qquad\gamma(t)\xrightarrow{\rm adm}\gamma(t^+).
\]
La relación de compatibilidad determina las transiciones; la parametrización
permite ordenarlas. En el espacio \(E=\mathcal T\times\widehat X\),
con \(p(t,x)=t\), la trayectoria determina la sección
\(\widetilde\gamma(t)=(t,\gamma(t))\), que satisface
\(p\circ\widetilde\gamma=\id_{\mathcal T}\).
Esta precisión conserva la diferencia entre un estado, su trayectoria y
la coordenada utilizada para describir la evolución.

La firma trítica \(\tau\in\{+1,0,-1\}\) tipa los generadores mediante
\[
 J_\tau^2=-\tau I.
\]
En el lector temporal del corpus, \(+1\) representa retorno y memoria;
\(0\), actualización de umbral; y \(-1\), apertura de prolongaciones
conjugadas. Su interpretación como pasado, presente y futuro se refiere
a las funciones del antecedente conservado, la transición efectiva y la
continuación admisible. Los tres tiempos verbales adquieren así un
contenido operatorio sobre historias.

Si \(\mathcal P_n(x)\) es el conjunto de continuaciones admisibles de
longitud \(n\), la supresión del último paso define
\(r_n:\mathcal P_{n+1}(x)\to\mathcal P_n(x)\). Una continuación
ilimitada es una colección \((\gamma_n)\) con
\(r_n(\gamma_{n+1})=\gamma_n\). La restricción reconstruye antecedentes;
la actualización \(U_t\) actúa en la transición; la prolongación mantiene
la frontera del prefijo. La existencia de estas colecciones se utiliza en
los dominios supervivientes fijados por el teorema de prolongación del
corpus.

Sobre una firma finita \(\sigma=(\tau_1,\ldots,\tau_n)\), la inversión
orientada es
\[
 C_{\rm rev}\sigma=(-\tau_n,\ldots,-\tau_1),
 \qquad C_{\rm rev}^2\sigma=\sigma.
\]
La trayectoria conjugada conserva también las transformaciones de hoja,
posición y frontera de su transporte. Por ello la bidireccionalidad se
expresa en la historia completa, mientras esta fórmula expresa su firma.
El retorno de fase de \(\Gamma_9\) admite un incremento de memoria:
estados con fase idéntica pueden ocupar profundidades diferentes.

\subsection{Álgebras cuadráticas y actividad del umbral parabólico}

En un régimen fijo, la conjugación de \(aI+bJ_\tau\) cambia el signo
de \(b\), y su norma algebraica es
\[
 (aI+bJ_\tau)(aI-bJ_\tau)=(a^2+\tau b^2)I.
\]
Los tipos elíptico, parabólico e hiperbólico corresponden a las tres
relaciones cuadráticas. Sus grupos locales se representan por
\[
 e^{tJ_{+1}}=\cos t\,I+\sin t\,J_{+1},\quad
 e^{tJ_0}=I+tJ_0,\quad
 e^{tJ_{-1}}=\cosh t\,I+\sinh t\,J_{-1}.
\]
La prueba separa términos pares e impares en la serie exponencial; en el
tipo nilpotente sólo permanecen los grados cero y uno. Cada colección
satisface la ley de composición en su parámetro y la inversión
\(e^{-tJ_\tau}=(e^{tJ_\tau})^{-1}\).

El valor trítico cero permite, en particular, un generador distinto de
cero con cuadrado nulo. Para
\(N=\left(\begin{smallmatrix}0&1\\0&0\end{smallmatrix}\right)\),
\[
 e^{tN}(q,p)=(q+tp,p).
\]
El umbral tiene una acción lineal efectiva sobre su fibra. Su parámetro,
su desplazamiento y la memoria de la transición permanecen determinados
aunque la firma de régimen sea cero. Éste es el contenido preciso que
permite desarrollar narrativamente el presente como actualización activa.

La forma elíptica positiva tiene niveles circulares en coordenadas
ortonormales y elípticos tras una transformación lineal; la forma hiperbólica
posee dos direcciones isotrópicas y niveles hiperbólicos; el tipo parabólico
posee una forma degenerada y transporte unipotente. Una geometría esférica
requiere además el espacio homogéneo o la métrica que la realiza. La
clasificación anterior fija los generadores que pueden intervenir en esa
realización y mantiene separados el tipo algebraico y la curvatura de un
espacio posterior.

El calendario de modos concreta la conexión con la autoescala. La regla
\(+1:m\mapsto\Sigma\), \(-1:m\mapsto\Pi\), \(0:m\mapsto m\)
aplicada al bloque cíclico \((+1,-1,0)\) produce
\((\Sigma,\Pi,\Pi)\). Contar sus tres transiciones, en el orden
\((\Sigma,\Pi)\), determina
\begin{equation}
 F_{\rm av}=\begin{pmatrix}0&1\\1&1\end{pmatrix}.
 \label{eq:geo-fav}
\end{equation}
Su rayo propio positivo, normalizado como \((1,x)^\mathsf T\), satisface
\(x=\lambda\), \(1+x=\lambda x\), y por ello
\(x^2-x-1=0\). La raíz positiva es la coordenada de autoescala expresada
como \(\varphi\). La matriz nace del calendario de operaciones; la
expresión radical y las identidades exponenciales siguientes son sus
realizaciones posteriores.

% Articulación posterior a la generación de pi, phi y e.
% Contenido acumulativo; no reemplaza el desarrollo previo del TRIT.
\subsection{Realizaciones exponenciales de las coordenadas generadas}
\label{geo-euler-trit-articulacion}

Las coordenadas regionales ya determinadas permiten reconocer las
realizaciones concretas de la colección cuadrática del TRIT. El orden es
sustantivo: la construcción de las regiones precede a la evaluación de una
exponencial en los parámetros \(\pi\) o \(2\log\varphi\). Estas expresiones
identifican la función de las coordenadas en operadores construidos desde
APP y TPK; no intervienen en la selección de los prefijos que las producen.

\subsubsection{Ciclo ternario, transporte nilpotente e incidencia de hojas}

La multiplicación \(a(r)=4r\) en \(\mathbb Z/9\mathbb Z\) tiene orden
\(3\). Sobre las unidades determina las órbitas \(\{1,4,7\}\) y
\(\{2,8,5\}\). En el módulo de aumento de cualquiera de ellas, formado por
las combinaciones enteras cuyos coeficientes suman cero, la base
\((e_r-e_{4r},e_{4r}-e_{7r})\) da
\[
 C=\begin{pmatrix}0&-1\\1&-1\end{pmatrix},
 \qquad C^2+C+I=0.
\]
La forma restringida tiene matriz de Gram
\(\bigl(\begin{smallmatrix}2&-1\\-1&2\end{smallmatrix}\bigr)\).
En las evaluaciones residuales módulo \(9\), la acción distingue las hojas:
\(S(ax,ay)=aS(x,y)\), mientras \(P(ax,ay)=a^{-1}P(x,y)\).
La realización cíclica conserva, por tanto, una orientación contragrediente
entre suma y producto. Los cocientes del levantamiento entero se conservan
en el registro aritmético correspondiente.

El transporte parabólico elemental se representa mediante
\[
 N=\begin{pmatrix}0&1\\0&0\end{pmatrix}.
\]
La incidencia areal--volumétrica ya construida en \eqref{eq:geo-fav} proporciona
\[
 F_{\rm av}=\begin{pmatrix}0&1\\1&1\end{pmatrix},
 \qquad A=F_{\rm av}^2=\begin{pmatrix}1&1\\1&2\end{pmatrix}.
\]
Los tres operadores actúan en espacios reales declarados: el módulo de
aumento tensorizado con \(\mathbb R\), el plano del transporte nilpotente y
el plano de incidencia modal. Compartir dimensión matricial no identifica
estos espacios ni sus coordenadas.

\begin{proposition}[Generadores concretos de los tres tipos]
\label{geo-euler-trit-generadores}
Los operadores
\[
 J_{\mathrm e}=\frac{2C+I}{\sqrt3},\qquad
 J_{\mathrm p}=N,\qquad
 J_{\mathrm h}=\frac{2A-3I}{\sqrt5}
\]
satisfacen, respectivamente,
\[
 J_{\mathrm e}^{\,2}=-I,\qquad
 J_{\mathrm p}^{\,2}=0,\qquad
 J_{\mathrm h}^{\,2}=I.
\]
\end{proposition}
\begin{proof}
La relación de \(C\) implica \((2C+I)^2=-3I\).
La multiplicación directa da \(N^2=0\).
Finalmente, \(A\) tiene traza \(3\) y determinante \(1\), por lo que
\(A^2-3A+I=0\) y \((2A-3I)^2=5I\).
\end{proof}

\begin{proposition}[Tres evaluaciones de la identidad exponencial]
\label{geo-euler-trit-evaluaciones}
En las realizaciones anteriores,
\begin{equation}
 \boxed{
 \exp(\pi J_{\mathrm e})+I=0,\qquad
 \exp(J_{\mathrm p})=I+J_{\mathrm p},\qquad
 \exp(2\log\varphi\,J_{\mathrm h})=F_{\rm av}^{\,2}.}
 \label{geo-euler-trit-tres-evaluaciones}
\end{equation}
\end{proposition}
\begin{proof}
La especialización elíptica de la exponencial trítica da
\(\cos\pi\,I+\sin\pi\,J_{\mathrm e}=-I\).
La nilpotencia elimina exactamente los términos de grado al menos \(2\)
en la segunda expresión. Para la tercera, \(\varphi^2=\varphi+1\) implica
\[
 \cosh(2\log\varphi)=\frac32,\qquad
 \sinh(2\log\varphi)=\frac{\sqrt5}{2}.
\]
Por tanto, la exponencial vale
\(\frac32I+\frac{\sqrt5}{2}J_{\mathrm h}=A\).
\end{proof}

La primera igualdad representa la media vuelta elíptica; la vuelta completa
satisface \(\exp(2\pi J_{\mathrm e})=I\). La segunda mantiene un término
lineal no nulo: el régimen parabólico expresa transporte, no ausencia de
evolución. La tercera realiza un avance hiperbólico unimodular. Sus
autovalores son \(\varphi^2\) y \(\varphi^{-2}\); una dirección se expande y
la otra se contrae conservando el determinante. Las tres evaluaciones
emplean parámetros distintos de una colección exponencial común.

En esa colección, \(e\) interviene mediante la evaluación escalar
\(\exp(I)=eI\) y la ley \(\exp(sI)\exp(tI)=\exp((s+t)I)\).
Su papel no se restringe al tipo parabólico. La coordenada \(e\), obtenida
previamente de su región, es reconocida aquí como la evaluación unitaria
de la exponencial; esta identidad no sustituye su genealogía coinductiva.

\subsubsection{Resolución polinómica de los regímenes}

Las tres realizaciones admiten una presentación compacta sobre la suma
directa de sus espacios. Defínanse
\[
 \mathbb J=J_{\mathrm e}\oplus J_{\mathrm p}\oplus J_{\mathrm h},
 \qquad Q=\mathbb J^2.
\]
Las relaciones anteriores implican
\[
 \mathbb J^6=\mathbb J^2,\qquad Q^3=Q.
\]
El uso de \(Q\) mantiene separado este cuadrado operatorial del registro
dodecafásico \(K\).

\begin{proposition}[Proyectores polinómicos de régimen]
\label{geo-euler-trit-proyectores}
Las aplicaciones
\[
 p_{\mathrm e}=\frac{Q^2-Q}{2},\qquad
 p_{\mathrm p}=I-Q^2,\qquad
 p_{\mathrm h}=\frac{Q^2+Q}{2}
\]
son idempotentes mutuamente ortogonales, suman \(I\) y recuperan los tres
sumandos de la representación.
\end{proposition}
\begin{proof}
Los tres polinomios toman los valores \((1,0,0)\), \((0,1,0)\) y
\((0,0,1)\) en los puntos \(-1,0,+1\), respectivamente. Las diferencias
entre sus cuadrados y ellos mismos, así como sus productos cruzados, son
múltiplos de \(X^3-X\). Sustituir \(Q\), que anula ese polinomio, demuestra
las identidades. En los tres bloques, \(Q\) actúa como \(-I,0,I\).
\end{proof}

Se obtiene así
\begin{align}
 \exp(t\mathbb J)
 ={}&p_{\mathrm e}(\cos t\,I+\sin t\,\mathbb J)
   +p_{\mathrm p}(I+t\mathbb J)\notag\\
   &+p_{\mathrm h}(\cosh t\,I+\sinh t\,\mathbb J).
 \label{geo-euler-trit-exponencial-total}
\end{align}
Cada sumando conserva su relación cuadrática y la conjugación invierte
su evolución. La representación total reúne los tres regímenes, pero no
define por sí sola un operador de transición entre ellos. Tal transición
requiere los mapas de transporte del TPK con sus dominios, orientación y
memoria. La articulación exponencial permite reconocer sus tipos sin
reemplazar esa dinámica por una identificación de las fibras.

% PROCEDENCIA FOCAL (no impresa)
% Resultado recuperado: tres generadores y tres evaluaciones exponenciales.
% Propietario completo:
% BASE/manuscrito/sections/hmt/09b_algebras_cuadraticas.tex:289-516.
% Antecedente causal de C y acción contragrediente:
% BASE/manuscrito/sections/hmt/03d_esqueleto_ciclotomico_app.tex:13-47,145-246.
% Los tres tipos y la orientación están en:
% BASE/manuscrito/sections/hmt/02_trit_geometrias.tex:83-196.
% Resolución polinómica recuperada:
% output/INVESTIGACION_HOLONOMIA_NONADICA_CRISTAL_APERIODICO_20260903/
% ALGEBRA_HOLONOMICA_UNIVERSAL_Y_EULER_TRITICO.md, secciones 8-11 y 22.
% Los cuatro propietarios se leyeron completos. BASE designa:
% /Users/ruben/Documents/HMT2/
% EL_CIERRE_HOLOGRAFICO_DEL_INFINITO_SUCESOR_102_CAPITULOS_2026-08-28_EN_TRABAJO/
% No se incorpora la identidad P9(alpha)=delta del documento especializado:
% esa identificación no pertenece a este bloque.
% Tampoco se identifica una dirección nilpotente con la vacancia electrónica,
% ni se presume una representación global del transporte por la suma directa.
% COMPROBACIONES: Python estándar, enteros y Fraction.
% C^2+C+I=0; (2C+I)^2=-3I; N^2=0; (2Fav^2-3I)^2=5I.
% 81 pares de residuos verifican las dos acciones contragredientes.
% Los 3 proyectores se verificaron en Q=-1,0,+1; prueba general en texto.
% 6 labels únicos y entornos equilibrados. Sin compilación.


\section{Transporte del bloque modular al cociente incidencial}
\label{sec:incidence-transport-new}
El bloque central de APP y la incidencia de caras y vértices proporcionan
dos realizaciones de una misma involución. Su relación se expresa mediante
un entrelazador rectangular, antes de efectuar cualquier interpretación
métrica. La congruencia diagonal adyacente
\[
 k^2\equiv(k+1)^2\pmod9
 \quad\Longleftrightarrow\quad 2k+1\equiv0\pmod9
\]
selecciona \(k=4\) en la ventana declarada. El bloque correspondiente es
\[
 B_c=\begin{pmatrix}7&2\\2&7\end{pmatrix}
 =7I_2+2R_c=9P_++5P_-,\qquad
 R_c=\begin{pmatrix}0&1\\1&0\end{pmatrix},\quad
 P_\pm=\frac{I_2\pm R_c}{2}.
\]
La diferencia entre los dos valores propios es cuatro. El transporte
incidencial determina a continuación cómo se conserva esa diferencia
cuando cambian las multiplicidades de los sectores.

Las caras del cubo se indexan por \((i,\epsilon)\), con
\(i\in\{1,2,3\}\) y \(\epsilon\in\{-1,1\}\); sus vértices se
indexan por \(\nu\in\{-1,1\}^3\). Definimos
\[
 M_{(i,\epsilon),\nu}=\mathbf1_{\{\nu_i=\epsilon\}}.
\]
Sean \(R_F\) la oposición de caras y \(R_V\) la inversión de vértices
\(\nu\mapsto-\nu\). En cada entrada se verifica
\[
 \mathbf1_{\{\nu_i=-\epsilon\}}
 =\mathbf1_{\{(-\nu)_i=\epsilon\}},\qquad R_FM=MR_V.
\]
La matriz \(M\) transporta, por tanto, la involución completa.

\begin{theorem}[Descenso incidencial del bloque central]
Sean \(B_F=7I_6+2R_F\) y \(B_V=7I_8+2R_V\). Entonces
\[
 B_FM=MB_V.
\]
Si \(u=(1,1,1,1,1,1)^{\mathsf T}\),
\(P_u=uu^{\mathsf T}/6\) y \(P_-=(I_6-R_F)/2\), el cociente
\[
 Q_\square=\RR^6/\ker(MM^{\mathsf T})
 \simeq\RR u\oplus E_-
\]
tiene dimensión \(1+3\), y el operador inducido es
\[
 \overline B_F=9P_u+5P_-.
\]
\end{theorem}
\begin{proof}
La primera identidad resulta de multiplicar \(R_FM=MR_V\) por dos
y sumar \(7M\). Cada cara contiene cuatro vértices; dos caras adyacentes
comparten dos y dos caras opuestas no comparten ninguno. De esta cuenta se
obtiene, entrada por entrada,
\[
 MM^{\mathsf T}=12P_u+4P_-.
\]
El primer proyector tiene rango uno y el segundo rango tres; sus imágenes
son ortogonales. Los valores propios son \(12,4,0\), con multiplicidades
\(1,3,2\). La oposición actúa como \(+1\) en \(\RR u\) y como
\(-1\) en \(E_-\), y conserva el núcleo de la matriz de Gram. Así,
\(B_F\) desciende al cociente y actúa por nueve y cinco en esos dos
sumandos. Su polinomio característico es
\((\lambda-9)(\lambda-5)^3\).
\end{proof}

La asignación temporal al eje \(\RR u\) fija la forma firmada
\(J_L=P_--P_u\). Sobre el cociente, \(I_Q=P_u+P_-\), de modo que
\[
 J_L=\frac{7I_Q-\overline B_F}{2}.
\]
Esta igualdad relaciona la forma de signatura \((3,1)\), con la
convención temporal indicada, y el operador modular transportado. La
matriz de Gram inicial conserva su positividad; la firma se introduce
en la realización posterior mediante el signo asignado a cada sector.
Las dilataciones nueve y cinco siguen siendo las del bloque: la identidad
afín determina la forma y no convierte al bloque en una isometría.

\subsection{Agregación aritmética, norma reticular y funcional de ruta}
La conservación del entero previo a la reducción se hace visible al sumar
las dos hojas. Las matrices de agregación modular son
\[
 M_\Pi=\begin{pmatrix}4&5&6\\5&4&6\\6&6&9\end{pmatrix},\qquad
 M_\Sigma=\begin{pmatrix}5&6&4\\6&4&5\\4&5&6\end{pmatrix}.
\]
Sus sumas son cincuenta y uno y cuarenta y cinco. Al reconstruir los
nueve elementos de cada bloque, el entero y su acarreo satisfacen
\[
 2025-810=(459-405)+9(174-45)=54+9\cdot129.
\]
El agregado residual y el cociente constituyen dos coordenadas de esta
igualdad. En la construcción reticular posterior, el valor cincuenta y
cuatro aparece como norma de la cámara marcada
\(v_\alpha=4\rho_1+\rho_2+\cdots+\rho_{12}\), para vectores
ortogonales \(\rho_j^2=2\). En la realización material aparece como
evaluación del funcional de ruta \(D_1\) del perfil \((80,54,6)\).
Cada aparición conserva su aplicación de origen: agregación con acarreo,
forma reticular y evaluación orientada de ruta. La coincidencia de sus
valores adquiere contenido estructural mediante esas aplicaciones, cuyos
dominios y pruebas se desarrollan en sus secciones respectivas.


\section{Extracción dodecafásica y proyección isotípica de la dirección}
\label{sec:register}
Antes de evaluar números, los lectores distinguen regiones dinámicas.
La simetría especular del corpus actúa sobre el bloque
\(B=(u^\pi,u^e,u^\varphi)\in(\mathbb F_3^6)^3\) por
\[
 \mathfrak cB=(u^e,u^\pi,u^\varphi).
\]
Con \(q=u^\pi+u^e+u^\varphi\),
\(a=u^\pi+u^e-u^\varphi\) y \(d=u^\pi-u^e\), las fórmulas
\[
 u^\varphi=2(q-a),\quad s=q-u^\varphi,\quad
 u^\pi=2(s+d),\quad u^e=2(s-d)
\]
son identidades componente a componente en \(\mathbb F_3\). El espejo
conserva el eje y el agregado en una fibra e invierte el reparto
antisimétrico. El transporte nonádico hace evolucionar la fibra misma.
La quinta fase tiene dato de frontera \(111111\), firma residual
\(000\), una solución local y una extensión: la marca de
\(\varphi\) en cinco identifica el eje regional de esa operación.
En las fases seis, siete y ocho los agregados regionales son
\(012100\), \(101001\) y \(112000\), respectivamente. Estos datos
pertenecen al dominio ternario; los lectores de historias compatibles
producen después sus evaluaciones reales.

El registro de doce ventanas y los bloques regionales conservan así
sus dominios y su enlace en el mismo estado. La realización positiva
que sigue utiliza el registro expresado, con su orden y su escala.
Para transportar el espejo regional a una deformación geométrica debe
componerse primero su acción en ese registro; la igualdad de genealogía
por sí sola expresa menos que la conmutatividad de ese cuadrado. Esta
precisión permite conservar la dinámica regional completa al estudiar
una de sus realizaciones finitas.

El calendario divide las ciento ocho posiciones en doce ventanas nonádicas.
La extracción terminal tiene la composición
\[
 x_{\rm term}\longmapsto\mathcal R_{12}(x_{\rm term})
 \longmapsto(b^{90},b^{120},Q)\longmapsto U_{\rm sgn}\longmapsto K.
\]
Las dos primeras flechas conservan la historia orientada en los lectores
del corpus. Para mostrar íntegramente la reconstrucción finita que utiliza
este artículo, sus salidas enteras son
\begin{align*}
 b^{90}&=(-495,-116,-684,108,601,-90,-173,-88,313,560,-397,461),\\
 b^{120}&=(-425,-281,-481,671,-255,30,475,-543,680,251,6,-128),
 \qquad Q=6263.
\end{align*}
Se fija el orden cíclico de doce posiciones. Sea
\(B_r=\sum_{j=0}^3 b^{120}_{r+3j}\) y
\(d_{r,j}=b^{90}_{r+3j}\), para \(r=1,2,3\). Entonces
\[
 s_1=\frac{Q+2B_1+B_2}{3},\quad s_2=s_1-B_1,\quad s_3=s_2-B_2,
 \qquad
 U^{(r)}=(s_r,d_{r,1}-d_{r,3},d_{r,0}+d_{r,2},d_{r,0}-d_{r,2}).
\]
La evaluación da
\[
 U^{(1)}=(2378,-452,-668,-322),\quad
 U^{(2)}=(1406,998,-204,-28),\quad
 U^{(3)}=(2479,-551,-371,-997).
\]
La matriz
\[
 H_4=\begin{pmatrix}1&1&1&1\\1&1&-1&-1\\1&-1&1&-1\\1&-1&-1&1\end{pmatrix},
 \qquad H_4^2=4I_4,
\]
reconstruye cada órbita mediante \(H_4U^{(r)}/4\). Restituir el orden
de las posiciones produce
\begin{equation}
 K=(234,543,140,729,659,824,621,58,914,794,146,601),
 \qquad Q=\sum_{j=1}^{12}K_j=6263.
 \label{eq:K-main}
\end{equation}
Cada componente es positiva. La división por tres procede de las tres
órbitas y la división por cuatro de la inversa de Hadamard. El control
directo adicional es \((I-S_3)K=b^{90}\),
\((I-S_4)K=b^{120}\), donde \((S_aK)_j=K_{j+a}\) con índices cíclicos.
Estas identidades hacen reproducible la reconstrucción sin seleccionar
coeficientes mediante una geometría objetivo.

La coordenatización incidencial del corpus realiza las doce posiciones como cosets
de \(A_5/C_5\). Si \(\rho\) es esa representación y \(\chi_3\) su
carácter tridimensional real, su proyector es
\[
 P_3=\frac3{60}\sum_{a\in A_5}\chi_3(a^{-1})\rho(a),\qquad
 P_{11}=I-\frac1{12}\mathbf1\mathbf1^{\mathsf T}.
\]
La enumeración y la orientación de cosets se conservan en la fuente de
esta coordenatización. Para que la extensión geométrica sea evaluable directamente,
se escribe su salida exacta \(u=P_3P_{11}K\), con \(r=\sqrt5\):
\begin{equation}
\begin{split}
u=\bigl(&-495/4-233r/10,\quad403/4+169r/10,\\
&-403/4-169r/10,\quad495/4+233r/10,\\
&601/4+1031r/10,\quad203/4+211r/5,\\
&-203/4-211r/5,\quad-601/4-1031r/10,\\
&313/4+569r/10,\quad162+219r/20,\\
&-162-219r/20,\quad-313/4-569r/10\bigr).
\end{split}
\label{eq:u-main}
\end{equation}
Se comprueba por suma y producto en \(\QQ(\sqrt5)\) que
\[
 \mathbf1^{\mathsf T}u=0,\qquad
 \|u\|^2=\frac{6638585+2275584\sqrt5}{20}>0.
\]
En particular, sus coordenadas son no constantes. La dirección \(u\),
el registro íntegro \(K\) y el proyector
\(P_{10}=P_{11}-uu^{\mathsf T}/\|u\|^2\) conservan diferentes cantidades
de información. Esta distinción permite utilizar la dirección para un
flujo sin sustituir por ella el registro que emplean los lectores de acción.

\section{Reconstrucción de la geometría positiva desde la memoria nonádica}
\label{sec:memory}

La extracción dodecafásica conserva más información que una dirección de
deformación. La dirección selecciona un movimiento; el registro completo
determina también la partición sobre la que ese movimiento actúa. La
recuperación del registro a partir de sus memorias permite formular esta
distinción como un teorema de reconstrucción. El resultado procede de la
composición de las compresiones nonádicas del corpus \cite{hmtX}; su aplicación
geométrica conserva la carga uniforme y los complementos de ambas etapas.

En esta sección se numeran las coordenadas de cero a once. Sea
\(\mathsf S:\RR^{12}\to\RR^{12}\) el desplazamiento
\((\mathsf Sx)_i=x_{i+1}\), con índices módulo doce. En la coordenatización euclídea
posterior al estado enriquecido, definimos
\[
 T_j=\frac{8I+\mathsf S^j}{9},\qquad
 D_j=\frac{\sqrt8}{9}(\mathsf S^j-I),\qquad j\in\{3,4\}.
\]
El coeficiente ocho cuenta las posiciones ordinarias de la coordenatización nonádica;
la novena posición contiene el transporte. La ortogonalidad del desplazamiento
implica por multiplicación directa
\[
 T_j^*T_j+D_j^*D_j=I.
\]
La primera memoria es el complemento de la primera compresión; la segunda
se toma sobre el estado que ya atravesó esa compresión. Por tanto,
\begin{equation}
 q=\mathbf1^*K,\qquad m_0=D_3K,\qquad m_1=D_4T_3K,
 \qquad MK=(m_0,m_1).
 \label{eq:mem-data}
\end{equation}
El orden de la composición forma parte del objeto. La aplicación
\(M:\RR^{12}\to\RR^{24}\) conserva dos observaciones correlacionadas,
producidas por una misma trayectoria lineal.

\begin{theorem}[Inversión exacta de la representación de memoria]
\label{thm:memory-inversion}
La aplicación \(K\mapsto(q,MK)\) es un isomorfismo lineal entre
\(\RR^{12}\) y \(\RR\times\operatorname{im}M\). Su inversa se obtiene
mediante
\begin{align}
 T_3^{-1}&=\frac{512I-64\mathsf S^3+8\mathsf S^6-\mathsf S^9}{455},
 \label{eq:memory-T-inverse}\\
 b&=-\frac9{\sqrt8}m_0,\qquad
 c=-\frac9{\sqrt8}T_3^{-1}m_1,\qquad w_i=c_i-b_{i+1},\notag\\
 B_0&=0,\qquad B_i=\sum_{j=0}^{i-1}w_j,\qquad
 K_i=\frac{q+\sum_{j=0}^{11}B_j}{12}-B_i.
 \label{eq:memory-inverse}
\end{align}
\end{theorem}
\begin{proof}
Escribiendo \(X=\mathsf S^3\), se tiene \(X^4=I\) y
\((8I+X)(512I-64X+8X^2-X^3)=4095I\). Dividir por nueve
prueba \eqref{eq:memory-T-inverse}. Todos los operadores utilizados son
polinomios del mismo desplazamiento y conmutan. Así,
\[
 b=(I-\mathsf S^3)K,\qquad c=(I-\mathsf S^4)K,
 \qquad w_i=K_i-K_{i+1}.
\]
La suma telescópica da \(B_i=K_0-K_i\); sumar esta igualdad sobre
las doce posiciones proporciona \eqref{eq:memory-inverse}. Además,
\(MK=0\) equivale a invariancia bajo \(\mathsf S^3\) y
\(\mathsf S^4\). Como \(\mathsf S=\mathsf S^4(\mathsf S^3)^{-1}\),
el núcleo es exactamente \(\RR\mathbf1\). La carga recupera esa
componente y demuestra la inyectividad. La descomposición
\(K=(q/12)\mathbf1+K^\perp\) prueba la sobreyectividad sobre el
codominio declarado.
\end{proof}

La forma geométrica queda determinada por la información que una lectura
comprimida parecía haber desplazado. El complemento vectorial contiene las
diferencias entre posiciones; la carga contiene la componente común. Su
reunión reconstruye las doce coordenadas antes de normalizarlas como
longitudes. Esta afirmación relaciona operaciones efectivas: compresión,
archivo de complementos e inversión. La conservación de una norma expresa
una propiedad distinta y, por sí sola, más débil que esa recuperabilidad.

\subsection{Cono de memoria compatible y estabilidad de los menores}

Sea \(L\) la pseudoinversa de \(M\), con imagen
\(\mathbf1^\perp\). Equivalentemente,
\[
 L=\left(M^*M\big|_{\mathbf1^\perp}\right)^{-1}M^*,
 \qquad LM=P_{11},\qquad K=\frac q{12}\mathbf1+Lm.
\]
La notación de inversa restringida se interpreta como un operador que toma
valores en \(\mathbf1^\perp\). La positividad estricta del registro
equivale a pertenecer al cono abierto
\begin{equation}
 \mathcal M_+=\left\{(q,m)\in\RR\times\operatorname{im}M:
 \frac q{12}\mathbf1+Lm>0\right\}.
 \label{eq:memory-positive-cone}
\end{equation}
Las desigualdades son coordenadas. El cono tiene dimensión doce; su sección
de carga positiva fija tiene dimensión once. En particular, las veinticuatro
entradas de las dos memorias satisfacen relaciones de compatibilidad.

\begin{proposition}[Constante óptima de reconstrucción en norma euclídea]
\label{prop:memory-stability}
Se cumple \(\|L\|=9/4\). Para perturbaciones arbitrarias de las observaciones,
\begin{equation}
 \|\delta K\|^2\leq\frac{|\delta q|^2}{12}
 +\frac{81}{16}\bigl(\|\delta m_0\|^2+\|\delta m_1\|^2\bigr).
 \label{eq:memory-stability}
\end{equation}
\end{proposition}
\begin{proof}
La base de Fourier diagonaliza \(\mathsf S\). Si \(z^{12}=1\),
el autovalor de \(M^*M\) en el modo correspondiente es
\[
 \lambda(z)=\frac8{81}|z^3-1|^2
 +\frac8{81}|z^4-1|^2\left|\frac{8+z^3}{9}\right|^2.
\]
Evaluar los doce modos da
\[
 \operatorname{Spec}(M^*M)=
 \left\{0^{[1]},(16/81)^{[2]},(8/27)^{[2]},(32/81)^{[1]},
 (952/2187)^{[4]},(1256/2187)^{[2]}\right\}.
\]
Los exponentes entre corchetes indican multiplicidad. El menor autovalor
positivo es \(16/81\), de donde \(\|L\|=9/4\). La componente uniforme
\((\delta q/12)\mathbf1\) es ortogonal a \(L\delta m\); la identidad
pitagórica y la norma de \(L\) producen \eqref{eq:memory-stability}.
\end{proof}

Para pasar a la geometría intervalar recuperamos la numeración de uno a doce
y ponemos \(d_j=K_j/q\), \(t_0=0\),
\(t_j=\sum_{r=1}^j d_r\). El cono \eqref{eq:memory-positive-cone}
determina exactamente \(0=t_0<t_1<\cdots<t_{12}=1\), y por tanto
\[
 \Delta_{ij}=t_j-t_i=\sum_{i<r\leq j}d_r>0\qquad(i<j).
\]
Son los menores de la matriz de caminos que se construye en la sección
\ref{app:network-cluster}. Su dependencia respecto de la memoria es
explícita. Los puntos \((t_j,t_j^2)\) quedan igualmente determinados y
producen la realización amplituhedral de rango uno desarrollada después.

Si la raíz del miembro derecho de \eqref{eq:memory-stability} es menor que
\(\min_j K_j\), la reconstrucción perturbada permanece en el cono positivo.
Para el registro expresado, \(\min_jK_j=58\); a carga fija basta
\[
 \sqrt{\|\delta m_0\|^2+\|\delta m_1\|^2}<\frac{232}{9}.
\]
La cota conserva el orden de los vértices y el signo de todos los menores
ordenados. Describe estabilidad de esta coordenatización matemática, con la
normalización indicada. Para observaciones fuera de \(\operatorname{im}M\),
la reconstrucción utiliza la proyección \(MLm\); el residuo
\((I-ML)m\) se conserva separadamente cuando interesa recuperar la
observación íntegra.

\subsection{Recuperación nodal y reducción a la frontera}

La partición obtenida tiene también una realización cuadrática:
\[
 \mathcal E_d(f)=\sum_{j=1}^{12}\frac{|f_j-f_{j-1}|^2}{d_j}
 =f^*L_df.
\]
El laplaciano nodal recupera cada longitud por
\(d_j=-1/(L_d)_{j-1,j}\). Si se conservan únicamente los extremos
globales, la minimización elimina los nodos interiores y produce, en
cambio, la matriz \(\left(\begin{smallmatrix}1&-1\\-1&1\end{smallmatrix}\right)\),
puesto que \(\sum_jd_j=1\). El lector nodal distingue la forma intervalar;
el lector de dos extremos sólo observa su longitud total.

La demostración de esta eliminación se formula íntegramente en la sección
energética: la prolongación afín minimizante \(Hf\) y el detalle \(z\),
nulo en los nodos antiguos, satisfacen
\[
 \mathcal E_{\widetilde d}(Hf+z)
 =\mathcal E_d(f)+\mathcal E_{\widetilde d}(z).
\]
El segundo término es no negativo y se anula exactamente para el detalle
nulo. La lectura efectiva conserva la energía mínima; el par \((f,z)\)
conserva además la configuración refinada. Así, positividad de menores,
positividad energética y recuperabilidad de memoria se relacionan mediante
mapas definidos sobre la misma partición, manteniendo sus respectivos
espacios de llegada.

\section{Reconstrucción íntegra del registro e isometría de incidencia}
\label{rev4:ki:section}

La reconstrucción de la sección~\ref{sec:memory} permite conservar en
una misma composición la geometría positiva, la incidencia excepcional
y la dirección de deformación. La información inicial de esta
composición es el registro producido por el estado enriquecido
APP--TRIT--TPK: las dos hojas aritméticas conservan sus operaciones;
la orientación trítica determina el régimen; el transporte, la
actualización y la prolongación nonádica conservan las diferencias
entre posiciones y la componente uniforme. Las realizaciones lineales
que siguen actúan sobre ese registro ya determinado. Su función es
transportar información y explicitar sus invariantes.

La distinción fundamental se refiere al dominio conservado. El
registro \(K\) contiene doce coordenadas ordenadas; la tétrada
\(B_K\) conserva cuatro posiciones seleccionadas por la incidencia;
una dirección geométrica conserva una recta en un espacio de
representación. Cada operación tiene una imagen y una fibra propias.
La construcción siguiente conserva primero las doce coordenadas y
deriva después las reducciones que intervienen en la geometría. Así
queda disponible una reconstrucción exacta antes de efectuar una
extracción de soporte o una reducción dimensional.

\subsection{Forma de Gram del diseño dodecafásico}
\label{rev4:ki:gram-subsection}

Adoptamos la numeración de uno a doce y denotamos por
\(\Omega=\{1,\ldots,12\}\) el conjunto de posiciones. La
codificación Paley--Witt del estado enriquecido determina el conjunto
\(\mathcal H\) de sus \(132\) hexadas. Una hexada es un soporte
de seis posiciones de un vector de peso seis del código ternario
extendido. La construcción explícita del código y la demostración de
que cada quintupla pertenece a una única hexada se desarrollan en
el teorema~\ref{thm:lf-code} y la
proposición~\ref{prop:lf-witt}. Utilizamos aquí precisamente esa
incidencia, con las posiciones transportadas en el mismo orden que
las coordenadas del registro.

Sean \(V=\RR^{12}\) y \(W=\RR^{\mathcal H}\), equipados con
sus productos interiores euclídeos. La matriz de incidencia define
la aplicación
\begin{equation}
 \mathcal I:V\longrightarrow W,
 \qquad (\mathcal I k)_H=\sum_{p\in H}k_p,
 \qquad \mathcal I_{H,p}=\mathbf1_{\{p\in H\}}.
 \label{rev4:ki:incidence-map}
\end{equation}
El símbolo \(\mathbf1\in V\) designa el vector uniforme de doce
componentes. Definimos
\[
 \mu(k)=\frac1{12}\mathbf1^{\mathsf T}k,\qquad
 P_0=I_{12}-\frac1{12}\mathbf1\mathbf1^{\mathsf T},\qquad
 V_0=\mathbf1^\perp.
\]
El proyector \(P_0\) coincide con \(P_{11}\) de la
sección~\ref{sec:memory}; el subíndice cero señala en esta
presentación la condición de media nula. La descomposición ortogonal
es \(k=\mu(k)\mathbf1+P_0k\).

\begin{lemma}[Identidad de Gram y normalización incidencial]
\label{rev4:ki:gram}
La incidencia satisface
\begin{equation}
 \mathcal I^{\mathsf T}\mathcal I
       =36I_{12}+30\mathbf1\mathbf1^{\mathsf T}.
 \label{rev4:ki:gram-formula}
\end{equation}
Por consiguiente,
\[
 U_W=\frac16\mathcal I\big|_{V_0}:
 V_0\longrightarrow E_{\rm inc}:=\mathcal I(V_0)
\]
es una isometría sobre un subespacio de dimensión once.
\end{lemma}
\begin{proof}
Fijado un punto, las quintuplas que lo contienen son
\(\binom{11}{4}\). Cada hexada que contiene ese punto comprende
\(\binom54\) de ellas. La unicidad de la extensión de cada
quintupla da \(66\) hexadas por punto. Fijada una pareja de
puntos, el mismo argumento da
\(\binom{10}{3}/\binom43=30\) hexadas por pareja.
La entrada \((p,r)\) de
\(\mathcal I^{\mathsf T}\mathcal I\) cuenta exactamente las
hexadas que contienen simultáneamente \(p\) y \(r\). Sus
entradas diagonales son, por tanto, \(66\), y las restantes,
\(30\), lo que prueba \eqref{rev4:ki:gram-formula}.

Para \(z\in V_0\), el término uniforme se anula y
\[
 \left\|\frac16\mathcal I z\right\|^2
 =\frac1{36}z^{\mathsf T}
       \mathcal I^{\mathsf T}\mathcal I z
 =\|z\|^2.
\]
La inyectividad y la dimensión de la imagen se deducen de esta
igualdad. El factor seis queda así fijado por el Gram del diseño.
\end{proof}

El conteo de incidencia determina la escala de la isometría. Las
sumas sobre hexadas son observables lineales de las posiciones; sus
productos interiores contienen la multiplicidad con que cada
posición y cada pareja aparecen en el diseño. La eliminación del
modo uniforme separa esos dos niveles de conteo y deja el factor
cuadrático \(36\). Esta operación proporciona un enlace explícito
entre la representación de doce posiciones y su representación
sobre hexadas.

\subsection{Conservación de la media e inversión de las doce coordenadas}
\label{rev4:ki:full-subsection}

La componente uniforme se conserva mediante una coordenada escalar
adicional. Equipamos \(\RR\oplus E_{\rm inc}\) con el producto
interior
\[
 \langle(\mu,y),(\nu,z)\rangle_*=
 12\mu\nu+\langle y,z\rangle.
\]
El peso doce es la norma cuadrada de \(\mathbf1\), fijada por el
número de posiciones. Definimos
\begin{equation}
 F:V\longrightarrow\RR\oplus E_{\rm inc},\qquad
 F(k)=\left(\mu(k),\frac16\mathcal I P_0k\right).
 \label{rev4:ki:F}
\end{equation}

\begin{theorem}[Isomorfismo isométrico del registro completo]
\label{rev4:ki:full-isometry}
La aplicación \(F\) es un isomorfismo lineal isométrico. Su inversa
es
\begin{equation}
 F^{-1}(\mu,y)=\mu\mathbf1+
                  \frac16P_0\mathcal I^{\mathsf T}y,
 \qquad y\in E_{\rm inc}.
 \label{rev4:ki:F-inverse}
\end{equation}
En particular,
\begin{equation}
 12|\mu(k)|^2+
 \left\|\frac16\mathcal I P_0k\right\|^2=\|k\|^2.
 \label{rev4:ki:F-norm}
\end{equation}
\end{theorem}
\begin{proof}
La descomposición uniforme--centrada es ortogonal y el
lema~\ref{rev4:ki:gram} conserva la norma de su segundo sumando.
Esto demuestra \eqref{rev4:ki:F-norm} y la inyectividad.
Además,
\[
 \frac1{36}P_0\mathcal I^{\mathsf T}\mathcal I P_0=P_0.
\]
Al sustituir \(F(k)\) en \eqref{rev4:ki:F-inverse}, se obtiene
\(\mu(k)\mathbf1+P_0k=k\). Recíprocamente, si
\(y\in E_{\rm inc}\), existe un único \(z\in V_0\) con
\(y=\mathcal I z/6\). La fórmula propuesta devuelve
\(\mu\mathbf1+z\), cuya imagen es \((\mu,y)\). Así se
prueban la sobreyectividad sobre el codominio declarado y las dos
identidades de inversión.
\end{proof}

La restricción \(y\in E_{\rm inc}\) expresa una compatibilidad
lineal efectiva. El operador
\begin{equation}
 \Pi_{\rm inc}=\frac1{36}\mathcal I P_0
                         \mathcal I^{\mathsf T}:W\longrightarrow W
 \label{rev4:ki:inc-projector}
\end{equation}
es el proyector ortogonal sobre \(E_{\rm inc}\). En efecto, es
simétrico y su cuadrado se reduce a él mediante
\eqref{rev4:ki:gram-formula}; actúa como la identidad en la imagen
de \(\mathcal I P_0\). Por tanto, la compatibilidad se comprueba
mediante \(\Pi_{\rm inc}y=y\). Para un dato arbitrario de
\(W\), la reconstrucción seguida de incidencia devuelve
\(\Pi_{\rm inc}y\); el componente ortogonal
\((I_W-\Pi_{\rm inc})y\) cuantifica la información externa a
esta representación del registro.

\begin{proposition}[Covariancia de la reconstrucción]
\label{rev4:ki:equivariance}
Sea \(g\) un automorfismo del diseño y sean \(\rho(g)\) y
\(\sigma(g)\) sus acciones por permutación sobre posiciones y
hexadas. Entonces
\[
 F\rho(g)=(1\oplus\sigma(g))F.
\]
La media se conserva y la reconstrucción conmuta con el transporte
simultáneo de las posiciones y de sus hexadas.
\end{proposition}
\begin{proof}
La equivalencia \(p\in H\Longleftrightarrow g(p)\in g(H)\)
implica \(\mathcal I\rho(g)=\sigma(g)\mathcal I\).
Las permutaciones fijan \(\mathbf1\), conservan la media y
conmutan con \(P_0\). Sustituir estas identidades en
\eqref{rev4:ki:F} demuestra la afirmación. La covariancia de la
inversa se obtiene componiendo con \(F^{-1}\).
\end{proof}

La normalización isométrica y la inversión aritmética cumplen
funciones distintas. En la coordenatización Hadamard del registro,
\(U=\mathcal H_{12}K\) y
\(K=\mathcal H_{12}U/4\), con
\(\mathcal H_{12}^{\mathsf T}=\mathcal H_{12}\) y
\(\mathcal H_{12}^2=4I_{12}\).
El operador \(J_H=\mathcal H_{12}/2\) es ortogonal. Por ello
\(FJ_H\) conserva la norma del registro firmado normalizado
\(U/2\) y lo representa en el espacio incidencial ampliado.
La inversión integral mediante \(\mathcal H_{12}/4\) conserva,
además, la subred de congruencias correspondiente. En ambos casos se
transporta el registro generado, con el orden de fase fijado antes
de la transformación \cite{hmtX}.

\subsection{Composición exacta desde las dos memorias nonádicas}
\label{rev4:ki:memory-subsection}

Retomamos \(M\), \(q\), \(m_0\) y \(m_1\)
de \eqref{eq:mem-data}, junto con el operador de reconstrucción
\(L=(M^*M|_{V_0})^{-1}M^*\) de la sección~\ref{sec:memory}.
El dominio compatible es
\(\mathcal D=\RR\times\operatorname{im}M\), con
\(m=(m_0,m_1)\). El teorema~\ref{thm:memory-inversion}
establece la reconstrucción
\[
 \mathcal R(q,m)=\frac q{12}\mathbf1+Lm,
 \qquad LM=P_0.
\]
Como \(L\) toma valores en \(V_0\), la composición con la
incidencia tiene una expresión directa:
\begin{equation}
 \mathcal C=F\mathcal R,\qquad
 \mathcal C(q,m)=\left(\frac q{12},\frac16\mathcal I Lm\right).
 \label{rev4:ki:memory-incidence}
\end{equation}

\begin{corollary}[Equivalencia entre memoria compatible e incidencia íntegra]
\label{rev4:ki:composition}
La aplicación \(\mathcal C\) es un isomorfismo entre
\(\mathcal D\) y \(\RR\oplus E_{\rm inc}\), con inversa
\begin{equation}
 \mathcal C^{-1}(\mu,y)=
 \left(12\mu,\frac16MP_0\mathcal I^{\mathsf T}y\right).
 \label{rev4:ki:C-inverse}
\end{equation}
Su forma cuadrática conservada es
\[
 \|\mathcal C(q,m)\|_*^2=
 \frac{|q|^2}{12}+\|Lm\|^2=\|\mathcal R(q,m)\|^2.
\]
\end{corollary}
\begin{proof}
La fórmula \eqref{rev4:ki:memory-incidence} resulta de
\(P_0\mathbf1=0\) y \(P_0L=L\). La inversa se obtiene
aplicando \(k\mapsto(\mathbf1^{\mathsf T}k,Mk)\) a
\eqref{rev4:ki:F-inverse}; se usa \(M\mathbf1=0\).
Las identidades de composición proceden de \(LM=P_0\) y de
\(MLm=m\) para \(m\in\operatorname{im}M\). La igualdad de
normas es \eqref{rev4:ki:F-norm} aplicada a \(\mathcal R(q,m)\).
\end{proof}

\begin{figure}[htbp]
\centering
\begin{tikzpicture}[>=Latex, every node/.style={font=\small},
  box/.style={draw,rounded corners=2pt,align=center,
              inner sep=6pt,minimum height=12mm}]
 \node[box] (memory) at (0,0) {Carga y memorias\\$(q,m_0,m_1)$};
 \node[box] (register) at (4.1,0) {Registro íntegro\\$K$};
 \node[box] (incidence) at (8.2,0) {Media e incidencia\\$(\mu,y)$};
 \node[box] (direction) at (4.1,-1.75)
    {Dirección geométrica\\$\ell_K=\RR P_3K$};
 \draw[<->] (memory) -- node[above] {$\mathcal R$} (register);
 \draw[<->] (register) -- node[above] {$F$} (incidence);
 \draw[->] (register) -- node[right] {proyección y paso a la recta}
    (direction);
\end{tikzpicture}
\caption{La carga y las dos memorias compatibles reconstruyen las doce
coordenadas. La media y el componente incidencial proporcionan otra
representación reversible del mismo registro. La dirección geométrica
es una reducción posterior con fibras de dimensión positiva.}
\label{rev4:ki:figure}
\end{figure}

Esta composición determina también la partición positiva de
\ref{sec:memory}. Si \(q>0\), sus longitudes se recuperan mediante
\begin{equation}
 d_j=\frac1{12}+\frac{1}{6q}
             \bigl(P_0\mathcal I^{\mathsf T}y\bigr)_j,
 \qquad y=\frac16\mathcal I P_0K.
 \label{rev4:ki:positive-lengths}
\end{equation}
Las doce condiciones \(d_j>0\) son las ecuaciones de la coordenatización
positiva expresadas en coordenadas de incidencia. Con
\(t_j=\sum_{r\leq j}d_r\), los menores ordenados
\(t_j-t_i=\sum_{i<r\leq j}d_r\) se reconstruyen desde
\((\mu,y)\) o desde \((q,m_0,m_1)\). De este modo, la
positividad de los menores y la conservación isométrica de la
incidencia quedan relacionadas por aplicaciones explícitas. La
primera es una desigualdad en la coordenatización ordenada; la segunda es una
identidad de productos interiores.

\subsection{Recuperación de la polarización y alcance de las reducciones}
\label{rev4:ki:direction-subsection}

La realización de las doce posiciones sobre las clases de
\(A_5/C_5\) determina el proyector tridimensional
\[
 P_3=\frac3{60}\sum_{a\in A_5}
               \chi_3(a^{-1})\varrho(a).
\]
Aquí \(\varrho\) es la representación por permutaciones y
\(\chi_3\) el carácter tridimensional elegido en la coordenatización
marcada. Sus representantes, valores de carácter y matriz efectiva
se fijan en \eqref{eq:leech-projectors}; el
lema~\ref{lem:leech-u} demuestra
\(P_3=P_3^{\mathsf T}=P_3^2\),
\(\operatorname{rank}P_3=3\) y \(P_3\mathbf1=0\).
En esa coordenatización, la dirección del registro se obtiene por
\begin{equation}
 u_K=P_3P_0K=P_3Lm
     =\frac16P_3P_0\mathcal I^{\mathsf T}y.
 \label{rev4:ki:direction-reconstruction}
\end{equation}
Las igualdades siguen de las dos inversas ya demostradas. Por tanto,
la misma dirección se recupera desde las memorias o desde el
componente incidencial; la media se conserva para reconstruir
\(K\), aunque el proyector tridimensional la anule.

Para el registro de \eqref{eq:leech-K}, la evaluación exacta
\eqref{eq:leech-u-norm} da
\[
 \|u_K\|^2=
       \frac{6638585+2275584\sqrt5}{20}>0.
\]
La recta \(\ell_K=\RR u_K\) y el complemento transversal
\(V_{10}(K)=V_0\cap u_K^\perp\) quedan así determinados.
Su proyector es
\begin{equation}
 P_{10}(K)=P_0-
       \frac{u_Ku_K^{\mathsf T}}{\|u_K\|^2}.
 \label{rev4:ki:rank-ten}
\end{equation}
En efecto, el segundo sumando es un proyector ortogonal de rango uno
contenido en \(V_0\); su producto con \(P_0\), en cualquiera
de los dos órdenes, coincide consigo mismo. Al elevar
\eqref{rev4:ki:rank-ten} al cuadrado se recupera el mismo
operador, y su traza vale \(11-1=10\). Se obtiene
\(V_0=\ell_K\oplus^\perp V_{10}(K)\).
Esta reducción es la que interviene después en la polarización de
los doce planos reticulares.

La incidencia transporta asimismo estos operadores. Sobre
\(E_{\rm inc}\), el operador
\(\widehat P_3=U_WP_3U_W^*\) satisface
\(\widehat P_3U_W=U_WP_3\); la identidad procede de
\(U_W^*U_W=I_{V_0}\). Transportar de igual modo
\(P_{10}(K)\) conserva su rango y su producto interior.
La correspondencia entre las dos representaciones se establece,
por tanto, sobre vectores y operadores, además de sus dimensiones.

La necesidad de conservar los distintos componentes admite
comprobaciones precisas. La observación centrada de \(k\) coincide
con la de \(k+c\mathbf1\) para todo \(c\), mientras sus
cargas difieren en \(12c\). En particular, \(k=0\) y
\(k=\mathbf1\) tienen el mismo componente incidencial centrado
y medias distintas. La coordenada \(\mu\) es exactamente la
información que separa esos dos registros. Del mismo modo,
\(u_{ak+b\mathbf1}=a u_k\) para \(a\ne0\); la recta
resultante coincide con \(\ell_k\), aunque cambien la amplitud
y la media. Estas identidades describen las fibras de las
aplicaciones sobre el espacio ambiente declarado.

La extracción \(k\mapsto\{j:k_j\geq729\}\) posee otras
fibras. Para el registro concreto \(K\), tanto \(K\) como
\(K+e_1\) tienen el soporte \(\{4,6,9,10\}\), donde
\(e_1\) es el primer vector de la base canónica. Son dos vectores
ambiente diferentes con la misma tétrada. Esta comprobación
caracteriza la pérdida de información del mapa de soporte; la
selección terminal del registro permanece fijada por su
construcción previa. La bandera
\(B_K\subset H^-_\alpha\) conserva además la hexada determinada
por los signos anteriores al acarreo, cuya coincidencia con la
selección incidencial se demuestra en la
proposición~\ref{prop:lf-flag}.

El resultado común de estas composiciones es una continuidad
matemática explícita. La memoria compatible reconstruye el registro;
su descomposición uniforme--centrada determina una representación
incidencial isométrica; la reconstrucción de sus longitudes determina
la coordenatización positiva; y la proyección tridimensional determina la
dirección de deformación. La tétrada y la bandera conservan las
incidencias que utilizan las construcciones excepcionales, mientras
\(F\) conserva la totalidad del registro dodecafásico. La historia
enriquecida mantiene, a su vez, la orientación, el recorrido y la
prolongación que sitúan ese registro dentro de la estructura
discreta conjunta del continuo.

% Procedencia HMT: output/EXTREMA_MEDIA_RAZON_NARRACION_Y_REVISION_20260916/
% ARTICULO/ES/source/libro/01_acoplamiento.tex, lib01:cadena-registro y lib01:K.
% Dirección: ARTICULO/ES/source/sections/k_direccion_dimensional.tex.
% Estatuto de las redes, cartas y flujos de este apéndice: FORMALIZACION_NUEVA.
% La extracción de K y el proyector P_3 son RESULTADO_RECUPERADO.
\section{Medición planar de frontera, coordenadas de Plücker y mutaciones positivas}
\label{app:network-cluster}

La representación dodecafásica del estado enriquecido admite una realización
positiva cuya red, coordenadas y cambios de coordenatización pueden escribirse
completamente. La composición de partida es
\[
 \mathrm{APP}\longrightarrow\mathrm{TRIT}\longrightarrow\mathrm{TPK}
 \longrightarrow x_{\rm term}\longrightarrow\mathcal R_{12}
 \longrightarrow(b^{90},b^{120},Q_{\rm TPK})
 \longrightarrow U_{\rm sgn}\longrightarrow K.
\]
Esta lectura actúa sobre el estado enriquecido que conserva las dos hojas,
la orientación, el acarreo, la memoria y la frontera. El registro ya generado
es
\[
 K=(234,543,140,729,659,824,621,58,914,794,146,601),
 \qquad Q=\sum_{r=1}^{12}K_r=6263.
\]
Su extracción se recupera exactamente de los dos canales y la carga total:
\[
 b^{90}=(-495,-116,-684,108,601,-90,-173,-88,313,560,-397,461),
\]
\[
 b^{120}=(-425,-281,-481,671,-255,30,475,-543,680,251,6,-128).
\]
En las órbitas de posiciones \((r,r+3,r+6,r+9)\), para \(r=1,2,3\),
sean \(B_r=\sum_{j=0}^3b^{120}_{r+3j}\) y
\(d_{r,j}=b^{90}_{r+3j}\). Se obtienen
\[
 s_1=(Q+2B_1+B_2)/3,\quad s_2=s_1-B_1,\quad s_3=s_2-B_2,
\]
\[
 U^{(r)}=(s_r,d_{r,1}-d_{r,3},d_{r,0}+d_{r,2},d_{r,0}-d_{r,2}),
 \qquad K|_r=\tfrac14H_4U^{(r)},
\]
\[
 H_4=\begin{pmatrix}1&1&1&1\\1&1&-1&-1\\1&-1&1&-1\\1&-1&-1&1\end{pmatrix},
 \qquad H_4^2=4I_4.
\]
La red construida a continuación realiza los valores positivos producidos
por esa extracción. Su forma combinatoria es una formalización posterior
del registro, con origen y orientación declarados.

\subsection{Red dirigida acíclica y determinantes de caminos no intersectantes}

Para un registro positivo de longitud \(L\), ponemos
\[
 d_r=K_r/Q>0,\qquad t_0=0,\qquad
 t_j=\sum_{r=1}^jd_r\quad(1\le j\le L).
\]
Así \(t_L=1\). Las \(L\) separaciones producen \(L+1\) puntos marcados;
en el caso dodecafásico son trece puntos. La matriz que los realiza es
\begin{equation}
 C(d)=\begin{pmatrix}1&1&\cdots&1\\t_0&t_1&\cdots&t_L\end{pmatrix}.
 \label{eq:nc-path-matrix}
\end{equation}
El Grassmanniano \(\operatorname{Gr}(2,L+1)\) identifica matrices de rango
dos por cambios invertibles de las dos filas. Su parte positiva está formada
por las clases que admiten todos sus menores ordenados de orden dos
estrictamente positivos.

Construimos una red dirigida con vértices \(a_j,b_j\), para
\(0\le j\le L\), y sumideros \(q_j\). Hay aristas horizontales
\(a_{j-1}\to a_j\) y \(b_{j-1}\to b_j\), de peso uno; aristas verticales
\(a_j\to b_j\), de peso \(d_j\), para \(j\ge1\); y aristas
\(b_j\to q_j\), de peso uno. Las fuentes ordenadas son \(b_0,a_0\).
Las dos filas se dibujan horizontalmente, las aristas verticales descienden
y los sumideros salen por debajo de la fila inferior. Este dibujo es planar
y la orientación es acíclica. La medición de frontera \(M_{ij}\) es la
suma de los productos de pesos de todos los caminos desde la fuente \(i\)
al sumidero \(q_j\).
La red de caminos tiene dos fuentes y \(L+1\) sumideros, es decir,
\(L+3\) conexiones exteriores. La matriz que realiza tiene \(L+1\)
columnas. La semilla plábica introducida después tiene precisamente
\(L+1\) vértices de frontera y constituye otra presentación de sus
coordenadas; estas dos cuentas de frontera se mantienen explícitas.

\begin{theorem}[Medición de frontera y menores del registro]
\label{thm:nc-network}
La medición de frontera de esta red es \(M=C(d)\). Para
\(0\le i<j\le L\),
\begin{equation}
 \Delta_{ij}(C(d))=t_j-t_i=\sum_{r=i+1}^{j}d_r>0.
 \label{eq:nc-minor-gap}
\end{equation}
El mismo número es la suma de pesos de las parejas de caminos
vértice-disjuntos \(b_0\to q_i\), \(a_0\to q_j\).
\end{theorem}
\begin{proof}
La fuente inferior tiene un único camino a cada sumidero, de peso uno.
Un camino desde \(a_0\) a \(q_j\) desciende exactamente una vez, por una
arista de índice \(1\le r\le j\); su peso es \(d_r\). Esto prueba la
fórmula de la matriz. El camino inferior hasta \(q_i\) ocupa todos los
vértices \(b_0,\ldots,b_i\). Por tanto, un camino de la otra fuente a
\(q_j\) es disjunto de él precisamente cuando desciende en
\(i<r\le j\). Su suma de pesos es el miembro derecho de
\eqref{eq:nc-minor-gap}. La pareja con destinos intercambiados siempre
comparte un vértice. El desarrollo del determinante cancela todas las
parejas que se cruzan y conserva exactamente las disjuntas. Aquí se ha
verificado directamente, para la red concreta, el mecanismo determinantal
de caminos no intersectantes.
\end{proof}

La identificación de las mediciones de redes planas con las celdas del
Grassmanniano no negativo proporciona su reconocimiento geométrico posterior;
véase A. Postnikov,
\href{https://arxiv.org/abs/math/0609764}{\emph{Total positivity,
Grassmannians, and networks}}. La prueba anterior da, además, una fórmula
finita para cada menor sin utilizar un valor metrológico.

La colección \(d_r>0\), \(\sum_rd_r=1\), tiene dimensión \(L-1\).
Su aplicación al Grassmanniano es inyectiva: si las filas de \(C(d)\) y
\(C(d')\) generan el mismo plano, las segundas filas difieren por una
transformación afín; las condiciones \(t_0=t'_0=0\) y
\(t_L=t'_L=1\) fijan esa transformación. En particular, para \(L=12\)
se obtiene una colección de dimensión once dentro de la celda positiva de
dimensión veintidós. El punto está en la celda superior; la colección conserva
su dimensión propia. Los trece puntos y las doce separaciones tienen así
papeles matemáticos distintos.

Permitir \(d_r\ge0\) conserva la no negatividad. Cuando \(d_r=0\), las
columnas consecutivas \(r-1\) y \(r\) coinciden y forman elementos
paralelos del matroide. El rango sigue siendo dos porque \(t_0=0\) y
\(t_L=1\). Las bases son exactamente los pares \(i<j\) que atraviesan
alguna separación positiva. Esta degeneración produce estratos de incidencia
consecutiva. Una columna nula, que daría un lazo, corresponde a una operación
distinta de la coincidencia de dos columnas no nulas.

\subsection{Semilla de triangulación y relación de intercambio}

Sea \(N=L+1\) y considérense los vértices \(0,\ldots,L\) de un polígono
convexo en su orden cíclico. A cada lado o diagonal \((i,j)\) se asigna
la coordenada \(A_{ij}=\Delta_{ij}(C(d))\), con índices ordenados.
La triangulación en abanico desde el vértice cero tiene diagonales
\((0,j)\), \(2\le j\le L-1\). Sus lados son variables congeladas y
sus diagonales variables mutables. El quiver se obtiene orientando
cíclicamente los tres lados de cada triángulo y cancelando los pares de
flechas opuestas. Esta es una semilla de tipo \(A_{N-3}\); para el registro
dodecafásico, su tipo es \(A_{10}\).

La semilla tiene una realización plábica explícita. Se coloca un vértice
negro dentro de cada triángulo; en los vértices del polígono se coloca un
vértice blanco si incide en alguna diagonal, y negro en caso contrario.
Se unen los centros a las tres esquinas de sus triángulos, se retiran los
lados originales y se contraen las aristas entre vértices del mismo color.
Se añade una hoja de frontera a cada vértice marcado. Esta construcción
recoge el algoritmo de triangulación de Kodama y Williams,
\href{https://people.math.harvard.edu/~williams/papers/KW-ArxivVersion.pdf}
{\emph{KP solitons and total positivity for the Grassmannian}, algoritmo 12.1}.
Las etiquetas de coordenadas de esta semilla se evalúan en los menores ya
producidos por la red de dos filas. Para una cara mutable, la coordenada
de tipo \(\mathcal X\) es el cociente de los productos de las
coordenadas \(\mathcal A\) incidentes que prescribe su columna del
quiver; en un cuadrilátero se escribirá explícitamente como una razón
doble. Cambiar el grafo cambia la coordenatización de esos menores, manteniendo su
procedencia.

\begin{theorem}[Intercambio positivo de la realización del registro]
\label{thm:nc-ptolemy}
Para cuatro vértices \(a<b<c<d\),
\begin{equation}
 A_{ac}A_{bd}=A_{ab}A_{cd}+A_{ad}A_{bc}.
 \label{eq:nc-ptolemy}
\end{equation}
En consecuencia, el cambio de diagonal \((a,c)\leftrightarrow(b,d)\)
se evalúa mediante
\[
 A_{bd}=\frac{A_{ab}A_{cd}+A_{ad}A_{bc}}{A_{ac}}>0.
\]
Todas las triangulaciones proporcionan coordenatizaciones positivas compatibles sobre
la misma realización \(C(d)\).
\end{theorem}
\begin{proof}
Al sustituir \(A_{ij}=t_j-t_i\), la igualdad se reduce a
\[
 (t_c-t_a)(t_d-t_b)
 =(t_b-t_a)(t_d-t_c)+(t_d-t_a)(t_c-t_b).
\]
La expansión de ambos miembros prueba la identidad. Las diferencias son
positivas; por tanto, cada flip conserva una evaluación positiva y el
cociente está definido. La conexión de las triangulaciones por flips
permite pasar entre ellas mediante estos intercambios. La interpretación
cluster de la relación de Plücker es el reconocimiento posterior desarrollado
por J. Scott en
\href{https://arxiv.org/abs/math/0311148}{\emph{Grassmannians and Cluster
Algebras}}. La prueba algebraica de la evaluación concreta ya está contenida
en la identidad mostrada.
\end{proof}

\subsection{Refinamiento de las separaciones y naturalidad de los intercambios}

Sea \(B\ge2\). En profundidad \(n\), cada separación \(d_r\) se divide
en \(B^n\) separaciones consecutivas iguales:
\[
 d^{(n)}_{(r-1)B^n+c}=d_r/B^n,
 \qquad 1\le c\le B^n.
\]
Esto define \(LB^n+1\) puntos y una matriz \(C_n\). Si \(\iota_n(j)=Bj\)
identifica un punto con su descendiente de igual posición, entonces
\begin{equation}
 C_{n+1}|_{\iota_n(\{0,\ldots,LB^n\})}=C_n,
 \qquad
 \Delta_{\iota_n(i),\iota_n(j)}(C_{n+1})=\Delta_{ij}(C_n).
 \label{eq:nc-refinement}
\end{equation}
La red refinada reemplaza cada arista vertical de peso \(d\) por una
secuencia de \(B\) posiciones de descenso, con pesos \(d/B\). La suma
de caminos entre extremos heredados sigue siendo \(d\). La composición de
refinamientos conserva esta igualdad porque las sumas finitas se agrupan
asociativamente.

\begin{proposition}[Naturalidad sobre extremos heredados]
\label{prop:nc-inherited}
Todo intercambio \eqref{eq:nc-ptolemy} entre cuatro extremos heredados
conmuta exactamente con el refinamiento. Toda triangulación antigua se
extiende a una del polígono refinado preservando sus triángulos interiores:
los polígonos añadidos entre dos extremos consecutivos se triangulan por
fuera de esos triángulos. Los flips entre diagonales antiguas conservan
entonces sus relaciones de intercambio.
\end{proposition}
\begin{proof}
La primera conclusión resulta al sustituir cada menor por su igualdad en
\eqref{eq:nc-refinement}. Para la segunda, se conservan las cuerdas que
unen extremos antiguos consecutivos y se triangula cada bolsillo exterior
formado por las nuevas marcas. Estas regiones tienen interiores disjuntos.
Los dos triángulos incidentes en una diagonal antigua permanecen iguales,
y, por tanto, también su cuadrilátero y su identidad de Ptolomeo.
\end{proof}

Esta naturalidad especifica los extremos sobre los que actúa. Un lado
antiguo, al recibir vértices nuevos en su arco de frontera, pasa a ser una
diagonal del polígono mayor. Por ello una variable antiguamente congelada
puede ser mutable en el problema ampliado. La compatibilidad anterior es
una compatibilidad exacta del sistema de evaluaciones y de los flips
heredados. Una inclusión de patrones de semillas con igual régimen de
coeficientes se obtiene manteniendo congeladas esas cuerdas y las nuevas
diagonales auxiliares. Esto declara la operación que conserva la semilla;
el incremento de profundidad y una mutación siguen siendo operaciones
distintas.

\subsection{Acción diagonal proyectiva e invariancia de las razones cruzadas}

Si \(\lambda_j>0\), la acción por columnas
\(C\mapsto C\operatorname{diag}(\lambda_0,\ldots,\lambda_L)\) da
\[
 A_{ij}\longmapsto\lambda_i\lambda_j A_{ij}.
\]
Preserva la positividad y todos los estratos definidos por anulación de
menores. La razón doble asociada a un cuadrilátero es
\begin{equation}
 X_{abcd}=\frac{A_{ab}A_{cd}}{A_{ad}A_{bc}}>0.
 \label{eq:nc-crossratio}
\end{equation}
Los cuatro factores \(\lambda_a\lambda_b\lambda_c\lambda_d\) se
cancelan exactamente. En consecuencia, toda acción diagonal positiva por
columnas deja fijo \(X_{abcd}\), incluso si depende de un parámetro y
tiene determinante uno. La modificación del peso de las columnas conserva
el punto proyectivo de cada columna; sus razones dobles expresan esa
conservación. Para que el registro cambie la forma proyectiva debe actuar
sobre las separaciones ordenadas.

\subsection{Deformación positiva de las separaciones dirigida por el registro}

El corpus fija la representación de las doce posiciones y su proyector
ortogonal \(P_3\) sobre el sector tridimensional. Con
\(P_{11}=I_{12}-J_{12}/12\), la dirección recuperada es
\[
 u=P_3P_{11}K=P_3K\in\mathbf1^\perp.
\]
En la coordenatización declarada, poniendo \(r=\sqrt5\), su evaluación exacta es
\begin{align*}
u=(&-495/4-233r/10,\ 403/4+169r/10,\ -403/4-169r/10,\\
   &495/4+233r/10,\ 601/4+1031r/10,\ 203/4+211r/5,\\
   &-203/4-211r/5,\ -601/4-1031r/10,\ 313/4+569r/10,\\
   &162+219r/20,\ -162-219r/20,\ -313/4-569r/10).
\end{align*}
Se conservan \(\sum_ru_r=0\) y
\(\|u\|^2=(6638585+2275584\sqrt5)/20>0\).

Para \(s\in\mathbb R\), definimos la deformación de los pesos mediante
\begin{equation}
 d_r(s)=\frac{K_re^{su_r}}{\sum_{\ell=1}^{12}K_\ell e^{su_\ell}},
 \qquad t_j(s)=\sum_{r=1}^jd_r(s),\qquad C(s)=C(d(s)).
 \label{eq:nc-gap-flow}
\end{equation}
La fórmula utiliza el registro y su dirección ya generados. Es una nueva
realización analítica del lector; el parámetro \(s\) mide su deformación.
Cuando se usa la salida interna \(\rho_A=\pi_{\rm HMT}/\varphi_{\rm HMT}^2\),
la escritura \(s=\tau\log\rho_A\) expresa exactamente la misma colección.
Esa reparametrización actúa después de la generación de las constantes.
Una identificación de \(s\) o \(\tau\) con un tiempo físico requiere la
ley temporal del lector correspondiente.

\begin{theorem}[Cambio de forma, positividad y compatibilidad con la profundidad]
\label{thm:nc-shape-flow}
La colección \eqref{eq:nc-gap-flow} pertenece a
\(\operatorname{Gr}_{>0}(2,13)\) para todo \(s\) real y conserva
\(t_0=0\), \(t_{12}=1\). Para \(i<j\), sean
\[
 \mu_{ij}(s)=
 \frac{\sum_{r=i+1}^jK_ru_re^{su_r}}{\sum_{r=i+1}^jK_re^{su_r}},
 \qquad
 \mu(s)=\frac{\sum_rK_ru_re^{su_r}}{\sum_rK_re^{su_r}}.
\]
Entonces
\begin{equation}
 \frac{d}{ds}\log A_{ij}(s)=\mu_{ij}(s)-\mu(s),
 \quad
 \frac{d}{ds}\log X_{abcd}(s)
 =\mu_{ab}(s)+\mu_{cd}(s)-\mu_{ad}(s)-\mu_{bc}(s).
 \label{eq:nc-crossratio-derivative}
\end{equation}
La deformación conmuta con cada refinamiento que asigna a los sucesores de la
separación \(r\) el peso \(d_r/B\) y la misma velocidad \(u_r\).
\end{theorem}
\begin{proof}
Todos los exponenciales y todos los \(K_r\) son positivos. La normalización
hace que la suma de separaciones sea uno; los menores son sumas de
separaciones positivas por \eqref{eq:nc-minor-gap}. Derivar el logaritmo de
\[
 A_{ij}(s)=\frac{\sum_{r=i+1}^jK_re^{su_r}}
                      {\sum_rK_re^{su_r}}
\]
da la primera identidad. La combinación de cuatro logaritmos en
\eqref{eq:nc-crossratio} cancela las cuatro medias globales con sus signos
y da la segunda. Finalmente, reemplazar \(K_r\) por \(B\) copias de
\(K_r/B\) con idéntica \(u_r\) conserva el numerador y el denominador
cuando se agrupan los sucesores. La igualdad vale para todos los parámetros,
para cualquier profundidad y para toda composición de refinamientos.
\end{proof}

Hay un testigo exacto de que esta colección cambia la forma. Para los primeros
cuatro extremos,
\begin{equation}
 X_{0123}(0)=\frac{1560}{23711},\qquad
 \left.\frac{d}{ds}\log X_{0123}(s)\right|_{s=0}
 =-\frac{309899}{917}-\frac{268596}{4585}\sqrt5<0.
 \label{eq:nc-nontrivial-witness}
\end{equation}
La cuenta utiliza únicamente las primeras tres separaciones
\((234,543,140)\) y sus velocidades ya recuperadas. La desigualdad
estricta prueba un movimiento proyectivo efectivo. Las identidades de
Ptolomeo siguen siendo exactas durante ese movimiento, pues se aplican a
las columnas \((1,t_j(s))\) en cada parámetro. El registro controla así
una deformación positiva de la forma y, simultáneamente, un sistema de
coordenatizaciones que permanece compatible con ella.

El alcance demostrado reúne una red planar de rango dos, su medición
explícita, una semilla plábica de triangulación, todas sus mutaciones por
flips, la naturalidad de los extremos heredados y una deformación de
razones dobles dirigida por \(K\). Cada operación tiene un dominio y un
resultado verificables. La comparación de esta realización con la
dinámica temporal enriquecida y con la rama material utiliza sus respectivos
mapas de lectura; ninguna identidad de coordenadas sustituye esos mapas.

\section{Geometrías amplituhedrales de rango uno y compatibilidad de residuos}
\label{sec:k-poligono-amplituhedral}

El registro dodecafásico permite reunir explícitamente datos positivos,
dominio grassmanniano, imagen cinemática, cobertura, grado sobre las celdas
y forma canónica. El primer resultado concierne al amplituedro
bidimensional de parámetros \((n,k,m)=(13,1,2)\); después se prolonga la
construcción a los polítopos cíclicos de rango uno y dimensión superior.
Se trata de una realización geométrica posterior del registro producido
por APP--TRIT--TPK. La operación adicional declarada es la elevación por
potencias de una partición racional ordenada. Esta elección fija una
realización particular; las correspondencias físicas posteriores conservan
sus lectores y condiciones propios.

La procedencia del registro conserva la composición
\[
 x_{\mathrm{term}}\longmapsto\mathcal R_{12}(x_{\mathrm{term}})
 \longmapsto(b^{90},b^{120},Q_{\mathrm{TPK}})
 \longmapsto U_{\mathrm{sgn}}\longmapsto K.
\]
APP conserva ambas hojas y sus residuos y cocientes; TRIT orienta el
régimen de cada transición; TPK transporta y actualiza el estado con
acarreo, frontera, ruta y memoria. La representación \(K\) se lee en el mismo
sistema de estados enriquecidos que genera conjuntamente el continuo.
Su extracción firmada y su inversión
de Hadamard producen
\begin{equation}
 K=(234,543,140,729,659,824,621,58,914,794,146,601),
 \qquad S_K=\sum_{j=1}^{12}K_j=6263.
 \label{eq:polygon-K}
\end{equation}
La prueba geométrica utiliza la positividad de esas doce coordenadas, su
orden y su total. Sus coeficientes proceden de aquella extracción y no de
la selección retrospectiva de un valor físico.

\subsection{Partición racional y elevación positiva}

Definimos trece puntos racionales de partición:
\begin{equation}
 t_0=0,\qquad t_i=\frac{K_1+\cdots+K_i}{S_K}\quad(1\leq i\leq12),
 \qquad Z_i=(1,t_i,t_i^2)^\mathsf T,\qquad Z=(Z_0\ \cdots\ Z_{12}).
 \label{eq:polygon-Z}
\end{equation}
Sus numeradores, con denominador común \(6263\), son
\[
 (0,234,777,917,1646,2305,3129,3750,3808,4722,5516,5662,6263).
\]
Así, \(0=t_0<t_1<\cdots<t_{12}=1\). La partición etiquetada y el total
conservan el registro mediante la inversa exacta
\begin{equation}
 K_i=S_K(t_i-t_{i-1}).
 \label{eq:polygon-recover-K}
\end{equation}
La partición conserva las proporciones; retener \(S_K\) conserva también
la escala entera.

\begin{proposition}[Datos cinemáticos estrictamente positivos]
\label{prop:polygon-vandermonde}
La matriz \(Z\) tiene rango tres y todos sus menores ordenados de tamaño
tres son positivos. Para \(0\leq a<b<c\leq12\),
\begin{equation}
 \langle abc\rangle:=\det(Z_a,Z_b,Z_c)
 =(t_b-t_a)(t_c-t_a)(t_c-t_b)>0.
 \label{eq:polygon-minors}
\end{equation}
\end{proposition}
\begin{proof}
Se resta la primera columna de las dos últimas y se desarrolla por la
primera fila. Cada diferencia \(t_j^2-t_a^2\) factoriza como
\((t_j-t_a)(t_j+t_a)\), lo que da el producto indicado. Las sumas parciales
estrictamente crecientes prueban su positividad. Cualquiera de esos
menores acredita el rango tres.
\end{proof}

\subsection{Imagen, fibras y celdas de grado uno}

El Grassmanniano \(\operatorname{Gr}_{\geq0}(1,13)\) es el espacio
proyectivo de vectores \(c=(c_0,\ldots,c_{12})\) no negativos y no nulos,
módulo escala positiva. La coordenatización \(\sum c_j=1\) lo identifica con el
símplex cerrado de dimensión doce. Definimos
\begin{equation}
 \Phi_Z:\operatorname{Gr}_{\geq0}(1,13)\longrightarrow\mathbb P^2,
 \qquad[c]\longmapsto\left[\sum_{j=0}^{12}c_jZ_j\right].
 \label{eq:polygon-map}
\end{equation}
En la coordenatización \(Y_0=1\), sean \(v_j=(t_j,t_j^2)\) y
\(P_K=\operatorname{conv}(v_0,\ldots,v_{12})\).

\begin{theorem}[Amplituedro poligonal del registro]
\label{thm:polygon-amplituhedron}
La imagen cerrada de \(\Phi_Z\) es exactamente el polígono convexo
\(P_K=\mathcal A_{13,1,2}(Z)\), de dimensión dos y trece vértices. La imagen
del Grassmanniano estrictamente positivo es el interior de \(P_K\).
Las once celdas
\begin{equation}
 \mathcal S_i=\{[c]:c_0,c_i,c_{i+1}>0,\
 c_j=0\text{ si }j\notin\{0,i,i+1\}\},\qquad 1\leq i\leq11,
 \label{eq:polygon-cells}
\end{equation}
se aplican biyectivamente, con grado uno y orientación positiva, sobre los
interiores de \(T_i=[v_0,v_i,v_{i+1}]\). Sus clausuras cubren \(P_K\) y sus
interiores son disjuntos dos a dos.
\end{theorem}
\begin{proof}
En la normalización \(\sum c_j=1\), la aplicación es
\(c\mapsto\sum c_jv_j\), cuya imagen es la envolvente convexa. Los menores
de \eqref{eq:polygon-minors} muestran que cada triple ordenado tiene
orientación positiva. Cada lado consecutivo deja todos los otros vértices
estrictamente a su izquierda; también lo hace el lado final de \(v_{12}\)
a \(v_0\). Todos los puntos son, por tanto, vértices extremos de un polígono
convexo.

Toda combinación con \(c_j>0\) es interior, pues en cualquier recta de
soporte alguno de los vértices está estrictamente en el semiplano
interior. Recíprocamente, sean \(y\) interior y
\(b=13^{-1}\sum v_j\). Para \(\varepsilon>0\) suficientemente pequeño,
\(y'=(y-\varepsilon b)/(1-\varepsilon)\) permanece en el polígono.
Escribiendo \(y'=\sum a_jv_j\), \(a_j\geq0\), \(\sum a_j=1\), se obtiene
\(y=\sum((1-\varepsilon)a_j+\varepsilon/13)v_j\), con todos los
coeficientes positivos.

Las diagonales desde \(v_0\) dividen el polígono en los once triángulos.
Sobre cada soporte triple, las coordenadas baricéntricas son únicas,
porque \(\langle0,i,i+1\rangle>0\); su inversa explícita aparece en
\eqref{eq:polygon-barycentric}. Cada soporte define una celda positroidal
de rango uno con coordenadas \([1:a:b]\), \(a,b>0\), en esas posiciones.
Una red dirigida planar con una fuente y tres trayectorias de pesos
\(1,a,b\) hacia los tres destinos realiza esas coordenadas por medición
de frontera. El determinante positivo prueba orientación y grado uno.
\end{proof}

La fibra genérica de la aplicación completa tiene dimensión diez:
\(\sum c_j=1\), \(\sum c_jt_j=x\) y \(\sum c_jt_j^2=y\) son tres
ecuaciones independientes en trece coordenadas. Para \(y\) interior la
fibra contiene un punto estrictamente positivo, por lo que su intersección
con el símplex tiene dimensión \(13-3=10\). El grado uno concierne a las
celdas triangulares; la forma bidimensional se obtiene mediante esas
celdas.

\subsection{Forma canónica y residuos orientados}

Para \(Y=(1,x,y)^\mathsf T\), definimos
\begin{equation}
 \ell_{ab}(Y)=\det(Z_a,Z_b,Y)
 =(t_b-t_a)\{y-(t_a+t_b)x+t_at_b\}.
 \label{eq:polygon-lines}
\end{equation}
La expresión vale también cuando \(b<a\). Para \(a<b<c\), sean
\(D_{abc}=\langle abc\rangle\) y
\begin{equation}
 \lambda_a=\frac{\ell_{bc}}{D_{abc}},\qquad
 \lambda_b=\frac{\ell_{ca}}{D_{abc}},\qquad
 \lambda_c=\frac{\ell_{ab}}{D_{abc}},\qquad
 \lambda_a+\lambda_b+\lambda_c=1.
 \label{eq:polygon-barycentric}
\end{equation}
La orientación afín es \(dx\wedge dy\). Se fija la convención de frontera
\(\operatorname{Res}^{\partial}(\eta\wedge du/u)=\eta|_{u=0}\):
la diferencial normal ocupa la última posición. En dimensión dos, el
residuo de Poincaré que coloca \(du/u\) primero tiene el signo opuesto.
Las cancelaciones siguientes usan de forma uniforme la convención declarada.

\begin{theorem}[Forma racional y cancelación de las diagonales]
\label{thm:polygon-form}
La forma canónica del triángulo orientado \([v_a,v_b,v_c]\) es
\begin{equation}
 \Omega_{abc}=
 \frac{D_{abc}^2\,dx\wedge dy}
 {\ell_{ab}(Y)\ell_{bc}(Y)\ell_{ca}(Y)}.
 \label{eq:polygon-triangle-form}
\end{equation}
La forma canónica del polígono es
\begin{equation}
 \boxed{\displaystyle
 \Omega_{P_K}(x,y)=
 \sum_{i=1}^{11}
 \frac{t_it_{i+1}(t_{i+1}-t_i)\,dx\wedge dy}
 {(y-t_ix)\{y-(t_i+t_{i+1})x+t_it_{i+1}\}(t_{i+1}x-y)}.}
 \label{eq:polygon-full-form}
\end{equation}
Los polos de las diagonales interiores se cancelan. Los únicos divisores
polares son las trece rectas de los lados exteriores. Parametrizado cada
lado desde su vértice inicial al final por \(s\in[0,1]\), su residuo de
frontera es \(ds/[s(1-s)]\).
\end{theorem}
\begin{proof}
La forma del símplex baricéntrico es
\(d\lambda_b\wedge d\lambda_c/(\lambda_a\lambda_b\lambda_c)\).
El jacobiano satisface
\(dx\wedge dy=D_{abc}\,d\lambda_b\wedge d\lambda_c\).
Sustituir \eqref{eq:polygon-barycentric} da
\eqref{eq:polygon-triangle-form}. Para \(a=0\), cancelar los tres factores
escalares de las rectas da cada sumando de \eqref{eq:polygon-full-form}.

En el lado \(ab\), pongamos \(\lambda_c=u\),
\(\lambda_b=(1-u)s\), \(\lambda_a=(1-u)(1-s)\). Entonces
\[
 \Omega_{abc}=\frac{ds}{s(1-s)}\wedge\frac{du}{u(1-u)}.
\]
El residuo es la forma del intervalo orientado. Dos triángulos contiguos
inducen orientaciones contrarias sobre su diagonal y producen residuos
opuestos. Al ser simples todos los polos, anular el residuo elimina el
polo de la diagonal como divisor. Cada lado exterior aparece una vez.
La fórmula homogénea es de grado cero como forma proyectiva y no tiene
otros divisores polares. Si dos formas racionales proyectivas tuvieran los
mismos residuos y ningún otro polo, su diferencia sería una dos-forma
holomorfa sobre \(\mathbb P^2\), que es necesariamente cero. Así quedan
probadas caracterización y unicidad.
\end{proof}

Se emplea la noción convencional de geometría positiva mediante polos
logarítmicos y residuos canónicos de frontera, desarrollada por
N.~Arkani-Hamed, Y.~Bai y T.~Lam, \emph{Positive Geometries and Canonical
Forms}, JHEP 11 (2017), 039,
\url{https://arxiv.org/abs/1703.04541}, secciones 2, 3, 5.2, 6.1 y 7.2.
La selección de los coeficientes del presente polígono precede a ese
reconocimiento y procede de \(K\).

\subsection{Refinamiento por acarreo sobre la geometría fija}

En un triángulo \(T=[A,B,C]\) de orientación positiva, escribimos
\begin{equation}
 Y(u,r)=(1-u)A+u\{(1-r)B+rC\},\qquad 0<u,r<1.
 \label{eq:polygon-blowdown}
\end{equation}
El vértice \(A\) corresponde a \(u=0\) y el lado opuesto a \(u=1\).
La forma queda
\begin{equation}
 \Omega_T=\frac{du}{u(1-u)}\wedge\frac{dr}{r(1-r)}.
 \label{eq:polygon-product-form}
\end{equation}
Sea \(B_0\geq2\) una base entera declarada, distinta del vértice \(B\).
La coordenada \(r\) admite la actualización por residuo y acarreo
\begin{equation}
 c=\lfloor B_0r\rfloor,\qquad r'=B_0r-c,\qquad
 r=\frac{c+r'}{B_0}.
 \label{eq:polygon-carry}
\end{equation}
Los extremos conservan sus marcas de frontera. El dígito \(c\) permanece
como memoria y permite invertir la actualización. Conviene precisar las
coordenatizaciones de los dos transportes afines. El zoom de la subcelda actúa sobre
\(r\in[c/B_0,(c+1)/B_0]\) mediante \(r'=B_0r-c\).
En cambio, el transporte de la pareja levantada
\(L_c(x,y)=(B_0x-c,B_0y+c)\), con \(x+y=M\), cambia la masa total
a \(B_0M\) y expresa la coordenada
\[
 \frac{B_0x-c}{B_0M}=\frac{x}{M}-\frac{c}{B_0M}.
\]
Son aplicaciones de dominios y normalizaciones diferentes. Ambas conservan
el acarreo en el estado enriquecido y admiten recuperación con sus datos
de escala. La demostración siguiente utiliza el zoom de subcelda y la
partición geométrica que éste identifica.

\begin{theorem}[Invariancia exacta bajo subdivisión con memoria]
\label{thm:polygon-carry}
Sean \(s_j=j/B_0\) y \(W_j=(1-s_j)B+s_jC\).
Los triángulos \(T_j=[A,W_j,W_{j+1}]\), \(0\leq j<B_0\), subdividen
\(T\) y satisfacen
\begin{equation}
 \sum_{j=0}^{B_0-1}\Omega_{T_j}=\Omega_T.
 \label{eq:polygon-carry-additivity}
\end{equation}
La igualdad permanece válida a toda profundidad finita y para cualquier
partición racional ordenada del lado. Subdividir los once triángulos
de \(P_K\) deja \(\Omega_{P_K}\) exactamente invariante.
\end{theorem}
\begin{proof}
Las regiones \(r\in[s_j,s_{j+1}]\) cubren \(T\) con interiores disjuntos.
En las coordenadas comunes,
\[
 \Omega_{T_j}=\frac{du}{u(1-u)}\wedge
 \frac{(s_{j+1}-s_j)\,dr}{(r-s_j)(s_{j+1}-r)}.
\]
La identidad racional
\[
 \frac{s_{j+1}-s_j}{(r-s_j)(s_{j+1}-r)}
 =\frac1{r-s_j}+\frac1{s_{j+1}-r}
\]
telescopa: cada término interior cancela el del subintervalo contiguo.
Quedan \(1/r+1/(1-r)=1/[r(1-r)]\). Esto prueba igualdad de formas
racionales; iterarla proporciona toda profundidad finita.
La prueba sólo utiliza \(s_0=0<s_1<\cdots<s_N=1\), de modo que cubre
particiones no uniformes. El cambio
\(r'=(r-s_j)/(s_{j+1}-s_j)\) devuelve localmente la forma unitaria.
La palabra de acarreos identifica la subcelda y su inversa recupera la
coordenada anterior. Los nuevos polos interiores se cancelan mediante
la misma identidad telescópica.
\end{proof}

La suma anterior reúne la partición geométrica completa. Cuando un lector
de supervivencia selecciona sólo algunas historias, su imagen es la unión
de las subceldas seleccionadas y su forma es la suma correspondiente;
coincide con la forma del polígono completo precisamente cuando esa unión
lo cubre con multiplicidad orientada uno. Por tanto, el criterio de
cobertura para una selección concreta de historias queda expresado por
una condición comprobable sobre su cadena de celdas.

Este refinamiento conserva el polígono fijo. Insertar un nuevo punto
\((t,t^2)\) de la parábola entre dos parámetros consecutivos \(a<t<b\)
produce otra geometría: la altura de la cuerda menos \(t^2\) es
\((t-a)(b-t)>0\), y el nuevo punto queda fuera del polígono previo.
La invariancia demostrada corresponde a la subdivisión rectilínea de
celdas, cuya frontera exterior permanece fija.

\subsection{Cambios de triangulación y reparametrización positiva}

Para cuatro vértices ordenados \(A,B,C,D\), se cumple
\begin{equation}
 \Omega_{ABC}+\Omega_{ACD}=\Omega_{ABD}+\Omega_{BCD}.
 \label{eq:polygon-flip}
\end{equation}
Los dos miembros tienen los mismos residuos exteriores y cancelan su
diagonal interior; su igualdad procede de la unicidad demostrada.
También se comprueba multiplicando por sus denominadores lineales.
Todo cambio entre triangulaciones de un polígono convexo se descompone
en estos cambios de diagonal, por lo que la forma es independiente
de la triangulación. Una identificación con coordenadas cluster debe
conservar adicionalmente la variedad y sus relaciones de cambio.

Si \(D=\operatorname{diag}(d_0,\ldots,d_{12})\), \(d_j>0\), entonces
\begin{equation}
 \mathcal A_{13,1,2}(ZD)=\mathcal A_{13,1,2}(Z).
 \label{eq:polygon-rescale}
\end{equation}
El automorfismo \([c]\mapsto[Dc]\) reparametriza la aplicación cinemática.
En la coordenatización afín, \(c_j\) cambia a
\(c_jd_j/\sum_\ell c_\ell d_\ell\), que recorre nuevamente el símplex.
La imagen y su forma canónica permanecen iguales. El desplazamiento
de una historia parametrizada y el cambio del conjunto geométrico
completo son así operaciones distinguibles.

\subsection{Extensión de rango uno a dimensiones superiores}

Para \(2\leq m\leq12\), definimos
\begin{equation}
 Z_i^{(m)}=(1,t_i,t_i^2,\ldots,t_i^m)^\mathsf T,\qquad
 P_K^{(m)}=\operatorname{conv}
 \{(t_i,t_i^2,\ldots,t_i^m):0\leq i\leq12\}.
 \label{eq:polygon-higher-lift}
\end{equation}
Los menores máximos ordenados son los productos de Vandermonde
\(\prod_{a<b}(t_{i_b}-t_{i_a})>0\). La matriz tiene rango \(m+1\).
La imagen de \(\operatorname{Gr}_{\geq0}(1,13)\) es el polítopo cíclico
\(P_K^{(m)}=\mathcal A_{13,1,m}(Z^{(m)})\), de dimensión \(m\); su fibra
genérica tiene dimensión \(12-m\). Aquí se usa la extensión matemática
de la definición amplituhedral a cualquier \(m\); las aplicaciones
habituales a amplitudes imponen sus propios valores de \(m\).

La triangulación se determina por el orden fijo de etiquetas. Se elige
el menor vértice de un polítopo, se triangulan recursivamente todas sus
facetas que no lo contienen y se toman conos desde ese vértice. La
inducción comienza con los puntos. Una recta desde el vértice inicial
hasta un punto interior corta la frontera opuesta en una faceta; ello
prueba la cobertura. La unicidad del punto de intersección y la
compatibilidad de las restricciones inducidas por el mismo orden prueban
que los interiores son disjuntos y que las intersecciones son caras
comunes. En cada paso disminuye la dimensión, por lo que el procedimiento
termina. Esta es una triangulación por tracción determinada por las
etiquetas.

Para un símplex orientado
\(T=[p_0,\ldots,p_m]\) de esa triangulación, sean
\[
 A_T=(p_1-p_0\ \cdots\ p_m-p_0),\qquad
 (\lambda_1,\ldots,\lambda_m)^\mathsf T=A_T^{-1}(x-p_0),
 \qquad\lambda_0=1-\sum_{j=1}^m\lambda_j.
\]
La orientación de sus vértices se toma de modo que \(\det A_T>0\).
Su forma canónica es
\begin{equation}
 \Omega_T(x)=
 \frac{dx_1\wedge\cdots\wedge dx_m}
 {\det(A_T)\prod_{j=0}^m\lambda_j(x)},\qquad
 \Omega_{P_K^{(m)}}=\sum_{T}\Omega_T.
 \label{eq:polygon-higher-form}
\end{equation}
El cambio baricéntrico prueba la fórmula y que la restricción cinemática
sobre el soporte de \(T\) tiene grado uno. Los residuos son las formas
de las caras con los signos inducidos por la orientación. En una cara
interior, las dos orientaciones son contrarias y el residuo total se
anula. En una cara exterior queda la suma de las formas de la
triangulación de esa cara. La inducción sobre la dimensión demuestra
que dicha suma es su forma canónica. Repetir el residuo a lo largo de
una bandera de caras alcanza los vértices, donde la normalización es
\(\pm1\). Se obtiene así la correspondencia de residuos en todas las
codimensiones para esta colección de rango uno.

La misma demostración cubre cualquier subdivisión simplicial compatible
del polítopo fijo: sus caras interiores se cancelan y las exteriores
conservan los residuos. La inexistencia de formas holomorfas de grado
máximo sobre \(\mathbb P^m\) asegura la unicidad de la forma racional.
En particular, la estabilidad bajo refinamiento, la cobertura y los
residuos quedan cerrados en esta colección de dimensiones \(2,\ldots,12\).
Los grados grassmannianos \(k>1\) conservan sus aplicaciones y sus
obligaciones demostrativas específicas.

\section{Matrices de momentos y realizaciones amplituhedrales de rango superior}
\label{sec:higher-rank}
La partición del registro dodecafásico permite prolongar la construcción
de rango uno a puntos explícitos de rango superior. Los mismos trece
nodos ordenados proporcionan simultáneamente menores de Plücker positivos
y una forma de momentos definida positiva. Estas dos positividades se
obtienen de una sola lista de intervalos; la segunda prueba que la
composición conserva el rango requerido.

Sea \(d_j=K_j/Q\), donde \(Q=\sum_{j=1}^{12}K_j=6263\), y definamos
\[
 t_0=0,\qquad t_j=\sum_{r=1}^jd_r.
\]
La positividad de los doce bloques da
\(0=t_0<t_1<\cdots<t_{12}=1\). Para cada \(2\le k\le9\), ponemos
\[
 C^{(k)}_{ai}=t_i^a\quad(0\le a<k),\qquad
 Z^{(k)}_{bi}=t_i^b\quad(0\le b<k+4),\qquad 0\le i\le12.
\]
La convención \(t_i^0=1\) incluye el nodo cero. Así, \(C^{(k)}\)
es una matriz de tamaño \(k\times13\), mientras \(Z^{(k)}\) tiene
tamaño \((k+4)\times13\).

\begin{theorem}[Positividad y rango de la realización de momentos]
La matriz \(C^{(k)}\) representa un punto de
\(\Gr_{>0}(k,13)\), y todos los menores máximos ordenados de
\(Z^{(k)}\) son positivos. La matriz
\[
 Y^{(k)}=C^{(k)}(Z^{(k)})^{\mathsf T},\qquad
 Y^{(k)}_{ab}=\sum_{i=0}^{12}t_i^{a+b},
\]
tiene rango \(k\). En particular, para el dato externo positivo
\(Z^{(k)}\) ya construido a partir del registro,
\[
 [Y^{(k)}]\in\mathcal A_{13,k,4}(Z^{(k)}),\qquad 2\le k\le9,
\]
donde
\[
 \mathcal A_{13,k,4}(Z^{(k)})
 =\{[C(Z^{(k)})^{\mathsf T}]:[C]\in\Gr_{\ge0}(k,13)\}.
\]
\end{theorem}
\begin{proof}
Un menor máximo de cualquiera de las dos matrices iniciales tiene la forma
\[
 \det(t_{i_j}^{a})_{a=0}^{r-1}
 =\prod_{p<q}(t_{i_q}-t_{i_p})>0,
\]
para índices crecientes. La desigualdad prueba las dos afirmaciones sobre
Plücker. El bloque inicial \(k\times k\) de \(Y^{(k)}\) es la matriz
de Hankel \(H^{(k)}\). Para \(x\ne0\),
\[
 x^{\mathsf T}H^{(k)}x
 =\sum_{i=0}^{12}\left(\sum_{a=0}^{k-1}x_at_i^a\right)^2>0.
\]
En efecto, un polinomio no nulo de grado inferior a \(k\le9\) no se
anula en trece puntos distintos. Luego \(H^{(k)}\) es definida positiva,
su determinante es estrictamente positivo y \(Y^{(k)}\) tiene rango
\(k\). La pertenencia amplituhedral es entonces la composición de las
dos matrices positivas anteriores, con el rango ya comprobado.
\end{proof}

En rango dos, la prueba adopta la forma particularmente explícita
\[
 \det H^{(2)}=13\sum_i t_i^2-\left(\sum_i t_i\right)^2
 =\sum_{i<j}(t_j-t_i)^2>0.
\]
La separación de los nodos, producida por la positividad del registro,
queda así convertida en coercividad de una forma cuadrática. A cada
registro fijo corresponde un punto determinado para cada rango. La
cobertura completa, las triangulaciones y los residuos demostrados en la
sección anterior mantienen su dominio de rango uno; el presente teorema
añade la colección explícita de puntos de rangos dos a nueve y sus matrices
de momentos. La distinción permite reutilizar cada demostración con su
alcance exacto.

% Sección aditiva. No modifica la edición de 211 páginas.
% Procedencia y comprobaciones: documentos internos de este mismo directorio.
\section{Refinamiento nonádico, banderas positivas y compatibilidad de los momentos}
\label{sec:dim-cofinal}

La longitud de una representación finita determina cuántos menores pueden
formarse con sus columnas. El refinamiento conserva las columnas heredadas
y añade otras con procedencia conocida. Por ello, la dimensión de una
realización matricial puede crecer sin sustituir el registro que fija sus
separaciones. La construcción siguiente reúne este crecimiento con dos
nociones precisas de compatibilidad: restricción de evaluaciones a las
marcas anteriores y promedio de observables sobre las celdas descendientes.
La primera preserva las banderas polinomiales; la segunda conserva una
medida y separa el incremento de información en detalles ortogonales.

El dato de partida es la salida dodecafásica ya producida por
APP--TRIT--TPK. Las dos hojas de APP conservan sus divisiones completas;
TRIT orienta el régimen y el operador
\(U_t=\operatorname{Upd}_t\circ\operatorname{Tra}_t\circ
\operatorname{Sel}_t\) transporta el estado enriquecido. La prolongación
\(w_6\to w_{12}\to w_{18}\to w_{24}\to w_{30}\to R_{36}\to G_9\)
conserva retorno de fase y avance de memoria. La estructura discreta conjunta
del continuo conserva correlativamente las representaciones de ese estado,
con sus hojas, orientación, memoria y fronteras. De ese mismo estado proceden
las representaciones regionales y el registro firmado cuya inversión da
\[
 K=(234,543,140,729,659,824,621,58,914,794,146,601),
 \qquad Q=\sum_{j=1}^{12}K_j=6263.
\]
Se utiliza esta salida con sus marcas y orientación; las matrices de
momentos son realizaciones posteriores. El índice de refinamiento que se
introduce abajo cuenta subdivisiones del lector geométrico. Su identificación
con una duración física o con un número de microtransiciones del TPK
requiere conservar el reloj de esa otra realización.

\subsection{Particiones marcadas y crecimiento cofinal}

Escribamos \(d_j=K_j/Q\), \(a_0=0\) y
\(a_j=\sum_{h=1}^j d_h\). Para \(r\geq0\), definimos
\[
 L_r=12\,9^r,\qquad N_r=L_r+1,\qquad
 t_{r,(j-1)9^r+c}=a_{j-1}+\frac{c}{9^r}d_j,
 \quad 0\leq c\leq9^r.
\]
Los extremos comunes se cuentan una sola vez. Las \(L_r\) celdas se
indexan por \((j,c)\), con \(1\leq j\leq12\) y
\(0\leq c<9^r\), y tienen longitud
\(\ell_{r,j,c}=d_j/9^r\). Se conserva también la escritura de \(c\)
en \(r\) cifras nonádicas, incluidos los ceros iniciales. Un sucesor tiene
índice \(9c+h\), con \(0\leq h\leq8\); la división euclídea recupera
su predecesor y la última cifra. Los extremos compartidos retienen su marca
de frontera; la convención de intervalos semiabiertos sólo organiza
las celdas y no elimina ninguna marca.

\begin{proposition}[Cofinalidad y recuperación del registro]
\label{prop:dim-grid}
Las marcas satisfacen
\[
 0=t_{r,0}<t_{r,1}<\cdots<t_{r,L_r}=1,
 \qquad t_{r+1,9i}=t_{r,i}.
\]
El máximo de las longitudes es
\(h_r=(\max_jK_j)/(Q9^r)\), que tiende a cero. Para todo par de
enteros \(k\geq1\), \(m\geq1\), existe \(r_0\) tal que
\(k+m\leq N_r\) para \(r\geq r_0\). El registro se recupera en
cualquier nivel por
\[
 K_j=Q\sum_{c=0}^{9^r-1}\ell_{r,j,c}.
\]
\end{proposition}
\begin{proof}
Cada diferencia consecutiva es \(d_j/9^r>0\). Al reemplazar \(c\)
por \(9c\), la fórmula del nivel siguiente da exactamente la marca
anterior. Las nueve longitudes de los sucesores suman la longitud del predecesor;
la iteración prueba la recuperación. Finalmente \(N_r=12\,9^r+1\)
crece sin cota. Una elección explícita es cualquier \(r_0\geq0\) con
\(9^{r_0}\geq(k+m-1)/12\). Todo el argumento conserva \(Q\), el
índice de bloque y las cifras de descendencia.
\end{proof}

La normalización determina las razones \(K_j/Q\). La recuperación de
los enteros utiliza, además, la carga \(Q\) conservada por el registro.
Así se distinguen la geometría normalizada y la información de escala
que acompaña a su representación.

\subsection{Banderas de evaluaciones y realizaciones en toda dimensión finita}

Para \(1\leq q\leq N_r\), sea
\[
 V_r^{(q)}=(t_{r,i}^{\,a})_{0\leq a<q,\,0\leq i<N_r},
 \qquad F_r^q=\operatorname{fila}V_r^{(q)}.
\]
La potencia de exponente cero vale uno también en el nodo cero.

\begin{theorem}[Bandera positiva y composición de rango arbitrario]
\label{thm:dim-all-ranks}
Los espacios \(F_r^q\) forman una bandera
\(F_r^1\subset F_r^2\subset\cdots\subset F_r^{N_r}\), con
\(\dim F_r^q=q\), y cada matriz \(V_r^{(q)}\) tiene todos sus
menores máximos ordenados estrictamente positivos. Para
\(k\geq1\), \(m\geq1\), \(k+m\leq N_r\), las matrices
\[
 C_r=V_r^{(k)},\quad Z_r=V_r^{(k+m)},\quad
 Y_r=C_rZ_r^{\mathsf T}
\]
satisfacen
\[
 (Y_r)_{ab}=\sum_{i=0}^{L_r}t_{r,i}^{a+b},\qquad
 \operatorname{rank}Y_r=k,\qquad
 [Y_r]\in\mathcal A_{N_r,k,m}(Z_r).
\]
La construcción alcanza todo par finito \((k,m)\) mediante un nivel
suficientemente profundo de la misma partición marcada.
\end{theorem}
\begin{proof}
Para \(i_1<\cdots<i_q\), el menor correspondiente es
\[
 \det(t_{r,i_b}^{a-1})_{a,b=1}^{q}
   =\prod_{a<b}(t_{r,i_b}-t_{r,i_a})>0.
\]
Esto prueba rango, positividad e inclusiones de la bandera. El bloque
inicial \(k\times k\) de \(Y_r\) es el Gram de las evaluaciones
de \(1,t,\ldots,t^{k-1}\). Para \(x\neq0\),
\[
 x^{\mathsf T}(Y_r)_{[0,k-1],[0,k-1]}x
 =\sum_i\left(\sum_{a=0}^{k-1}x_at_{r,i}^{a}\right)^2>0,
\]
porque un polinomio no nulo de grado inferior a \(k\) no se anula
en los \(N_r\geq k+m\) nodos distintos. El bloque es invertible,
y \(Y_r\), que tiene \(k\) filas, tiene rango \(k\).
La pertenencia indicada es la aplicación de la matriz externa positiva
\(Z_r\) al punto positivo representado por \(C_r\).
\end{proof}

Se emplea la convención matemática
\[
 \mathcal A_{N,k,m}(Z)=
 \{[CZ^{\mathsf T}]:[C]\in\operatorname{Gr}_{\geq0}(k,N)\}.
\]
La definición para \(m\) general puede consultarse en S.~N.~Karp,
L.~K.~Williams y Y.~X.~Zhang,
\href{https://people.math.harvard.edu/~williams/papers/amplituhedron-plus.pdf}
{\emph{Decompositions of amplituhedra}}, definición~1.1.
Aquí el dato externo ya procede de la partición expresada por HMT.
La conclusión construye una colección determinada de puntos y banderas;
las dimensiones del Grassmanniano ambiente y de un atlas de su imagen
son objetos adicionales a los rangos probados.

Sea \(E_r:\mathbb R^{N_{r+1}}\to\mathbb R^{N_r}\) la restricción
a las posiciones \(9i\). Para todo \(q\leq N_r\),
\[
 E_r\bigl((V_{r+1}^{(q)})^{\mathsf T}x\bigr)
       =(V_r^{(q)})^{\mathsf T}x.
\]
Por tanto, \(E_r:F_{r+1}^q\to F_r^q\) es un isomorfismo sobre
estos espacios de evaluaciones: sus dos realizaciones identifican el
mismo polinomio de grado inferior a \(q\). La composición de las
restricciones coincide con la restricción directa. Los menores sobre
índices heredados conservan exactamente su valor.

\begin{proposition}[Recuperación desde una bandera calibrada]
Si se conservan el orden de las columnas, las marcas extremas y la
calibración \(t_{r,0}=0\), \(t_{r,L_r}=1\), el plano \(F_r^2\)
determina los nodos y las separaciones. Junto con \(Q\) y las
marcas de descendencia determina \(K\).
\end{proposition}
\begin{proof}
El plano contiene el vector constante y el vector de nodos. Cualquier
otra coordenada que, junto con el vector constante, genere el mismo
plano es \(a+bt\), con \(b\neq0\). Las dos calibraciones fijan
\(a=0\), \(b=1\). Las diferencias consecutivas recuperan las
longitudes, y la proposición~\ref{prop:dim-grid} recupera el registro.
\end{proof}

\subsection{Cuerpos convexos de rango uno y crecimiento del dominio}

Para \(m\geq1\) y \(m+1\leq N_r\), definimos
\[
 P_r^{(m)}=\operatorname{conv}
 \{(t_{r,i},t_{r,i}^2,\ldots,t_{r,i}^m):0\leq i\leq L_r\}.
\]
La normalización de una fila positiva de \(C\) como coeficientes
baricéntricos prueba
\(\mathcal A_{N_r,1,m}(Z_r)=P_r^{(m)}\). El determinante de
Vandermonde proporciona \(m+1\) puntos afínmente independientes,
de modo que el cuerpo tiene dimensión \(m\).

\begin{proposition}[Cobertura simplicial a dimensión finita y límite convexo]
\label{prop:dim-polytopes}
Cada \(P_r^{(m)}\) admite una triangulación por sus vértices, con
grado uno sobre el interior de cada celda simplicial. La suma de sus
formas canónicas tiene residuos interiores cancelados. Además,
\[
 P_r^{(m)}\subseteq P_{r+1}^{(m)},\qquad
 \overline{\bigcup_{r\geq r_0}P_r^{(m)}}=
 \operatorname{conv}\{(t,t^2,\ldots,t^m):0\leq t\leq1\}.
\]
\end{proposition}
\begin{proof}
Se ordenan los vértices por sus marcas. En cada cara se elige su menor
vértice; se triangulan recursivamente las facetas que no lo contienen
y se forman los conos correspondientes. La inducción sobre la dimensión
prueba cobertura e interiores disjuntos, y la misma regla en caras
comunes garantiza compatibilidad. En cada simplex orientado
\([v_0,\ldots,v_m]\), con matriz
\(A=(v_1-v_0,\ldots,v_m-v_0)\) de determinante positivo y
coordenadas baricéntricas \(\lambda_0,\ldots,\lambda_m\), la
aplicación baricéntrica es afín e invertible. Su forma es
\[
 \Omega_{[v_0,\ldots,v_m]}=
 \frac{dx_1\wedge\cdots\wedge dx_m}
 {\det A\prod_{j=0}^{m}\lambda_j}.
\]
La restricción residual a cada cara es su forma orientada. Dos simplices
adyacentes inducen orientaciones opuestas sobre su cara común, de modo
que los residuos interiores se cancelan. Esta es la misma operación de
triangulación y residuo aplicada en la realización de rango uno anterior,
con el número de vértices ahora aumentado por el refinamiento.

La inclusión de nodos heredados prueba la inclusión de cuerpos. La malla
tiende a cero, por lo que las curvas muestreadas son densas en la curva
compacta de momentos. Toda combinación convexa finita de puntos de la
curva se aproxima por combinaciones con nodos de un nivel común.
Recíprocamente, todos los cuerpos están contenidos en la envolvente
convexa compacta de la curva. Ambas inclusiones prueban la igualdad.
\end{proof}

La cancelación de residuos opera sobre una triangulación o subdivisión
de un cuerpo fijo. El paso de \(P_r^{(m)}\) a \(P_{r+1}^{(m)}\)
añade vértices y puede ampliar el cuerpo. Sus formas pertenecen a esos
dominios respectivos. El límite convexo precedente especifica la
convergencia de los cuerpos, sin imponer una identidad entre sus formas.

\subsection{Medida del calendario y convergencia de los momentos nodales}

Sea \(F_K:[0,1]\to[0,1]\) la función creciente afín por tramos
que lleva \((j-1)/12\) a \(a_{j-1}\) y \(j/12\) a \(a_j\).
En el intervalo correspondiente su pendiente es \(12d_j>0\).
La medida del calendario transportado es
\[
 \mu_K=(F_K)_*(ds),\qquad
 \frac{d\mu_K}{dt}=\frac1{12d_j}\quad
 (a_{j-1}<t<a_j).
\]
Cada intervalo inicial tiene masa \(1/12\), y cada celda de nivel
\(r\) tiene masa \(1/L_r\). Esta elección asigna el mismo peso
a las posiciones del calendario antes de realizar sus longitudes.
La medida de longitud \(dt\) es otro lector; ambas coinciden cuando
todas las separaciones iniciales son iguales.

Sus momentos son racionales y están determinados por \(K\):
\begin{equation}
 M_p=\int_0^1t^p\,d\mu_K(t)
   =\frac1{12(p+1)}\sum_{j=1}^{12}
       \frac{a_j^{p+1}-a_{j-1}^{p+1}}{d_j},\qquad p\geq0.
 \label{eq:dim-moments}
\end{equation}
\begin{proposition}[Recuperación por momentos y cuantiles]
\label{prop:dim-moment-recovery}
La sucesión completa \((M_p)_{p\geq0}\) determina \(\mu_K\).
Si \(H_K(t)=\mu_K([0,t])\) es su función de distribución, entonces
\[
 a_j=H_K^{-1}(j/12),\qquad
 K_j=Q\bigl(H_K^{-1}(j/12)-H_K^{-1}((j-1)/12)\bigr).
\]
En consecuencia, los momentos completos, la carga \(Q\) y el orden
del calendario recuperan el registro inicial.
\end{proposition}
\begin{proof}
Dos probabilidades sobre \([0,1]\) con los mismos momentos tienen las
mismas integrales para todos los polinomios. La aproximación uniforme
de funciones continuas por polinomios y la masa finita extienden esa
igualdad a todas las funciones continuas; éstas determinan la medida
de Borel. La densidad de \(\mu_K\) es positiva y constante en cada
intervalo inicial. Su distribución es, por tanto, continua y
estrictamente creciente, y cumple \(H_K(a_j)=j/12\). La inversión
de esta igualdad y la diferencia de extremos prueban las fórmulas.
\end{proof}

La probabilidad nodal correspondiente al teorema
\ref{thm:dim-all-ranks} es
\(\nu_r=N_r^{-1}\sum_{i=0}^{L_r}\delta_{t_{r,i}}\).

\begin{theorem}[Límite de la representación nodal]
\label{thm:dim-nodal-limit}
Se cumple \(\nu_r\rightharpoonup\mu_K\). Si \(f\) es continua,
\[
 \left|\int f\,d\nu_r-\int f\,d\mu_K\right|
 \leq \omega_f(h_r)+\frac{\operatorname{osc}(f)}{N_r},
\]
donde \(\omega_f\) es su módulo de continuidad. Para \(p\geq1\),
el error del momento de orden \(p\) está acotado por
\(p h_r+1/N_r\). Para cada par fijo \((k,m)\),
\[
 \frac1{N_r}Y_r\longrightarrow (M_{a+b})_{
     0\leq a<k,\,0\leq b<k+m}=:Y_\infty^{(k,m)},
 \qquad\operatorname{rank}Y_\infty^{(k,m)}=k.
\]
\end{theorem}
\begin{proof}
La medida que asigna masa \(1/L_r\) a cada extremo izquierdo de
celda aproxima su integral con error a lo sumo \(\omega_f(h_r)\).
En efecto, cada celda tiene esa masa y diámetro no mayor que \(h_r\).
La medida \(\nu_r\) es \(L_r/N_r\) veces dicha medida más
\(\delta_1/N_r\), lo que añade a lo sumo
\(\operatorname{osc}(f)/N_r\). Para \(f(t)=t^p\), la constante
de Lipschitz es \(p\). La convergencia de cada entrada se sigue de
estas cotas. El bloque inicial de la matriz límite es definido positivo:
la integral del cuadrado de un polinomio no nulo es positiva, porque
\(\mu_K\) tiene densidad estrictamente positiva en todo intervalo
abierto de \([0,1]\). Así el límite conserva rango \(k\).
\end{proof}

Los momentos nodales varían exactamente por la incorporación de nodos.
Si \(\mathcal J_r\) es el conjunto de las nuevas posiciones y
\(z_q(t)=(1,t,\ldots,t^{q-1})^{\mathsf T}\), entonces
\[
 V_{r+1}^{(q)}(V_{r+1}^{(q)})^{\mathsf T}
 =V_r^{(q)}(V_r^{(q)})^{\mathsf T}
  +\sum_{i\in\mathcal J_r}z_q(t_{r+1,i})z_q(t_{r+1,i})^{\mathsf T}.
\]
Cada sumando es positivo semidefinido. Para las matrices rectangulares
\(Y_r\) vale la misma actualización con
\(z_kz_{k+m}^{\mathsf T}\). La igualdad expresa el cambio de la
medida de conteo; la normalización y el teorema anterior describen su
límite. El objeto límite de rango \(k\) tiene coordenadas en un
Grassmanniano finito, mientras los datos externos \(Z_r\) pertenecen
a las distintas representaciones de cada nivel.

\subsection{Promedios celulares, operadores compatibles y detalle ortogonal}

Para una celda \(I_{r,i}=[b_{r,i},b_{r,i}+\ell_{r,i}]\), ponemos
\[
 a_{r,i}^{(p)}=
 \frac{(b_{r,i}+\ell_{r,i})^{p+1}-b_{r,i}^{p+1}}
 {(p+1)\ell_{r,i}},\qquad
 A_r^{(q)}=(a_{r,i}^{(p)})_{0\leq i<L_r,\,0\leq p<q}.
\]
Son las medias de los monomios respecto de la probabilidad condicional
de \(\mu_K\) en cada celda. Esta probabilidad es uniforme en la
celda, aunque \(\mu_K\) no sea la medida de longitud global.

Sea \(\mathcal H_r=\mathbb R^{L_r}\) con producto
\(\langle x,y\rangle_r=L_r^{-1}\sum_i x_i y_i\). La aplicación
\(J_r:\mathcal H_r\to\mathcal H_{r+1}\) repite el valor de un
predecesor en sus nueve sucesores. Su adjunta \(R_r=J_r^*\) promedia esos
nueve valores. Se cumple \(R_rJ_r=I\).

\begin{theorem}[Compatibilidad exacta y balance de los momentos]
\label{thm:dim-cell-balance}
Para todo \(p\geq0\),
\[
 R_r a_{r+1}^{(p)}=a_r^{(p)},\qquad
 \frac1{L_r}\sum_i a_{r,i}^{(p)}=M_p.
\]
Definamos
\[
 D_r^{(q)}=A_{r+1}^{(q)}-J_rA_r^{(q)},\quad
 G_r^{(q)}=\frac1{L_r}(A_r^{(q)})^{\mathsf T}A_r^{(q)}.
\]
Entonces
\[
 R_rD_r^{(q)}=0,\qquad
 G_{r+1}^{(q)}=G_r^{(q)}+
       \frac1{L_{r+1}}(D_r^{(q)})^{\mathsf T}D_r^{(q)},
\]
y
\[
 G_r^{(q)}\longrightarrow H_K^{(q)}=(M_{a+b})_{0\leq a,b<q}.
\]
Los detalles y el nivel inicial reconstruyen todos los niveles. Para
\(q\leq L_r\), la matriz \(A_r^{(q)}\) tiene rango \(q\) y
todos los menores ordenados de orden \(q\) son positivos.
\end{theorem}
\begin{proof}
Las masas de los nueve sucesores son iguales y sus intervalos particionan
el intervalo predecesor. La aditividad de la integral da la primera identidad;
la suma sobre todas las celdas da la segunda. La descomposición
\(A_{r+1}=J_rA_r+D_r\) es ortogonal en cada columna porque
\(R_rD_r=0\). Al expandir su Gram se anulan los términos cruzados,
y \(J_r\) es una isometría. Esto prueba el balance matricial.

El observable constante por celdas que toma el valor \(a_{r,i}^{(p)}\)
aproxima uniformemente \(t^p\), con error no mayor que \(p h_r\)
para \(p\geq1\). Por tanto converge en \(L^2(\mu_K)\), y los
productos internos convergen a \(M_{a+b}\). La reconstrucción se
obtiene iterando \(A_{r+1}=J_rA_r+D_r\).

Para probar los menores, elijamos \(q\) celdas ordenadas. Por
multilinealidad y Fubini, el determinante de sus promedios es la media
sobre el producto de esas celdas de
\(\prod_{a<b}(x_b-x_a)\). Este integrando es positivo en el
interior del producto porque las celdas son disjuntas y están ordenadas.
Sus fronteras tienen medida nula. La media es estrictamente positiva.
\end{proof}

La identidad se prolonga al par de grados \(k\) y \(k+m\):
\[
 \overline Y_r=\frac1{L_r}(A_r^{(k)})^{\mathsf T}A_r^{(k+m)},
 \qquad
 \overline Y_{r+1}-\overline Y_r
 =\frac1{L_{r+1}}(D_r^{(k)})^{\mathsf T}D_r^{(k+m)}.
\]
Para \(k+m\leq L_r\), las matrices de promedios definen una
segunda realización amplituhedral positiva, de \(L_r\) columnas,
y \(\overline Y_r\) tiene rango \(k\). Son promedios de celdas,
mientras \(Y_r\) usa evaluaciones en \(N_r=L_r+1\) nodos. Ambas
realizaciones convergen a \(Y_\infty^{(k,m)}\); su coincidencia
en el límite procede de las demostraciones anteriores.

La versión operatoria es igualmente explícita. Para una función continua
\(f\), sea \(\Psi_r(f)\) el operador diagonal de sus medias
condicionales en \(\mathcal H_r\). Es lineal, positivo y unital, y
\[
 R_r\Psi_{r+1}(f)J_r=\Psi_r(f).
\]
Esta igualdad se prueba aplicando ambos miembros a cada vector base.
Para \(X(t)=t\), los dos primeros operadores determinan, celda a
celda,
\[
 \Psi_r(X)_i=b_{r,i}+\frac{\ell_{r,i}}2,\qquad
 \Psi_r(X^2)_i-\Psi_r(X)_i^2=\frac{\ell_{r,i}^2}{12}.
\]
Si \(v_{r,i}\) denota esa varianza, la fórmula
\(\ell_{r,i}=\sqrt{12v_{r,i}}\) recupera la longitud; la
primera identidad recupera el extremo inicial. Con orientación, marcas y \(Q\),
se recupera nuevamente \(K\). El promedio es una compresión positiva:
la varianza escrita cuantifica por qué \(\Psi_r(X^2)\) y
\(\Psi_r(X)^2\) son operaciones distintas.

\subsection{Transporte de la medida y alcance del crecimiento dimensional}

Dos listas positivas de separaciones \(d\) y \(d'\), con las
mismas marcas, determinan un homeomorfismo creciente afín por intervalos
\(F=F_{K'}\circ F_K^{-1}\). Lleva cada celda a su celda homóloga
y satisface \(F_*\mu_K=\mu_{K'}\). Por ello
\[
 \mathcal U_F:L^2(\mu_K)\longrightarrow L^2(\mu_{K'}),
 \qquad\mathcal U_F f=f\circ F^{-1}
\]
es una isometría sobreyectiva. Si \(P_r\) y \(P'_r\) son las
proyecciones a funciones constantes por celdas, entonces
\(P'_r\mathcal U_F=\mathcal U_FP_r\): las medias condicionales
se transportan porque las masas de las celdas correspondientes coinciden.
Esto aplica en particular al flujo positivo de separaciones dirigido por
\(K\), conservando la dirección de cada bloque en sus sucesores. Los
momentos de la coordenada espacial pueden cambiar; la isometría transporta
el observable junto con la medida y sus proyecciones.

El crecimiento demostrado es cofinal en toda dimensión externa finita.
Conserva un sistema de banderas positivas, una colección explícita de
puntos amplituhedrales, cuerpos convexos completos en rango uno y una
medida cuyos operadores de promedio son exactamente compatibles. La
memoria de los detalles permite invertir la compresión entre niveles.
Las pruebas de cobertura y residuos de rango uno conservan su dominio;
para rangos mayores, el resultado presente fija matrices, positividad,
recuperación, compatibilidad y límite con tipos explícitos. Esta
separación sitúa el alcance de cada construcción en el mismo desarrollo
sin confundir crecimiento de dimensión, parametrización de un punto y
clasificación de todas las celdas de una imagen.

\section{Momentos genealógicos, operadores de Jacobi y residuos espectrales}
\label{sec:jacobi-generated}

La medida inducida por el calendario proporciona un enlace entre la
geometría de las celdas y una realización espectral positiva. Este enlace
utiliza toda la partición marcada: cada uno de sus doce intervalos conserva
la misma masa cronológica, mientras su longitud procede de la coordenada
correspondiente del registro. La densidad compensa exactamente esa
longitud. Por consiguiente, la distribución espacial de la medida retiene
la anisotropía del registro aunque la distribución entre ventanas sea
uniforme.

Sean \(K_j>0\), \(Q=\sum_{j=1}^{12}K_j\), \(d_j=K_j/Q\) y
\(t_j=\sum_{i\leq j}d_i\), con \(t_0=0\). La medida de probabilidad
utilizada en esta sección es
\[
 \dd\mu_K(t)=\sum_{j=1}^{12}
 \frac{\mathbf1_{[t_{j-1},t_j]}(t)}{12d_j}\,\dd t.
\]
Los extremos compartidos tienen medida nula. Sus momentos son racionales
para todo registro entero:
\begin{equation}
 m_r=\int_0^1t^r\,\dd\mu_K(t)
 =\frac1{12(r+1)}\sum_{j=1}^{12}
   \frac{t_j^{r+1}-t_{j-1}^{r+1}}{d_j},\qquad r\geq0.
 \label{eq:jacobi-moments}
\end{equation}
La igualdad \(m_0=1\) expresa conservación de la masa total. Para cada
\(n\geq1\), la matriz de Hankel \(H_n=(m_{a+b})_{0\leq a,b<n}\)
es definida positiva: su forma cuadrática es la integral del cuadrado
de un polinomio no nulo sobre un intervalo donde la densidad es positiva.
Esta observación permite pasar de los momentos a un operador sin elegir
un espectro objetivo.

\begin{theorem}[Recurrencia tridiagonal determinada por la partición]
\label{thm:jacobi-recurrence}
Existe una única sucesión de polinomios mónicos \(p_n\), de grado \(n\),
ortogonales para \(\mu_K\). Si
\[
 h_n=\int p_n^2\,\dd\mu_K>0,\qquad
 a_n=\frac{\int t p_n^2\,\dd\mu_K}{h_n},\qquad
 \beta_n=\frac{h_n}{h_{n-1}}>0\quad(n\geq1),
\]
entonces \(p_0=1\), \(p_1=t-a_0\) y
\begin{equation}
 p_{n+1}(t)=(t-a_n)p_n(t)-\beta_np_{n-1}(t),\qquad n\geq1.
 \label{eq:jacobi-recurrence}
\end{equation}
Todos los coeficientes \(a_n,\beta_n\) son racionales para el registro
entero declarado. En la base ortonormal \(\psi_n=p_n/\sqrt{h_n}\),
la compresión de la multiplicación por \(t\) a los polinomios de grado
menor que \(n\) tiene matriz simétrica
\[
 J_n=\begin{pmatrix}
 a_0&\sqrt{\beta_1}&&0\\
 \sqrt{\beta_1}&a_1&\ddots&\\
 &\ddots&\ddots&\sqrt{\beta_{n-1}}\\
 0&&\sqrt{\beta_{n-1}}&a_{n-1}
 \end{pmatrix}.
\]
\end{theorem}
\begin{proof}
La positividad de \(H_n\) garantiza que el sistema lineal que impone
ortogonalidad a un polinomio mónico tiene solución única. Como sus
coeficientes son momentos racionales, la solución es racional. Se
descompone \(tp_n\) en la base mónica hasta grado \(n+1\). Para
\(j\leq n-2\), la simetría del producto interno da
\(\langle tp_n,p_j\rangle=\langle p_n,tp_j\rangle=0\).
Quedan únicamente los términos de grados \(n+1,n,n-1\). Sus
coeficientes son, respectivamente, \(1,a_n,h_n/h_{n-1}\), lo que
demuestra la recurrencia. La normalización por \(\sqrt{h_n}\)
convierte los dos coeficientes adyacentes en \(\sqrt{\beta_n}\).
\end{proof}

Para el registro expresado en la sección~\ref{sec:register}, la primera
posición diagonal y el cuadrado del primer acoplamiento son
\[
 a_0=m_1=\frac{71195}{150312},\qquad
 \beta_1=m_2-m_1^2=\frac{2206226167}{22593697344}.
\]
La primera igualdad calcula el centro de masa de la partición cronológica;
la segunda calcula su dispersión y determina el primer enlace de la
matriz tridiagonal. Ambas cifras se obtienen integrando las longitudes
generadas, y su significado estructural antecede a cualquier evaluación
decimal.

\begin{theorem}[Espectro simple, cuadratura y residuos positivos]
\label{thm:jacobi-residues}
Los autovalores \(\lambda_1<\cdots<\lambda_n\) de \(J_n\) son
simples y pertenecen a \((0,1)\). Existen pesos únicos
\(w_i>0\), de suma uno, tales que
\begin{equation}
 \int f(t)\,\dd\mu_K(t)=\sum_{i=1}^{n}w_i f(\lambda_i)
 \qquad(\deg f\leq2n-1).
 \label{eq:jacobi-quadrature}
\end{equation}
El resolvente observado en la constante unitaria \(e_0\) satisface
\begin{equation}
 g_n(z)=\langle e_0,(zI-J_n)^{-1}e_0\rangle
       =\sum_{i=1}^{n}\frac{w_i}{z-\lambda_i}
       =\cfrac1{z-a_0-\cfrac{\beta_1}{z-a_1-
          \cfrac{\beta_2}{\ddots-\cfrac{\beta_{n-1}}{z-a_{n-1}}}}}.
 \label{eq:jacobi-resolvent}
\end{equation}
En particular, sus polos y residuos se obtienen del registro por las
operaciones de momento, ortogonalización y compresión.
\end{theorem}
\begin{proof}
Para cualquier polinomio no nulo \(p\) de grado menor que \(n\),
\[
 0<\int_0^1 t|p(t)|^2\,\dd\mu_K(t)
   <\int_0^1 |p(t)|^2\,\dd\mu_K(t).
\]

La caracterización variacional coloca el espectro de \(J_n\) en
\((0,1)\). Cada subdiagonal es positiva. Las ecuaciones de un autovector
determinan recursivamente todas sus coordenadas a partir de la primera;
si ésta se anulase, se anularía el vector entero. Cada autoespacio tiene
así dimensión uno. La simetría de la matriz da diagonalización ortogonal
y simplicidad del espectro. Para autovectores unitarios \(v_i\), se toma
\(w_i=|\langle e_0,v_i\rangle|^2>0\). La completitud de la base
propia da \(\sum_iw_i=1\) y la primera expresión del resolvente.

La recurrencia tridiagonal muestra por expansión del determinante que
\(\det(tI-J_n)=p_n(t)\). Para \(\deg f\leq2n-1\), la división
euclídea escribe \(f=qp_n+r\), con \(\deg q,\deg r<n\).
La ortogonalidad anula la integral de \(qp_n\), y su evaluación en los
\(\lambda_i\) también se anula. Para \(r\) de grado menor que
\(n\), las sucesivas multiplicaciones de la constante por \(t\)
permanecen dentro del espacio de polinomios hasta el grado requerido;
por ello \(\int r\,\dd\mu_K=\langle e_0,r(J_n)e_0\rangle\).
La resolución espectral prueba la fórmula de cuadratura. Su unicidad
se deduce de la invertibilidad de la matriz de Vandermonde evaluada
en las raíces distintas. Finalmente, el complemento de Schur de la
primera fila y columna de \(zI-J_n\), iterado sobre las submatrices
terminales, produce la fracción continua.
\end{proof}

\begin{figure}[htbp]
\centering
\begin{tikzpicture}[>=Latex, node distance=8mm and 8mm,
 every node/.style={font=\small}, box/.style={draw,rounded corners=1pt,
 text width=37mm,minimum height=13mm,align=center}]
\node[box] (k) {Registro marcado\\\(K,Q\)};
\node[box,right=of k] (m) {Medida cronológica\\\(\mu_K\)};
\node[box,right=of m] (h) {Momentos positivos\\\(H_n\)};
\node[box,below=of h] (j) {Operador tridiagonal\\\(J_n\)};
\node[box,left=of j] (g) {Resolvente\\\(g_n(z)\)};
\node[box,left=of g] (w) {Polos y residuos\\\(\lambda_i,w_i>0\)};
\draw[->] (k)--(m);\draw[->] (m)--(h);\draw[->] (h)--(j);
\draw[->] (j)--(g);\draw[->] (g)--(w);
\end{tikzpicture}
\caption{Realización espectral de la partición HMT. La geometría de las
celdas determina la medida; la positividad de sus momentos produce el
operador y los residuos positivos de su resolvente.}
\label{fig:jacobi-generated}
\end{figure}

La fórmula del resolvente aporta una lectura espectral de la memoria
espacial. Los pesos positivos describen cuánto observa la coordenada
constante de cada modo; los polos son las frecuencias algebraicas del
operador de multiplicación comprimido. Los momentos de orden creciente
recuperan detalles sucesivos de la medida, y los coeficientes de la
fracción continua conservan esa información mediante una recurrencia
local de tres términos. La apariencia clásica de una matriz de Jacobi
se reconoce después de haber construido sus coeficientes desde las
celdas HMT. La cadena causal queda fijada por
\[
 K\longmapsto(d,t)\longmapsto\mu_K\longmapsto(m_r)_{r\geq0}
 \longmapsto\bigl((a_n)_{n\geq0},(\beta_n)_{n\geq1}\bigr)
 \longmapsto(J_n,g_n).
\]

\begin{proposition}[Compatibilidad de dimensión y de refinamiento]
\label{prop:jacobi-refinement}
Las matrices \(J_n\) son compresiones principales compatibles de la
misma multiplicación acotada por \(t\) en \(L^2(\mu_K)\). Sus
medidas espectrales \(\sum_iw_i\delta_{\lambda_i}\) convergen
débilmente a \(\mu_K\). Si cada subintervalo nonádico se representa
por su punto medio y por su masa \(1/(12\cdot9^r)\), la medida
atómica resultante \(\mu_{K,r}\) aproxima los momentos de
\(\mu_K\). Para un grado fijo \(n\), la construcción desde esos
momentos recupera los coeficientes \(a_j,\beta_j\) de la misma
matriz cuando el refinamiento tiende a profundidad infinita.
\end{proposition}
\begin{proof}
La ortogonalización de la sucesión completa utiliza las mismas fórmulas
al aumentar \(n\), de modo que cada bloque principal permanece fijo.
La cuadratura reproduce exactamente cualquier polinomio cuando
\(2n-1\) supera su grado. La aproximación uniforme de funciones
continuas en \([0,1]\) por polinomios, junto con la masa total uno,
prueba la convergencia débil. Las sumas de puntos medios convergen
en cada uno de los doce intervalos con su densidad constante, y por
tanto aproximan todos los momentos. En dimensión fija, los coeficientes de
ortogonalización son funciones racionales de un número finito de
momentos, con denominadores formados por determinantes positivos de
Gram. La convergencia de esos momentos y la positividad del límite
aseguran la convergencia de los coeficientes.
\end{proof}

\begin{falsifier}
Una longitud nula invalida la densidad declarada y puede destruir el
rango de los momentos. La sustitución de la medida fija por conteos
no normalizados en mallas distintas cambia el operador. Los residuos
de \(g_n\) son pesos espectrales del resolvente, mientras los residuos
de una forma canónica se definen sobre divisores geométricos: una
identificación entre ambas nociones exige especificar la forma y el
mapa correspondiente. Del mismo modo, \(J_n\) es adimensional; su
identificación con un Hamiltoniano físico requiere la representación
y la escala de acción y reloj que utilice ese sector.
\end{falsifier}

\begin{theorem}[Fidelidad de la realización espectral marcada]
\label{thm:jacobi-marked-faithfulness}
El operador infinito de Jacobi \(J\), junto con su vector cíclico
unitario \(e_0\) y la carga \(Q\), determina el registro \(K\).
Dos realizaciones de este tipo son unitariamente equivalentes mediante
una aplicación que conserva el vector marcado y la carga si y sólo
si sus registros coinciden. El espectro como conjunto es, en cambio,
\([0,1]\) para todo registro positivo del dominio.
\end{theorem}
\begin{proof}
Los polinomios son densos en \(L^2(\mu_K)\): primero se aproximan
funciones continuas y después se usa su densidad para la medida finita.
La base ortonormal construida define una aplicación unitaria que lleva
\(e_n\) a \(\psi_n\). Bajo ella, \(J\) es la multiplicación por
\(t\), un operador autoadjunto acotado, y \(e_0\) representa la
función constante uno. Por tanto,
\[
 \langle e_0,J^re_0\rangle=m_r\qquad(r\geq0).
\]
Estos momentos determinan la medida: dos medidas con los mismos
momentos coinciden en polinomios, después en funciones continuas por
aproximación uniforme y finalmente como medidas de Borel.
Su distribución \(F(t)=\mu_K([0,t])\) es continua y estrictamente
creciente. Cada intervalo inicial tiene masa \(1/12\), de modo que
\[
 t_j=F^{-1}(j/12),\qquad
 K_j=Q\bigl(F^{-1}(j/12)-F^{-1}((j-1)/12)\bigr).
\]
Una equivalencia unitaria que conserva \(e_0\) conserva los momentos
y por ello el registro. La recíproca se sigue de la construcción única.

Para el último enunciado, la densidad de \(\mu_K\) es positiva
en todo subintervalo. Cada \(\lambda\in[0,1]\) admite funciones
unitarias concentradas en intervalos cada vez menores alrededor de
\(\lambda\), sobre las que la norma de \((J-\lambda I)f\)
tiende a cero. Así todo el intervalo pertenece al espectro. Fuera de
él, la multiplicación por \(1/(t-\lambda)\) es acotada y define
el resolvente. Esto prueba la igualdad del espectro.
\end{proof}

El teorema distingue la extensión del espectro y la información que
éste conserva respecto de una observación marcada. El mismo intervalo
espectral admite las distintas distribuciones de la colección; la medida
observada desde la unidad recupera las separaciones originales. El
enlace con la representación grassmanniana es ahora reversible:
los menores marcados reconstruyen \(K\), el registro construye
\((J,e_0,Q)\), y sus momentos y cuantiles recuperan el mismo registro.
Cada lector que factoriza por \((K,\xi)\) posee, por consiguiente,
una expresión espectral marcada una vez conservadas las restantes
coordenadas \(\xi\). La realización espectral contiene estructura,
además de posiciones de polos o valores de autovalores.

\begin{figure}[htbp]
\centering
\begin{tikzpicture}[x=1.9cm,y=1.05cm,>=Stealth,font=\small]
\foreach \j/\lab in {0/0,1/1,2/2,3/3,5/12}{
  \node[circle,fill=black,inner sep=1.6pt,label=above:{$a_{\lab}$}] (a\j) at (\j,1.4) {};
  \node[circle,fill=black,inner sep=1.6pt,label=below right:{$b_{\lab}$}] (b\j) at (\j,0) {};
  \node (q\j) at (\j,-1) {$q_{\lab}$};
  \draw[->] (b\j) -- (q\j);
}
\foreach \j/\prev in {1/0,2/1,3/2}{
 \draw[->] (a\prev)--(a\j);\draw[->] (b\prev)--(b\j);
 \draw[->,blue!65!black,thick] (a\j)--node[right]{$d_{\j}$}(b\j);
}
\draw[->,blue!65!black,thick] (a5)--node[right]{$d_{12}$}(b5);
\draw[dashed,->] (a3)--(a5);\draw[dashed,->] (b3)--(b5);
\node at (4,.7) {$\cdots$};
\node[left] at (a0) {$s_2$};\node[left] at (b0) {$s_1$};
\end{tikzpicture}
\caption{Red planar de dos filas. Las aristas horizontales y de salida
tienen peso uno. Cada camino desde la fuente superior desciende una sola
vez. Para dos destinos \(i<j\), las parejas disjuntas descienden después
de \(i\) y antes de \(j\); su peso total es
\(\Delta_{ij}=d_{i+1}+\cdots+d_j\). Los trazos discontinuos conservan
las ocho posiciones intermedias omitidas únicamente en el dibujo.}
\end{figure}

\begin{figure}[htbp]
\centering
\begin{tikzpicture}[x=9cm,y=6cm,font=\small]
\foreach \j/\num in {0/0,1/234,2/777,3/917,4/1646,5/2305,6/3129,7/3750,8/3808,9/4722,10/5516,11/5662,12/6263}{
 \pgfmathsetmacro{\tx}{\num/6263}
 \coordinate (v\j) at (\tx,\tx*\tx);
}
\fill[blue!5] (v0)--(v1)--(v2)--(v3)--(v4)--(v5)--(v6)--(v7)--(v8)--(v9)--(v10)--(v11)--(v12)--cycle;
\foreach \j in {2,...,11}{\draw[blue!45,thin] (v0)--(v\j);}
\draw[blue!80!black,thick] (v0)--(v1)--(v2)--(v3)--(v4)--(v5)--(v6)--(v7)--(v8)--(v9)--(v10)--(v11)--(v12)--cycle;
\foreach \j in {0,...,12}{\fill (v\j) circle[radius=.035cm];}
\node[below left] at (v0) {$v_0$};
\node[below right] at (v3) {$v_3$};
\node[below right] at (v6) {$v_6$};
\node[below right] at (v9) {$v_9$};
\node[above] at (v12) {$v_{12}$};
\end{tikzpicture}
\caption{Polígono racional determinado por los trece parámetros
\(t_j\) de \eqref{eq:polygon-Z}, dibujado con sus coordenadas efectivas
\((t_j,t_j^2)\). La triangulación desde \(v_0\) tiene once celdas.
Cada diagonal interior recibe dos residuos opuestos. Se rotulan cinco
vértices para facilitar la lectura; los trece están representados.}
\end{figure}

\section{Flujo de pesos, geometría de información y dualidad de Legendre}
\label{sec:convex-refinement}

La deformación de los intervalos admite una descripción variacional
exacta. Sean \(d_j=K_j/Q>0\), \(\sum_jd_j=1\), y \(u_j\) las
coordenadas de la dirección \eqref{eq:u-main}. Para \(s\in\RR\),
se definen
\begin{equation}
 Z(s)=\sum_{j=1}^{12}d_j e^{su_j},\qquad
 \psi(s)=\log Z(s),\qquad
 d_j(s)=\frac{d_j e^{su_j}}{Z(s)}.
 \label{eq:convex-partition}
\end{equation}
Todos los argumentos de la exponencial son adimensionales. El parámetro
\(s\) es la coordenada de esta deformación geométrica. La asignación de
una duración física pertenece a un lector temporal declarado.

\begin{theorem}[Convexidad estricta y coordenada dual]
La función \(\psi\) es analítica y estrictamente convexa sobre \(\RR\).
Si \(u_- =\min_j u_j\) y \(u_+=\max_j u_j\), la aplicación
\(v=\psi'(s)\) es un difeomorfismo de \(\RR\) sobre \((u_-,u_+)\).
Para cada \(v\) en ese intervalo existe una única coordenada \(s(v)\), y
\[
 \psi^*(v)=s(v)v-\psi(s(v)),\qquad
 (\psi^*)'(v)=s(v),\qquad
 (\psi^*)''(v)=\frac1{\psi''(s(v))}>0.
\]
\end{theorem}
\begin{proof}
Las sumas finitas permiten diferenciar término a término. Se obtiene
\begin{align}
 \psi'(s)&=\sum_jd_j(s)u_j=:\overline u(s),\notag\\
 \psi''(s)&=\sum_jd_j(s)(u_j-\overline u(s))^2.
 \label{eq:variance}
\end{align}
Los pesos son estrictamente positivos y los \(u_j\) no son todos iguales;
la varianza es por ello estrictamente positiva. Cuando \(s\) tiende a
\(+\infty\), dividir numerador y denominador de \(\psi'\) por
\(e^{su_+}\) muestra que sólo sobreviven los índices con \(u_j=u_+\).
Así \(\psi'(s)\to u_+\). El argumento con \(u_-\) da el límite
opuesto. La monotonía y el teorema de la función inversa proporcionan la
biyectividad suave. La función \(sv-\psi(s)\) tiene su máximo único
en \(s(v)\); diferenciar el valor máximo demuestra las dos identidades
duales.
\end{proof}

La evolución de cada intervalo es
\begin{equation}
 \dot d_j(s)=d_j(s)\bigl(u_j-\overline u(s)\bigr).
 \label{eq:replicator}
\end{equation}
La suma de las derivadas es cero. Por tanto la longitud total permanece
unitaria mientras se redistribuyen las longitudes internas. Esta conservación
de escala global permite atribuir el cambio de las razones dobles a una
deformación de posiciones y no a una dilatación común.

En el simplex abierto de pesos, la métrica
\(g_d(\xi,\zeta)=\sum_j\xi_j\zeta_j/d_j\), sobre vectores tangentes
de suma cero, es positiva. La función lineal
\(F(d)=\sum_jd_ju_j\) tiene como gradiente para esta métrica
precisamente el campo de \eqref{eq:replicator}, pues
\[
 g_d\bigl(d_j(u_j-\overline u),\xi\bigr)
 =\sum_j(u_j-\overline u)\xi_j=\sum_ju_j\xi_j=dF(\xi).
\]
En consecuencia \(dF/ds=\psi''(s)>0\) a lo largo del flujo. El
carácter ascendente corresponde a esta función y a esta métrica concretas.

\begin{theorem}[Invariancia variacional por subdivisión heredada]
Para cada intervalo \(j\), sean \(r_{j,c}>0\),
\(\sum_c r_{j,c}=1\), los pesos relativos de una subdivisión fija.
Asigne a los sucesores
\[
 d_{j,c}=r_{j,c}d_j,\qquad u_{j,c}=u_j.
\]
La función de partición, su logaritmo, todas sus derivadas y la transformada
de Legendre coinciden exactamente con los de la partición madre.
Además, el flujo conmuta con la suma de sucesores:
\[
 \sum_c d_{j,c}(s)=d_j(s),\qquad
 d_{j,c}(s)=r_{j,c}d_j(s).
\]
\end{theorem}
\begin{proof}
Por agrupación de la suma finita,
\[
 Z_{\rm ref}(s)=\sum_{j,c}r_{j,c}d_j e^{su_j}
 =\sum_j d_j e^{su_j}=Z(s).
\]
Todas las conclusiones se deducen de esta identidad y de dividir cada
peso por el mismo \(Z(s)\). El mismo argumento vale después de cualquier
número finito de subdivisiones. Para una subdivisión numerable, la suma de
términos positivos converge al mismo valor mediante su agrupación por predecesores.
\end{proof}

La hipótesis de dirección heredada tiene contenido verificable. Si un
intervalo con dirección \(a\) se divide en dos mitades de direcciones
\(a+b\) y \(a-b\), su contribución cambia de
\(d e^{sa}\) a \(d e^{sa}\cosh(sb)\). Para \(b\ne0\) la igualdad
de funciones de partición falla aunque el peso total en \(s=0\) sea
el mismo. Así, la conservación geométrica y variacional exige conservar
también la información direccional pertinente. La identidad demostrada
se aplica a la subdivisión heredada; su transporte a una actualización
TPK que cambie esas direcciones se decide componiendo el lector de dicha
actualización.

Existe una formulación hamiltoniana regular de la coordenada dual. En el
plano cotangente con coordenadas \((q,p)\), tome
\(H(q,p)=\psi(p)\). Las ecuaciones son
\(\dot q=\psi'(p)\), \(\dot p=0\). Para velocidades
\(\dot q\in(u_-,u_+)\), la transformación de Legendre da
\[
 L(q,\dot q)=\psi^*(\dot q),\qquad
 \frac{\partial^2L}{\partial\dot q^2}>0.
\]
Este sistema realiza variacionalmente el potencial del flujo de intervalos.
Su ecuación \(\dot p=0\) muestra que el tiempo hamiltoniano es distinto
del parámetro \(s\) que recorre la colección \eqref{eq:convex-partition}.
La evaluación de \(H\) y \(L\) se conserva bajo la subdivisión del
teorema, porque \(\psi\) y \(\psi^*\) permanecen idénticas.

Por su parte, cualquier flujo \(\dot d=F(d)\) admite en una coordenatización local
del simplex una elevación cotangente con
\(\mathscr H(d,p)=p\cdot F(d)\). Esta segunda construcción reproduce
el propio flujo \eqref{eq:replicator} en su base, pero su Hessiano respecto
de \(p\) es nulo. Las dos realizaciones responden a preguntas distintas:
una produce una dualidad convexa regular y la otra transporta una dinámica
de base dada. La distinción preserva el significado de Hamiltoniano y
Lagrangiano al comparar esta geometría con los lectores físicos del corpus.

% Desarrollo acotado: lector energético intervalar del registro K.
% Estatuto: formalización añadida sobre K y sobre las identidades de Schur
% ya presentes en el corpus. No se atribuye novedad histórica a Schur.
\section{Formas de Dirichlet y equivalencia variacional del refinamiento}
\label{sec:energia-refinamiento-k}

La representación positiva del registro dodecafásico define una partición
ordenada mediante
\[
 K=(234,543,140,729,659,824,621,58,914,794,146,601),
 \qquad d_j=K_j/6263,
 \qquad t_i=\sum_{j=1}^id_j.
\]
El registro se recibe de la composición
\(x_{\rm term}\mapsto\mathcal R_{12}\mapsto
(b^{90},b^{120},Q_{\rm TPK})\mapsto U_{\rm sgn}\mapsto K\),
posterior a APP--TRIT--TPK y a su estado enriquecido, conforme a los lectores
reconstruidos en \cite{hmtX}. Esta sección aplica
un lector cuadrático a la partición resultante. La comparación entre
profundidades se efectuará por restricciones a extremos heredados, de modo
que cada coeficiente de la energía conserve su origen en una separación
positiva ya producida.

El corpus contiene una energía ponderada de diferencias para el grafo de
las dos memorias y una reducción de Schur de las formas completas. Aquí se
reúnen esas operaciones sobre la red intervalar del registro. La
especialización a conductancias \(1/d_j\), el operador de interpolación
y su compatibilidad con la deformación de las separaciones constituyen
la formalización desarrollada en esta sección.

\subsection{Reconstrucción de las separaciones a partir de la matriz de rigidez}

Sea \(d=(d_1,\ldots,d_m)\), con \(d_j>0\), y sean
\(f=(f_0,\ldots,f_m)\in\mathbb C^{m+1}\) los valores de un observable
en los extremos de la partición. Definimos
\begin{equation}
 \mathcal E_d(f)=\sum_{j=1}^{m}\frac{|f_j-f_{j-1}|^2}{d_j},
 \qquad L_d=D^*\operatorname{diag}(d_1^{-1},\ldots,d_m^{-1})D,
 \label{eq:energia-intervalar}
\end{equation}
donde \(D:\mathbb C^{m+1}\to\mathbb C^m\) es el operador de
incidencia orientada, \((Df)_j=f_j-f_{j-1}\), y el asterisco denota la
adjunción hermítica. Así \(\mathcal E_d(f)=f^*L_df\).
La matriz de rigidez tiene entradas
\[
 (L_d)_{00}=d_1^{-1},\quad (L_d)_{mm}=d_m^{-1},\quad
 (L_d)_{ii}=d_i^{-1}+d_{i+1}^{-1}\quad(1\le i<m),
\]
\[
 (L_d)_{i-1,i}=(L_d)_{i,i-1}=-d_i^{-1},
\]
y las demás son nulas. Es hermítica positiva semidefinida, su núcleo
consiste en las funciones constantes y su forma es positiva definida
en el cociente por ese núcleo. La matriz completa recupera las
separaciones mediante
\begin{equation}
 d_j=-\frac1{(L_d)_{j-1,j}}.
 \label{eq:recuperacion-d-rigidez}
\end{equation}

\begin{proposition}[Equivalencia entre separaciones y lector nodal completo]
\label{prop:energia-bidireccional}
La aplicación \(d\mapsto L_d\) es una biyección entre
\((0,\infty)^m\) y el conjunto de matrices reales simétricas
tridiagonales cuyas entradas inmediatamente adyacentes a la diagonal son
negativas y cuya suma en cada fila es cero. Su inversa es
\eqref{eq:recuperacion-d-rigidez}. La normalización \(\sum_jd_j=1\)
equivale a \(\sum_{j=1}^m-1/(L_d)_{j-1,j}=1\).
\end{proposition}
\begin{proof}
Las fórmulas explícitas de \(L_d\) verifican todas las condiciones.
Recíprocamente, las entradas adyacentes negativas determinan separaciones
positivas únicas por \eqref{eq:recuperacion-d-rigidez}. La suma nula de
cada fila determina sus entradas diagonales como la suma de las
conductancias incidentes. La tridiagonalidad fija en cero las restantes
entradas. Se recupera exactamente la matriz de
\eqref{eq:energia-intervalar}; su semidefinitud positiva sigue entonces
de la factorización \(D^*\operatorname{diag}(1/d_j)D\).
\end{proof}

La derivada discreta ponderada
\(J_j=(f_j-f_{j-1})/d_j\) es el flujo orientado de la arista \(j\).
En los vértices interiores, \((L_df)_i=J_i-J_{i+1}\); en los extremos,
\((L_df)_0=-J_1\) y \((L_df)_m=J_m\). La ecuación interior
\(L_df=0\) expresa que el flujo tiene el mismo valor en cada arista.
Esta definición es la de un observable cuadrático del lector intervalar;
su realización como magnitud física requiere la coordenatización dimensional pertinente.

\subsection{Interpolación armónica, complemento de Schur y reconstrucción del detalle}

Refinemos cada separación como
\(d_j=\sum_{c=1}^{r_j}\delta_{j,c}\), con \(\delta_{j,c}>0\).
Sea \(\widetilde d\) la concatenación de esas separaciones, \(R\) la
restricción a los extremos antiguos y \(H\) la interpolación definida por
\begin{equation}
 (Hf)_{j,c}=(1-\theta_{j,c})f_{j-1}+\theta_{j,c}f_j,
 \qquad
 \theta_{j,c}=\frac{\sum_{a=1}^c\delta_{j,a}}{d_j}.
 \label{eq:interpolacion-armonica-k}
\end{equation}
Los extremos compartidos se cuentan una sola vez. Se cumple \(RH=I\).

\begin{theorem}[Refinamiento isométrico de la energía]
\label{thm:energia-schur-refinamiento}
Para todo vector de valores antiguos \(f\),
\begin{equation}
 H^*L_{\widetilde d}H=L_d,
 \qquad L_{\widetilde d}H=R^*L_d,
 \qquad
 \min_{Rg=f}\mathcal E_{\widetilde d}(g)=\mathcal E_d(f).
 \label{eq:isometria-energia-k}
\end{equation}
El minimizador es único y vale \(Hf\). Todo \(g\) admite la
descomposición única
\begin{equation}
 g=H(Rg)+z,\qquad Rz=0,
 \qquad
 \mathcal E_{\widetilde d}(g)
 =\mathcal E_d(Rg)+\mathcal E_{\widetilde d}(z).
 \label{eq:detalle-energia-k}
\end{equation}
\end{theorem}
\begin{proof}
En una arista antigua escribamos \(h_c=g_{j,c}-g_{j,c-1}\), de modo
que \(\sum_ch_c=f_j-f_{j-1}=a\). Si
\(h_c=\delta_{j,c}a/d_j+r_c\), entonces \(\sum_cr_c=0\) y
\[
 \sum_c\frac{|h_c|^2}{\delta_{j,c}}
 =\frac{|a|^2}{d_j}+\sum_c\frac{|r_c|^2}{\delta_{j,c}}.
\]
La igualdad se obtiene expandiendo; el término cruzado es proporcional a
\(\sum_cr_c\) y se anula. El mínimo se alcanza exactamente cuando
todos los \(r_c\) son cero, condición que determina
\eqref{eq:interpolacion-armonica-k}. Al sumar las aristas se obtienen
el mínimo y la descomposición energética. El flujo de \(Hf\) es
constante dentro de cada arista antigua, por lo que sus componentes
interiores de \(L_{\widetilde d}Hf\) son cero y sus componentes heredadas
coinciden con \(L_df\). Esto prueba la segunda identidad matricial; al
multiplicar por \(H^*\), usando \(RH=I\), se obtiene la primera.
\end{proof}

El término \(\mathcal E_{\widetilde d}(z)\) es no negativo y se anula
exactamente para \(z=0\): un detalle de energía nula sería constante y,
al anularse en los extremos heredados, se anularía en toda la red.
La pareja \((Rg,z)\) conserva el vector refinado completo y permite
recuperarlo por \(g=H(Rg)+z\). La minimización selecciona su
componente armónica; el archivo del detalle \(z\) mantiene la
información que la sola forma reducida no distingue. La conservación
de la energía observada y la conservación del estado refinado quedan
así expresadas por mapas diferentes y explícitos.

Ordenando los vértices como extremos heredados \(E\) e interiores \(I\),
escribimos
\[
 L_{\widetilde d}=\begin{pmatrix}A&B^*\\B&C\end{pmatrix}.
\]
Para un refinamiento con vértices interiores, el bloque \(C\) es
positivo definido: una función nula en los extremos
que tenga energía cero es constante en cada subintervalo y, por tanto,
idénticamente nula. Se obtiene
\begin{equation}
 H=\begin{pmatrix}I\\-C^{-1}B\end{pmatrix},\qquad
 \operatorname{Schur}_{I}(L_{\widetilde d})
 =A-B^*C^{-1}B=L_d.
 \label{eq:schur-intervalar-k}
\end{equation}
Si el refinamiento conserva todos los vértices y no introduce ninguno,
la misma afirmación se reduce a \(H=I\) y \(L_{\widetilde d}=L_d\).

\begin{proposition}[Naturalidad de las prolongaciones y de la eliminación]
\label{prop:energia-naturalidad-sucesiva}
Para tres particiones anidadas, las interpolaciones armónicas y los
complementos de Schur satisfacen
\begin{equation}
 H_{n+2\leftarrow n}
 =H_{n+2\leftarrow n+1}H_{n+1\leftarrow n},
 \qquad
 \operatorname{Schur}_{n+1\to n}
 \bigl(\operatorname{Schur}_{n+2\to n+1}(L_{n+2})\bigr)=L_n.
 \label{eq:energia-naturalidad-sucesiva}
\end{equation}
\end{proposition}
\begin{proof}
La interpolación de valores afines sobre un intervalo sigue siendo la
misma función afín después de cualquier subdivisión de ese intervalo.
Esto prueba la primera identidad. Minimizar primero sobre los vértices
incorporados en el nivel último y después sobre los del nivel intermedio
es minimizar sobre todos ellos, con los mismos extremos heredados fijados.
La identidad variacional de \eqref{eq:isometria-energia-k} y la unicidad
del minimizador prueban la segunda identidad. La restricción del mínimo
se realiza, por tanto, mediante la misma composición que la prolongación
armónica de los datos.
\end{proof}

\subsection{Respuesta de Dirichlet a Neumann y defecto energético de una extensión}

Si sólo se conservan los extremos inicial y final, sea
\(D_0=\sum_{j=1}^md_j\). Para valores \((a,b)\), la solución
armónica es
\[
 f_i=a+(b-a)\frac{\sum_{j=1}^id_j}{D_0},\qquad
 J_j=(b-a)/D_0.
\]
Su operador de Dirichlet a Neumann, que transforma valores de frontera en
flujos salientes, es
\begin{equation}
 \Lambda_d=\frac1{D_0}
 \begin{pmatrix}1&-1\\-1&1\end{pmatrix},\qquad
 \min\mathcal E_d=\frac{|b-a|^2}{D_0}.
 \label{eq:dtn-intervalar-k}
\end{equation}
Cuando se prescribe el valor en todos los vértices heredados, su operador
de respuesta multipuntual es, por el mismo argumento, \(\Lambda=L_d\).

Una extensión aproximada \(g=(f,y)\), donde \(y\) reúne los valores
en los vértices interiores, tiene residual interior
\(r=Bf+Cy\). La identidad exacta
\begin{equation}
 \mathcal E_{\widetilde d}(g)-f^*\Lambda f
 =r^*C^{-1}r\ge0
 \label{eq:residual-dtn-k}
\end{equation}
se deduce de \(y+C^{-1}Bf=C^{-1}r\) y de la descomposición
\[
 g^*L_{\widetilde d}g
 =f^*(A-B^*C^{-1}B)f
 +(y+C^{-1}Bf)^*C(y+C^{-1}Bf).
\]
Así, el valor efectivo y el detalle mantienen una descomposición
hermítica completa. En particular, si una
extensión lineal es \(\widetilde H=\begin{pmatrix}I\\Z\end{pmatrix}\), su matriz de
residuales es \(\mathscr R=B+CZ\) y
\[
 \widetilde H^*L_{\widetilde d}\widetilde H-\Lambda
 =\mathscr R^*C^{-1}\mathscr R\succeq0.
\]
El residual nulo caracteriza la interpolación armónica exacta. Este
residual energético vive en el espacio de vértices interiores; la
cancelación de residuos de formas meromorfas pertenece a otro mapa de
frontera y conserva su definición propia.

\subsection{Conmutación del refinamiento y del flujo de separaciones}

Sea \(u=P_3P_{11}K\) la dirección algebraica ya recuperada de los
lectores del registro \cite{hmtX}. Para el parámetro real \(s\), definamos
\[
 d_j(s)=\frac{K_je^{su_j}}{\sum_\ell K_\ell e^{su_\ell}}.
\]
Fijadas fracciones positivas \(\vartheta_{j,c}\) de suma uno para cada
\(j\), el refinamiento heredado es
\[
 \delta_{j,c}(s)=\vartheta_{j,c}d_j(s),\qquad
 u_{j,c}=u_j.
\]
Producir primero los pesos refinados \(K_j\vartheta_{j,c}\) y
deformarlos da la misma colección: los sucesores tienen idéntico exponencial y
sus coeficientes suman \(K_j\). Para todo \(s\),
\begin{equation}
 \operatorname{Schur}_{I}(L_{\widetilde d(s)})=L_{d(s)},
 \qquad H^*L_{\widetilde d(s)}H=L_{d(s)}.
 \label{eq:schur-flujo-k}
\end{equation}
En esta comparación padre--hijos, \(H\) es constante: sus coeficientes
son las sumas parciales de \(\vartheta_{j,c}\). En cambio, la
interpolación desde los dos extremos globales usa
\(t_i(s)=\sum_{j\le i}d_j(s)\) y cambia con el parámetro. Ambas
afirmaciones corresponden a restricciones diferentes de la misma red.

La normalización \(\sum_jd_j(s)=1\) hace que el operador global de dos
extremos permanezca igual a
\(\left(\begin{smallmatrix}1&-1\\-1&1\end{smallmatrix}\right)\).
La matriz de todos los nodos \(L_{d(s)}\) sí conserva la deformación y
la recupera mediante \eqref{eq:recuperacion-d-rigidez}. La pérdida de
información en el lector de dos extremos queda identificada exactamente:
retiene la longitud total y sus flujos, mientras el lector nodal retiene
todas las separaciones. El archivo de detalles proporciona el complemento
necesario para reconstruir el observable refinado.

\section{Reconstrucción analítica de la energía y memoria del refinamiento}
\label{sec:limite-energetico-k}

Las separaciones positivas construidas a partir del registro dodecafásico
determinan una partición ordenada y, por la sección anterior, una forma
de Dirichlet sobre sus valores nodales. Su refinamiento posee dos
operaciones complementarias: la interpolación armónica conserva la energía
de los datos heredados, mientras cada detalle registra la variación
incorporada entre esos datos. El paso a profundidad ilimitada consiste en
reunir esta colección compatible mediante una norma que controle la suma de
sus detalles. El resultado concierne a la realización analítica
unidimensional de la energía; conserva como antecedente la construcción
HMT del registro y de la estructura discreta conjunta del continuo.

\subsection{Particiones anidadas y espacios de interpolación}

Fijemos la partición inicial
\(P_0=\{0=t_0<t_1<\cdots<t_{12}=1\}\), cuyas longitudes son
\(d_j=t_j-t_{j-1}=K_j/\sum_iK_i>0\). Sea
\((P_n)_{n\geq0}\) una sucesión de particiones finitas anidadas de
\([0,1]\), obtenidas por subdivisión de los intervalos anteriores.
Escribiremos
\[
 |P_n|=\max\{b-a:[a,b]\text{ es un intervalo de }P_n\},
 \qquad \lim_{n\to\infty}|P_n|=0.
\]
La subdivisión reiterada de cada intervalo en un número fijo de sucesores
iguales, al menos dos, satisface esta condición. Para otras reglas, la
contracción de la malla es una hipótesis explícita de este lector.
Las etiquetas de los sucesores y sus datos de acarreo acompañan la
subdivisión y conservan la procedencia del refinamiento.

Sea \(V_n\) el espacio complejo de funciones continuas sobre \([0,1]\)
que son afines en cada intervalo de \(P_n\). La energía y la norma
energética se definen por
\begin{equation}
 \mathcal E_n(f)
 =\sum_{[a,b]\in P_n}\frac{|f(b)-f(a)|^2}{b-a}
 =\int_0^1|f'(x)|^2\,dx,
 \qquad
 \|f\|_{\mathcal E}^2=|f(0)|^2+\mathcal E_n(f).
 \label{eq:energia-afines-k}
\end{equation}
La igualdad integral resulta de que la derivada de una función afín es
constante en cada intervalo. La energía coincide así exactamente con
la forma nodal ya construida. La inclusión \(V_n\subset V_{n+1}\)
interpreta la misma función afín en una partición más fina; conserva
tanto \(f(0)\) como la energía y es isométrica para
\(\|\cdot\|_{\mathcal E}\).

El valor inicial tiene una función precisa: fija la componente constante
que la forma de Dirichlet identifica en su núcleo. La energía controla
las diferencias, y el par formado por el valor inicial y esas diferencias
determina el observable. En el espacio de Sobolev
\(H^1([0,1];\mathbb C)\), cuyos elementos tienen representante
absolutamente continuo con derivada débil en \(L^2\), esta norma es
equivalente a la norma usual. En efecto,
\(f(x)=f(0)+\int_0^xf'(t)\,dt\) controla la norma de \(f\) a
partir de \(|f(0)|\) y \(\|f'\|_{L^2}\); la integración de la
misma identidad controla a su vez \(|f(0)|\) mediante
\(\|f\|_{L^2}+\|f'\|_{L^2}\).

\subsection{Reconstrucción de una función a partir de sus valores compatibles}

\begin{theorem}[Reconstrucción energética y completación de las interpolaciones]
\label{thm:limite-energia-h1-k}
Sea \(f_n\in V_n\) una colección compatible por restricción a los
nodos heredados:
\[
 f_{n+1}(a)=f_n(a)
 \quad\text{para todo nodo }a\text{ de }P_n.
\]
Si \(\sup_n\mathcal E_n(f_n)<\infty\), existe una única función
\(f\in H^1([0,1];\mathbb C)\), considerada en su representante
continuo, cuyos valores en todos esos nodos son los prescritos. Además,
\(f_n\to f\) uniformemente y en \(H^1\), y
\begin{equation}
 \int_0^1|f'(x)|^2\,dx
 =\lim_{n\to\infty}\mathcal E_n(f_n).
 \label{eq:limite-energia-h1-k}
\end{equation}
Recíprocamente, la interpolación nodal de cualquier elemento de \(H^1\)
produce una colección de esta clase. Por tanto, la completación de
\(\bigcup_n V_n\) en la norma energética es \(H^1([0,1];\mathbb C)\).
\end{theorem}

\begin{proof}
Pongamos \(g_n=f_n'\in L^2([0,1];\mathbb C)\). Para
\(m\geq n\), la compatibilidad de los valores en los extremos de cada
intervalo \([a,b]\) de \(P_n\) proporciona
\[
 \frac1{b-a}\int_a^bg_m(x)\,dx
 =\frac{f_m(b)-f_m(a)}{b-a}
 =\frac{f_n(b)-f_n(a)}{b-a}
 =g_n\big|_{(a,b)}.
\]
Sea \(W_n\) el espacio de funciones de \(L^2\) constantes en cada
intervalo de \(P_n\), y sea \(Q_n:L^2\to W_n\) la proyección
ortogonal dada por los promedios sobre dichos intervalos. La identidad
anterior afirma \(Q_ng_m=g_n\). En consecuencia,
\(g_m-g_n\) es ortogonal a \(g_n\) y
\begin{equation}
 \|g_m-g_n\|_{L^2}^2
 =\mathcal E_m(f_m)-\mathcal E_n(f_n).
 \label{eq:pitagoras-gradientes-k}
\end{equation}
Las energías son crecientes y acotadas. La identidad de Pitágoras
prueba que \((g_n)\) es una sucesión de Cauchy en \(L^2\), con
límite \(g\). Definamos
\[
 f(x)=f_0(0)+\int_0^xg(t)\,dt.
\]
Esta función pertenece a \(H^1\), tiene derivada débil \(g\) y
satisface, por la desigualdad de Cauchy--Schwarz,
\begin{equation}
 \sup_{x\in[0,1]}|f_n(x)-f(x)|
 \leq\|g_n-g\|_{L^2}\longrightarrow0.
 \label{eq:convergencia-uniforme-energia-k}
\end{equation}
La convergencia uniforme conserva cada valor nodal prescrito; junto con
la convergencia de las derivadas en \(L^2\), da la convergencia en
\(H^1\). La convergencia de las normas de las derivadas prueba
\eqref{eq:limite-energia-h1-k}. La unión de los nodos es densa porque
\(|P_n|\to0\); dos representantes continuos con los mismos valores
en esa unión coinciden. Se obtiene la unicidad.

Para la recíproca, tomemos \(f\in H^1\) en su representante
continuo y sea \(f_n\) su interpolación en los nodos de \(P_n\).
Su derivada es \(g_n=Q_nf'\). Las proyecciones \(Q_n\) son
contractivas en \(L^2\). Para una función continua \(v\), el
error \(|Q_nv-v|\) está acotado por su módulo de continuidad
evaluado en \(|P_n|\), que tiende a cero. La densidad de las
funciones continuas en \(L^2\) y la contractividad implican
\(Q_ng\to g\) para todo \(g\in L^2\): dado un aproximante
continuo \(v\), se emplea
\[
 \|Q_ng-g\|_{L^2}
 \leq 2\|g-v\|_{L^2}+\|Q_nv-v\|_{L^2}.
\]
Aplicando esta convergencia a \(g=f'\), se obtienen la convergencia
de las derivadas y, mediante \eqref{eq:convergencia-uniforme-energia-k},
las de las funciones. La contractividad acota sus energías por
\(\|f'\|_{L^2}^2\), y la convergencia de las normas prueba de
nuevo la identidad del límite. Toda función de \(H^1\) es, por
tanto, límite energético de elementos de \(\bigcup_nV_n\).
La equivalencia de normas y la completitud de \(H^1\) identifican
la completación enunciada.
\end{proof}

El argumento describe concretamente qué se conserva al pasar de la red
a la función. Cada derivada finita registra el promedio de todas las
derivadas posteriores en su intervalo; al refinar, se incorporan
componentes ortogonales que mantienen ese promedio. La energía acotada
controla la suma total de esas incorporaciones. El valor continuo se
recupera entonces integrando la derivada límite y conservando el valor
inicial. La realización analítica reúne, por esta vía explícita, la
información compatible de los niveles finitos.

\subsection{Descomposición ortogonal de la memoria por escalas}

Sea \(H_n\) la interpolación armónica de los valores de \(P_n\)
sobre los nodos de \(P_{n+1}\), interpretada como función de
\(V_{n+1}\), y definamos el detalle
\[
 z_{n+1}=f_{n+1}-H_nf_n.
\]
Por \eqref{eq:detalle-energia-k}, el detalle se anula en los nodos
heredados. Su derivada tiene promedio cero en cada intervalo antiguo
y es ortogonal a \(W_n\). Para niveles distintos, la derivada de
un detalle anterior pertenece a ese espacio de funciones constantes;
las derivadas de los detalles son, por tanto, ortogonales entre sí.
La misma ortogonalidad vale frente a \(f_0'\).

\begin{proposition}[Identidad de energía y recuperación de los detalles]
\label{prop:memoria-energetica-escalas-k}
Para toda profundidad finita \(N\),
\begin{equation}
 \mathcal E_N(f_N)
 =\mathcal E_0(f_0)
 +\sum_{n=0}^{N-1}\mathcal E_{n+1}(z_{n+1}).
 \label{eq:energia-telescopica-detalles-k}
\end{equation}
Si la colección satisface las hipótesis del
Teorema~\ref{thm:limite-energia-h1-k}, la serie de la derecha
converge y su suma es \(\int_0^1|f'|^2\). Los datos \(f_0\)
y todos los detalles determinan la colección compatible completa y su
realización analítica.
\end{proposition}
\begin{proof}
La identidad finita resulta de sumar
\(\mathcal E_{n+1}(f_{n+1})
=\mathcal E_n(f_n)+\mathcal E_{n+1}(z_{n+1})\), que es la
descomposición ortogonal de la sección anterior. Los sumandos son
no negativos y las sumas parciales convergen por
\eqref{eq:limite-energia-h1-k}. La reconstrucción finita se efectúa
recursivamente mediante \(f_{n+1}=H_nf_n+z_{n+1}\); el teorema
precedente determina después la función límite. Recíprocamente, una
función de \(H^1\) determina sus interpolaciones y, por diferencia,
todos estos detalles. Cada flecha conserva así sus datos propios.
\end{proof}

La eliminación variacional de nodos interiores selecciona el detalle
nulo en el nivel eliminado. La conservación de los detalles permite
recuperar cualquier configuración compatible de energía finita. La
respuesta efectiva y la configuración completa poseen, de esta manera,
una relación precisa: la primera registra el mínimo sobre las variables
interiores; la segunda añade las componentes ortogonales y su energía.

\subsection{Transporte de la realización analítica bajo el flujo positivo}

La dirección \(u=P_3P_{11}K\), recuperada antes de la evaluación
energética \cite{hmtX}, determina la colección de separaciones de
\eqref{eq:schur-flujo-k}. Su parámetro \(s\in\mathbb R\)
describe esta deformación analítica. Escribamos
\[
 d_j(s)=\frac{K_je^{su_j}}{\sum_\ell K_\ell e^{su_\ell}},
 \qquad
 t_i(s)=\sum_{j=1}^id_j(s),
 \qquad
 a_j(s)=\frac{d_j(s)}{d_j(0)}>0.
\]
Sea \(F_s:[0,1]\to[0,1]\) la aplicación que envía cada
\(t_i(0)\) a \(t_i(s)\) y es afín en los intervalos iniciales.
Es un homeomorfismo creciente, de pendiente \(a_j(s)\) en el
intervalo \([t_{j-1}(0),t_j(0)]\). Su inversa es también afín
por intervalos y ambas aplicaciones son Lipschitz.

Supongamos que cada subdivisión hereda las fracciones de longitud y la
dirección de su intervalo predecesor, como en la sección anterior. Todos
los nodos refinados se transportan entonces mediante el mismo \(F_s\).
Denotemos sus particiones por \(P_n(s)=F_s(P_n(0))\) y sus
espacios de interpolación por \(V_n(s)\). La aplicación
\[
 U_s f=f\circ F_s^{-1}
\]
transporta \(V_n(0)\) a \(V_n(s)\), preserva los valores
nodales correspondientes y prolonga esta acción a \(H^1\).

\begin{theorem}[Equivalencia energética y conmutación con el refinamiento]
\label{thm:transporte-h1-flujo-k}
Para todo \(f\in H^1([0,1];\mathbb C)\),
\begin{equation}
 \int_0^1|(U_sf)'(y)|^2\,dy
 =\sum_{j=1}^{12}\frac1{a_j(s)}
 \int_{t_{j-1}(0)}^{t_j(0)}|f'(x)|^2\,dx.
 \label{eq:transporte-energia-h1-k}
\end{equation}
Para \(s\) en cualquier intervalo compacto, las normas energéticas
antes y después del transporte son uniformemente equivalentes y
\(|P_n(s)|\to0\) uniformemente en ese compacto. La interpolación
armónica y la restricción a nodos heredados conmutan con \(U_s\);
la reconstrucción y la minimización de los detalles se transportan
tanto a profundidad finita como en la realización \(H^1\).
\end{theorem}

\begin{proof}
En el intervalo inicial \(j\), la regla de la cadena y el cambio
de variable \(y=F_s(x)\) proporcionan
\((U_sf)'(F_s(x))=f'(x)/a_j(s)\) y \(dy=a_j(s)\,dx\).
La integración y la suma sobre los intervalos prueban
\eqref{eq:transporte-energia-h1-k}. La fórmula vale para los
representantes absolutamente continuos, y su lado derecho es finito.

Fijado un intervalo compacto de parámetros, la continuidad y positividad
de los doce factores dan constantes \(0<a_-\leq a_+<\infty\)
tales que \(a_-\leq a_j(s)\leq a_+\) para todo \(j\) y
todo \(s\) de ese compacto. Por tanto,
\[
 \frac1{a_+}\mathcal E(f)
 \leq\mathcal E(U_sf)
 \leq\frac1{a_-}\mathcal E(f),
 \qquad
 |P_n(s)|\leq a_+|P_n(0)|.
\]
Como \((U_sf)(0)=f(0)\), se obtiene además
\[
 \min\{1,a_+^{-1}\}\|f\|_{\mathcal E}^2
 \leq\|U_sf\|_{\mathcal E}^2
 \leq\max\{1,a_-^{-1}\}\|f\|_{\mathcal E}^2.
\]
Estas cotas prueban la equivalencia uniforme y la conservación de la
condición de malla decreciente.

En cada arista predecesor, los coeficientes de interpolación son las sumas
parciales de las fracciones heredadas, independientes de \(s\).
Si \(H_n(s)\) denota su interpolación, se tiene
\(U_sH_n(0)=H_n(s)U_s\) en los espacios correspondientes.
La restricción conmuta porque los nodos se transportan con sus valores.
La descomposición en función armónica y detalle se conserva, y la
identidad de Schur de \eqref{eq:schur-flujo-k} preserva su
caracterización variacional. Finalmente, la continuidad acotada de
\(U_s\) permite pasar al límite \(H^1\). En particular, una
colección reconstruida antes del transporte converge después a \(U_sf\),
con los mismos valores nodales transportados.
\end{proof}

La ley \eqref{eq:transporte-energia-h1-k} identifica los factores
métricos exactos de la deformación: el cambio de longitud de cada
intervalo modifica inversamente su contribución energética. La
equivalencia de normas conserva las configuraciones de energía finita,
y el archivo de detalles conserva su reconstrucción. Se obtiene así una
realización analítica completa de este lector del registro positivo.
Su compatibilidad con las formas canónicas se articula sobre las
subdivisiones de un dominio fijo: la suma de formas conserva el objeto
meromorfo, mientras la reducción variacional y los detalles conservan
la respuesta energética y el observable refinado. La identificación
con formas amplituhedrales de rango superior conserva las condiciones
específicas de cobertura, grado y correspondencia de residuos
establecidas para aquella realización.

\section{Transportes nonádicos y conservación de la memoria}
\label{sec:nonadic}

La reconstrucción de la sección~\ref{sec:memory} proporciona un enlace
efectivo entre la memoria del transporte y la geometría del registro
dodecafásico. El registro generado por APP--TRIT--TPK determina sus
lecturas de incidencia; las dos compresiones conservan, en componentes
separadas, la representación terminal y sus complementos; la reunión de
las memorias con la carga restituye las doce coordenadas. La dirección
geométrica y el proyector dimensional se obtienen después de esa
restitución. El desarrollo siguiente establece el balance operatorio
que sostiene esta cadena y su comportamiento bajo refinamiento,
realización espectral y cambio de métrica
\cite{HMT-VIII,hmtX}.

\subsection{Calendario de nueve fases y descomposición ortogonal}

Sea \(\mathcal H\) un espacio de Hilbert complejo y sea
\(C\in\mathcal B(\mathcal H)\). La suma ortogonal de nueve copias,
\(\mathcal H^{\oplus9}\), conserva por separado las nueve posiciones
del calendario. Definimos
\begin{equation}
 \begin{split}
 S_C(v_0,\ldots,v_8)&=(Cv_8,v_0,\ldots,v_7),\\
 J:\mathcal H\longrightarrow\mathcal H^{\oplus9},
 \qquad Jv&=\frac13(v,\ldots,v).
 \end{split}
 \label{nt:calendar}
\end{equation}
La normalización \(1/3=9^{-1/2}\) da \(J^*J=I\). El operador
\(P=JJ^*\) es la proyección ortogonal sobre la sección uniforme;
su acción sustituye las nueve componentes por su media. Cada
componente atraviesa una vez la unión que contiene \(C\) durante una
vuelta, de modo que
\begin{equation}
 S_C^9=I_9\otimes C.
 \label{nt:full-return}
\end{equation}
El retorno de la posición del calendario realiza así una acción de
\(C\) sobre la fibra de memoria.

\begin{definition}
La compresión uniforme del transporte y su complemento son
\begin{equation}
 T_C=J^*S_CJ=\frac{8I+C}{9},
 \qquad
 \eta_C=(I-P)S_CJ.
 \label{nt:compression}
\end{equation}
La primera aplicación tiene codominio \(\mathcal H\); la segunda
tiene codominio \(J(\mathcal H)^\perp\subset\mathcal H^{\oplus9}\).
\end{definition}

Para \(z=(C-I)v\), el complemento se escribe
\begin{equation}
 \eta_Cv=\frac1{27}(8z,-z,-z,-z,-z,-z,-z,-z,-z).
 \label{nt:complement}
\end{equation}
La suma de sus componentes es cero. Los coeficientes de la
descomposición proceden, por tanto, del promedio de ocho
continuaciones idénticas y una continuación transportada.

\begin{proposition}[Balance de la compresión y de la unión de memoria]
Para todo \(C\in\mathcal B(\mathcal H)\),
\begin{align}
 \eta_C^*\eta_C
   &=\frac8{81}(I-C)^*(I-C),\label{nt:eta-gram}\\
 I-T_C^*T_C
   &=\eta_C^*\eta_C+\frac19(I-C^*C).\label{nt:general-balance}
\end{align}
Si \(C^*C=I\), la aplicación
\(v\mapsto(T_Cv,\eta_Cv)\) es isométrica.
\end{proposition}

\begin{proof}
La expresión~\eqref{nt:complement} da
\[
 \|\eta_Cv\|^2
 =\frac{8^2+8}{27^2}\|(C-I)v\|^2
 =\frac8{81}\|(C-I)v\|^2.
\]
La polarización prueba~\eqref{nt:eta-gram}. Por otra parte,
\[
 T_C^*T_C
 =\frac{64I+8C+8C^*+C^*C}{81}.
\]
Sumando \eqref{nt:eta-gram} se obtiene
\[
 T_C^*T_C+\eta_C^*\eta_C=\frac{8I+C^*C}{9}.
\]
Restar este resultado de \(I\) da~\eqref{nt:general-balance}.
Cuando \(C^*C=I\), la suma de los dos productos de Gram es la
identidad, lo que conserva todos los productos interiores.
\end{proof}

El término \(\eta_C^*\eta_C\) mide la dispersión entre las nueve
componentes. El término \((I-C^*C)/9\) registra, por separado,
la variación de la norma en la unión de memoria. La realización
unitaria sitúa toda la variación visible en el primer término y
conserva la norma en la suma ortogonal completa.

\begin{proposition}[Conservación de un observable]
Sean \(C\) unitario y \(A=A^*\in\mathcal B(\mathcal H)\) tales que
\(C^*AC=A\). Entonces
\begin{equation}
 A=T_C^*AT_C+
   \eta_C^*(I_9\otimes A)\eta_C
  =T_C^*AT_C+\frac8{81}(I-C)^*A(I-C).
 \label{nt:observable-balance}
\end{equation}
Para \(A\succeq0\), ambos términos de la suma son positivos.
\end{proposition}

\begin{proof}
La fórmula del complemento, aplicada al producto sesquilineal
definido por \(A\), prueba la segunda igualdad. La expansión de
los dos términos da
\[
 \frac1{81}\bigl(
 64A+8AC+8C^*A+C^*AC
 +8A-8AC-8C^*A+8C^*AC\bigr)
 =\frac{8A+C^*AC}{9}=A.
\]
La positividad se conserva bajo una aplicación de la forma
\(B\mapsto L^*BL\).
\end{proof}

\subsection{Iteración de la compresión y conservación del sector fijo}

Fijamos en esta subsección un operador unitario \(C\) y escribimos
\(T=T_C\), \(\eta=\eta_C\). El transporte \(S_C^N\) mantiene todas
las fases durante \(N\) pasos. La potencia \(T^N\), en cambio,
aplica sucesivamente la compresión uniforme; en cada etapa
\(\eta T^jv\) conserva el complemento de la representación anterior.
La distinción queda expresada por los siguientes operadores.

\begin{proposition}[Registro isométrico a horizonte finito]
Para todo entero \(N\geq1\), la aplicación
\begin{equation}
 V_Nv=(T^Nv,\eta v,\eta Tv,\ldots,\eta T^{N-1}v)
 \label{nt:finite-record}
\end{equation}
de \(\mathcal H\) a
\(\mathcal H\oplus(\mathcal H^{\oplus9})^{\oplus N}\)
es isométrica. Si \(A\) satisface las hipótesis de
\eqref{nt:observable-balance}, entonces
\begin{equation}
 A=T^{*N}AT^N+
 \sum_{j=0}^{N-1}T^{*j}\eta^*(I_9\otimes A)\eta T^j.
 \label{nt:telescopy}
\end{equation}
\end{proposition}

\begin{proof}
Multiplicar~\eqref{nt:observable-balance} por \(T^{*j}\) y
\(T^j\) da
\[
 T^{*j}AT^j-T^{*(j+1)}AT^{j+1}
 =T^{*j}\eta^*(I_9\otimes A)\eta T^j.
\]
La suma telescópica prueba~\eqref{nt:telescopy}. Para \(A=I\),
el miembro derecho es el producto de Gram de \(V_N\).
El paso de \(N\) a \(N+1\) sustituye únicamente \(T^Nv\) por
\((T^{N+1}v,\eta T^Nv)\), conservando todas las componentes
de memoria previamente producidas.
\end{proof}

\begin{theorem}[Límite de la compresión y memoria ilimitada]
Sea \(P_{\mathrm{fix}}\) la proyección ortogonal sobre
\(\ker(I-C)\). Entonces
\begin{align}
 \operatorname*{s\!-\!lim}_{N\to\infty}T^N
   &=P_{\mathrm{fix}},\label{nt:strong-limit}\\
 I&=P_{\mathrm{fix}}+
 \sum_{j=0}^{\infty}T^{*j}\eta^*\eta T^j.\label{nt:infinite-record}
\end{align}
La serie converge fuertemente. Para todo observable acotado
autoadjunto que satisfaga \(C^*AC=A\), se conserva además
\begin{equation}
 A=P_{\mathrm{fix}}AP_{\mathrm{fix}}+
 \sum_{j=0}^{\infty}
 T^{*j}\eta^*(I_9\otimes A)\eta T^j.
 \label{nt:infinite-observable}
\end{equation}
\end{theorem}

\begin{proof}
El teorema espectral del operador unitario realiza \(C\) como
la función coordenada \(z\) sobre la circunferencia unidad.
La compresión corresponde a \(t(z)=(8+z)/9\), y
\begin{equation}
 1-|t(z)|^2=\frac8{81}|1-z|^2,
 \qquad |z|=1.
 \label{nt:circular-defect}
\end{equation}
Así, \(t(z)^N\) converge a cero en cada \(z\ne1\) y conserva
el valor uno en \(z=1\). La cota \(|t(z)^N|\leq1\) permite
aplicar convergencia dominada a cada medida espectral vectorial.
Se obtiene~\eqref{nt:strong-limit} y, por el mismo argumento,
\(T^{*N}\to P_{\mathrm{fix}}\) fuertemente. La acotación de
\(A\), junto con \(\|T^N\|\leq1\), da
\(T^{*N}AT^N\to P_{\mathrm{fix}}AP_{\mathrm{fix}}\).
El límite de~\eqref{nt:telescopy} prueba las dos identidades.
\end{proof}

El registro
\[
 V_\infty v=(P_{\mathrm{fix}}v,\eta v,\eta Tv,\eta T^2v,\ldots)
\]
es, por consiguiente, una isometría. El sector fijo conserva la
parte que persiste en la representación terminal; las demás
componentes permanecen en la memoria de los complementos.

Para el desplazamiento bilateral
\(C e_n=e_{n+1}\) en \(\ell^2(\mathbb Z)\), una sucesión fija
es constante. La condición de cuadrado-sumabilidad fuerza su
anulación, por lo que \(P_{\mathrm{fix}}=0\). Toda la norma
queda entonces representada en los complementos. El espectro de
este desplazamiento es la circunferencia completa, y el máximo
espectral de \(|t(z)|^N\) es uno; de ahí
\(\|T^N\|=1\) para cada \(N\). La convergencia fuerte conserva
este comportamiento no uniforme.

Para una permutación cíclica de \(d\) posiciones, el sector fijo
es la recta de los vectores uniformes. Si la permutación tiene
varias órbitas, cada órbita aporta su propia dirección uniforme.
Esta dimensión determina la parte terminal que debe conservarse.

\begin{proposition}[Compresión de un transporte de orden tres]
Sea \(C\) unitario y \(C^3=I\). Entonces
\begin{equation}
 P_{\mathrm{fix}}=\frac{I+C+C^2}{3},
 \qquad
 T_C^*T_C=P_{\mathrm{fix}}+
 \frac{19}{27}(I-P_{\mathrm{fix}}).
 \label{nt:order-three}
\end{equation}
\end{proposition}

\begin{proof}
El promedio de las tres potencias es el proyector ortogonal
sobre los vectores fijos. Como \(C^*=C^2\),
\[
 T_C^*T_C=\frac{65I+8(C+C^2)}{81}
 =\frac{57I+24P_{\mathrm{fix}}}{81},
\]
que coincide con~\eqref{nt:order-three}.
\end{proof}

En la realización de doce posiciones de
\(\mathbb R^{12}\), \((Sx)_i=x_{i+1}\) con índices módulo
doce, \(S^3\) tiene orden cuatro y \(S^4\) tiene orden tres.
Las cuatro órbitas de \(S^4\) dan
\(\operatorname{rank}P_{\mathrm{fix}}=4\).
La realización excepcional \(g_c\) del artículo~X satisface
\(g_c^3=I\) y \(\ker(I-g_c)=0\) en su propio portador;
en ella~\eqref{nt:order-three} se reduce a
\(T_{g_c}^*T_{g_c}=19I/27\)
\cite{hmtX}. El coeficiente y el sector fijo se transportan
conjuntamente con el dominio.

La contracción radial \(M=qC\), \(0<q<1\), posee otro registro
exacto:
\begin{equation}
 L_Nv=\bigl(\sqrt{1-q^2}\,v,
 \sqrt{1-q^2}\,qCv,\ldots,
 \sqrt{1-q^2}\,q^{N-1}C^{N-1}v,q^NC^Nv\bigr).
 \label{nt:radial-record}
\end{equation}
Su producto de Gram es la identidad porque
\((1-q^2)\sum_{j=0}^{N-1}q^{2j}+q^{2N}=1\).
El terminal tiene norma \(q^N\|v\|\), y el registro infinito
conserva el producto interior. El factor \(q=5/9\) del modo
diferencial central del artículo~VIII se aplica a una vuelta
completa; su reparto uniforme entre nueve pasos utiliza
\(q^{1/9}S_C\), cuya novena potencia es
\(q(I_9\otimes C)\) \cite{HMT-VIII}.

\subsection{Reducción del registro y conservación de la traza}

La dirección fija del complemento en las nueve componentes
permite una realización mínima en dos sumandos. Definamos
\[
 D_C=\frac{\sqrt8}{9}(I-C).
\]
Las identidades anteriores dan
\(T_C^*T_C+D_C^*D_C=I\) para \(C\) unitario.
El cambio de signo respecto de la expresión de
\(\eta_C\) corresponde a una orientación de la componente
complementaria y conserva su producto de Gram.

\begin{proposition}[Canal positivo de la compresión con memoria]
Para \(C\) unitario, la aplicación
\begin{equation}
 \mathcal Q_C(X)=T_CXT_C^*+D_CXD_C^*
              =\frac{8X+CXC^*}{9}
 \label{nt:channel}
\end{equation}
es completamente positiva en \(\mathcal B(\mathcal H)\) y
conserva la identidad. Su restricción a los operadores de
clase traza conserva la traza.
\end{proposition}

\begin{proof}
Al expandir los dos términos de la primera expresión,
los productos \(CX\) y \(XC^*\) se cancelan. Los coeficientes
restantes son \(72/81=8/9\) para \(X\) y \(9/81=1/9\)
para \(CXC^*\), lo que prueba~\eqref{nt:channel}.
Para cualquier entero \(m\geq1\), la ampliación a matrices
de orden \(m\) tiene la forma
\[
 X\longmapsto
 (I_m\otimes T_C)X(I_m\otimes T_C)^*
 +(I_m\otimes D_C)X(I_m\otimes D_C)^*.
\]
Cada sumando conserva positividad; por ello la aplicación es
completamente positiva. La segunda expresión de
\eqref{nt:channel} da \(\mathcal Q_C(I)=I\).
Si \(X\) es de clase traza, la invariancia de la traza bajo
conjugación unitaria implica
\(\operatorname{Tr}\mathcal Q_C(X)=\operatorname{Tr}X\).
\end{proof}

La aplicación \eqref{nt:channel} se obtiene igualmente
preparando la isometría
\(v\mapsto T_Cv\otimes e_0+D_Cv\otimes e_1\) y efectuando
la traza parcial sobre \(\mathbb C^2\).
La conservación del vector completo pertenece a la isometría;
la conservación de la traza del estado reducido pertenece a
\(\mathcal Q_C\). Ambas afirmaciones proceden del mismo
balance ortogonal.

En la prolongación nonádica de álgebras finitas aparece una
conservación complementaria. Para
\(\mathcal A_d=M_{9^d}(\mathbb C)\), con
\(\tau_d(a)=9^{-d}\operatorname{Tr}a\), las aplicaciones
\[
 \iota_d(a)=a\otimes I_9,\qquad
 \jmath_{d,j}(a)=a\otimes p_j,
 \quad p_j=|e_j\rangle\langle e_j|
\]
satisfacen
\begin{align}
 \tau_{d+1}(\iota_d(a))&=\tau_d(a),
 &\iota_d(I)&=I,\label{nt:unital-trace}\\
 \tau_{d+1}(\jmath_{d,j}(a))&=\frac{\tau_d(a)}9,
 &\jmath_{d,j}(I)&=I\otimes p_j.\label{nt:corner-trace}
\end{align}
La demostración usa
\(\operatorname{Tr}(a\otimes b)=
\operatorname{Tr}(a)\operatorname{Tr}(b)\),
\(\operatorname{Tr}I_9=9\) y \(\operatorname{Tr}p_j=1\).
La primera inclusión conserva toda la fibra y la unidad;
la segunda conserva una continuación marcada y su soporte.

Esta distinción también se manifiesta en el árbol de
historias enriquecidas. Sean \(\mathcal H_n\) los prefijos
supervivientes de longitud \(n\), y sea \(d(h)\geq1\) el número de
emisiones salientes conservadas de \(h\). Cada emisión con multiplicidad
se representa por sucesores marcados distintos; la suma siguiente recorre
ese conjunto de \(d(h)\) sucesores.
La ley de pesos
\[
 \mu_{n+1}(h\epsilon)=\mu_n(h)\,p_h(\epsilon),
 \qquad p_h(\epsilon)=d(h)^{-1}
\]
conserva
\(\sum_{\epsilon}\mu_{n+1}(h\epsilon)=\mu_n(h)\).
Por tanto, en
\(\mathscr L_n=L^2(\mathcal H_n,\mu_n)\), el pullback
\[
 R_nf(h\epsilon)=f(h)
\]
es isométrico y conserva la función unidad. Su adjunto es
\[
 E_ng(h)=\sum_{\epsilon}p_h(\epsilon)g(h\epsilon),
 \qquad E_nR_n=I,\quad E_n1=1.
\]
En efecto, agrupar la suma de \(\|R_nf\|^2\) por predecesor
produce \(\sum_h\mu_n(h)|f(h)|^2\); el mismo agrupamiento
en el producto interior determina \(E_n\).
La memoria y las multiplicidades quedan presentes en el
dominio de las historias. Estos mapas constituyen la
realización funcional del refinamiento descrito en el
artículo~X \cite{hmtX}.

\subsection{Coordenatización espectral y reciprocidad del radio}

El transporte anterior admite realizaciones espectrales
con dominios propios. La coordenatización de Möbius de VIII actúa,
en coordenadas finitas, por
\begin{equation}
 \tau(z)=\frac{z-1}{z},
 \qquad z\in\mathbb C\setminus\{0\}.
 \label{nt:mobius}
\end{equation}
La inversa es \(z=(1-\tau)^{-1}\), para \(\tau\ne1\).
La extensión a la esfera de Riemann envía \(0\) a
\(\infty\) e \(\infty\) a \(1\).

\begin{proposition}[Espejo espectral y radio recíproco]
Sea \(z^\sharp=1-\overline z\). Para
\(z\notin\{0,1\}\),
\begin{equation}
 \tau(z^\sharp)=\frac1{\overline{\tau(z)}}.
 \label{nt:mirror}
\end{equation}
Si \(z=\sigma+it\), entonces
\begin{equation}
 |\tau(z)|^2
 =\frac{(\sigma-1)^2+t^2}{\sigma^2+t^2};
 \qquad
 |\tau(z)|=1\ \Longleftrightarrow\ \sigma=\frac12.
 \label{nt:critical-circle}
\end{equation}
\end{proposition}

\begin{proof}
La sustitución directa da
\[
 \tau(1-\overline z)
 =\frac{-\overline z}{1-\overline z}
 =\frac{\overline z}{\overline z-1}
 =\frac1{\overline{(z-1)/z}}.
\]
El cociente de módulos cuadrados produce
\eqref{nt:critical-circle}; la igualdad entre numerador y
denominador equivale a \(-2\sigma+1=0\).
\end{proof}

Escribiendo \(\tau(z)=r e^{i\theta}\), con \(r>0\), el
espejo preserva \(\theta\) y transforma \(r\) en \(r^{-1}\).
Su conjunto fijo en esta coordenatización es la circunferencia unidad.
La conjugación \(z\mapsto\overline z\), por su parte,
produce \(\tau\mapsto\overline\tau\) y cambia la orientación
angular. Las dos involuciones conservan, respectivamente,
la coordenada angular y la coordenada radial.

Sobre \(z=1/2+i\gamma\), con \(\gamma\in\mathbb R\),
\begin{equation}
 \tau_\gamma=\frac{\gamma+i/2}{\gamma-i/2},
 \qquad
 \gamma=\frac{i}{2}\frac{\tau_\gamma+1}{\tau_\gamma-1}
       =\frac12\cot\frac{\theta}{2},
 \quad \tau_\gamma=e^{i\theta}.
 \label{nt:spectral-phase}
\end{equation}
La compactificación conserva la altura mediante la inversa;
el punto \(\tau=1\) corresponde al infinito.

Sea \(\mathcal D=C_c^\infty(\mathbb R;\mathbb C)\), con su topología
usual de funciones de prueba. Para la realización positiva de la forma
de Weil se utilizan expresamente su positividad semidefinida, su
invariancia por traslación y su continuidad en esa topología, en el
dominio demostrado por el lector espectral correspondiente. Si
\(\mathscr W(f,g)=\langle Ef,Eg\rangle\), su núcleo es
\(\mathcal N=\{f\in\mathcal D:\mathscr W(f,f)=0\}\), y se define
\[
 \mathcal H_{\mathscr W}
 =\overline{\mathcal D/\mathcal N}.
\]
Las traslaciones inducen un grupo unitario fuertemente
continuo \(U_t[f]=[f(\,\cdot-t)]\): la invariancia preserva la norma,
y la continuidad de traslaciones en \(\mathcal D\) y de la forma
prueba la continuidad fuerte sobre un subespacio denso y después en
toda la completación. Con la convención
\(U_t=e^{itH}\), su generador autoadjunto actúa sobre
las clases de pruebas por \(H[f]=[if']\). El Cayley de
esta realización es el operador acotado
\begin{equation}
 C_{\mathscr W}
 =(H+iI/2)(H-iI/2)^{-1}
 =I+i(H-iI/2)^{-1}.
 \label{nt:cayley-hilbert}
\end{equation}
El cálculo funcional lo hace unitario, pues el cociente
en~\eqref{nt:spectral-phase} tiene módulo uno para
\(\gamma\) real. La identidad de compresión se aplica
a \(C_{\mathscr W}\) en \(\mathcal H_{\mathscr W}\).
El espacio, el dominio del generador y la positividad
que permite su construcción son las premisas de esta
realización espectral \cite{HMT-VIII}.

La memoria bilateral actúa, en cambio, en
\(\ell^2(\mathbb Z)\), y las permutaciones \(S^3,S^4\)
actúan en \(\mathbb R^{12}\). Cada realización determina
su espectro, sus sectores fijos y la acción de una vuelta.
Los respectivos mapas de transporte especifican las
composiciones entre estos espacios.

\subsection{Cayley sobre pruebas y conservación de la forma aritmética}

La misma representación aritmética admite un lector de Cayley
definido directamente sobre funciones de prueba. Fijemos
\(a>b>1/2\) y
\[
 \mathscr E_b=
 \left\{f\in C^\infty(\mathbb R):
   \sup_x e^{b|x|}|f^{(j)}(x)|<\infty
   \text{ para todo }j\geq0\right\}.
\]
Las funciones suaves de soporte compacto pertenecen a
\(\mathscr E_b\). Definimos
\begin{align}
 (R_a^+f)(x)&=\int_0^\infty e^{-at}f(x+t)\,dt,
 &(R_a^-f)(x)&=\int_0^\infty e^{-at}f(x-t)\,dt,\label{nt:resolvents}\\
 C_af&=f-2aR_a^+f,
 &C_a^{-1}f&=f-2aR_a^-f.\label{nt:test-cayley}
\end{align}
Cada derivada de \(R_a^\pm f\) satisface la cota del
correspondiente seminorma de \(f\), multiplicada por
\((a-b)^{-1}\). Esto sigue de
\(e^{b|x|}|f(x\pm t)|\leq M_0e^{bt}\).
Ambas aplicaciones preservan \(\mathscr E_b\).

Con la transformada
\(\widehat f(\xi)=\int_{\mathbb R}f(x)e^{-i\xi x}\,dx\),
la aplicación de Fubini da
\begin{equation}
 \widehat{C_af}(\xi)
 =c_a(\xi)\widehat f(\xi),\qquad
 c_a(\xi)=\frac{\xi-ia}{\xi+ia}.
 \label{nt:cayley-multiplier}
\end{equation}
El multiplicador de la segunda fórmula
de~\eqref{nt:test-cayley} es \(c_a^{-1}\), lo que prueba
la inversa. Para \(\xi\in\mathbb R\),
\(|c_a(\xi)|=1\); Plancherel proporciona su extensión
unitaria a \(L^2(\mathbb R)\).

Para precisar la conservación aritmética, sean
\[
 \mathcal C_{f,g}(u)
 =\int_{\mathbb R}f(x)\overline{g(x-u)}\,dx,
 \qquad
 P_f^\pm=\int_{\mathbb R}f(x)e^{\pm x/2}\,dx.
\]
Las longitudes \(\ell_{p,k}=k\log p\) y los pesos
\(w_{p,k}=(\log p)p^{-k/2}\) son las representaciones
primo--potencia del lector aritmético. La forma completa
del artículo~VIII se escribe
\begin{equation}
 \begin{split}
 \mathscr W(f,g)={}&c_\Gamma\mathcal C_{f,g}(0)\\
 &+\int_0^\infty\frac{e^{-u/2}}{1-e^{-2u}}
 \bigl[2\mathcal C_{f,g}(0)
       -\mathcal C_{f,g}(u)-\mathcal C_{f,g}(-u)\bigr]\,du\\
 &-\sum_{p,k}w_{p,k}
   \bigl[\mathcal C_{f,g}(\ell_{p,k})
        +\mathcal C_{f,g}(-\ell_{p,k})\bigr]
 +P_f^+\overline{P_g^-}+P_f^-\overline{P_g^+},
 \end{split}
 \label{nt:weil-complete}
\end{equation}
donde \(c_\Gamma=\psi(1/4)-\log\pi\), y
\(\psi=\Gamma'/\Gamma\) es la derivada logarítmica de la
función gamma. Estas funciones y constantes intervienen
en la realización analítica posterior del lector
aritmético del corpus.

La cota
\[
 |\mathcal C_{f,g}(u)|
 \leq M_fM_g\bigl(|u|+b^{-1}\bigr)e^{-b|u|}
\]
controla la suma prima por una serie convergente del tipo
\(\sum_{n\geq2}(\log n)(1+\log n)n^{-b-1/2}\).
La diferencia simétrica de correlaciones es \(O(u^2)\)
en el origen, mientras que el peso gamma es \(O(u^{-1})\).
En infinito ese peso decrece exponencialmente.
El margen \(b>1/2\) asegura asimismo la convergencia de
los dos lectores polares. Así queda fijado el dominio
de todos los términos de~\eqref{nt:weil-complete}.

\begin{proposition}[Invariancia aritmética y balance nonádico]
Para \(f,g\in\mathscr E_b\),
\begin{equation}
 \mathscr W(C_af,C_ag)=\mathscr W(f,g).
 \label{nt:weil-invariance}
\end{equation}
Con \(\mathcal T_a=(8I+C_a)/9\), se cumple
\begin{equation}
 \mathscr W(f,g)
 =\mathscr W(\mathcal T_af,\mathcal T_ag)
 +\frac8{81}\mathscr W((I-C_a)f,(I-C_a)g).
 \label{nt:weil-balance}
\end{equation}
\end{proposition}

\begin{proof}
El multiplicador \(c_a\) es unitario en la recta real y
conmuta con las traslaciones. En consecuencia,
\(\mathcal C_{C_af,C_ag}(u)=\mathcal C_{f,g}(u)\)
para todo \(u\). Los lectores polares se transforman
por Fubini según
\[
 P_{C_af}^+
 =-\frac{a-1/2}{a+1/2}P_f^+,\qquad
 P_{C_af}^-
 =-\frac{a+1/2}{a-1/2}P_f^-.
\]
Sus factores son reales y recíprocos, de modo que
cada producto polar cruzado permanece inalterado.
La reunión de los términos de
\eqref{nt:weil-complete} prueba~\eqref{nt:weil-invariance}.
La expansión sesquilineal del segundo miembro de
\eqref{nt:weil-balance} cancela los términos mixtos y
produce
\(\bigl(8\mathscr W(f,g)+
\mathscr W(C_af,C_ag)\bigr)/9\), igual a
\(\mathscr W(f,g)\).
\end{proof}

La iteración finita conserva la forma completa:
\[
 \mathscr W(f,g)
 =\mathscr W(\mathcal T_a^Nf,\mathcal T_a^Ng)
 +\frac8{81}\sum_{j=0}^{N-1}
 \mathscr W((I-C_a)\mathcal T_a^jf,
            (I-C_a)\mathcal T_a^jg).
\]
Esta identidad sesquilineal conserva simultáneamente
gamma, primos y términos polares. Su evaluación en una
realización positiva permite interpretarla como reparto
de energía; su validez algebraica precede a esa
interpretación. El margen \(a>b>1/2\) pertenece al
dominio exponencial de pruebas de esta construcción.
El Cayley~\eqref{nt:cayley-hilbert} con parámetro \(1/2\)
pertenece al dominio hilbertiano de su generador.

\subsection{Covariancia métrica del flujo dodecafásico}

La reconstrucción de la sección~\ref{sec:memory} determina
\(u_K\in\mathbb R^{12}\), con
\(\mathbf1^{\mathsf T}u_K=0\) y \(u_K\ne0\).
Escribimos \(A_K=\operatorname{diag}(u_K)\), con la dirección y la
normalización declaradas en \eqref{eq:u-main}, y
\[
 g_s=e^{sA_K}\in\operatorname{GL}(12,\mathbb R).
\]
La traza nula de \(A_K\) implica \(\det g_s=1\).
Cada \(g_s\) transporta los pesos y las coordenadas
geométricas según la realización correspondiente.
Para conservar también el balance de memoria entre
coordenatizaciones hay que transportar el producto interior.

\begin{proposition}[Transporte métrico de compresión y complemento]
Sean \(V\) un espacio hermítico finito, \(C_0\) unitario
y \(g:V\to V\) invertible. Escribimos
\(g^{-*}=(g^{-1})^*\) y definimos
\begin{equation}
 C_g=gC_0g^{-1},\qquad
 G_g=g^{-*}g^{-1},\qquad
 \langle v,w\rangle_g=\langle G_gv,w\rangle.
 \label{nt:metric-chart}
\end{equation}
Entonces \(C_g^*G_gC_g=G_g\), y
\begin{align}
 T_{C_g}g&=gT_{C_0},\label{nt:metric-naturality}\\
 \eta_{C_g}g&=g^{\oplus9}\eta_{C_0},\label{nt:eta-naturality}\\
 G_g&=T_{C_g}^*G_gT_{C_g}
 +\frac8{81}(I-C_g)^*G_g(I-C_g).\label{nt:metric-balance}
\end{align}
\end{proposition}

\begin{proof}
La positividad de \(G_g\) se obtiene de
\(\langle G_gv,v\rangle=\|g^{-1}v\|^2\).
Sustituyendo~\eqref{nt:metric-chart},
\[
 C_g^*G_gC_g
 =g^{-*}C_0^*g^*g^{-*}g^{-1}gC_0g^{-1}
 =g^{-*}C_0^*C_0g^{-1}=G_g.
\]
La identidad \(C_gg=gC_0\) da
\eqref{nt:metric-naturality} al sumar \(8g\) y dividir
por nueve. También da
\((C_g-I)g=g(C_0-I)\); su sustitución en
\eqref{nt:complement} prueba~\eqref{nt:eta-naturality}.
Finalmente, la expansión de
\eqref{nt:observable-balance}, con \(A=G_g\) y
\(C_g^*G_gC_g=G_g\), prueba
\eqref{nt:metric-balance}.
\end{proof}

Aplicando la proposición a \(g=g_s\) y a los operadores
\(C_0=S^3,S^4\) de la realización dodecafásica, las
compresiones, las memorias y sus productos interiores
se transportan por cuadrados conmutativos explícitos.
La identificación \(g_s:(V,\langle\cdot,\cdot\rangle)
\to(V,\langle\cdot,\cdot\rangle_{g_s})\) es isométrica.
Por tanto, la fórmula transporta también los sectores
fijos y los registros a toda profundidad:
\[
 P_{\mathrm{fix},s}
 =g_sP_{\mathrm{fix},0}g_s^{-1}.
\]
El proyector del miembro derecho es ortogonal respecto
de \(G_{g_s}\).

La condición de volumen y la condición métrica tienen
efectos diferenciados. Para la matriz racional
\[
 g=\begin{pmatrix}2&0\\0&1/2\end{pmatrix},
 \qquad \det g=1,
\]
la sustitución directa en la compresión da
\[
 \frac{8I+g}{9}
 =\begin{pmatrix}10/9&0\\0&17/18\end{pmatrix},
 \qquad
 T_g^*T_g+\frac8{81}(I-g)^*(I-g)
 =\frac{8I+g^*g}{9}.
\]
La primera coordenada se expande en la métrica
euclídea. El balance~\eqref{nt:metric-balance}
conserva en cambio la unidad métrica cuando se
transporta un operador unitario \(C_0\) junto con
su producto interior. El determinante controla
el volumen; \(G_g\) controla las normas y los
productos cruzados.

El flujo satisface \(g_{-s}=g_s^{-1}\).
La reciprocidad del radio de
\eqref{nt:mirror} pertenece a la coordenatización espectral,
y la inversión de \(g_s\) pertenece a la
realización dimensional del registro. El enlace
constructivo utilizado aquí conserva las
aplicaciones explícitas de memoria a \(K\) y de
\(K\) a su dirección geométrica, seguido por
el transporte métrico de cada realización.
Esta composición mantiene, en un mismo desarrollo,
la reconstrucción exacta, el refinamiento positivo,
el registro de fronteras y la conservación
de los productos interiores.

\section{Incidencia ternaria, pegado de discriminantes y construcción de Leech}
\label{sec:lf-foundations}

La construcción reticular utiliza la información de incidencia que
acompaña al registro dodecafásico. Su dominio conserva las palabras
ternarias, el marco en que se representan, la orientación previa al
acarreo y la bandera que distingue una de las doce posiciones. El
paso a un retículo transmite esos datos mediante operaciones
sucesivas: codificación lineal, selección incidencial, pegado de
discriminantes y formación de un vecino. La clasificación reticular
interviene después de probar la positividad, paridad, autodualidad
y ausencia de raíces de la construcción obtenida.

Estas operaciones continúan la misma generación
APP--TRIT--TPK y el mismo refinamiento compatible del continuo
conjunto que anteceden a la lectura numérica \cite{HMT-IV}.
La matriz finita, el código y la métrica que se definen a
continuación son realizaciones con dominios precisos. La prueba
reticular recibe el estado marcado; el escalar final de una
constante no desempeña el papel de código de pegado ni de longitud.

\subsection{Coordenatización elíptica y transporte de la incidencia}
\label{subsec:lf-incidence}

Las seis unidades \(\{1,2,4,5,7,8\}\) módulo nueve, conservadas
por el transporte modular, admiten una coordenatización de seis posiciones.
La identificamos con
\(\mathbb P^1(\mathbb F_5)=\{\infty,0,1,2,3,4\}\), en ese orden.
Esta identificación fija coordenadas para realizar la relación
cuadrática del régimen elíptico del TRIT.

Sea \(\chi:\mathbb F_5\to\{0,1,-1\}\) el carácter cuadrático:
\(\chi(0)=0\), \(\chi(1)=\chi(4)=1\) y
\(\chi(2)=\chi(3)=-1\). Definimos la matriz entera
\(C_{\mathrm P}\) con diagonal cero,
\((C_{\mathrm P})_{\infty,a}=(C_{\mathrm P})_{a,\infty}=1\),
y entradas finitas \((C_{\mathrm P})_{a,b}=\chi(b-a)\).
Su cuadrado es \(5I_6\). En efecto, cada entrada diagonal del
cuadrado suma cinco unidades. Las entradas que involucran
\(\infty\) fuera de la diagonal se anulan porque la suma de
\(\chi\) es cero. Para dos índices finitos distintos, la
contribución uno de la posición \(\infty\) se cancela con
\(\sum_c\chi(c-a)\chi(c-b)=-1\). Esta última identidad se verifica
directamente sobre los cinco residuos, o mediante la sustitución
afín que lleva \((a,b)\) a \((0,1)\).

La reducción módulo tres produce la matriz simétrica
\begin{equation}
 A_W=
 \begin{pmatrix}
 0&1&1&1&1&1\\
 1&0&1&2&2&1\\
 1&1&0&1&2&2\\
 1&2&1&0&1&2\\
 1&2&2&1&0&1\\
 1&1&2&2&1&0
 \end{pmatrix},
 \qquad A_W^2=2I_6=-I_6
 \quad\text{en }\mathbb F_3.
 \label{eq:lf-aw}
\end{equation}
La igualdad conserva el tipo elíptico en este codominio finito.
El cuerpo finito y la posterior realización arquimediana de la
unidad imaginaria tienen dominios diferentes.

Para una matriz \(A\in M_6(\mathbb F_3)\) y una palabra fila
\(w\in\mathbb F_3^6\), escribimos
\[
 \gamma_A(w)=(w,wA)\in\mathbb F_3^{12}.
\]
Si \(f\in\operatorname{GL}_6(\mathbb F_3)\), entonces
\begin{equation}
 \gamma_{fAf^{-1}}(wf^{-1})
   =\gamma_A(w)(f^{-1}\oplus f^{-1}).
 \label{eq:lf-frame-covariance}
\end{equation}
La primera mitad es \(wf^{-1}\); la segunda es
\(wf^{-1}fAf^{-1}=wAf^{-1}\), lo que prueba la identidad.

Sea \(X_n^{\mathrm{mar}}\) el espacio de prefijos enriquecidos de
longitud \(n\), con marco inicial declarado, y sean
\(p_{n,m}:X_n^{\mathrm{mar}}\to X_m^{\mathrm{mar}}\) sus
truncaciones. En el paso \(j\), el prefijo contiene una palabra
visible \(w_j\) y un marco \(f_j\). La actualización del marco
tiene la forma \(f_{j+1}=g_j(x)f_j\), donde \(g_j(x)\) depende
únicamente de información ya contenida en el prefijo. Una
prolongación compatible conserva, por ello, todos los marcos
anteriores.

\begin{proposition}[Naturalidad del transporte incidencial]
\label{prop:lf-incidence-naturality}
La aplicación
\begin{equation}
 \mathcal Q_n(x)=
 \bigl((f_j,\gamma_{f_jA_Wf_j^{-1}}(w_j))\bigr)_{0\leq j<n}
 \label{eq:lf-Qn}
\end{equation}
satisface \(r_{n,m}\mathcal Q_n=\mathcal Q_mp_{n,m}\), donde \(r_{n,m}\) trunca
las primeras \(m\) componentes codificadas. Por tanto induce una
aplicación entre los límites inversos de los sistemas de
prefijos y de incidencias marcadas.
\end{proposition}
\begin{proof}
Dos prefijos compatibles tienen las mismas palabras y los mismos
operadores de transporte hasta el paso \(m-1\). Partiendo del
mismo marco inicial, la relación recursiva da por inducción los
mismos \(f_j\) hasta ese paso. Coinciden entonces las primeras
\(m\) componentes de \eqref{eq:lf-Qn}. La propiedad universal
del límite inverso produce la aplicación de prolongación.
\end{proof}

En cada marco, la primera proyección de \(\gamma_A\) recupera
exactamente la palabra visible. Este alcance es local: la
memoria profunda se conserva en el estado enriquecido y su
recuperación exige las coordenadas correspondientes. Al pasar
de palabras a soportes se retiene todavía menos información.
Asimismo, un cambio general de base en
\(\operatorname{GL}_6(\mathbb F_3)\) transporta también la métrica
de Hamming y la incidencia de la coordenatización. Los cambios monomiales
las preservan directamente. Los soportes utilizados en esta
sección se calculan en la coordenatización explícita \eqref{eq:lf-aw}.

La naturalidad anterior concreta la lectura incidencial que
acompaña a los cilindros numéricos. El encajamiento de cilindros
y la restricción de \eqref{eq:lf-Qn} actúan sobre un mismo
prefijo compatible, conservando respectivamente su evaluación
posicional y sus relaciones marcadas. El ambiente
\(\mathbb F_3^6\), de \(729\) palabras, proporciona el dominio
lineal de la codificación; conserva su distinción respecto de
las \(243\) palabras visibles y las \(468\) emisiones del censo
admisible de APP--TRIT--TPK.

\subsection{Código ternario autodual y enumeración de sus pesos}
\label{subsec:lf-code}

\begin{definition}[Código de la coordenatización dodecafásica]
El código lineal asociado a \eqref{eq:lf-aw} es
\[
 C_W=\{(w,wA_W):w\in\mathbb F_3^6\}\subset\mathbb F_3^{12}.
\]
El peso \(\operatorname{wt}(c)\) cuenta las coordenadas no nulas
de \(c\), y \(\operatorname{supp}(c)\subset\{1,\ldots,12\}\)
es su conjunto de posiciones no nulas. Una hexada es el soporte
de una palabra de peso seis; su colección se denota \(\mathcal H\).
\end{definition}

\begin{theorem}[Autodualidad, distancia mínima y enumerador]
\label{thm:lf-code}
El código \(C_W\) es autodual de parámetros \([12,6,6]_3\) y
tiene enumerador
\begin{equation}
 W_{C_W}(z)=1+264z^6+440z^9+24z^{12}.
 \label{eq:lf-enumerator}
\end{equation}
\end{theorem}
\begin{proof}
La matriz generadora \([I_6\mid A_W]\) tiene rango seis.
Por la simetría de \(A_W\) y \eqref{eq:lf-aw},
\[
 [I_6\mid A_W][I_6\mid A_W]^{\mathsf T}
       =I_6+A_W^2=0.
\]
Así, el código es autoortogonal de dimensión mitad de la
longitud y, en consecuencia, autodual.

La enumeración restante es una operación entera finita sobre
la matriz explícita. Para \(r\in\{0,\ldots,6\}\), agrupamos
las palabras por \(\operatorname{wt}(w)=r\). La tabla siguiente
da el número de palabras con cada peso total:
\[
\begin{array}{c|rrrr|r}
r&0&6&9&12&\text{total}\\ \hline
0&1&0&0&0&1\\
1&0&12&0&0&12\\
2&0&60&0&0&60\\
3&0&120&40&0&160\\
4&0&60&180&0&240\\
5&0&12&180&0&192\\
6&0&0&40&24&64\\ \hline
\text{total}&1&264&440&24&729
\end{array}
\]
Cada fila contiene \(\binom6r2^r\) palabras. Para verificar
todas sus entradas, sea \(q\) un entero de \(0\) a \(728\), y
definamos sus seis cifras ternarias
\[
 d_i(q)=\left\lfloor\frac{q}{3^{6-i}}\right\rfloor\bmod3,
 \qquad
 e_j(q)=\left(\sum_{i=1}^{6}d_i(q)(A_W)_{ij}\right)\bmod3.
\]
La fila y la columna de la contribución de \(q\) son,
respectivamente,
\[
 r(q)=\sum_i\mathbf1_{\{d_i(q)\ne0\}},\qquad
 h(q)=r(q)+\sum_j\mathbf1_{\{e_j(q)\ne0\}}.
\]
La recurrencia de acumulación empieza con
\(T^{(0)}_{r,h}=0\) y aplica
\[
 T^{(q+1)}_{r,h}
 =T^{(q)}_{r,h}
   +\mathbf1_{\{r=r(q)\}}\mathbf1_{\{h=h(q)\}}
 \quad(0\leq q<729).
\]
Produce exactamente la tabla mostrada. La expansión ternaria
da una biyección entre esos \(729\) enteros y todas las palabras
de \(\mathbb F_3^6\), por lo que la operación agota el código.
La suma de sus filas prueba \eqref{eq:lf-enumerator}; el primer
peso no nulo es seis.
\end{proof}

\begin{remark}[Evaluación reproducible de la recurrencia finita]
El siguiente bucle implementa literalmente las sumas de la prueba.
Sus únicos datos son los enteros de la matriz
\eqref{eq:lf-aw}; calcula el enumerador antes de compararlo con
la igualdad obtenida.
\end{remark}
{\small
\begin{verbatim}
from collections import Counter
from itertools import product
A = ((0,1,1,1,1,1), (1,0,1,2,2,1),
     (1,1,0,1,2,2), (1,2,1,0,1,2),
     (1,2,2,1,0,1), (1,1,2,2,1,0))
rows = [Counter() for _ in range(7)]
for w in product(range(3), repeat=6):
    v = tuple(sum(w[i]*A[i][j] for i in range(6)) % 3
              for j in range(6))
    r = sum(x != 0 for x in w)
    h = r + sum(x != 0 for x in v)
    rows[r][h] += 1
total = sum(rows, Counter())
if dict(total) != {0:1, 6:264, 9:440, 12:24}:
    raise ArithmeticError("Enumerador incompatible")
\end{verbatim}
}

\begin{proposition}[Hexadas y diseño de Witt]
\label{prop:lf-witt}
La colección \(\mathcal H\) contiene \(132\) hexadas y cada
quintupla de posiciones pertenece a una única hexada.
\end{proposition}
\begin{proof}
Dos palabras de peso seis con el mismo soporte difieren por
un escalar. En efecto, entre sus seis cocientes de coordenadas
no nulas, uno de los signos \(1,-1\) aparece al menos tres
veces. Restando el múltiplo correspondiente, si las palabras
no fueran proporcionales, se obtendría una palabra no nula
de peso a lo sumo tres, en contradicción con
el teorema~\ref{thm:lf-code}. Cada soporte corresponde así a
la pareja \(\{c,-c\}\), y hay \(264/2=132\) soportes.

Si dos soportes distintos compartieran cinco posiciones,
en esa intersección uno de los dos cocientes aparecería al
menos tres veces. La combinación lineal que cancela dichas
coordenadas tendría soporte en una unión de tamaño siete
y peso a lo sumo cuatro. La distancia mínima lo impide.
Una quintupla pertenece, por tanto, a lo sumo a una hexada.
Como
\[
 132\binom65=792=\binom{12}{5},
\]
pertenece a exactamente una.
\end{proof}

El objeto construido es el código ternario extendido de Golay
con su pequeño diseño de Witt. El reconocimiento se efectúa
después de la codificación, la autodualidad y el cómputo de
distancia. Para la construcción reticular sólo se necesitan
las propiedades ya demostradas y la bandera marcada que se
obtiene a continuación.

\subsection{Bandera dodecafásica y selección del origen}
\label{subsec:lf-flag}

En el marco del registro generado se conservan las palabras
\[
 w_\pi=010211,\qquad
 w_e=201101,\qquad
 w_\varphi=121200,\qquad
 w_{\mathrm{APP}}=022020.
\]
Los subíndices \(\pi,e,\varphi\) identifican las correspondientes
regiones de lectura de las constantes; \(w_e\) conserva aquí
la región de Euler. Aplicar \(\gamma_{A_W}\) y tomar soportes da
\[
\begin{aligned}
 P_\pi&=\{2,4,5,6,7,8,9,10,11\},\\
 H_e&=\{1,3,4,6,9,10\},\\
 H_\varphi&=\{1,2,3,4,7,9\},\\
 H_{\mathrm{APP}}&=\{2,3,5,10,11,12\}.
\end{aligned}
\]
El primer soporte tiene peso nueve; los tres restantes son
hexadas. El registro
\[
 K=(234,543,140,729,659,824,621,58,914,794,146,601)
\]
selecciona
\[
 B_K=\{j:K_j\geq729\}=\{4,6,9,10\}=H_e\cap P_\pi.
\]
El umbral \(729\) pertenece al formato de bloques de la
lectura declarada. En este registro concreto, los umbrales
enteros de \(660\) a \(729\) producen el mismo soporte, de
modo que el soporte conserva la incidencia seleccionada,
mientras el registro completo conserva también los bloques
y el umbral de su lectura.

El registro firmado anterior al acarreo es
\[
\begin{split}
 Z_0={}&(7,297,353,-430,-715,-1199,\\
       &\hspace{9mm}-3,286,-894,-528,380,663).
\end{split}
\]
Su soporte negativo es la hexada
\[
 H_\alpha^-=\{4,5,6,7,9,10\}.
\]
Los signos se reciben del estado previo al acarreo y no se
extraen de la expansión del escalar final.

\begin{proposition}[Coincidencia de selección firmada e incidencial]
\label{prop:lf-flag}
La hexada \(H_\alpha^-\) es la única \(H\in\mathcal H\) que cumple
\[
\begin{gathered}
 B_K\subset H\subset P_\pi,\\
 (|H\cap H_e|,|H\cap H_\varphi|,
                   |H\cap H_{\mathrm{APP}}|)=(4,3,2).
\end{gathered}
\]
\end{proposition}
\begin{proof}
La codificación de las palabras de peso seis, o la recurrencia
finita anterior aplicada a sus soportes, da las cuatro
hexadas sobre \(B_K\). Sus pares complementarios son
\[
 \{1,3\},\qquad\{2,11\},\qquad\{5,7\},\qquad\{8,12\}.
\]
Sólo el segundo y el tercero están contenidos en \(P_\pi\).
La hexada del segundo tiene intersecciones de tamaño tres
con \(H_\varphi\) y con \(H_{\mathrm{APP}}\), y por tanto
incumple la última condición. La del tercero tiene
intersecciones \(4,3,2\), y coincide con \(H_\alpha^-\).
\end{proof}

\begin{lemma}[Origen de la bandera marcada]
\label{lem:lf-origin}
El origen seleccionado por la incidencia es
\[
 o_*=(H_e\cap H_\varphi)
                 \setminus(H_\alpha^-\cup H_{\mathrm{APP}})
     =\{1\}.
\]
La selección conmuta con toda biyección aplicada
simultáneamente a los soportes.
\end{lemma}
\begin{proof}
La intersección \(H_e\cap H_\varphi\) es
\(\{1,3,4,9\}\). La unión sustraída contiene \(3,4,9\)
y excluye \(1\). La equivariancia procede de que las
biyecciones conmutan con unión, intersección y diferencia.
\end{proof}

La bandera distingue una componente entre doce antes del
paso métrico. El vector radial se determina después,
dentro de una cámara positiva declarada. Así quedan
separadas y articuladas la selección combinatoria y la
operación reticular.

\subsection{Pegado ternario de las doce componentes}
\label{subsec:lf-glue}

Sea \(A_2=\mathbb Za_1\oplus\mathbb Za_2\), con forma
bilineal de matriz
\[
 B_2=\begin{pmatrix}2&-1\\-1&2\end{pmatrix},
 \qquad \rho=a_1+a_2,\qquad \rho^2=2.
\]
Para \(x=ua_1+va_2\),
\[
 x^2=2(u^2-uv+v^2).
\]
La identidad
\(4(u^2-uv+v^2)=(2u-v)^2+3v^2\) muestra
la positividad. El determinante es tres. Las tres
clases de \(A_2^*/A_2\) se representan mediante
\[
 \varpi_0=0,\qquad
 \varpi_1=\frac{2a_1+a_2}{3},\qquad
 \varpi_2=\frac{a_1+2a_2}{3}.
\]
El emparejamiento con \(a_1,a_2\) verifica que estos
vectores pertenecen a \(A_2^*\), y
\(2\varpi_1-\varpi_2=a_1\). Por ello
\(j\mapsto\varpi_j+A_2\) identifica aditivamente
\(\mathbb F_3\) con el discriminante.

Para \(c\in C_W\), definimos
\(\phi(c)=(\varpi_{c_1},\ldots,\varpi_{c_{12}})\)
y
\begin{equation}
 N=\bigcup_{c\in C_W}\bigl(A_2^{12}+\phi(c)\bigr).
 \label{eq:lf-glue}
\end{equation}

\begin{theorem}[Retículo de pegado par y unimodular]
\label{thm:lf-niemeier}
El conjunto \(N\) es un retículo positivo, par y
unimodular de rango \(24\). Sus vectores de norma dos
son exactamente las raíces de \(A_2^{12}\).
\end{theorem}
\begin{proof}
La linealidad del código y la aditividad del discriminante
hacen que la unión sea un subgrupo de
\((A_2^*)^{12}\) que contiene \(A_2^{12}\).
Es, por tanto, un retículo de rango \(24\).

El emparejamiento de las clases discriminantes es
\(2cd/3\) módulo \(\mathbb Z\) en cada componente.
La autoortogonalidad de \(C_W\) anula su suma y prueba
la integralidad de \(N\). La norma de una clase
representada por \(\phi(c)\) es
\(2\operatorname{wt}(c)/3\) módulo \(2\mathbb Z\).
Todos los pesos de \eqref{eq:lf-enumerator} son
múltiplos de tres, luego \(N\) es par.

Su índice sobre \(A_2^{12}\) es \(|C_W|=3^6\).
Con el determinante de la matriz de Gram,
\[
 \det N=\frac{\det(A_2^{12})}{[N:A_2^{12}]^2}
       =\frac{3^{12}}{(3^6)^2}=1.
\]
La forma ambiente es la suma ortogonal de doce formas
positivas; queda demostrada la positividad.
La integralidad y el determinante uno implican
\(N^*=N\).

En una clase no nula del discriminante local, un
vector tiene la forma \((ua_1+va_2)/3\), con residuos
\((u,v)\equiv(2,1)\) o \((1,2)\pmod3\).
El entero \(u^2-uv+v^2\) es positivo y divisible
por tres, luego es al menos tres. Los representantes
\(\varpi_1,\varpi_2\) alcanzan ese valor: la norma
mínima de cada clase no nula es \(2/3\).
Una palabra no nula del código tiene por lo menos
seis coordenadas no nulas. Todo vector de una clase
de pegado no trivial tiene, por tanto, norma
al menos \(6(2/3)=4\).

En \(A_2^{12}\), una norma dos sólo puede proceder
de una única componente, porque cada componente
tiene norma par positiva mínima dos. Las soluciones
de \(u^2-uv+v^2=1\) son
\((\pm1,0),(0,\pm1),(1,1),(-1,-1)\).
Son las seis raíces de \(A_2\), y completan el censo
de raíces de \(N\).
\end{proof}

Las propiedades anteriores identifican \(N\) como el
retículo de Niemeier de sistema de raíces \(A_2^{12}\).
La descripción de todas sus raíces es necesaria para
controlar cuáles sobreviven al paso al vecino.

\subsection{Cámara radial y vecino sin raíces}
\label{subsec:lf-neighbour}

Se fija la cámara radial
\[
 v(a)=\sum_{j=1}^{12}a_j\rho_j,\qquad
 a_j=1+3b_j,\qquad b_j\in\mathbb Z_{\geq0},
\]
donde \(\rho_j\) es la copia de \(\rho\) en la
componente \(j\). La paridad de un vecino de índice
tres requiere \(v(a)^2\equiv0\pmod{18}\).
Como
\[
 v(a)^2=2\sum_j(1+3b_j)^2
       =24+12\sum_jb_j+18\sum_jb_j^2,
\]
la congruencia equivale a \(\sum_jb_j\equiv1\pmod3\).
La norma mínima dentro de la cámara se alcanza con
un único \(b_j=1\) y los demás cero, y vale \(54\).
El origen \(o_*=1\) del lema~\ref{lem:lf-origin}
selecciona
\begin{equation}
 v_\alpha=4\rho_1+\rho_2+\cdots+\rho_{12},
 \qquad v_\alpha^2=54.
 \label{eq:lf-valpha}
\end{equation}
El calificativo de mínimo se refiere a la cámara
declarada. En la clase residual completa,
\((a_1-2a_2,\rho,\ldots,\rho)\) representa la
misma clase módulo tres y tiene norma
\(14+11\cdot2=36\). Esta distinción preserva
la condición que determina el coeficiente cuatro.

Definimos
\begin{equation}
\begin{split}
 N_0&=\{x\in N:\langle x,v_\alpha\rangle\equiv0\pmod3\},\\
 \Lambda&=N_0+\mathbb Z\,v_\alpha/3.
\end{split}
\label{eq:lf-lambda}
\end{equation}

\begin{theorem}[Vecino par autodual sin raíces]
\label{thm:lf-leech}
La construcción \eqref{eq:lf-lambda} es un
retículo positivo, par y autodual de rango \(24\)
sin vectores de norma dos. En consecuencia, es
isométrica al retículo de Leech.
\end{theorem}
\begin{proof}
El carácter \(x\mapsto\langle x,v_\alpha\rangle\bmod3\)
es no nulo. Una raíz simple de la componente \(j\)
tiene emparejamiento \(a_j\equiv1\pmod3\) con
\(a_j\rho_j\). Por tanto \([N:N_0]=3\), de donde
\(\det N_0=9\).

Para \(x\in N_0\), el emparejamiento
\(\langle x,v_\alpha/3\rangle\) es entero, y así
\(v_\alpha/3\in N_0^*\). Además,
\(\langle v_\alpha,v_\alpha\rangle=54\) implica
\(v_\alpha\in N_0\). Si \(v_\alpha/3\in N\),
la integralidad de \(N\) obligaría a que todos los
emparejamientos de \(N\) con \(v_\alpha\) fuesen
divisibles por tres, contradiciendo el carácter
no nulo. La clase de \(v_\alpha/3\) sobre \(N_0\)
tiene entonces orden exactamente tres. Resulta
\[
 [\Lambda:N_0]=3,\qquad
 \det\Lambda=\frac9{3^2}=1.
\]
Los emparejamientos de \(N_0\) con el nuevo
generador son enteros y la norma de este generador
es seis. La extensión es integral. Para
\(x\in N_0\) y \(m\in\mathbb Z\),
\[
 \left(x+m\frac{v_\alpha}{3}\right)^2
 =x^2+2m\frac{\langle x,v_\alpha\rangle}{3}
       +6m^2\in2\mathbb Z.
\]
Por ello \(\Lambda\) es par; su integralidad y
determinante uno prueban la autodualidad.
La positividad y el rango se heredan del espacio
euclídeo ambiente.

Para excluir raíces, primero consideramos la
clase \(N_0\). Todas las raíces de \(N\) pertenecen
a \(A_2^{12}\) por el teorema~\ref{thm:lf-niemeier}.
Los emparejamientos de sus seis raíces locales
con \(a_j\rho_j\) son \(\pm a_j,\pm2a_j\),
congruentes con \(\pm1,\pm2\pmod3\).
Ninguna raíz pertenece a \(N_0\).

Las otras dos clases de \(\Lambda/N_0\) pueden
representarse con \(\pm v_\alpha/3\).
En cada componente, sus vectores pertenecen a
una de las clases
\[
 A_2+\varpi_j\pm\rho/3,\qquad j=0,1,2.
\]
La componente distinguida tiene la misma
descripción porque \(4\rho/3\equiv\rho/3\pmod{A_2}\).
Para los numeradores \((u,v)\) en
\((ua_1+va_2)/3\), los seis residuos son
\[
 (1,1),(0,2),(2,0),(2,2),(1,0),(0,1)\pmod3.
\]
Todos son no nulos. La positividad de
\(u^2-uv+v^2\) obliga a que sea al menos uno,
y los representantes
\((1,1),(0,-1),(-1,0),(-1,-1),(1,0),(0,1)\)
alcanzan ese mínimo en las seis clases.
Cada componente tiene norma al menos \(2/9\).
Un vector de cualquiera de las dos clases
nuevas tiene norma al menos
\[
 12\cdot\frac29=\frac83>2.
\]
Se han excluido raíces en las tres clases.
El teorema de clasificación de los retículos
positivos pares unimodulares de rango veinticuatro
identifica el único caso sin raíces con el
retículo de Leech \cite{ConwaySloane}.
\end{proof}

La identificación final utiliza una clasificación
externa sobre un objeto cuyas propiedades ya han
sido demostradas. La selección de la componente
distinguida conserva la bandera procedente del
registro firmado; el subíndice \(\alpha\) expresa
esa procedencia. La norma y los coeficientes del
vector radial se obtienen de la métrica y de
la cámara mediante la congruencia indicada.

El retículo \(\Lambda\), su forma bilineal, sus
doce planos ambientes y el vector \(v_\alpha/3\)
quedan así determinados con todas las propiedades
utilizadas por la deformación posterior. El
código contiene asimismo la palabra
\[
 c_*=(1,1,1,1,1,1,2,1,1,1,1,1),
\]
pues \(q_*=(1,1,1,1,1,1)\) satisface
\(q_*A_W=(2,1,1,1,1,1)\). Éste es el dato
de soporte completo que permite construir y
probar la estabilidad de la isometría de orden
tres en la sección~\ref{sec:leech-dualidad}.

\section{Polarización dodecafásica, dualidad reticular y reciprocidad theta}
\label{sec:leech-dualidad}

La incidencia excepcional y la polarización del registro dodecafásico
proceden de dos realizaciones del mismo estado enriquecido. La primera
conserva las relaciones de frontera, el código de pegado y el origen
marcado que determinan el retículo; la segunda hace actuar una
representación sobre las doce coordenadas enteras ya obtenidas por los
lectores del Topological Phase Kernel. Su composición permite estudiar
una deformación métrica cuyo parámetro modifica las longitudes relativas
de doce planos y conserva simultáneamente el volumen total y una
simetría ternaria explícita.

El punto de partida causal sigue siendo la construcción
APP--TRIT--TPK del estado enriquecido y del continuo conjunto.
El registro \(K\) se recibe después del registro dodecafásico y de sus
operaciones de incidencia y orientación. La representación empleada
actúa sobre esa salida y conserva sus doce valores como datos ya
determinados. La realización reticular se recibe con su forma bilineal,
sus clases de pegado y su vector marcado. Estas estructuras, desarrolladas
respectivamente en \cite{hmtX,HMT-IV}, se mantienen al componerlas.
La colección que se construye a continuación es una realización geométrica
de esos datos con una coordenatización común de doce posiciones.

\subsection{Polarización del registro y marcación de los doce planos}
\label{subsec:leech-polarizacion}

El registro dodecafásico y su suma son
\begin{equation}
 \begin{split}
 K={}&(234,543,140,729,659,824,\\
     &\hspace{13mm}621,58,914,794,146,601),\\
 \sum_{j=1}^{12}K_j={}&6263.
 \end{split}
 \label{eq:leech-K}
\end{equation}
Sea \(A_5\) el grupo de permutaciones pares de cinco elementos y sea
\(C_5=\langle(1\,2\,3\,4\,5)\rangle\). Sus doce clases izquierdas
determinan una representación por permutaciones
\(\varrho:A_5\to O(\mathbb R^{12})\). La coordenatización marcada se fija mediante
los representantes siguientes, escritos como listas de imágenes:
\[
\begin{array}{c|c@{\qquad}c|c}
j&r_j&j&r_j\\ \hline
1&(4,3,5,2,1)&7&(5,4,3,2,1)\\
2&(4,2,1,3,5)&8&(1,3,4,2,5)\\
3&(1,2,4,5,3)&9&(4,2,3,5,1)\\
4&(2,5,3,4,1)&10&(3,2,5,4,1)\\
5&(4,3,1,5,2)&11&(4,5,2,3,1)\\
6&(5,1,2,3,4)&12&(5,3,2,4,1).
\end{array}
\]
Concretamente, \(\varrho(a)e_j=e_\ell\) si
\(ar_jC_5=r_\ell C_5\). El carácter tridimensional \(\chi_3\) toma
los valores
\[
\begin{array}{c|ccccc}
\text{clase}&1&2^2&3&5_a&5_b\\ \hline
\chi_3&3&-1&0&(1+\sqrt5)/2&(1-\sqrt5)/2.
\end{array}
\]
La clase \(5_a\) contiene el generador de \(C_5\) indicado y \(5_b\)
contiene su cuadrado. Esta elección distingue los dos caracteres
tridimensionales conjugados. Si \(I_{12}\) es la identidad y
\(\mathbf1\) es el vector de doce unos, se definen
\begin{equation}
 P_3=\frac3{60}\sum_{a\in A_5}\chi_3(a^{-1})\varrho(a),
 \qquad
 P_{11}=I_{12}-\frac1{12}\mathbf1\mathbf1^{\mathsf T}.
 \label{eq:leech-projectors}
\end{equation}
El campo \(\mathbb Q(\sqrt5)\) aparece aquí en la realización de la
representación que actúa sobre el registro generado.

\begin{lemma}[Polarización de suma nula]
\label{lem:leech-u}
Los operadores de \eqref{eq:leech-projectors} satisfacen
\[
 P_3=P_3^{\mathsf T}=P_3^2,\qquad
 \operatorname{rank}P_3=3,\qquad P_3\mathbf1=0,\qquad
 P_3P_{11}=P_{11}P_3=P_3.
\]
Por consiguiente,
\begin{equation}
 u=P_3P_{11}K=P_3K,\qquad
 \sum_{j=1}^{12}u_j=0,\qquad
 \|u\|^2=\frac{6638585+2275584\sqrt5}{20}>0.
 \label{eq:leech-u-norm}
\end{equation}
\end{lemma}
\begin{proof}
La suma del carácter es el proyector isótipo de la representación.
La realidad del carácter y la sustitución \(a\mapsto a^{-1}\)
dan su simetría. La multiplicidad del carácter tridimensional en la
representación de \(A_5/C_5\) es la dimensión de sus invariantes por
\(C_5\), que vale
\[
 \frac15\left(
 3+2\frac{1+\sqrt5}{2}+2\frac{1-\sqrt5}{2}
 \right)=1.
\]
El proyector tiene por ello rango tres. Su ortogonalidad al carácter
trivial implica \(P_3\mathbf1=0\) y las identidades con \(P_{11}\).
La aplicación de la suma finita al registro
\eqref{eq:leech-K} produce las doce coordenadas
\begin{equation}
\begin{aligned}
u={}&(-495/4-233\sqrt5/10,\quad403/4+169\sqrt5/10,\\
&-403/4-169\sqrt5/10,\quad495/4+233\sqrt5/10,\\
&601/4+1031\sqrt5/10,\quad203/4+211\sqrt5/5,\\
&-203/4-211\sqrt5/5,\quad-601/4-1031\sqrt5/10,\\
&313/4+569\sqrt5/10,\quad162+219\sqrt5/20,\\
&-162-219\sqrt5/20,\quad-313/4-569\sqrt5/10).
\end{aligned}
\label{eq:leech-u-coordinates}
\end{equation}
Sumarlas y elevarlas al cuadrado en \(\mathbb Q(\sqrt5)\) da
\eqref{eq:leech-u-norm}. Equivalentemente,
\[
 \|u\|^2=K^{\mathsf T}P_3K
 =\frac1{20}\sum_{a\in A_5}
          \chi_3(a^{-1})K^{\mathsf T}\varrho(a)K.
\]
Los coeficientes positivos de su evaluación prueban que \(u\ne0\).
\end{proof}

La suma nula tiene una función geométrica concreta: las expansiones de
unos planos podrán compensar las contracciones de los restantes en el
determinante total. La no nulidad de \(u\), en cambio, garantiza que esa
compensación contiene una deformación efectiva.

El espacio euclídeo ambiente de la realización reticular es
\begin{equation}
 E=\bigoplus_{j=1}^{12}E_j,\qquad
 B=\operatorname{diag}(B_2,\ldots,B_2),\qquad
 B_2=\begin{pmatrix}2&-1\\-1&2\end{pmatrix}.
 \label{eq:leech-ambient}
\end{equation}
Cada \(E_j\) es el plano real generado por el retículo de raíces
\(A_2\), expresado en su base de raíces simples. La descomposición
ortogonal pertenece al espacio ambiente. El retículo de Leech se
obtiene mediante pegado y paso a un vecino, operaciones que relacionan
los doce factores.

Sea \(C_W\subset\mathbb F_3^{12}\) el código ternario de dimensión seis
de la construcción Paley--Witt. Identificando el discriminante de cada
factor \(A_2\) con \(\mathbb F_3\), el pegado es
\[
 N=N(C_W)=
 \bigcup_{c\in C_W}\bigl(A_2^{12}+\phi(c)\bigr),
\]
donde \(\phi(c)\) elige representantes de las clases discriminantes.
La autoortogonalidad y autodualidad del código transmiten paridad e
integralidad a \(N\); su determinante es
\[
 \det N=\frac{3^{12}}{(3^6)^2}=1.
\]
La construcción excepcional fija el origen marcado \(j=1\) y el
vector
\begin{equation}
 v_\alpha=4\rho_1+\rho_2+\cdots+\rho_{12},
 \qquad \rho_j=(1,1)_j,\qquad
 \|v_\alpha\|_B^2=54.
 \label{eq:leech-marked-vector}
\end{equation}
Con él se forman
\begin{equation}
\begin{split}
 N_{\alpha,0}
  &=\{x\in N:\langle x,v_\alpha\rangle_B\equiv0\pmod3\},\\
 \Lambda&=N_{\alpha,0}+\mathbb Z\,v_\alpha/3.
\end{split}
\label{eq:leech-neighbour}
\end{equation}
La realización excepcional establece que \(\Lambda\) es un retículo
positivo, par, unimodular, sin raíces y de rango \(24\), de modo que es
isométrico al retículo de Leech \cite{HMT-IV}. En particular,
\(\Lambda^*=\Lambda\), con dual respecto de \(B\).
Estas propiedades y la pertenencia \(v_\alpha/3\in\Lambda\) son las
hipótesis reticulares precisas empleadas en los teoremas siguientes.

\subsection{Isometría ternaria y estabilidad del vecino}
\label{subsec:leech-c3}

Definimos las matrices bidimensionales
\[
 C=\begin{pmatrix}0&-1\\1&-1\end{pmatrix},
 \qquad
 R=\begin{pmatrix}0&1\\1&0\end{pmatrix},
 \qquad
 \varpi=\frac13\begin{pmatrix}2\\1\end{pmatrix}.
\]
La multiplicación directa verifica
\[
 C^3=I_2,\quad C^2+C+I_2=0,\quad
 C^{\mathsf T}B_2C=B_2,\quad
 RCR=C^{-1},\quad R^{\mathsf T}B_2R=B_2.
\]
Además, \(C\varpi-\varpi=(-1,0)^{\mathsf T}\), por lo que \(C\)
actúa trivialmente en \(A_2^*/A_2\). La reflexión \(R\) cambia el signo
de la clase discriminante. El código \(C_W\) contiene la palabra de
soporte completo
\[
 c_*=(1,1,1,1,1,1,2,1,1,1,1,1).
\]
En la presentación \(C_W=\{(q,qA_W):q\in\mathbb F_3^6\}\), dicha
palabra procede de \(q_*=(1,1,1,1,1,1)\) y
\(q_*A_W=(2,1,1,1,1,1)\).

Para cada plano se fija el marco
\begin{equation}
 Q_j=\begin{cases}
 R,&(c_*)_j=1,\\
 I_2,&(c_*)_j=2,
 \end{cases}
 \qquad
 g_c=\bigoplus_{j=1}^{12}Q_jCQ_j^{-1}.
 \label{eq:leech-c3}
\end{equation}
Los símbolos \(Q_j\) designan marcos de planos; el símbolo \(P_3\)
queda reservado para el proyector sobre el sector tridimensional.

\begin{proposition}[Simetría de orden tres del retículo marcado]
\label{prop:leech-c3}
La transformación \(g_c\) es una isometría de orden tres de \(E\),
tiene como único vector fijo el cero y satisface
\[
 g_cN=N,\qquad g_cN_{\alpha,0}=N_{\alpha,0},
 \qquad g_c\Lambda=\Lambda.
\]
\end{proposition}
\begin{proof}
Cada bloque es \(C\) o \(C^{-1}\), cuyo polinomio mínimo es
\(X^2+X+1\). Esto prueba el orden y la afirmación sobre vectores fijos,
así como la conservación de \(B\). Cada bloque actúa trivialmente en
el discriminante y conserva, por tanto, todas las clases de pegado de
\(N\).

Escribamos \(v=v_\alpha=\sum_j a_j\rho_j\), con \(a_1=4\) y
\(a_j=1\) para \(j>1\). Todos los \(a_j\) son congruentes con uno
módulo tres. Las identidades
\[
 (C-I_2)\rho=(-2,-1)^{\mathsf T},
 \qquad (C^{-1}-I_2)\rho=(-1,-2)^{\mathsf T}
\]
muestran que
\[
 y=\frac{g_cv-v}{3},\qquad z=\frac{g_c^{-1}v-v}{3}
\]
tienen clases discriminantes \(c_*\) y \(-c_*\), respectivamente.
Ambas pertenecen a \(C_W\), luego \(y,z\in N\). En cada factor,
\[
 \left\langle
 \frac{(C^{\pm1}-I_2)a_j\rho}{3},a_j\rho
 \right\rangle_{B_2}=-a_j^2.
\]
Por suma, \(\langle y,v\rangle_B=\langle z,v\rangle_B=-27\);
por ello \(y,z\in N_{\alpha,0}\).
Si \(x\in N_{\alpha,0}\), entonces
\[
 \langle g_cx,v\rangle_B
  =\langle x,g_c^{-1}v\rangle_B
  =\langle x,v\rangle_B+3\langle x,z\rangle_B
  \equiv0\pmod3.
\]
La misma operación con \(g_c^{-1}\) demuestra la igualdad
\(g_cN_{\alpha,0}=N_{\alpha,0}\). Finalmente,
\(g_c(v/3)=v/3+y\) conserva la clase que genera
\(\Lambda=N_{\alpha,0}+\mathbb Z(v/3)\).
\end{proof}

La simetría ternaria actúa, por consiguiente, sobre el retículo
pegado completo. Su conservación utiliza la palabra del código y
la congruencia del vector marcado, además de la rotación de cada
plano. Esta precisión permite transportar la simetría al flujo sin
perder la información aritmética que hizo posible su construcción.

\subsection{Flujo anisótropo y conservación del covolumen}
\label{subsec:leech-flow}

\begin{definition}[Deformación determinada por la polarización]
Para \(s\in\mathbb R\), parámetro adimensional, se define
\begin{equation}
 H_u=\bigoplus_{j=1}^{12}u_jI_2,\qquad
 D_s=\exp(sH_u)=\bigoplus_{j=1}^{12}e^{su_j}I_2,\qquad
 \Lambda_s=D_s\Lambda.
 \label{eq:leech-flow}
\end{equation}
La coordenada \(u_j\) se asigna al plano \(E_j\) de la misma coordenatización
marcada.
\end{definition}

El factor positivo \(e^{su_j}\) conserva la orientación de cada
plano. Las doce etiquetas transmiten la correspondencia entre la
polarización y la realización reticular. Esta correspondencia
forma parte de la estructura de partida; una permutación de las
etiquetas se trata mediante su transporte simultáneo en ambos objetos.

\begin{theorem}[Covolumen y simetría del flujo]
\label{thm:leech-flow}
Para todos \(s,t\in\mathbb R\) se cumple
\[
\begin{gathered}
 D_{s+t}=D_sD_t,\qquad D_0=I,\qquad D_s^{-1}=D_{-s},\\
 D_s^\dagger=D_s,\qquad D_sg_c=g_cD_s,\qquad \det D_s=1,
\end{gathered}
\]
donde \(A^\dagger=B^{-1}A^{\mathsf T}B\).
Cada \(\Lambda_s\) es una red discreta de rango \(24\) y covolumen
uno, preservada por la isometría \(g_c\) de orden tres.
\end{theorem}
\begin{proof}
Las identidades de grupo se verifican en cada bloque exponencial.
Los bloques escalares son autoadjuntos para \(B_2\), y conmutan
con los bloques de \(g_c\). La suma nula de
\eqref{eq:leech-u-norm} da
\[
 \det D_s=\prod_{j=1}^{12}e^{2su_j}
 =\exp\left(2s\sum_{j=1}^{12}u_j\right)=1.
\]
La imagen de una red por un operador invertible es una red del
mismo rango; su covolumen se multiplica por el valor absoluto del
determinante. El covolumen de \(\Lambda\) es uno, de modo que el
de \(\Lambda_s\) también lo es. La conmutación y
la proposición~\ref{prop:leech-c3} dan
\[
 g_c\Lambda_s=g_cD_s\Lambda=D_sg_c\Lambda=D_s\Lambda.
\]
El orden y el espacio de puntos fijos de \(g_c\) permanecen
inalterados.
\end{proof}

\begin{remark}[Parámetro y simetría conservada]
La dirección \(u\) queda determinada por el registro; \(s\) recorre
una colección de realizaciones. Reemplazar \(u\) por \(u/\|u\|\)
reparametriza la misma colección. Una interpretación física que
seleccione un valor particular de \(s\) añade su regla de realización.
El teorema conserva la simetría explícita \(g_c\); el transporte de
otros automorfismos requiere sus correspondientes relaciones con
\(H_u\).
\end{remark}

\subsection{Dual euclídeo y deformación efectiva}
\label{subsec:leech-dual}

\begin{definition}[Dual euclídeo]
Para una red \(L\subset E\), su dual respecto de \(B\) es
\[
 L^*=\{y\in E:\langle y,x\rangle_B\in\mathbb Z
                         \text{ para todo }x\in L\}.
\]
\end{definition}

\begin{theorem}[Reciprocidad del dual]
\label{thm:leech-dual}
La colección \eqref{eq:leech-flow} satisface
\begin{equation}
 \Lambda_s^*=\Lambda_{-s}\qquad(s\in\mathbb R).
 \label{eq:leech-dual}
\end{equation}
Su autodualidad en un parámetro \(s\) equivale a
\(D_{2s}\Lambda=\Lambda\).
\end{theorem}
\begin{proof}
Para un operador invertible \(A\),
\[
 (AL)^*=A^{-\dagger}L^*.
\]
En efecto, \(\langle y,Ax\rangle_B=\langle A^\dagger y,x\rangle_B\)
es entero para todo \(x\in L\) precisamente cuando
\(A^\dagger y\in L^*\). Tomando \(A=D_s\), la autoadjunción,
la identidad \(D_s^{-1}=D_{-s}\) y \(\Lambda^*=\Lambda\) prueban
\eqref{eq:leech-dual}. La igualdad
\(D_s\Lambda=D_{-s}\Lambda\) equivale, multiplicando por \(D_s\),
a \(D_{2s}\Lambda=\Lambda\).
\end{proof}

El covolumen controla el volumen de un dominio fundamental; la
autodualidad controla todos los productos enteros con la red.
La primera propiedad es continua a lo largo de esta colección. La
segunda impone la condición aritmética exacta del teorema anterior.
Se empleará por ello la expresión \emph{covolumen uno} para el
flujo euclídeo y se especificará la integralidad cuando intervenga
la unimodularidad aritmética.

\begin{proposition}[Testigo de pérdida local de integralidad]
\label{prop:leech-nonintegral}
Existe \(\delta>0\) tal que, si \(0<|s|<\delta\), la red
\(\Lambda_s\) es no integral y no es isométrica al retículo de Leech.
\end{proposition}
\begin{proof}
El vector \(v_\alpha/3\) pertenece a \(\Lambda\). Su norma
transportada es
\begin{equation}
 f(s)=\left\|D_s\frac{v_\alpha}{3}\right\|_B^2
 =\frac29\left(16e^{2su_1}+
                       \sum_{j=2}^{12}e^{2su_j}\right).
 \label{eq:leech-witness}
\end{equation}
Por \(\sum_j u_j=0\) y por la primera coordenada de
\eqref{eq:leech-u-coordinates},
\begin{equation}
\begin{split}
 f(0)&=6,\\
 f'(0)&=\frac49\left(16u_1+\sum_{j=2}^{12}u_j\right)
       =\frac{20}{3}u_1
       =-825-\frac{466}{3}\sqrt5\ne0.
\end{split}
\label{eq:leech-witness-derivative}
\end{equation}
La continuidad de \(f'\) permite elegir un intervalo
\((-\delta,\delta)\) en el que \(f\) sea estrictamente monótona.
Reduciendo \(\delta\), se obtiene además
\(11/2<f(s)<13/2\). Para \(0<|s|<\delta\), resulta \(f(s)\ne6\);
en ese intervalo, seis es el único entero. Por tanto, \(\Lambda_s\)
contiene un vector de norma no entera. Una red integral tiene
norma entera en todos sus vectores. Puesto que Leech es integral
y la norma se conserva por isometrías, queda demostrada también
la segunda afirmación.
\end{proof}

La deformación preserva volumen y simetría ternaria al tiempo que
modifica efectivamente la estructura métrica aritmética. Los pares
opuestos de coordenadas de \(u\) definen una permutación del espacio
ambiente que intercambia dilataciones y contracciones. Su promoción
a un automorfismo del retículo pegado exigiría conservar el código
y el vecino. La reciprocidad del dual ya demostrada utiliza
únicamente la autoadjunción del flujo y la autodualidad inicial.

\subsection{Convergencia y reciprocidad de las funciones theta}
\label{subsec:leech-theta}

Para \(s\in\mathbb R\) y \(t>0\), definimos
\begin{equation}
 \Theta_s(t)=\sum_{\lambda\in\Lambda}
                \exp\bigl(-\pi t\|D_s\lambda\|_B^2\bigr).
 \label{eq:leech-theta-real}
\end{equation}
Es la suma gaussiana de todas las longitudes de la red deformada,
con sus multiplicidades. Si \(M=\max_j|u_j|\), se tiene
\begin{equation}
 e^{-2|s|M}\|x\|_B^2
 \leq\|D_sx\|_B^2
 \leq e^{2|s|M}\|x\|_B^2.
 \label{eq:leech-quadratic-bounds}
\end{equation}
Sobre un compacto de \(\mathbb R\times(0,\infty)\), la cota
inferior proporciona una gaussiana dominante sobre la red fija.
La serie converge absoluta y uniformemente. Las derivadas de
orden finito añaden factores polinomiales en las coordenadas de
\(\lambda\); dichos factores quedan dominados por otra gaussiana.
Las derivaciones término a término son, por consiguiente,
legítimas en esos compactos.

\begin{theorem}[Reciprocidad gaussiana]
\label{thm:leech-theta}
Con el volumen euclídeo asociado a \(B\),
\begin{equation}
 \Theta_s(t)=t^{-12}\Theta_{-s}(t^{-1}),
 \qquad s\in\mathbb R,\quad t>0.
 \label{eq:leech-theta-reciprocity}
\end{equation}
En particular, \(\Theta_s(1)=\Theta_{-s}(1)\).
\end{theorem}
\begin{proof}
Fijamos la transformada de Fourier
\[
 \widehat f(y)=\int_E
      f(x)e^{-2\pi i\langle x,y\rangle_B}\,dx.
\]
Para una red \(L\subset E\) y
\(f_t(x)=e^{-\pi t\|x\|_B^2}\), periodizamos
\[
 F_t(x)=\sum_{\ell\in L}f_t(x+\ell).
\]
La convergencia normal de la serie y sus derivadas da una función
suave sobre \(E/L\). Si \(\xi\in L^*\), su coeficiente de Fourier es
\[
\begin{split}
 c_\xi
 &=\frac1{\operatorname{covol}(L)}
   \int_{E/L}F_t(x)e^{-2\pi i\langle x,\xi\rangle_B}\,dx\\
 &=\frac{\widehat f_t(\xi)}{\operatorname{covol}(L)}.
\end{split}
\]
La segunda igualdad descompone \(E\) en traslaciones de un dominio
fundamental; el carácter tiene valor uno sobre \(L\). El producto
de las integrales gaussianas unidimensionales, en una base
ortonormal para \(B\), da
\[
 \widehat f_t(\xi)=
 t^{-12}\exp\bigl(-\pi\|\xi\|_B^2/t\bigr).
\]
La serie de Fourier es absolutamente convergente. Evaluándola en
\(x=0\), obtenemos la suma de Poisson
\[
 \sum_{\ell\in L}e^{-\pi t\|\ell\|_B^2}
 =\frac{t^{-12}}{\operatorname{covol}(L)}
      \sum_{\xi\in L^*}e^{-\pi\|\xi\|_B^2/t}.
\]
Ahora \(L=\Lambda_s\) tiene covolumen uno y
\(L^*=\Lambda_{-s}\) por el teorema~\ref{thm:leech-dual}.
La sustitución prueba \eqref{eq:leech-theta-reciprocity}.
El caso \(t=1\) es inmediato.
\end{proof}

La periodización anterior es la demostración gaussiana de Poisson;
puede cotejarse con \cite[teorema 2, pp.~10--11]{Elkies}.
La convención con signo positivo en Fourier conduce a la misma
transformada de la gaussiana, por su paridad. La igualdad de theta
afecta a las series completas y conserva el dominio analítico de
convergencia.

\begin{proposition}[Extensión holomorfa]
\label{prop:leech-theta-complex}
Para \(\tau\) en el semiplano superior
\(\mathbb H=\{\tau\in\mathbb C:\operatorname{Im}\tau>0\}\), la serie
\begin{equation}
 \theta_s(\tau)=\sum_{\lambda\in\Lambda}
             e^{\pi i\tau\|D_s\lambda\|_B^2}
 \label{eq:leech-theta-complex}
\end{equation}
converge normalmente y satisface
\begin{equation}
 \theta_s(-1/\tau)=(-i\tau)^{12}\theta_{-s}(\tau).
 \label{eq:leech-theta-modular}
\end{equation}
\end{proposition}
\begin{proof}
El valor absoluto de cada sumando es
\(e^{-\pi\operatorname{Im}\tau\|D_s\lambda\|_B^2}\).
Sobre compactos del semiplano superior, la parte imaginaria tiene
una cota inferior positiva, de modo que
\eqref{eq:leech-quadratic-bounds} prueba convergencia normal y
holomorfía. Para \(\tau=it\), con \(t>0\),
\eqref{eq:leech-theta-reciprocity} aplicada en \(t^{-1}\) da
\(\theta_s(i/t)=t^{12}\theta_{-s}(it)\).
Ambos miembros de \eqref{eq:leech-theta-modular} son holomorfos
en \(\mathbb H\); el teorema de identidad prolonga la igualdad
desde el eje imaginario positivo.
\end{proof}

La inversión modular relaciona los parámetros \(s\) y \(-s\).
La periodicidad adicional bajo \(\tau\mapsto\tau+1\) utiliza la
paridad de las normas. La proposición~\ref{prop:leech-nonintegral}
exhibe parámetros arbitrariamente próximos a cero para los que
la integralidad euclídea se pierde. Por ello la colección conserva
la reciprocidad analítica demostrada, mientras una realización
mediante formas modulares escalares o álgebras de vértices
reticulares incorpora además las hipótesis aritméticas
correspondientes. En la realización excepcional inicial, esas
hipótesis sostienen la continuación hacia las álgebras de vértices
y Monstrous Moonshine \cite{HMT-IV}.

\subsection{Forma partida y autodualidad en dimensión cuarenta y ocho}
\label{subsec:leech-split}

La pérdida de integralidad de una proyección euclídea conduce a
considerar ambas proyecciones de manera simultánea. En el espacio
\(E\oplus E\) se introduce la forma bilineal
\begin{equation}
 \mathcal B((x,p),(x',p'))=
       \langle x,p'\rangle_B+\langle p,x'\rangle_B.
 \label{eq:leech-split-form}
\end{equation}
Sus productos emparejan las dos coordenadas. Definimos
\begin{equation}
\begin{split}
 \mathcal D_s&=D_s\oplus D_{-s},\\
 \Gamma_s&=\mathcal D_s(\Lambda\oplus\Lambda)\\
 &=\{(D_s\lambda,D_{-s}\mu):\lambda,\mu\in\Lambda\}.
\end{split}
\label{eq:leech-split-lattice}
\end{equation}

\begin{theorem}[Autodualidad integral para la forma partida]
\label{thm:leech-split}
La forma \(\mathcal B\) tiene signatura \((24,24)\).
Para todo \(s\in\mathbb R\), \(\mathcal D_s\) es una isometría de
\(\mathcal B\), y \(\Gamma_s\) es una red integral, par y autodual
respecto de esta forma.
\end{theorem}
\begin{proof}
Introduciendo
\[
 p_L=\frac{x+p}{\sqrt2},\qquad
 p_R=\frac{x-p}{\sqrt2},
\]
se obtiene
\[
 \mathcal B((x,p),(x,p))=\|p_L\|_B^2-\|p_R\|_B^2.
\]
Como \(B\) es positiva en dimensión \(24\), la signatura es
\((24,24)\). La autoadjunción y reciprocidad de los bloques dan
\[
 \langle D_sx,D_{-s}p'\rangle_B=\langle x,p'\rangle_B,
\]
y análogamente para el segundo producto cruzado. Por ello
\(\mathcal B(\mathcal D_s\xi,\mathcal D_s\eta)=\mathcal B(\xi,\eta)\).

En \(\Gamma_0=\Lambda\oplus\Lambda\), todos los productos cruzados
son enteros y
\[
 \mathcal B((\lambda,\mu),(\lambda,\mu))
   =2\langle\lambda,\mu\rangle_B\in2\mathbb Z.
\]
Esto prueba integralidad y paridad. Si \((x,p)\) pertenece al dual
de \(\Gamma_0\) para \(\mathcal B\), el emparejamiento con
\((\lambda,0)\) y \((0,\mu)\) exige respectivamente
\(p,x\in\Lambda^*=\Lambda\). El dual es, por tanto, \(\Gamma_0\).
La isometría \(\mathcal D_s\) transporta estas tres propiedades
a \(\Gamma_s\).
\end{proof}

La métrica especifica el contenido de la conclusión. La colección
\(\Gamma_s\) conserva la integralidad para los productos cruzados
de \(\mathcal B\), incluso cuando \(\Lambda_s\) pierde integralidad
para \(B\). Como retículos abstractos con forma partida, los
\(\Gamma_s\) son isométricos. Sus dos proyecciones euclídeas,
consideradas con las métricas positivas respectivas, varían con
el parámetro.

\begin{proposition}[Intercambio de las proyecciones recíprocas]
\label{prop:leech-split-involution}
La involución \(\mathcal J(x,p)=(p,x)\) satisface
\[
 \mathcal J\Gamma_s=\Gamma_{-s},\qquad
 \mathcal J\mathcal D_s\mathcal J=\mathcal D_{-s}.
\]
En las coordenadas \((p_L,p_R)\), fija \(p_L\) y cambia el signo
de \(p_R\).
\end{proposition}
\begin{proof}
El intercambio convierte
\((D_s\lambda,D_{-s}\mu)\) en
\((D_{-s}\mu,D_s\lambda)\), que recorre \(\Gamma_{-s}\).
La conjugación de los dos bloques da la segunda identidad.
Finalmente, \(x+p\) se conserva y \(x-p\) cambia de signo.
\end{proof}

La involución expresa una dualidad de la construcción reticular
completa. Su eventual representación física conserva, además de
la forma partida, los operadores y los datos de identificación
propios de esa realización.

\subsection{Energía reticular y equivalencia unitaria}
\label{subsec:leech-energy}

La energía de radio recíproco
\(m^2/R^2+w^2R^2\) intercambia sus dos términos bajo
\((m,w,R)\mapsto(w,m,R^{-1})\), con una conjugación unitaria
que preserva los dominios operatorios \cite{HMT-IV}.
La realización sobre los doce planos marcados emplea
\begin{equation}
 \mathcal E_s(\lambda,\mu)=
       \|D_s\lambda\|_B^2+\|D_{-s}\mu\|_B^2,
 \qquad \lambda,\mu\in\Lambda.
 \label{eq:leech-energy}
\end{equation}
La energía es positiva y cumple
\(\mathcal E_s(\lambda,\mu)=\mathcal E_{-s}(\mu,\lambda)\).

\begin{theorem}[Operador positivo y conjugación de dominios]
\label{thm:leech-energy}
En el espacio de Hilbert \(\ell^2(\Lambda\oplus\Lambda)\), sea
\[
 (H_s\psi)(\lambda,\mu)=
       \mathcal E_s(\lambda,\mu)\psi(\lambda,\mu)
\]
con dominio
\begin{equation}
 \mathcal D(H_s)=
 \left\{\psi\in\ell^2(\Lambda\oplus\Lambda):
 \sum_{\lambda,\mu}\mathcal E_s(\lambda,\mu)^2
             |\psi(\lambda,\mu)|^2<\infty\right\}.
 \label{eq:leech-energy-domain}
\end{equation}
Entonces \(H_s\) es positivo, autoadjunto y de resolvente compacto.
La aplicación \(U\psi(\lambda,\mu)=\psi(\mu,\lambda)\) es unitaria y
\begin{equation}
 U\mathcal D(H_s)=\mathcal D(H_{-s}),\qquad
 UH_sU^{-1}=H_{-s}.
 \label{eq:leech-unitary}
\end{equation}
\end{theorem}
\begin{proof}
La multiplicación por una función real en un conjunto discreto,
con el dominio máximo \eqref{eq:leech-energy-domain}, es
autoadjunta. Puede comprobarse directamente mediante la base
ortonormal de los vectores de soporte unitario: la pertenencia
al dominio del adjunto exige que la sucesión obtenida al
multiplicar por \(\mathcal E_s\) sea cuadrado sumable, exactamente
la misma condición. Además,
\[
 \langle\psi,H_s\psi\rangle
   =\sum_{\lambda,\mu}\mathcal E_s(\lambda,\mu)
                         |\psi(\lambda,\mu)|^2\geq0.
\]
Por \eqref{eq:leech-quadratic-bounds},
\[
 \mathcal E_s(\lambda,\mu)\geq
 e^{-2|s|M}\bigl(\|\lambda\|_B^2+\|\mu\|_B^2\bigr).
\]
Una red tiene finitos puntos en cualquier conjunto acotado; por
ello cada subnivel de \(\mathcal E_s\) es finito. Los coeficientes
diagonales \((1+\mathcal E_s)^{-1}\) tienden a cero fuera de
conjuntos finitos, de modo que \((H_s+I)^{-1}\) es límite en norma
de operadores de rango finito y es compacto.

El intercambio de los dos índices conserva la norma de
\(\ell^2(\Lambda\oplus\Lambda)\), por lo que \(U\) es unitario
y \(U^{-1}=U\). La igualdad de energías al invertir \(s\) transforma
la condición de dominio de \(H_s\) en la de \(H_{-s}\), y la
evaluación sobre cada par \((\lambda,\mu)\) prueba
\eqref{eq:leech-unitary}.
\end{proof}

Quedan así articuladas tres conservaciones de distinto dominio.
En los planos euclídeos, el flujo mantiene covolumen y simetría
ternaria y relaciona cada red con el dual del parámetro opuesto.
En el espacio duplicado, los productos cruzados conservan una
red integral, par y autodual. En la realización operatoria, el
intercambio recíproco conserva la energía mediante una equivalencia
unitaria de operadores con sus dominios. Las fórmulas theta
expresan la misma reciprocidad en la suma convergente de las
longitudes de toda la red.

% Procedencia primaria (lectura de 18-09-2026):
% XROOT = output/EXTREMA_MEDIA_RAZON_NARRACION_Y_REVISION_20260916/ARTICULO/ES/source
% XROOT/libro/01_acoplamiento.tex:51-147,727-749.
% XROOT/sections/k_direccion_dimensional.tex:54-155.
% XROOT/nuclear/12.tex:9-88.
% VIROOT = output/AMPLIACION_FORMAL_LEAN_SERIE_20260916/delivery/USB_ES_EN_STAGE10_20260917/ES/PAQUETES/06_PARTICULAS_Y_MASAS_ES/documentacion_original/payload/source_es
% VIROOT/sections/05_registro_operador_masa.tex:11-119,193-294.
% VIROOT/sections/02_particula_persistencia.tex:210-235.
% XROOT/nuclear/07.tex:5-25,74-154: Hamiltoniano de transporte eliptico.
% XROOT/nuclear/11.tex:878-979: elevacion cotangente y accion discreta.
% Estatutos: cadena K->alpha->accion->masa = RESULTADO_RECUPERADO;
% carta 1+1+10 de eventos y criterio de entrelazamiento = FORMALIZACION_NUEVA;
% congruencia, espectro generalizado y control racional = pruebas completas aqui.

\section{Congruencias positivas, lectores de acción y dinámica hamiltoniana}
\label{sec:k-mass-bridge}

La relación entre el registro dodecafásico y la Ley General de Masas tiene
una composición efectiva. APP produce las dos hojas con sus residuos y
cocientes; TRIT orienta la transición y determina su régimen; TPK compone
selección, transporte y actualización conservando ruta, acarreo, frontera,
supervivencia y memoria. Del estado enriquecido y de su prolongación
compatible proceden conjuntamente las representaciones regionales y el
registro de doce posiciones. La clausura numérica y la lectura material
comparten ese origen y conservan sus operaciones específicas.

Con origen de fase, orientación y base mil declarados, el registro es
\[
 K=(234,543,140,729,659,824,621,58,914,794,146,601).
\]
El lector posicional utiliza
\begin{equation}
 \kappa_{\rm ag}(K)=\frac{\sum_{j=1}^{12}K_j1000^{12-j}}{1000^{12}-1},
 \qquad
 \alpha=\pi_{\rm HMT}+e_{\rm HMT}-\varphi_{\rm HMT}-4-
              \kappa_{\rm ag}(K).
 \label{eq:mass-k-alpha}
\end{equation}
Los valores de la derecha son representaciones internas del mismo proceso
coinductivo. La fórmula expresa el lector de clausura infinita compatible,
con la frontera de acarreo que lo define.

La acción recibe esas representaciones mediante los lectores
\begin{align}
 H_5={}&\frac12\sqrt{1000\alpha/\varphi}-\alpha+
       \frac9{16}\alpha^2-\frac59\alpha^3+
       \frac7{48}\alpha^4-\frac1{54}\alpha^5,\notag\\
 D_A={}&\exp(-100\pi\alpha/9),\qquad
 \mathcal C_\pi=169\alpha^6+\frac{D_A\alpha^7}{1-D_A\alpha},\notag\\
 \eta_{\rm ret}={}&H_5-\frac{90}{\pi}\mathcal C_\pi,
 \qquad R_{\rm act}=\frac{H_5}{\eta_{\rm ret}}>0.
 \label{eq:mass-k-action}
\end{align}
Aquí y en las fórmulas siguientes, \(\pi,\varphi,e\) abrevian las
representaciones HMT ya construidas. La norma determinantal local
\(\Sigma_H=\varphi\alpha^{16}\) fija la década \(-34\), y la
sección \(\hbar_{\rm ret}=\eta_{\rm ret}10^{-34}\mathcal U_S\)
pertenece a la recta de acción declarada. La razón \(R_{\rm act}\)
conserva la relación entre las dos secciones de acción al cambiar su base
dimensional común.

La ruta electrónica central proporciona, mediante sus lectores de
desplazamiento y autocorrelación, el triple \((80,54,6)\). Con
\[
 \Delta_4=\frac{\pi-e\log\pi}{270}\left(1+\frac\pi{729}\right),
 \qquad
 b_e^{\rm alg}=\frac{\sqrt3}{4}
       \bigl[\exp(\varphi/\pi^2)-22\alpha^3\bigr](1+15\Delta_4),
\]
su representación material contiene
\begin{equation}
 \kappa_e=80-54\alpha+6\Delta_4,
 \qquad \varepsilon_e^\beta=b_e^{\rm alg}R_{\rm act}^{\kappa_e}.
 \label{eq:mass-k-electron}
\end{equation}
La composición \eqref{eq:mass-k-alpha}--\eqref{eq:mass-k-electron}
identifica la intervención concreta de \(K\): su lectura de clausura
participa en \(\alpha\), y esa salida actúa después en el lector de
acción y en el exponente de retorno. La elección central del electrón
pertenece a la involución de la coordenatización APP, cuyo único punto fijo es
\((5,5)\). La centralidad de esa celda y su evaluación escalar son
dos resultados enlazados por la ruta y su registro.

\subsection{Coordenadas reversibles del registro de eventos materiales}

El registro material asigna a una ruta \(\gamma\) doce eventos
orientados \(\mathbf e(\gamma)=(e_0,\ldots,e_{11})\), uno por cada
ventana nonádica del calendario de ciento ocho posiciones. El mismo
orden temporal permite realizar estos eventos en el portador de doce
coordenadas utilizado por \(K\). Esta identificación requiere conservar
el origen y la orientación del calendario.

Sean \(\mathbf1\) el vector uniforme,
\(P_{11}=I-\mathbf1\mathbf1^*/12\), y \(P_3\) el proyector
tridimensional de la coordenatización incidencial declarada. La construcción
geométrica del registro determina
\[
 u_K=P_3P_{11}K,\qquad
 \|u_K\|^2=\frac{6638585+2275584\sqrt5}{20}>0,
 \qquad u_K\perp\mathbf1,
\]
y el proyector transversal
\[
 P_{10}(K)=P_{11}-\frac{u_Ku_K^*}{\|u_K\|^2}.
\]

\begin{proposition}[Descomposición material relativa a la dirección de \(K\)]
Para todo vector de eventos \(\mathbf e\), la aplicación
\begin{equation}
 \mathbf e\longmapsto
 \left(s_C=\mathbf1^*\mathbf e,\quad
       a_K=\frac{u_K^*\mathbf e}{\|u_K\|^2},\quad
       v_K=P_{10}(K)\mathbf e\right)
 \label{eq:mass-k-event-chart}
\end{equation}
es invertible sobre su imagen y satisface
\begin{equation}
 \mathbf e=\frac{s_C}{12}\mathbf1+a_Ku_K+v_K,
 \qquad
 \|\mathbf e\|^2=\frac{|s_C|^2}{12}
          +|a_K|^2\|u_K\|^2+\|v_K\|^2.
 \label{eq:mass-k-event-inverse}
\end{equation}
\end{proposition}
\begin{proof}
Los proyectores \(\mathbf1\mathbf1^*/12\),
\(u_Ku_K^*/\|u_K\|^2\) y \(P_{10}(K)\) son ortogonales dos a
dos y suman la identidad. Aplicar esa igualdad a \(\mathbf e\)
demuestra la inversión. La ortogonalidad de los tres sumandos da la
identidad de normas.
\end{proof}

Esta coordenatización conserva el vector de eventos completo. Por ello conserva
también los funcionales que dependen de sus signos, sus sumas y sus
correlaciones, una vez realizada la reconstrucción
\eqref{eq:mass-k-event-inverse}. El registro de predecesores de los nueves,
la cuenta de acción y las coordenadas de torre mantienen su residencia
en la trayectoria enriquecida: la coordenatización anterior modifica únicamente
las coordenadas de los doce eventos. El carácter se evalúa sobre la misma
firma recuperada. La construcción proporciona así una representación
geométrica de la información utilizada por la masa, con una dirección
distinguida por \(K\) y diez componentes transversales.

La cantidad de información que retiene cada representación puede probarse
exactamente. En el dominio algebraico de registros, \(K'=K+\mathbf1\)
satisface
\[
 u_{K'}=u_K,\qquad P_{10}(K')=P_{10}(K),\qquad
 \kappa_{\rm ag}(K')-\kappa_{\rm ag}(K)=\frac1{999}.
\]
La última identidad resulta de sumar la progresión geométrica de los
doce pesos posicionales. Si las representaciones regionales se mantienen
fijas, la lectura de \(\alpha\) cambia en \(-1/999\).
Este control algebraico determina el alcance exacto de la dirección:
su proyector conserva una polarización, mientras la clausura escalar
utiliza además la información uniforme y posicional del registro.
La admisibilidad de \(K'\) como otra historia terminal requiere su
propia construcción APP--TRIT--TPK.

\subsection{Congruencia positiva y espectro del operador material}

Para una multisección \(M\), sea \(\mathscr H_M\) el espacio finito
con base ortonormal etiquetada por sus rutas. El lector de retorno
\(r\in\{D,\Omega\}\) produce un operador autoadjunto
\(B_M^r e_\gamma=\kappa_\gamma^r e_\gamma\). El carácter refinado,
evaluado sobre diferencias respecto del electrón, da el operador basal
\[
 M_M^{(0)}e_\gamma=
 m_e^{(0)}\chi_{\rm ref}(\sigma_\gamma-\sigma_e,
                  \eta_\gamma-\eta_e)e_\gamma.
\]
Los coeficientes de torre y su convergencia pertenecen al lector de cada
ruta. Su representación beta es la congruencia
\begin{equation}
 D_M^r=R_{\rm act}^{B_M^r/2},\qquad
 M_M^{\beta,r}=D_M^rM_M^{(0)}D_M^r.
 \label{eq:mass-k-congruence}
\end{equation}

\begin{proposition}[Invariantes exactos de la congruencia material]
El operador \(D_M^r\) es positivo e invertible. La congruencia
\eqref{eq:mass-k-congruence} conserva rango e inercia; en particular,
conserva la positividad. Si \(M_M^{(0)}\) y \(B_M^r\) son diagonales
en la base de rutas, entonces
\[
 m_\gamma^\beta=m_\gamma^{(0)}R_{\rm act}^{\kappa_\gamma^r}.
\]
Para dos rutas en la misma coordenatización dimensional se obtiene
\[
 \frac{m_Y^\beta}{m_X^\beta}=
 \chi_{\rm ref}(\sigma_Y-\sigma_X,\eta_Y-\eta_X)
 R_{\rm act}^{\kappa_Y^r-\kappa_X^r}.
\]
\end{proposition}
\begin{proof}
El cálculo funcional da
\(D_M^r=\exp((\log R_{\rm act})B_M^r/2)\), con espectro
estrictamente positivo. Su inversa es la exponencial del opuesto.
Para todo \(v\), la forma cuadrática transformada vale
\((D_M^rv)^*M_M^{(0)}(D_M^rv)\). El cambio invertible de variable
preserva las dimensiones máximas de subespacios positivos, negativos
y nulos; éstas determinan la inercia y el rango. En una base común
diagonal, los dos factores aportan conjuntamente
\(R_{\rm act}^{\kappa_\gamma^r}\). Dividir dos valores cancela
la sección dimensional común y utiliza la ley multiplicativa del
carácter.
\end{proof}

La distinción entre transporte de una forma y modificación de un
operador respecto de una métrica fija es comprobable sin aproximaciones.
Para
\[
 M_0=\begin{pmatrix}1&0\\0&4\end{pmatrix},\qquad
 D=\begin{pmatrix}2&0\\0&1/2\end{pmatrix},\qquad
 M'=D^*M_0D=\begin{pmatrix}4&0\\0&1\end{pmatrix},
\]
el determinante de \(D\) es uno y los valores asociados a las dos
rutas se intercambian. Incluso el espectro no ordenado cambia si
\(M_0=I\): se transforma de \(\{1,1\}\) en \(\{4,1/4\}\).
Ambos controles preservan positividad, rango e inercia.

\begin{proposition}[Invarianza espectral con métrica transportada]
Sean \(G_0>0\) una métrica y \(D\) invertible. Si
\[
 M'=D^*M_0D,\qquad G'=D^*G_0D,
\]
los pares \((M',G')\) y \((M_0,G_0)\) tienen el mismo espectro
generalizado, con las mismas multiplicidades.
\end{proposition}
\begin{proof}
Se tiene
\[
 \det(M'-\lambda G')=|\det D|^2\det(M_0-\lambda G_0),
 \qquad (G')^{-1}M'=D^{-1}G_0^{-1}M_0D.
\]
La primera igualdad preserva las raíces y sus multiplicidades; la
segunda exhibe una semejanza de los operadores asociados.
\end{proof}

El flujo diagonal positivo inducido por \(K\) posee, por tanto, dos
realizaciones matemáticas precisas sobre una forma material. Cuando
transporta simultáneamente la forma y su métrica describe el mismo
objeto en otra coordenatización. Cuando actúa sobre la forma conservando la métrica
de rutas, produce una colección operatoria que puede cambiar las masas
relativas. Su identificación física se decide por el lector de rutas y
por la acción declarada.

\subsection{Compatibilidad del flujo dodecafásico con el retorno material}

Sea
\[
 A_K=(\log\rho_A)\operatorname{diag}(\widehat u_{K,1},\ldots,
       \widehat u_{K,12}),\quad
 \widehat u_K=u_K/\|u_K\|,\quad \rho_A=\pi/\varphi^2,
\]
de modo que \(g_K(t)=e^{tA_K}\). La colección de retorno material
en la multisección \(M\) tiene generador
\[
 C_M^r=\tfrac12(\log R_{\rm act})B_M^r,\qquad
 S_M^r(t)=e^{tC_M^r}.
\]

\begin{proposition}[Transporte de la congruencia por un lector entrelazante]
Una aplicación lineal declarada
\(J_M:\mathbb C^{12}\to\mathscr H_M\) que cumple
\begin{equation}
 C_M^rJ_M=J_MA_K
 \label{eq:mass-k-intertwiner}
\end{equation}
satisface
\[
 S_M^r(t)J_M=J_Mg_K(t),
\]
y, para \(M_M(t)=S_M^r(t)M_M^{(0)}S_M^r(t)\),
\begin{equation}
 J_M^*M_M(t)J_M
 =g_K(t)^*\bigl(J_M^*M_M^{(0)}J_M\bigr)g_K(t).
 \label{eq:mass-k-pullback}
\end{equation}
\end{proposition}
\begin{proof}
La identidad \eqref{eq:mass-k-intertwiner} se itera a cada potencia
y se suma en la serie exponencial convergente. Tomar adjuntos y
sustituir en la forma material demuestra
\eqref{eq:mass-k-pullback}.
\end{proof}

El criterio es finito y admite una prueba de esfuerzo directamente
operatoria. En las bases diagonales declaradas, si \(a_j\) y
\(c_\gamma\) son los autovalores de \(A_K\) y \(C_M^r\),
respectivamente, se debe cumplir
\[
 (c_\gamma-a_j)(J_M)_{\gamma j}=0
 \qquad\text{para cada }(\gamma,j).
\]
Cada coeficiente no nulo del lector exige la igualdad correspondiente
de autovalores. La composición recuperada de
\eqref{eq:mass-k-alpha}--\eqref{eq:mass-k-electron} y la coordenatización reversible
\eqref{eq:mass-k-event-chart} son resultados efectivos. La identificación
adicional de sus dos flujos requiere un lector \(J_M\) con la propiedad
\eqref{eq:mass-k-intertwiner}; esta proposición determina exactamente
qué debe comprobar una realización que la sostenga.

\subsection{Transporte elíptico y transformación regular de Legendre}

La realización variacional del transporte tiene un antecedente explícito
en la elipse de acción. El par angular ya expresado determina
\(\eta=\operatorname{artanh}(C^*/A)\), en la cámara
\(A\ne0\), \(|C^*/A|<1\). Las coordenadas \((Q,P)\), ambas de
dimensión raíz de acción, realizan la forma
\[
 \mathcal I_\eta(Q,P)=\frac12\bigl(e^{-\eta}Q^2+e^\eta P^2\bigr),
 \qquad \lambda=P\,dQ,\quad \Omega=dQ\wedge dP.
\]
El transporte entre formas, con
\(M_\eta=\operatorname{diag}(e^{\eta/2},e^{-\eta/2})\) y
\(J_2=\left(\begin{smallmatrix}0&-1\\1&0\end{smallmatrix}\right)\),
es \(M_{\eta'}e^{-\vartheta J_2}M_\eta^{-1}\). Una
interpolación \(\eta(t)\), \(\vartheta(t)\), con
\(\dot\vartheta=\omega(t)>0\), tiene el Hamiltoniano recuperado
del artículo X:
\begin{equation}
 H_{\rm tr}(Q,P,t)=\omega(t)\mathcal I_{\eta(t)}(Q,P)
                 +\frac{\dot\eta(t)}2QP.
 \label{eq:mass-k-hamiltonian}
\end{equation}
El segundo término procede de \(\dot M_\eta M_\eta^{-1}\).
Las ecuaciones
\[
 \dot Q=\omega e^\eta P+\frac{\dot\eta}2Q,
 \qquad
 \dot P=-\omega e^{-\eta}Q-\frac{\dot\eta}2P
\]
conservan \(\mathcal I_{\eta(t)}(Q(t),P(t))\): la variación
explícita de \(\eta\) y la contribución del término de conexión
se cancelan. Este resultado y la identidad
\(P\dot Q-H_{\rm tr}=p\dot q-\omega(q^2+p^2)/2\), para
\(Q=e^{\eta/2}q\), \(P=e^{-\eta/2}p\), están demostrados
en la sección de conservación variacional de la obra de extrema y
media razón holográfica.

\begin{proposition}[Lagrangiano regular del transporte elíptico]
Supongamos \(\eta\in C^2\), \(\omega\in C^1\) y \(\omega>0\).
La transformación de Legendre respecto de \(P\) es invertible y
produce
\begin{equation}
 L_{\rm tr}(Q,\dot Q,t)=
 \frac{\bigl(\dot Q-\dot\eta Q/2\bigr)^2}{2\omega e^\eta}
 -\frac{\omega e^{-\eta}}2Q^2.
 \label{eq:mass-k-lagrangian}
\end{equation}
\end{proposition}
\begin{proof}
Escribamos \(a=\dot\eta/2\), \(b=\omega e^\eta>0\) y
\(c=\omega e^{-\eta}\). Entonces
\(H_{\rm tr}=bP^2/2+aQP+cQ^2/2\), y
\[
 \dot Q=\partial_PH_{\rm tr}=bP+aQ,
 \qquad P=\frac{\dot Q-aQ}{b}.
\]
Al sustituir esta expresión en \(P\dot Q-H_{\rm tr}\), los
términos que contienen \(P\) se reúnen en
\(P(\dot Q-aQ)-bP^2/2=(\dot Q-aQ)^2/(2b)\), lo cual
demuestra \eqref{eq:mass-k-lagrangian}. Además,
\[
 \partial_{\dot Q}L_{\rm tr}=P,
 \qquad
 \partial_P^2H_{\rm tr}=\omega e^\eta>0,
 \qquad
 \partial_{\dot Q}^2L_{\rm tr}=\frac1{\omega e^\eta}>0.
\]
La transformación es regular y sus ecuaciones de Euler--Lagrange
equivalen a las ecuaciones hamiltonianas anteriores.
\end{proof}

Esta eliminación algebraica reúne el Hamiltoniano ya construido con
su forma lagrangiana de segundo orden. La regularidad procede del
término cuadrático de la realización elíptica. Para un levantamiento
cotangente puramente lineal \(H_K(q,p)=p^{\mathsf T}A_Kq\), la
ecuación \(\dot q=A_Kq\) es independiente de \(p\) y el Hessiano
\(\partial_p^2H_K\) es nulo. Su formulación variacional utiliza la
acción de primer orden \(\int p^{\mathsf T}(\dot q-A_Kq)\,dt\).
El corpus contiene también su antecedente discreto
\(L_n=p_{n+1}^{\mathsf T}(q_{n+1}-U_nq_n)\), cuyas ecuaciones
producen el transporte \(q_{n+1}=U_nq_n\) y el transporte dual
\(p_{n+1}=U_n^{-\mathsf T}p_n\). Las dos formulaciones poseen
dominios y operadores determinados. La interpolación temporal y la
sección de acción se conservan como parte de la realización física
declarada.

% Incorporación aditiva. No modifica las fuentes ni el PDF de 211 páginas.
% Procedencia, dependencias y alcance en preservation_map.json.
\section{Representación grassmanniana de los lectores y formas con valores en la recta de acción}
\label{sec:compat-lectores-formas}

Las realizaciones positivas conservan información que puede volver a ser
utilizada por los lectores físicos. Esta conservación tiene una formulación
precisa: una representación marcada del registro permite reconstruir sus
coordenadas y transportar, por composición, cada función que las utiliza.
La geometría representa un resultado de la construcción; su elección de
coordenadas conserva la procedencia de ese resultado.

El dominio de partida está formado por los estados producidos por APP,
TRIT y TPK en las secciones precedentes. APP conserva sus hojas aditiva y
multiplicativa mediante residuo y cociente; TRIT selecciona el régimen
local con su orientación; la composición
\(U_t=\operatorname{Upd}_t\circ\operatorname{Tra}_t\circ
\operatorname{Sel}_t\) actualiza posición, registro y memoria.
La prolongación
\[
 w_6\longrightarrow w_{12}\longrightarrow w_{18}\longrightarrow
 w_{24}\longrightarrow w_{30}\longrightarrow R_{36}\longrightarrow G_9
\]
conserva el prefijo y distingue retorno de fase de retorno del estado.
Las realizaciones espectral, solenoidal, gauge, cohomológica y determinantal
permanecen como operaciones correlativas de la estructura discreta
conjunta del continuo. La extracción del registro y las representaciones
regionales se realizan en ese mismo dominio. En lo que sigue, todas ellas
son salidas HMT previamente construidas.

\subsection{La colección marcada y su inversa}

Sea \(L\geq2\), y sea \(K=(K_1,\ldots,K_L)\) un registro positivo
con carga \(Q=\sum_jK_j\). En la aplicación dodecafásica,
\(L=12\) y \(Q=6263\). Definimos
\[
 d_j=K_j/Q,\qquad t_0=0,\qquad
 t_j=\sum_{a=1}^{j}d_a,\qquad
 B(K)=\begin{pmatrix}1&\cdots&1\\t_0&\cdots&t_L\end{pmatrix}.
\]
La colección \(\mathcal G_L^{\rm aff}\) es la imagen de estos planos
de filas en \(\operatorname{Gr}_{>0}(2,L+1)\), conservando las
columnas ordenadas \(0,\ldots,L\), las marcas extremas y la carga
\(Q\). La normalización afín fija \(t_0=0,t_L=1\).
Escribimos \(\Delta_{ij}\) para el menor de las columnas \(i,j\)
de cualquier representante del plano.

\begin{proposition}[Inversa de la realización marcada]
\label{prop:compat-inversa-marcada}
La aplicación \(K\mapsto([B(K)],Q)\) es inyectiva. Su inversa
sobre la imagen está dada por
\begin{equation}
 K_j=Q\frac{\Delta_{j-1,j}}{\Delta_{0L}},
 \qquad
 t_j=\frac{\Delta_{0j}}{\Delta_{0L}},
 \quad 1\leq j\leq L.
 \label{eq:compat-inversa-plucker}
\end{equation}
Los cocientes son invariantes bajo cualquier cambio invertible de base
de las dos filas.
\end{proposition}
\begin{proof}
En el representante declarado,
\(\Delta_{ij}=t_j-t_i\), y \(\Delta_{0L}=1\).
Las fórmulas recuperan las diferencias consecutivas y sus sumas.
Si se sustituye \(B\) por \(AB\), con \(A\) invertible,
todos los menores se multiplican por \(\det A\); sus cocientes
permanecen iguales. La carga recupera la escala anterior a la
normalización. Por tanto, dos registros con iguales datos marcados
tienen las mismas coordenadas.
\end{proof}

La condición de pertenencia a \(\mathcal G_L^{\rm aff}\) importa.
Los cocientes consecutivos de un plano positivo arbitrario pueden
incumplir la identidad \(\sum_j\Delta_{j-1,j}=\Delta_{0L}\).
En la colección construida esa identidad es telescópica. La inversa actúa
en la imagen marcada que acaba de definirse.

\subsection{Álgebra de lectores y cambios de coordenatización}

Sea \(\mathcal D\) la imagen de una representación del estado enriquecido
de la forma \(x\mapsto(K(x),\xi(x))\). El símbolo \(\xi\)
reúne exactamente las restantes coordenadas utilizadas por el lector:
representaciones regionales, origen y orientación, ruta, firma, condiciones
de frontera y secciones dimensionales, según el objeto considerado.
No se supone que estas coordenadas sean independientes. Definimos
\[
 \mathcal F:\mathcal D\longrightarrow\widetilde{\mathcal G},
 \qquad (K,\xi)\longmapsto([B(K)],Q,\xi),
 \qquad \widetilde{\mathcal G}:=\mathcal F(\mathcal D).
\]
La proposición anterior proporciona la inversa \(\mathcal I\) de
\(\mathcal F\) en su imagen.

\begin{theorem}[Representación fiel del álgebra de lectores retenidos]
\label{thm:compat-algebra-lectores}
Para cada lector \(\mathcal R:\mathcal D\to V\), con codominio
declarado \(V\), existe un único lector geométrico
\(\mathcal R^{\rm geom}:\widetilde{\mathcal G}\to V\) tal que
\[
 \mathcal R^{\rm geom}=\mathcal R\circ\mathcal I,
 \qquad \mathcal R^{\rm geom}\circ\mathcal F=\mathcal R.
\]
Para lectores escalares, esta correspondencia preserva sumas, productos,
unidad y conjugación. Es un isomorfismo del álgebra de funciones sobre
\(\mathcal D\) con el álgebra correspondiente sobre
\(\widetilde{\mathcal G}\).
\end{theorem}
\begin{proof}
La composición está definida porque \(\mathcal I\mathcal F\) es
la identidad. Si dos lectores geométricos tienen la misma composición
con \(\mathcal F\), coinciden en su imagen, que es todo
\(\widetilde{\mathcal G}\). Las operaciones algebraicas se verifican
punto a punto; por ejemplo,
\((\mathcal R_1\mathcal R_2)\circ\mathcal I=
(\mathcal R_1\circ\mathcal I)(\mathcal R_2\circ\mathcal I)\).
La composición inversa es la restricción por \(\mathcal F\).
\end{proof}

La fidelidad se refiere a la representación \((K,\xi)\). Para un
observable \(\mathcal O\) del estado enriquecido completo, el criterio
de descenso es que \(\mathcal O(x)=\mathcal O(x')\) siempre que
\((K(x),\xi(x))=(K(x'),\xi(x'))\). La necesidad se obtiene por
composición; la suficiencia permite definir el valor en una fibra
mediante cualquiera de sus representantes. Este criterio conserva
explícitamente qué información debe acompañar al registro.

Sea ahora \(\mathcal T\) una triangulación del polígono etiquetado
que determina una coordenatización de Plücker positiva. Se incluyen sus coordenadas
de frontera y una normalización homogénea, por ejemplo
\(\Delta_{0L}=1\). Si se intercambia la diagonal \(ac\) por
\(bd\), con \(a<b<c<d\), el cambio de coordenatización es
\begin{equation}
 \Delta_{bd}=
 \frac{\Delta_{ab}\Delta_{cd}+\Delta_{ad}\Delta_{bc}}{\Delta_{ac}}.
 \label{eq:compat-cambio-carta}
\end{equation}
El denominador es positivo y la relación de Plücker prueba que ambas
coordenatizaciones evalúan el mismo punto.

\begin{corollary}[Compatibilidad de los lectores con las mutaciones]
\label{cor:compat-mutaciones}
Sean \(r_{\mathcal T}\) y \(r_{\mathcal T'}\) las expresiones
de \(\mathcal R^{\rm geom}\) en dos coordenatizaciones que conservan las marcas,
y \(\mu_{\mathcal T\mathcal T'}\) su cambio positivo. Entonces
\[
 r_{\mathcal T'}\circ\mu_{\mathcal T\mathcal T'}
 =r_{\mathcal T}
\]
en el dominio común. Una sucesión cerrada de cambios de coordenatización que
conserva el punto marcado conserva todos estos lectores.
\end{corollary}
\begin{proof}
Las dos reconstrucciones mediante
\eqref{eq:compat-inversa-plucker} recuperan el mismo \(K\).
Las coordenadas \(\xi\) se transportan con sus marcas. Aplicar
\(\mathcal R\) prueba la igualdad; iterarla prueba la afirmación
sobre una sucesión cerrada.
\end{proof}

Los lectores pueden expresarse en coordenatizaciones distintas y seguir siendo
funciones del mismo objeto. Una permutación de las etiquetas actúa
adicionalmente sobre la marca posicional; una transformación que cambia
el punto actúa sobre el argumento del lector. Ambas operaciones tienen
su ley de transporte específica.

La acción diagonal por columnas proporciona un control explícito.
Para \(\lambda_i>0\),
\(\Delta_{ij}\mapsto\lambda_i\lambda_j\Delta_{ij}\).
La razón doble \(\Delta_{ab}\Delta_{cd}/
(\Delta_{ad}\Delta_{bc})\) se conserva, mientras
\(\Delta_{j-1,j}/\Delta_{0L}\) recibe el factor
\(\lambda_{j-1}\lambda_j/(\lambda_0\lambda_L)\).
Con \(t=(0,1/4,1/2,1)\) y \(\lambda=(1,2,1,1)\), las
separaciones reconstruidas serían \((1/2,1/2,1/2)\), cuya suma
es \(3/2\). Las razones dobles solas conservan un cociente diferente
del requerido por la inversa marcada.

\subsection{Composición explícita con acción y respuesta constitutiva}

En la coordenatización dodecafásica, escribimos \(b=1000\) y
\begin{equation}
 \kappa^{\rm geom}=
 \frac{Q}{b^{12}-1}\sum_{j=1}^{12}
 b^{12-j}\frac{\Delta_{j-1,j}}{\Delta_{0,12}},
 \qquad
 \alpha^{\rm geom}=\pi_{\rm HMT}+e_{\rm HMT}
 -\varphi_{\rm HMT}-4-\kappa^{\rm geom}.
 \label{eq:compat-alpha-geometrica}
\end{equation}
Son las expresiones geométricas del lector de clausura
\eqref{eq:mass-k-alpha}. Las tres representaciones regionales y \(K\)
comparten generación coinductiva; la expresión utiliza su orden
constructivo interno. El lector de acción
\eqref{eq:mass-k-action}, evaluado en \(\alpha^{\rm geom}\),
determina \(\eta_{\rm ret}^{\rm geom}\), y por tanto
\[
 \hbar_{\rm ret}^{\rm geom}
 =\eta_{\rm ret}^{\rm geom}10^{-34}\mathcal U_S
 =\hbar_{\rm ret}.
\]
La igualdad tiene lugar en la recta de acción \(\mathcal L_S\)
con su sección \(\mathcal U_S\); conserva la década procedente
del lector determinantal \(\varphi_{\rm HMT}\alpha^{16}\).
La composición no modifica la sección ni la aplica una segunda vez.

En la coordenatización angular levantada que conserva el origen y la rama real
de la representación, el par angular posterior determina
\[
 x=\frac{\pi_{\rm HMT}}{180}(1000\alpha),\qquad
 y=\frac{\pi_{\rm HMT}}{90}(\eta_{\rm ret}+\alpha),\qquad
 q_\pm=e^{-x\mp y}.
\]
La rama forma parte de las marcas retenidas: la reducción de un ángulo
módulo \(360\) por sí sola conserva otra información. Las igualdades
que siguen emplean el levantamiento declarado, en el cual la primera
coordenada es \(1000\alpha\).
En la cámara \(x>|y|\), con proyectores de hoja fijados
\(P_\pm\), sea \(T=q_+P_++q_-P_-\). La respuesta del corpus
compara las incidencias \(90\) y \(120\) mediante
\[
 \mathcal V=(I-T^{90})(I-T^{120})^{-1}
 =r_+P_++r_-P_-,\qquad
 r_\pm=f(q_\pm^{30}),\quad
 f(s)=\frac{1+s+s^2}{1+s+s^2+s^3}.
\]
La regla satisface
\(\mathcal V\Sigma_4(T^{30})=\Sigma_3(T^{30})\), donde
\(\Sigma_j(S)=\sum_{a=0}^{j-1}S^a\). Su unicidad se deduce
de la invertibilidad de \(\Sigma_4(T^{30})\).
Las cuatro lecturas normalizadas son
\[
 \widehat\varepsilon=r_+^2,\quad
 \widehat\mu=r_-^2,\quad
 \widehat Z=r_-/r_+,\quad
 \widehat c=(r_+r_-)^{-1}.
\]
En las rectas dimensionales de acción, carga y velocidad se realizan como
\[
 Z=\widehat Z\,u_Su_Q^{-2},\quad c=\widehat c\,u_C,
 \qquad
 \mu=\widehat\mu\,u_Su_C^{-1}u_Q^{-2},\quad
 \varepsilon=\widehat\varepsilon\,u_S^{-1}u_C^{-1}u_Q^2.
\]
Las identidades \(\mu\varepsilon=c^{-2}\) y
\(\mu/\varepsilon=Z^2\) se obtienen por producto y cociente.
El corolario de cambios de coordenatización conserva esta respuesta completa,
con las marcas de hoja y las secciones dimensionales incluidas en
\(\xi\). El mismo argumento se aplica a un lector de masa con su
ruta, firma y operador de retorno retenidos, en el dominio de la
sección~\ref{sec:k-mass-bridge}.

\begin{theorem}[Recuperación del registro desde la respuesta etiquetada]
\label{thm:compat-recuperacion-constitutiva}
Fijadas las representaciones regionales, la orientación de las hojas, la
base \(b=1000\), el origen, la longitud doce y el levantamiento angular
declarado, la aplicación
\(K\mapsto(\widehat\varepsilon,\widehat\mu)\) es inyectiva
sobre los registros enteros de su imagen física que satisfacen
\(0\leq K_j<b\) y la cámara contractiva indicada. Su inversa exacta
es la composición de las siguientes operaciones:
\begin{align*}
 r_+&=\sqrt{\widehat\varepsilon},&
 r_-&=\sqrt{\widehat\mu},\\
 r_\pm s_\pm^3+(r_\pm-1)(1+s_\pm+s_\pm^2)&=0,
 &0<s_\pm<1,\\
 x&=-\frac1{60}\log(s_+s_-),&
 \alpha&=\frac{180x}{1000\pi_{\rm HMT}},\\
 \kappa&=\pi_{\rm HMT}+e_{\rm HMT}-\varphi_{\rm HMT}-4-\alpha,
 &N&=(b^{12}-1)\kappa,\\
 K_j&=\left\lfloor N/b^{12-j}\right\rfloor\bmod b.
\end{align*}
\end{theorem}
\begin{proof}
Las raíces positivas recuperan los coeficientes etiquetados.
La derivada
\[
 f'(s)=-\frac{s^2(3+2s+s^2)}{(1+s+s^2+s^3)^2}<0
 \qquad(0<s<1)
\]
y sus límites \(1\) y \(3/4\) prueban que la ecuación cúbica
posee exactamente una raíz en ese intervalo. Puesto que
\(s_\pm=e^{-30x\mp30y}\), su producto es \(e^{-60x}\).
Se recuperan \(x\), \(\alpha\) y \(\kappa\).
En la imagen considerada, \(N=\sum_jK_jb^{12-j}\) es un entero
entre cero y \(b^{12}-1\). La división euclídea sucesiva recupera
de manera única sus doce cifras, incluidos los ceros iniciales.
\end{proof}

La precisión necesaria para una recuperación numérica también puede
expresarse. Con representaciones regionales exactas, un error
\(|\widetilde\alpha-\alpha|<1/[2(b^{12}-1)]\) implica
\(|\widetilde N-N|<1/2\); el redondeo al entero más próximo
recupera \(N\). Si las representaciones regionales tienen errores, se
aplica la misma condición a la suma de errores de
\(\pi+e-\varphi-\alpha\). La inyectividad exacta y la resolución
metrológica son así propiedades diferenciadas.

El dominio entero es esencial. En una prolongación analítica de los
registros positivos, la perturbación de tres posiciones consecutivas
proporcional a \((1,-(b+1),b)\) conserva tanto la suma como la
lectura posicional, pues
\(1-(b+1)+b=0\) y \(b^2-b(b+1)+b=0\).
Para perturbaciones suficientemente pequeñas conserva la positividad
y cambia el punto grassmanniano. La inyectividad precedente utiliza
la condición de doce cifras enteras; la colección analítica continua
posee otras fibras.

\subsection{Refinamiento y lectores heredados}

Si cada separación \(d_j\) se subdivide en separaciones positivas
\(d_{j,a}\) con \(\sum_ad_{j,a}=d_j\), los extremos originales
conservan su posición. Sea \(\rho\) la operación de sumar los sucesores
de cada intervalo y \(H\) la inclusión de sus extremos en el nuevo
conjunto ordenado. Entonces
\[
 \mathcal I_{\rm heredada}(B(\widetilde d),Q)
 =\bigl(Q\sum_ad_{j,a}\bigr)_j=K,
 \qquad
 \mathcal R_{\rm heredada}=\mathcal R(K,\xi).
\]
La igualdad prueba la conmutatividad entre reconstrucción heredada y
subdivisión. En una sucesión de subdivisiones, las sumas se asocian;
el resultado persiste en cada composición finita. Esto conserva el
lector de longitud doce y sus marcas. Un lector definido sobre los
nuevos bloques tendrá su propio dominio y su ley de comparación.

\subsection{Formas canónicas con valores en una recta dimensional}

Sea \(P\) uno de los politopos de rango uno construidos en el
manuscrito y \(P=\bigcup_\sigma P_\sigma\) una subdivisión
orientada del mismo soporte, con multiplicidad uno y caras compatibles.
Escribimos \(\Omega_\sigma\) y \(\Omega_P\) para sus formas
canónicas. Sus identidades son
\(\Omega_P=\sum_\sigma\Omega_\sigma\) y
\(\operatorname{Res}_F\Omega_P=\Omega_F\), con la convención
de orientación fijada en el desarrollo amplituhedral.
La recta de acción \(\mathcal L_S\) se considera constante sobre
este politopo; su sección procede del estado HMT fijado.

\begin{theorem}[Transporte de la forma por una sección de acción]
\label{thm:compat-forma-accion}
Para \(a\in\mathcal L_S\), constante respecto de las coordenadas
del politopo,
\[
 \boldsymbol\Omega_P:=a\otimes\Omega_P
 =\sum_\sigma a\otimes\Omega_\sigma,
 \qquad
 \operatorname{Res}_F\boldsymbol\Omega_P=a\otimes\Omega_F.
\]
Las identidades se conservan bajo subdivisiones sucesivas y residuos
iterados sobre banderas de caras. Se aplican, en particular, a
\(a=\hbar_{\rm ret}\).
\end{theorem}
\begin{proof}
El producto tensorial con una recta es lineal. El residuo actúa sobre
el factor diferencial y conmuta con el coeficiente constante. La
compatibilidad de subdivisiones y de residuos iterados se transporta
por las mismas operaciones lineales. En un cambio de base dimensional
\(u'_S=\lambda\,u_S\), con \(\lambda>0\), el coeficiente numérico
cambia por \(\lambda^{-1}\);
el tensor y sus residuos permanecen iguales.
\end{proof}

\begin{theorem}[Condición de pegado de los coeficientes de frontera]
\label{thm:compat-pegado}
Sean \(a_\sigma\) secciones de la misma recta, regulares en un
entorno de las caras y meromorfas en una realización ambiental común.
Supóngase que sus productos con las formas canónicas tienen, a lo
largo de las caras consideradas, polos a lo sumo simples. En el
interior relativo de una cara común \(F\) de codimensión uno de dos celdas
\(\sigma,\tau\), con orientaciones opuestas,
\begin{equation}
 \operatorname{Res}_F\left(a_\sigma\Omega_\sigma+a_\tau\Omega_\tau\right)
 =\left(a_\sigma|_F-a_\tau|_F\right)\Omega_F.
 \label{eq:compat-defecto-frontera}
\end{equation}
Por tanto, el polo de esta pareja sobre el divisor de \(F\) desaparece
exactamente cuando las dos trazas coinciden. Si los coeficientes son
constantes y el grafo de adyacencia de las celdas es conexo, la cancelación
de cada pareja de celdas a través de su cara común equivale a la existencia
de un único coeficiente común \(a\).
\end{theorem}
\begin{proof}
Con coordenada local transversal \(z\), una forma logarítmica se
escribe \(dz/z\wedge\omega+\eta\), donde \(\eta\) es regular
respecto de ese divisor. La regularidad de \(a_\sigma\) da
\(\operatorname{Res}_F(a_\sigma\Omega_\sigma)
=a_\sigma|_F\operatorname{Res}_F\Omega_\sigma\).
Las dos orientaciones producen \(\Omega_F\) y \(-\Omega_F\).
La forma \(\Omega_F\) es no nula en el interior relativo de la
cara; la igualdad de residuos equivale a la igualdad de las trazas.
Para un polo simple, la anulación del residuo elimina el polo.
En el caso constante, la igualdad se propaga a lo largo de cada
arista del grafo conexo de adyacencia.
\end{proof}

Para la suma global deben contarse todos los términos con un polo en el
mismo divisor ambiental. Si \(H\) es el hiperplano que contiene \(F\),
y \(I_H\) reúne esos términos, la identidad general es
\[
 \operatorname{Res}_{H}\left(\sum_\nu a_\nu\Omega_\nu\right)
 =\sum_{\nu\in I_H}a_\nu|_H\operatorname{Res}_H\Omega_\nu,
\]
siempre que los coeficientes sean regulares sobre \(H\). Una faceta
coplanar, aunque no sea adyacente a la cara real \(F\), participa en esta
suma. La fórmula de dos términos representa el residuo global cuando
los restantes términos son regulares sobre ese divisor. La cancelación
por parejas de todas las facetas interiores elimina sus contribuciones;
las facetas exteriores conservan los residuos de frontera correspondientes.

En una cara de codimensión superior, se aplica esta fórmula sucesivamente
con las orientaciones inducidas. Los signos son los de la incidencia
celular y satisfacen la cancelación de borde de borde. La colección de
lectores locales define un pegado compatible por parejas cuando sus
trazas coinciden después del transporte a la misma recta de coeficientes.
El residuo global se obtiene sumando las contribuciones de cada divisor.
Este enunciado da un criterio efectivo para distribuir una lectura
física sobre coordenatizaciones o celdas sin introducir discontinuidades espurias.

\subsection{Cancelación de polos y normalización global}

La cancelación interior determina una condición local. La identificación
con una forma canónica global requiere asimismo el divisor de polos y
los residuos exteriores normalizados. En espacio proyectivo, dos
formas racionales de grado máximo con los mismos polos logarítmicos
y los mismos residuos difieren por una forma holomorfa de grado máximo;
\(H^0(\mathbb P^m,\Omega^m)=0\) demuestra su igualdad.
Ésta es la condición global utilizada en la prueba del politopo
\cite{positive}.

Un ejemplo intervalar muestra por qué los coeficientes variables
requieren dicha condición. Para \(0<c<1\), sean
\[
 \Omega_{0c}=\frac{c\,dx}{x(c-x)},\qquad
 \Omega_{c1}=\frac{(1-c)\,dx}{(x-c)(1-x)},\qquad
 \Omega_{01}=\frac{dx}{x(1-x)}.
\]
Con \(a_1(x)=1+x(c-x)\) y \(a_2(x)=1\), sus trazas coinciden
en todos los extremos pertinentes, y sin embargo
\[
 a_1\Omega_{0c}+a_2\Omega_{c1}=\Omega_{01}+c\,dx.
\]
El término adicional es regular en la coordenatización afín y tiene un polo en
el infinito proyectivo. Los residuos de los extremos finitos y la
cancelación en \(c\) se conservan, mientras la forma resultante
tiene otro comportamiento global. Multiplicar el ejemplo por una
sección de \(\mathcal L_S\) produce el mismo control dimensional.

La consecuencia de conjunto es precisa: la representación grassmanniana
marcada transporta los lectores retenidos; las mutaciones conservan sus
valores cuando sólo cambian la coordenatización; el refinamiento conserva sus
evaluaciones heredadas; y las formas con valores en la recta de acción
se ensamblan mediante igualdad de trazas y control de polos. Cada
igualdad expresa la operación que conserva el objeto correspondiente.

\section{Incidencia hexádico--octádica, operador de área e información sectorial}
\label{rev4:ae:seccion}

La recuperación geométrica del registro adquiere contenido físico cuando se
compone con observables cuya construcción conserva las incidencias y las
condiciones de estado utilizadas. El paso de una héxada a su óctada marcada
proporciona un caso preciso: la ampliación del soporte transforma dos
espacios de incidencia, determina sus multiplicidades y permite conservar
esa distinción en un operador de área y en el logaritmo de un estado
reducido. La exposición siguiente reúne esa construcción del segundo
artículo de la serie~\cite{hmtII} con la representación grassmanniana
marcada de la sección precedente.

El punto de partida permanece en el estado enriquecido producido por
APP--TRIT--TPK. Las dos operaciones de APP conservan sus divisiones
completas, TRIT determina régimen y orientación, y la composición
\(U_t=\operatorname{Upd}_t\circ\operatorname{Tra}_t\circ
\operatorname{Sel}_t\) transporta y actualiza la información de posición,
acarreo y memoria. La prolongación compatible y la holonomía nonádica
conservan conjuntamente las realizaciones aritmética e incidencial de la
estructura discreta del continuo. El registro \(K\), las coordenadas
regionales \(\pi_{\rm HMT},\varphi_{\rm HMT},e_{\rm HMT}\), la
coordenada \(\alpha_{\rm HMT}\) y la configuración excepcional proceden
de esa construcción común. Aquí se utilizan sus resultados y las marcas
que permiten prolongarlos hacia los observables de frontera.

\subsection{Elevación de incidencia y diferencia de multiplicidades}

Sea \(\mathcal H\) el sistema de héxadas de Witt sobre las doce posiciones
del registro. Fijemos la coordenatización de incidencia \(\ell:[12]\to\mathscr D\),
donde \(\mathscr D\) es una dodécada del código binario de Golay extendido
\(G_{24}\), de parámetros \([24,12,8]_2\). La coordenatización transporta las
héxadas de \(\mathcal H\) al diseño residual sobre \(\mathscr D\).
Este dato de realización conserva la configuración marcada determinada
en el desarrollo excepcional. Los soportes de peso ocho de \(G_{24}\)
constituyen las óctadas del sistema \(S(5,8,24)\).

\begin{proposition}[Elevación única de una héxada residual]
\label{rev4:ae:elevacion}
Los conjuntos \(O\cap\mathscr D\), con \(O\) una óctada y
\(|O\cap\mathscr D|=6\), forman un sistema \(S(5,6,12)\).
Cada héxada de ese sistema es la intersección de \(\mathscr D\) con
una única óctada. En particular, la héxada marcada \(H\subset\mathscr D\)
tiene una elevación determinada \(O_H\) y una inclusión \(H\subset O_H\).
\end{proposition}
\begin{proof}
Los pesos del código son \(0,8,12,16,24\). Si
\(j=|O\cap\mathscr D|\), la suma binaria de los vectores de soporte
\(O\) y \(\mathscr D\) tiene peso \(20-2j\); por tanto,
\(j\in\{2,4,6\}\). Toda quíntupla \(T\subset\mathscr D\)
pertenece a una única óctada. Como su intersección con \(\mathscr D\)
contiene \(T\), esa intersección tiene seis elementos. Se obtiene así
una única héxada residual sobre cada quíntupla. Si dos óctadas elevaran
la misma héxada, ambas contendrían cualquiera de sus quíntuplas y
coincidirían por la propiedad de Steiner.
\end{proof}

En lo que sigue, \(\binom H2\) denota el conjunto de subconjuntos de dos
elementos de \(H\). Conservamos ese segundo factor al ampliar el primero:
\begin{equation}
 \mathcal F_6=H\times\binom H2,
 \qquad
 \mathcal F_8=O_H\times\binom H2,
 \qquad
 \jmath(h,P)=(h,P).
 \label{rev4:ae:incidencias}
\end{equation}
Un elemento de estos espacios es una incidencia punto--pareja. El punto
distinguido recorre primero la héxada y después su óctada, mientras la
pareja sigue perteneciendo a la héxada de origen. Esta especificación
fija la operación responsable del incremento.

\begin{proposition}[Ampliación exacta del espacio de incidencias]
\label{rev4:ae:incremento}
La aplicación \(\jmath\) es inyectiva y su complemento es
\((O_H\setminus H)\times\binom H2\). En consecuencia,
\begin{equation}
 |\mathcal F_6|=6\binom62=90,
 \qquad |\mathcal F_8|=8\binom62=120,
 \qquad |\mathcal F_8\setminus\jmath(\mathcal F_6)|=30.
 \label{rev4:ae:cardinales}
\end{equation}
\end{proposition}
\begin{proof}
La inclusión \(H\subset O_H\) conserva ambas coordenadas y prueba la
inyectividad. La partición disjunta \(O_H=H\sqcup(O_H\setminus H)\)
induce la partición del producto cartesiano. Como
\(|O_H\setminus H|=2\) y \(\binom62=15\), el complemento tiene
\(2\cdot15=30\) elementos.
\end{proof}

Las dos multiplicidades se obtienen antes de cualquier distribución
probabilística sobre la frontera. La diferencia normalizada
\((90-120)/(24\binom62)=-1/12\) conserva el signo de la comparación
orientada. Su interpretación depende del espacio de incidencias en que
se ha efectuado la operación.

\subsection{Transporte colectivo y parámetro de Barbero--Immirzi}

La representación de frontera retiene las etiquetas de ambos sectores.
Para explicitar su soporte, escribamos
\(x_j=r_{2j-1}+3r_{2j}\), con \(r_k\in\{0,1,2\}\), en tres
coordenadas de \(\mathbb Z/9\mathbb Z\). El reagrupamiento posicional
\begin{equation}
 \beta_{729}(x_1,x_2,x_3)
 =\bigl(r_1+3r_2+9r_3,\ r_4+3r_5+9r_6\bigr)
 \label{rev4:ae:reagrupamiento}
\end{equation}
es una biyección hacia \((\mathbb Z/27\mathbb Z)^2\). La unicidad de
la expansión ternaria prueba que su inversa extrae los seis dígitos y
restituye las tres parejas iniciales. Se preserva el orden posicional
con su información de acarreo; las respectivas adiciones se mantienen
definidas en sus propios grupos.

En el bloque \(B_\partial=\{0,\ldots,8\}^2\) del toro de orden
veintisiete se distinguen los enlaces orientados
\begin{align}
 E_{90}&=\{((-1,j),(0,j)),\ ((8,j),(9,j)):0\le j\le8\},\\
 E_{120}&=\{((i,-1),(i,0)),\ ((i,8),(i,9)):0\le i\le8\}.
 \label{rev4:ae:enlaces}
\end{align}
Las coordenadas se interpretan módulo veintisiete. Los dos conjuntos
son disjuntos y contienen dieciocho enlaces cada uno: dos lados de
nueve enlaces por dirección. Las etiquetas \(90\) y \(120\) proceden
de los espacios de incidencia; los cardinales de los conjuntos de
enlaces son \(18\) y \(18\).

Sean \(u_m\) los vectores uniformes normalizados de
\(\ell^2(\mathcal F_6)\oplus\ell^2(\mathcal F_8)\), con
\(m=90,120\), y sean \(v_m\) los correspondientes vectores
uniformes normalizados de \(\ell^2(E_{90})\oplus\ell^2(E_{120})\).
Por ejemplo, \(u_{90}=90^{-1/2}\sum_{f\in\mathcal F_6}(\delta_f,0)\)
y \(v_{90}=18^{-1/2}\sum_{e\in E_{90}}(\delta_e,0)\).
Cada pareja es ortonormal.

\begin{proposition}[Conservación de la estructura espectral colectiva]
\label{rev4:ae:colectivo}
La aplicación lineal \(\mathcal J u_m=v_m\) es una isometría
sobreyectiva entre los dos planos colectivos. Si
\[
 D_{\rm inc}=90|u_{90}\rangle\langle u_{90}|
             +120|u_{120}\rangle\langle u_{120}|,
 \qquad
 D_\partial=90|v_{90}\rangle\langle v_{90}|
             +120|v_{120}\rangle\langle v_{120}|,
\]
entonces \(\mathcal J D_{\rm inc}=D_\partial\mathcal J\).
Los proyectores sectoriales y sus combinaciones orientadas se
transportan mediante la misma aplicación.
\end{proposition}
\begin{proof}
La aplicación lleva una base ortonormal a otra y queda determinada
por esas imágenes. La igualdad de operadores vale sobre \(u_m\),
pues ambos miembros tienen valor \(m v_m\). Los proyectores son
\((120I-D)/30\) y \((D-90I)/30\) en cada plano, de modo que
sus combinaciones se conservan por composición lineal.
\end{proof}

El transporte anterior actúa sobre los modos uniformes. Las fluctuaciones
intrínsecas de cada espacio de incidencias permanecen en sus complementos
ortogonales. La ponderación areal utiliza precisamente las dos etiquetas
que esta isometría conserva.

El parámetro de Barbero--Immirzi se determina a partir de esos datos y
de la orientación angular previamente construida. Sean \(A\) y \(C^*\)
las coordenadas angulares en grados, en el levantamiento fijado en la
sección precedente, y pongamos
\[
 x=\frac{\pi_{\rm HMT}A}{180},\qquad
 y=\frac{\pi_{\rm HMT}C^*}{180},\qquad
 q_\pm=e^{-x\mp y},\qquad
 S_m(q)=\frac{q^m}{1-q^{3m}}.
\]
Se conserva la cámara \(x>|y|\), por lo que \(0<q_\pm<1\).
La cadena orientada de una vuelta asigna doce sectores a
\(\mathcal F_6\) y una costura de retorno con orientación contraria
a \(\mathcal F_8\). Su evaluación es \(12S_{90}-S_{120}\).
El coeficiente de \((1-q)^{-1}\) queda determinado por
\[
 \lim_{q\uparrow1}(1-q)(12S_{90}(q)-S_{120}(q))
 =\frac{12}{270}-\frac1{360}=\frac1{24},
\]
porque \(1-q^{3m}=(1-q)\sum_{j=0}^{3m-1}q^j\).
La evaluación conjugada de la cadena, junto con la componente angular,
define el parámetro estructural
\begin{equation}
 \gamma_{\rm HMT}
 =\frac{\pi_{\rm HMT}AC^*}{180}
 +12\bigl[S_{90}(q_-)-S_{90}(q_+)\bigr]
 -\bigl[S_{120}(q_-)-S_{120}(q_+)\bigr].
 \label{rev4:ae:gamma}
\end{equation}
El origen de los coeficientes \(12:-1\) reside en la cadena orientada.
La factorización precedente determina su coeficiente de polo.
La inversión \(C^*\mapsto-C^*\) intercambia los canales y cambia
el signo de \(\gamma_{\rm HMT}\). Conservamos a continuación la
rama \(C^*>0\), \(\gamma_{\rm HMT}>0\), con la normalización
areal fijada por esa misma orientación.

\subsection{Retorno longitudinal y determinación de la escala areal}
\label{rev5:longitud}

La unidad dimensional del área admite una construcción explícita
desde el mismo estado enriquecido. Escribiremos \(L_\alpha\) para la
longitud física obtenida; el entero \(L\) que designa la longitud de un
registro en la sección~\ref{sec:compat-lectores-formas} conserva su
significado combinatorio. La distinción separa el cardinal de una
representación, la extensión de una arista y el radio de su realización.

La composición utiliza las coordenadas regionales y el registro ya
generados por APP--TRIT--TPK. APP conserva ambas hojas con sus residuos
y cocientes; TRIT orienta el régimen local; TPK prolonga el estado,
conservando acarreo, frontera, ruta y memoria. La cadena
\(w_6\to w_{12}\to w_{18}\to w_{24}\to w_{30}\to R_{36}\to G_9\)
y el orden \(U_t\to\Gamma_9\to K_{\rm ph}\) mantienen el
retorno de fase con incremento de memoria. La construcción conjunta
del continuo conserva sus aspectos espectral, solenoidal, gauge,
cohomológico y determinantal. La longitud que sigue es una lectura
dimensional posterior de esa estructura y de sus salidas correlacionadas
\(\pi_{\rm HMT},\varphi_{\rm HMT},e_{\rm HMT},\alpha_{\rm HMT}\).

Fijemos el registro marcado
\[
 \Omega=\{(j,k,q,u,v):j,k\in\mathbb Z/12\mathbb Z,\quad
 k-j\in\{0,3,6,9\},\quad1\le q\le8,\quad
 u<v,\quad u,v\in\{1,2,4,5,7,8\}\}.
\]
La resta se calcula en ese grupo cíclico. Su cardinal es
\(N_\Omega=12\cdot4\cdot8\binom62=5760\).
El factor \(120\) procede del espacio octádico de incidencias
punto--pareja ya construido; las doce posiciones y las cuatro
posiciones de cada órbita conservan sus marcas respectivas. Esta
resolución amplía el portador del retorno de Hadamard. En los espacios
hermitianos de partida y llegada, de dimensión \(N_\Omega\), sea
\(B\) el inversor transportado, con
\(B^*B=BB^*=I/4\). La identidad local
\(H_4^*H_4=4I\) produce esa normalización para el inversor;
su extensión ortogonal y el transporte de las marcas la conservan.

En el espacio de llegada, sean
\(u_\Omega=N_\Omega^{-1/2}\mathbf1\),
\(P=u_\Omega u_\Omega^*\), \(C=(I-P)B\) y \(M=PB\).
Se conserva el transporte completo \(C+M=B\). Los proyectores
atómicos \(P_i=e_ie_i^*\), determinados por la inscripción,
definen la expectativa diagonal
\(\Delta_\Omega(T)=\sum_iP_iTP_i\).

\begin{proposition}[Respuesta registrada y complemento conservado]
\label{rev5:retorno}
Entre las aplicaciones lineales hacia el álgebra diagonal que fijan
los \(P_i\) y son bimodulares respecto de esa álgebra,
\(\Delta_\Omega\) es única. Además,
\begin{equation}
 \Delta_\Omega(CC^*)=r_\Omega I,
 \qquad r_\Omega=\frac{N_\Omega-1}{4N_\Omega}
 =\frac{5759}{23040},
 \qquad \Delta_\Omega(MM^*)=\frac{I}{4N_\Omega}.
 \label{rev5:retorno-escalar}
\end{equation}
Las dos respuestas suman \(I/4\).
\end{proposition}
\begin{proof}
Para la unidad matricial \(E_{ij}\), la bimodularidad impone
\(E(E_{ij})=P_iE(E_{ij})P_j\). Como la imagen es diagonal,
este término vale cero si \(i\ne j\); si \(i=j\), fijar
\(P_i\) determina su imagen. Las unidades matriciales constituyen
una base y determinan el mapa. Por otra parte,
\(CC^*=(I-P)/4\), \(MM^*=P/4\), y cada entrada diagonal
de \(P\) es \(1/N_\Omega\). Aplicar la expectativa demuestra
las tres identidades.
\end{proof}

La inscripción isométrica \(We_i=e_i\otimes e_i\) realiza
\(W^*(T\otimes I)W=\Delta_\Omega(T)\). La inscripción
completa conserva la información del registro; la expectativa extrae
su respuesta diagonal. El complemento \(M\) permanece en el
transporte y posee la respuesta específica de
\eqref{rev5:retorno-escalar}. Esta separación identifica el observable
que interviene en la longitud y el componente complementario conservado.

La forma normal de Smith de \(H_4\) es
\(\operatorname{diag}(1,2,2,4)\); su conúcleo tiene orden
dieciséis. El producto del escalar generado
\(\alpha=\alpha_{\rm HMT}\), indexado por las clases de ese
conúcleo, define el lector multiplicativo finito
\[
 n_H:=\prod_{\bar v\in\operatorname{coker}_{\mathbb Z}(H_4)}
 \alpha=\alpha^{16}.
\]
Sea \(L_*>0\) una sección de referencia
radial de la recta dimensional de longitud. La composición longitudinal
actúa linealmente sobre una cocadena orientada \(\beta_*\):
\begin{equation}
 \mathscr T_L\beta_*
 =\alpha^{16}(\Delta_\Omega(CC^*)\otimes I)\beta_*
 =q_\alpha\beta_*,\qquad
 q_\alpha=\alpha^{16}r_\Omega,\qquad
 L_\alpha=q_\alpha L_*.
 \label{rev5:longitud-generada}
\end{equation}
La fórmula fija una composición concreta de lectores. La linealidad
conserva el tipo de longitud y su orientación; el coeficiente procede
del retorno registrado y de la norma local. Por ejemplo, la respuesta
diagonal de una fuente \(a e_i\) es \(a r_\Omega\), mientras
la norma \(\|C^*ae_i\|=|a|\sqrt{r_\Omega}\) constituye
otra lectura del mismo transporte.

Si \(\beta_*(\bar a)=-U(a)\beta_*(a)\), multiplicar por
\(q_\alpha\) conserva la inversión orientada. Si una arista se
subdivide en nueve tramos compatibles,
\(\beta_*(a)=\sum_jU_j^{-1}\beta_*(a_j)\), la misma
multiplicación conserva esa identidad y la concatenación de memoria.
En una ampliación auxiliar del registro, se transportan
\(P\), \(CC^*\) y \(\Delta_\Omega\) por tensorización
con la identidad. Así permanece \(r_\Omega\); el cardinal de
una malla refinada y el cardinal del registro marcado tienen funciones
distintas.

\subsection{Reciprocidad radial y reconocimiento de la longitud de Planck}
\label{rev5:radial}

La lectura circular del calendario de ciento ocho avances conserva
\(108\ell_0=2\pi_{\rm HMT}L_\alpha\), con
\(\ell_0=ct_0\). Por consiguiente,
\begin{equation}
 t_0=\frac{\pi_{\rm HMT}L_\alpha}{54c},\qquad
 \omega=\frac{c}{L_\alpha},\qquad
 Q_{\rm clk}=\hbar_{\rm ret}\omega
 =\frac{\hbar_{\rm ret}c}{L_\alpha}.
 \label{rev5:reloj}
\end{equation}
La acción retornada se recibe de su ecuación HMT completa en
\eqref{eq:mass-k-action}; la velocidad se recibe de los dos canales
constitutivos de la sección~\ref{sec:compat-lectores-formas}.
La relación circular convierte el avance levantado en radio. Si
\(c=\widehat c\,u_C\), \(t_0=u_T\) y
\(u_L=u_Cu_T\), resulta
\(L_*/u_L=54\widehat c/(\pi_{\rm HMT}q_\alpha)\).
La referencia radial \(L_*\) y la base de arista \(u_L\)
quedan relacionadas con sus tipos dimensionales explícitos.

En una fibra finita, sea \(X=X^*>0\) un carácter positivo
invertible, \(U=2B\) y \(Y=UXU^*\). El transporte
longitudinal es \(D_L=L_\alpha U\). Su respuesta radial positiva
se define mediante la métrica de llegada,
\(\|R_Xv\|^2=\|XD_L^*v\|^2\) para todo \(v\).
La energía y la longitud conjugada conservan el mismo carácter:
\[
 H_X=Q_{\rm clk}Y,\qquad M_X=H_X/c^2,\qquad
 \Lambda_X=L_\alpha Y^{-1}.
\]

\begin{theorem}[Coeficiente radial común y escala areal]
\label{rev5:rigidez-radial}
Las condiciones anteriores determinan una única respuesta positiva
\(R_X=L_\alpha Y\). Se cumplen
\begin{equation}
 R_X\Lambda_X=L_\alpha^2I,\qquad
 H_X\Lambda_X=\hbar_{\rm ret}cI,\qquad
 R_X=\frac{L_\alpha^2}{\hbar_{\rm ret}c}H_X.
\end{equation}
En la realización física que identifica \(R_X\) con el radio
gravitatorio reducido \(R_X=(G/c^2)M_X\), el coeficiente es
único y común a todos los caracteres positivos admisibles:
\begin{equation}
 G=\frac{c^3L_\alpha^2}{\hbar_{\rm ret}},\qquad
 \frac{\hbar_{\rm ret}G}{c^3}=L_\alpha^2.
 \label{rev5:G-area}
\end{equation}
\end{theorem}
\begin{proof}
La polarización de la igualdad de normas determina
\(R_X^2=D_LX^2D_L^*=L_\alpha^2UX^2U^*\).
El operador \(L_\alpha UXU^*\) es positivo y su cuadrado es
el indicado. La unicidad de la raíz positiva prueba la primera
afirmación. Las dos identidades de productos siguen por sustitución.
La convención física da
\(L_\alpha Y=(GQ_{\rm clk}/c^4)Y\). La invertibilidad
de \(Y\) permite cancelarlo, y la expresión de
\(Q_{\rm clk}\) proporciona \eqref{rev5:G-area}.
Si otro coeficiente conserva los mismos datos, su diferencia
multiplicada por \(M_X\) vale cero y, por tanto, es nula.
\end{proof}

El dominio de la prueba es la fibra positiva finita; también admite
operadores acotados uniformemente positivos con las mismas identidades.
La lectura de radio gravitatorio reducido constituye la convención
física declarada de esta realización. El reconocimiento convencional
posterior \(\ell_P^2=\hbar_{\rm ret}G/c^3\) identifica
\(\ell_P=L_\alpha\). La longitud se obtiene primero mediante
\eqref{rev5:longitud-generada}; esta identificación conserva su
constructor. La fórmula circular equivalente es
\[
 G=\left(\frac{54}{\pi_{\rm HMT}}\right)^2
 \frac{c^5t_0^2}{\hbar_{\rm ret}}.
\]
El factor circular ya está incluido cuando se utiliza \(L_\alpha\).
La longitud y la sección de acción determinan asimismo
\[
 t_P=\frac{L_\alpha}{c},\qquad
 m_P=\frac{\hbar_{\rm ret}}{cL_\alpha},\qquad
 E_P=\frac{\hbar_{\rm ret}c}{L_\alpha}.
\]
Estas identidades siguen al sustituir \eqref{rev5:G-area} en las
expresiones convencionales de las tres escalas; la positividad
determina sus raíces. Para un modo de carácter \(y>0\), las
lecturas son \(m(y)=m_Py\), \(r_g(y)=L_\alpha y\) y
\(\bar\lambda(y)=L_\alpha/y\). Su producto longitudinal
vale \(r_g(y)\bar\lambda(y)=L_\alpha^2\). La igualdad de
ambas longitudes equivale a \(y=y^{-1}\) y, por positividad,
caracteriza exactamente \(y=1\). El calendario conserva
\(t_0=(\pi_{\rm HMT}/54)t_P\) y
\(108t_0=2\pi_{\rm HMT}t_P\).

Un radio de horizonte \(R_h=L_\alpha\zeta_h\) conserva su
factor propio \(\zeta_h=R_h/L_\alpha\). Su área espacial
pertenece al horizonte declarado; la escala elemental del operador
sectorial siguiente es \(L_\alpha^2\). En la especialización
de Schwarzschild para la misma masa, el radio es \(2r_g(y)\)
y, por tanto, \(\zeta_h=2y\); un horizonte de otra realización
conserva su condición geométrica de formación.


\subsection{Ocupaciones de frontera y reconstrucción sectorial}

Cada enlace de \(E_{90}\sqcup E_{120}\) determina un factor
\(\mathbb C^3\), con base \(|0\rangle,|1\rangle,|2\rangle\)
y operador número \(N_e=\operatorname{diag}(0,1,2)\). El espacio de
estados de frontera es
\(\mathscr H_A=(\mathbb C^3)^{\otimes18}\otimes
(\mathbb C^3)^{\otimes18}\). Definimos \(N_m\) como la suma de
los operadores número de los dieciocho enlaces del sector \(m\).
Los dos operadores conmutan y cada uno tiene espectro entero entre
cero y treinta y seis.

La escala angular procede del reloj y del carácter normalizado del
triplete conjugado:
\begin{equation}
 \theta_{\rm clk}=\frac{C^*}{6}\frac{\pi_{\rm HMT}}{180},
 \quad
 \chi_{10}=\frac13\Tr\operatorname{diag}
 (1,e^{i\pi_{\rm HMT}/18},e^{-i\pi_{\rm HMT}/18})
 =\frac{1+2\cos(\pi_{\rm HMT}/18)}3,
 \quad a_0=\chi_{10}\theta_{\rm clk}.
 \label{rev4:ae:escala}
\end{equation}
La identidad del carácter se obtiene sumando las fases conjugadas.
En la rama considerada \(a_0>0\). La unidad areal es
\(L_\alpha^2=\ell_P^2\), construida y reconocida en
\eqref{rev5:longitud-generada}--\eqref{rev5:G-area}. Escribamos
\(a_\partial=\gamma_{\rm HMT}a_0>0\). El operador de área es
\begin{equation}
 \widehat{\mathcal A}=L_\alpha^2a_\partial N,
 \qquad N=N_{90}+N_{120},
 \qquad D=90N_{90}+120N_{120}.
 \label{rev4:ae:area}
\end{equation}
Su espectro es \(\{nL_\alpha^2a_\partial:0\le n\le72\}\).
La multiplicidad de \(n\) es el coeficiente de \(z^n\) en
\((1+z+z^2)^{36}\), pues cada factor tensorial aporta una de
las tres ocupaciones.

\begin{theorem}[Reconstrucción de las dos ocupaciones sectoriales]
\label{rev4:ae:recuperacion}
La pareja de observables \((N,D)\) determina
\begin{equation}
 N_{90}=4N-\frac{D}{30},
 \qquad N_{120}=\frac{D}{30}-3N.
 \label{rev4:ae:inversa}
\end{equation}
En el espectro conjunto, sus valores \((n,d)\) satisfacen
\(n\in\mathbb Z\), \(d\in30\mathbb Z\), \(90n\le d\le120n\) y
\(0\le4n-d/30\le36\), \(0\le d/30-3n\le36\).
Estas condiciones caracterizan la imagen de las ocupaciones enteras.
\end{theorem}
\begin{proof}
La transformación de \((N_{90},N_{120})\) en \((N,D)\) tiene matriz
\(\left(\begin{smallmatrix}1&1\\90&120\end{smallmatrix}\right)\),
cuyo determinante es \(30\). Su inversa da las identidades de
operadores. La conmutatividad permite aplicarlas en cada autoespacio
conjunto. Las desigualdades expresan la no negatividad y las capacidades
máximas de los sectores; la divisibilidad hace enteras ambas ocupaciones.
Recíprocamente, cada entero entre cero y treinta y seis se expresa como
suma de dieciocho términos de \(\{0,1,2\}\), y por tanto todo par
que satisface esas condiciones pertenece al espectro conjunto.
\end{proof}

El área conserva la suma de ocupaciones; el observable ponderado conserva
su composición sectorial. Así, \((N_{90},N_{120})=(1,0)\) y \((0,1)\)
tienen igual área \(L_\alpha^2a_\partial\) y pesos \(90\) y \(120\).
En sentido complementario, \((4,0)\) y \((0,3)\) tienen igual peso
\(D=360\), pero áreas diferentes. La pareja separa ambas situaciones.
Dentro de cada autoespacio conjunto se conservan las degeneraciones
microscópicas correspondientes a las distribuciones entre enlaces.

\begin{figure}[htbp]
\centering
\begin{tikzpicture}[>=Latex,every node/.style={font=\small},node distance=15mm]
\node[draw,rounded corners,align=center] (a) {Ocupación sectorial\\$(1,0)$};
\node[draw,rounded corners,align=center,right=54mm of a] (b) {Ocupación sectorial\\$(0,1)$};
\node[draw,rounded corners,align=center] (n) at ($(a)!0.5!(b)+(0,-1.8)$)
 {Área común\\$L_\alpha^2a_\partial$};
\draw[->] (a)--node[left] {$N=1$}(n);
\draw[->] (b)--node[right] {$N=1$}(n);
\node[above=5mm of a] {$D=90$};
\node[above=5mm of b] {$D=120$};
\end{tikzpicture}
\caption{La igualdad de área conserva la ocupación total. La ponderación
incidencial distingue las dos composiciones; su lectura conjunta permite
reconstruir las ocupaciones de ambos sectores.}
\label{rev4:ae:figura}
\end{figure}

\subsection{Estado reducido, purificación y entropía de entrelazamiento}

La interpretación informacional requiere fijar el estado. Para
\(\Delta\in\mathbb R\) finito, con la normalización del parámetro
modular de la construcción sectorial, definimos
\begin{equation}
 x_m=e^{-m\Delta},\qquad z_m=1+x_m+x_m^2,
 \qquad Z=z_{90}^{18}z_{120}^{18},
 \qquad \rho_A(\Delta)=Z^{-1}e^{-\Delta D}.
 \label{rev4:ae:rho}
\end{equation}
La suma de ocupaciones conmutantes y la descomposición tensorial prueban
la expresión de \(Z\). Todos los autovalores de \(\rho_A\) son
estrictamente positivos para \(\Delta\) real finito. Su Hamiltoniano
modular, denotado \(K_A\) para distinguirlo del registro aritmético
\(K\), es
\begin{equation}
 K_A=-\log\rho_A=(\log Z)I+\Delta D.
 \label{rev4:ae:modular}
\end{equation}
Este operador caracteriza la distribución espectral del estado reducido.
El generador de evolución temporal se define en el desarrollo dinámico
con sus propios parámetros y dominio. Para \(\Delta\ne0\), conocido
\(\Delta\) y conocida la normalización, la pareja
\((\widehat{\mathcal A},K_A)\) recupera \((N,D)\) y, mediante
\eqref{rev4:ae:inversa}, ambas ocupaciones. Para \(\Delta=0\),
el estado es uniforme y \(K_A=36\log3\,I\).

Construimos la bipartición de manera explícita. Sea \(\mathscr H_{\bar A}\)
una copia de \(\mathscr H_A\), con las bases conjugadas fijadas. Para
un enlace de peso \(w\), el vector
\[
 |\operatorname{TFD}(w)\rangle
 =\frac1{\sqrt{1+e^{-w}+e^{-2w}}}
 \sum_{n=0}^{2}e^{-wn/2}|n\rangle_A\otimes|n\rangle_{\bar A}
\]
tiene norma uno. El estado puro
\begin{equation}
 |\Psi_\Delta\rangle
 =|\operatorname{TFD}(90\Delta)\rangle^{\otimes18}
  \otimes|\operatorname{TFD}(120\Delta)\rangle^{\otimes18}
 \label{rev4:ae:tfd}
\end{equation}
determina la bipartición \(A\mid\bar A\).

\begin{proposition}[Reducción sectorial y entropía de la purificación]
\label{rev4:ae:entropia}
La traza parcial de \(|\Psi_\Delta\rangle\langle\Psi_\Delta|\)
sobre \(\mathscr H_{\bar A}\) es \(\rho_A(\Delta)\). Definidos
\[
 f_m=\frac{x_m+2x_m^2}{z_m},
 \qquad s(\Delta)=-\Tr(\rho_A(\Delta)\log\rho_A(\Delta)),
\]
se cumplen
\begin{align}
 \langle N\rangle_\Delta&=18(f_{90}+f_{120}),\\
 \langle D\rangle_\Delta&=18(90f_{90}+120f_{120}),\\
 \langle\widehat{\mathcal A}\rangle_\Delta
 &=18L_\alpha^2a_\partial(f_{90}+f_{120}),\\
 s(\Delta)&=18\log z_{90}+18\log z_{120}
             +\Delta\langle D\rangle_\Delta.
 \label{rev4:ae:entropia-formula}
\end{align}
La última expresión es la entropía de entrelazamiento de la
purificación \eqref{rev4:ae:tfd}, en unidades adimensionales.
\end{proposition}
\begin{proof}
En la traza parcial de un enlace, la ortogonalidad de
\(|n\rangle_{\bar A}\) elimina los términos de índices distintos y
deja los pesos \((1,x_m,x_m^2)/z_m\). La traza parcial de un producto
tensorial es el producto de las trazas parciales; se obtiene
\eqref{rev4:ae:rho}. El primer momento de cada enlace es \(f_m\).
La aditividad de los operadores número da las tres primeras identidades.
La última resulta de sustituir \eqref{rev4:ae:modular} en
\(s=\Tr(\rho_AK_A)\). Los coeficientes de la construcción anterior
son coeficientes de Schmidt de la bipartición declarada, por lo que
esa entropía coincide con su entropía de entrelazamiento.
\end{proof}

La entropía máxima de este espacio es \(36\log3\), alcanzada en
\(\Delta=0\). La cantidad \(36\log3-s(\Delta)\) es la
entropía relativa de \(\rho_A(\Delta)\) respecto de
\(I/3^{36}\). La incidencia fija los sectores y sus multiplicidades;
el parámetro del estado fija las probabilidades; la coordenatización dimensional
convierte la ocupación en área. La identificación física de
\(\mathscr H_{\bar A}\) conserva la bipartición y el álgebra de
observables que acaban de especificarse.

\subsection{Identidad modular finita y transporte de la información geométrica}

Fijemos el estado de referencia \(\rho_0=\rho_A(\Delta_0)\), con
\(\Delta_0\) real finito, y mantengamos constantes
\(\Delta_0,a_\partial,L_\alpha\). Las áreas sectoriales
normalizadas \(\mathcal A_m^{\rm n}=a_\partial N_m\) satisfacen
\begin{equation}
 K_0=-\log\rho_0=c_0I+
 \beta_{90}\mathcal A_{90}^{\rm n}
 +\beta_{120}\mathcal A_{120}^{\rm n},
 \qquad c_0=\log Z(\Delta_0),\qquad
 \beta_m=\frac{m\Delta_0}{a_\partial}.
 \label{rev4:ae:betas}
\end{equation}
La relación \(\beta_{120}=\tfrac43\beta_{90}\) conserva la razón
de incidencias. Para \(\Delta_0\ne0\), la razón de ambos coeficientes
está definida y vale \(4/3\).

\begin{theorem}[Identidad finita de área e información relativa]
\label{rev4:ae:ley-finita}
Para todo estado normalizado \(\sigma\) de \(\mathscr H_A\), sea
\(\delta_\sigma\langle B\rangle=\Tr[(\sigma-\rho_0)B]\).
Entonces
\begin{equation}
 s(\sigma)-s(\rho_0)
 =\sum_{m\in\{90,120\}}\beta_m
       \delta_\sigma\langle\mathcal A_m^{\rm n}\rangle
   -D(\sigma\Vert\rho_0),
 \label{rev4:ae:ley-finita-formula}
\end{equation}
donde \(D(\sigma\Vert\rho_0)=
\Tr[\sigma(\log\sigma-\log\rho_0)]\) es finita y no negativa,
con la convención \(0\log0=0\).
\end{theorem}
\begin{proof}
De la definición se obtiene
\(D(\sigma\Vert\rho_0)=-s(\sigma)+\Tr(\sigma K_0)\),
mientras \(s(\rho_0)=\Tr(\rho_0K_0)\). Restar las identidades
y sustituir \eqref{rev4:ae:betas} prueba
\eqref{rev4:ae:ley-finita-formula}; el término escalar desaparece
por la normalización de las trazas. Para la positividad, sean \(p_i\)
y \(u_i\) los autovalores y autovectores de \(\sigma\), y
\(q_j>0,v_j\) los de \(\rho_0\). Con
\(r_i=\langle u_i,\rho_0u_i\rangle\), la concavidad de
\(\log\) da
\(\langle u_i,\log\rho_0u_i\rangle\le\log r_i\).
Así, \(D(\sigma\Vert\rho_0)\ge\sum_i p_i\log(p_i/r_i)\ge0\),
por convexidad de \(t\log t\), ya que \(\sum_i r_i=1\).
El rango completo de \(\rho_0\) garantiza la finitud.
\end{proof}

Su versión infinitesimal mantiene fijo el estado de referencia. Si
\(\sigma(\varepsilon)\) es diferenciable, normalizado y
\(\sigma(0)=\rho_0\), entonces
\[
 \left.\frac{\dd s(\sigma(\varepsilon))}{\dd\varepsilon}\right|_0
 =\sum_m\beta_m
  \left.\frac{\dd\langle\mathcal A_m^{\rm n}\rangle_{
                          \sigma(\varepsilon)}}{\dd\varepsilon}\right|_0.
\]
En efecto, la derivada de la traza de \(-\sigma\log\sigma\) es
\(-\Tr[\dot\sigma(\log\rho_0+I)]\); la normalización elimina
\(\Tr\dot\sigma\) y la sustitución de \(K_0\) da la fórmula.
La identidad finita añade a esa respuesta tangente la corrección exacta
de entropía relativa. Los coeficientes de las áreas físicas son
\(\beta_m/L_\alpha^2\), conservando la misma coordenatización dimensional.

La composición con la representación grassmanniana precisa finalmente
qué información se conserva. En
\(\mathcal F(K,\xi)=([B(K)],Q,\xi)\), las coordenadas adicionales
\(\xi\) retienen las magnitudes regionales, la orientación angular,
la inclusión marcada, la identificación de los sectores, la base de
ocupaciones, el parámetro \(\Delta\) y la unidad areal pertinente.
La inversa de la proposición~\ref{prop:compat-inversa-marcada}
recupera \(K\); las otras condiciones se conservan mediante sus
marcas. Por el teorema~\ref{thm:compat-algebra-lectores}, la composición
con esa inversa transporta \(\gamma_{\rm HMT}\),
\(\widehat{\mathcal A}\), \(\rho_A\), \(K_A\) y sus evaluaciones
en el dominio declarado. Los cambios de coordenatización del mismo punto marcado
conservan esos objetos, y los refinamientos conservan las evaluaciones
heredadas cuando mantienen los mismos datos sectoriales.

Esta articulación conserva tres niveles distinguibles de información.
La representación positiva marcada recupera el registro con los datos
retenidos; la pareja área--Hamiltoniano modular recupera las dos
ocupaciones bajo las condiciones de \eqref{rev4:ae:modular}; la
purificación determina el entrelazamiento de una bipartición concreta.
La pérdida de información en una proyección se calcula por sus fibras:
el área identifica ocupaciones de igual suma, mientras la lectura
conjunta separa su composición y conserva la degeneración microscópica
restante. El vínculo entre geometría e información queda así expresado
por aplicaciones e identidades verificables sobre la misma construcción.
\subsection{Naturalidad de la escala y formas canónicas operatoriales}
\label{rev5:naturalidad-areal}

La longitud generada permite precisar la compatibilidad entre la
realización positiva, el operador de área y la información sectorial.
La representación marcada \(\mathcal F(K,\xi)\) conserva,
entre los datos pertinentes de \(\xi\), la resolución
\(\Omega\), sus proyectores, el retorno transportado y la sección
\(L_*\). Las coordenadas regionales generadas, la acción y los canales
constitutivos mantienen su genealogía anterior. La inversa marcada
restituye así todos los argumentos de
\eqref{rev5:longitud-generada} y \eqref{rev5:G-area}.
En particular, el lector geométrico ya construido de
\(\alpha^{\mathrm{geom}}\) determina las expresiones
\[
 L_\alpha^{\mathrm{geom}}
 =L_*(\alpha^{\mathrm{geom}})^{16}\frac{5759}{23040},
 \qquad
 G^{\mathrm{geom}}
 =\frac{(c^{\mathrm{geom}})^3
 (L_\alpha^{\mathrm{geom}})^2}{\hbar_{\rm ret}^{\mathrm{geom}}}.
\]
Los lectores de velocidad y acción se obtienen por la misma
composición con la inversa marcada. Estas expresiones conservan
los datos del retorno y la sección radial de referencia.

\begin{corollary}[Conservación por cambios de coordenatización y refinamiento heredado]
\label{rev5:conservacion}
Los cambios positivos de coordenatización que conservan el punto marcado y los
refinamientos que restituyen el registro inicial conservan
\(L_\alpha\), \(G\), \(\gamma_{\rm HMT}\) y
\(\widehat{\mathcal A}=L_\alpha^2a_\partial N\), cuando
transportan asimismo los datos sectoriales y dimensionales declarados.
El espectro normalizado
\(\widehat{\mathcal A}/L_\alpha^2=a_\partial N\), sus
multiplicidades y el estado \(\rho_A(\Delta)\) permanecen iguales.
\end{corollary}
\begin{proof}
El teorema~\ref{thm:compat-algebra-lectores} transporta cada lector
por composición con la inversa. Las mutaciones del mismo punto
reconstruyen el mismo registro. En el refinamiento heredado, sumar
las separaciones subdivididas reconstruye cada coordenada inicial.
La tensorización de la inscripción conserva \(r_\Omega\),
y sus restantes argumentos quedan fijados por las marcas. Por
sustitución se conservan la longitud, el coeficiente gravitatorio y
el operador de área. Los proyectores de ocupación conservan el
polinomio de multiplicidades \((1+z+z^2)^{36}\) y la expresión
normalizada del estado.
\end{proof}

El refinamiento heredado considerado aquí conserva el mismo espacio
de treinta y seis enlaces de frontera. Una ampliación que introduzca
un espacio físico de frontera diferente requiere entrelazadores
especificados para \(N_{90}\) y \(N_{120}\), junto con el
transporte de los datos sectoriales declarados.

Sea ahora \(P\) un politopo de rango uno de este manuscrito,
con una subdivisión orientada \(P=\bigcup_\sigma P_\sigma\)
del mismo soporte, multiplicidad uno y caras compatibles. Se fija
la misma fibra sectorial \(\mathscr H_A\) sobre todas sus
celdas. La forma de grado máximo con valores operatoriales y
dimensión de área se define por
\[
 \boldsymbol\Omega^{\mathcal A}_P
 =\widehat{\mathcal A}\otimes\Omega_P.
\]
Aquí \(\widehat{\mathcal A}\) es constante respecto de las
coordenadas diferenciales del politopo: procede del estado HMT
fijado, mientras \(\Omega_P\) describe su soporte geométrico.
El coeficiente \(\widehat{\mathcal A}\) es semidefinido positivo.
Su componente en el sector de vacío \(N=0\) es nula. Sobre un
autoespacio de ocupación \(n>0\), dividir por
\(nL_\alpha^2a_\partial\) recupera la forma canónica
\(\Omega_P\) multiplicada por la identidad de ese autoespacio.

\begin{proposition}[Aditividad areal y residuos de frontera]
\label{rev5:forma-areal}
En la trivialización sectorial común se cumplen
\begin{equation}
 \boldsymbol\Omega^{\mathcal A}_P
 =\sum_\sigma\widehat{\mathcal A}\otimes\Omega_{P_\sigma},
 \qquad
 \operatorname{Res}_F\boldsymbol\Omega^{\mathcal A}_P
 =\widehat{\mathcal A}\otimes\Omega_F.
\end{equation}
Las identidades persisten bajo subdivisiones sucesivas y residuos
iterados. Para un estado fijo \(\rho\), tomar la traza da
\(\Tr(\rho\widehat{\mathcal A})\Omega_P\) y conserva
las mismas identidades escalares.
\end{proposition}
\begin{proof}
El espacio de operadores de \(\mathscr H_A\) es finito
dimensional. En cualquier base, la afirmación es la identidad de
aditividad de las formas canónicas multiplicada por cada entrada
constante de \(\widehat{\mathcal A}\). El residuo y la traza
son lineales. La cancelación de facetas interiores conserva sus
orientaciones y la misma entrada operatorial en ambos lados.
La iteración aplica esas operaciones a cada subdivisión y a cada
bandera de caras.
\end{proof}

Si distintas celdas llevan coeficientes operatoriales regulares
\(A_\sigma\), el defecto de una pareja sobre su cara común es
\((A_\sigma|_F-A_\tau|_F)\otimes\Omega_F\).
La igualdad de restricciones operatoriales a la frontera proporciona
el pegado, bajo las condiciones de polos simples del
teorema~\ref{thm:compat-pegado}. La suma global conserva además
todas las contribuciones sobre cada divisor ambiental y el control
de polos exteriores. La normalización dimensional, la cancelación
de residuos y la cobertura del soporte tienen así criterios
separados y compatibles.

La identidad modular finita también conserva una expresión directa
en áreas físicas \(\mathcal A_m=L_\alpha^2a_\partial N_m\):
\begin{equation}
 s(\sigma)-s(\rho_0)
 =\sum_{m=90,120}\frac{m\Delta_0}{L_\alpha^2a_\partial}
   \Tr[(\sigma-\rho_0)\mathcal A_m]
  -D(\sigma\Vert\rho_0).
 \label{rev5:modular-fisica}
\end{equation}
Es \eqref{rev4:ae:ley-finita-formula} después de insertar la
longitud construida. El factor angular \(a_0\), contenido en
\(a_\partial=\gamma_{\rm HMT}a_0\), permanece íntegro.
Si se comparan secciones dimensionales con \(L'_\alpha=sL_\alpha\)
y se conserva el mismo estado y los mismos datos adimensionales,
las áreas se multiplican por \(s^2\) y sus coeficientes modulares
por \(s^{-2}\). La entropía y la entropía relativa permanecen
iguales. Esta comparación de secciones no constituye otra solución
para los mismos datos longitudinales fijados.

Para un horizonte esférico de radio \(R_h=L_\alpha\zeta_h\),
el área geométrica vale \(4\pi_{\rm HMT}L_\alpha^2\zeta_h^2\).
Su identificación con la esperanza del operador sectorial, si se
requiere para un mismo corte, exige la condición explícita
\[
 a_\partial\Tr(\rho N)=4\pi_{\rm HMT}\zeta_h^2.
\]
La geometría del horizonte y el estado de frontera deben satisfacer
esa relación con sus propias construcciones. El factor
\(\zeta_h\) caracteriza ese horizonte; la escala longitudinal
del acoplamiento permanece \(L_\alpha\).


\begin{figure}[htbp]
\centering
\begin{tikzpicture}[
 box/.style={draw=black!70,rounded corners=1pt,align=center,
             font=\small,inner sep=5pt,text width=40mm},
 arrow/.style={-{Stealth[length=2mm]},thick}]
\node[box] (memory) at (0,0) {Carga y memorias\\\((q,m_0,m_1)\)};
\node[box] (register) at (0,-1.65) {Registro reconstruido\\\(K=(q/12)\mathbf1+Lm\)};
\node[box] (weights) at (-3.8,-3.55) {Partición ordenada\\\(d_j=K_j/q\)};
\node[box] (direction) at (3.8,-3.55) {Dirección de suma nula\\\(u=P_3P_{11}K\)};
\node[box] (geometry) at (-4.6,-6) {Medición de frontera\\y forma canónica\\\(C(d),\ \Omega_{P_K}\)};
\node[box] (energy) at (0,-6) {Energía nodal\\y detalles\\\(L_d,\ \mathcal E(z)\)};
\node[box] (lattice) at (4.6,-6) {Deformación reticular\\y red dual\\\(\Lambda_s^*=\Lambda_{-s}\)};
\draw[arrow] (memory)--(register);
\draw[arrow] (register)--(weights);
\draw[arrow] (register)--(direction);
\draw[arrow] (weights)--(geometry);
\draw[arrow] (weights)--(energy);
\draw[arrow] (direction)--(lattice);
\draw[arrow] (direction) to[bend left=12] node[right,font=\scriptsize,align=left,fill=white,inner sep=1pt]{flujo\\heredado} (energy);
\end{tikzpicture}
\caption{Realizaciones posteriores del registro recuperable. La rama
intervalar determina los menores de caminos, la forma canónica del
polígono y el laplaciano nodal. La dirección algebraica determina el
flujo de intervalos y, junto con la realización reticular marcada, la
colección dual de veinticuatro dimensiones. Cada flecha conserva el dominio
y las hipótesis de su construcción; los codominios inferiores son
distintos.}
\label{fig:multimorphology}
\end{figure}


\clearpage
\part{Realización gauge, refinamiento y brecha espectral}
\section{Positividades, espacios físicos y transporte del espectro}
La geometría positiva y la teoría gauge reciben estructuras distintas de
un mismo proceso genealógico. La primera realiza los pesos y sus
incidencias mediante menores ordenados, celdas y formas canónicas.
La segunda conserva la incidencia gauge, la holonomía, las
representaciones, las rutas y la memoria dentro de una medida compatible.
Su forma de positividad relevante para la reconstrucción física es la
positividad por reflexión. El nexo conceptual adquiere contenido
matemático cuando se conservan los mapas que producen cada forma.

En la parte geométrica se ha demostrado que la subdivisión de una celda
conserva su forma canónica mediante cancelación de polos interiores.
Para las energías nodales, la reducción de Schur conserva una forma
efectiva y registra el detalle. En la parte gauge, la compatibilidad del
refinamiento debe conservar además el vacío, el espacio de invariantes,
el generador y el testigo espectral. Ésta es la cadena desarrollada a
continuación: escala aritmético--geométrica, medida común, semigrupo,
proyección física, refinamiento y reconstrucción continua.

La norma de un vector, la positividad de una matriz de momentos y la
existencia de una brecha de un Hamiltoniano poseen dominios diferentes.
Una congruencia con factor invertible conserva la positividad y la
inercia de una forma; un entrelazamiento isométrico del generador, con su
dominio, transporta su realización sobre la imagen reducida de la
inclusión. La cota uniforme y el testigo persistente fijan después la
brecha límite. Por ello el desarrollo mantiene visibles las
inclusiones y las identidades de semigrupo que justifican el paso al
límite. La escala de la brecha procede de la incidencia areal--volumétrica
y de su unidad dimensional declarada.

Los cuerpos siguientes reúnen la construcción proyectiva y su
desarrollo posterior dependiente del acoplamiento \(g\). Las primeras
comparaciones con una parametrización Gibbs se completan después mediante
la acción, la medida y el transporte explícitos. El orden expositivo
conserva las demostraciones del proceso y sitúa la formulación completa
tras sus dependencias. Las cuatro reflexiones corresponden a las cuatro
direcciones euclídeas de la reconstrucción. El desarrollo aquí reunido
tiene por objeto esa teoría gauge y su brecha, con el grupo compacto,
conexo y simple declarado en los teoremas.

El sector supersimétrico planar en el que se estudian amplitudes mediante
amplituedros conserva una pregunta física específica. Para enlazarlo con
la presente realización se requieren los operadores de supercarga, sus
dominios y la representación correspondiente. La dimensión cuatro de la
coordenatización euclídea y la diferencia cuatro del bloque modular permanecen
resultados ya definidos de sus respectivos operadores.

\section{Transición areal--volumétrica y escala espectral}
\label{sec:ym-areal-volumetric-scale}

La constante que fija el semigrupo gauge se obtiene antes de introducir el
Hamiltoniano. Su procedencia combina dos lectores independientes del mismo
estado genealógico: la frecuencia de retorno que expresa
\(\pi_{\rm HMT}\) y la distribución estacionaria de las dos hojas
areal--volumétrica.

\subsection{Dinámica de las dos hojas}

Sea
\begin{equation}
 F_{\rm av}=\begin{pmatrix}0&1\\1&1\end{pmatrix},
 \qquad e_{\rm ar}=\binom10,
 \qquad e_{\rm vol}=\binom01.
 \label{eq:ym-av-transition-matrix}
\end{equation}
La recurrencia de Fibonacci da, por inducción,
\begin{equation}
 F_{\rm av}^{n}
 =\begin{pmatrix}F_{n-1}&F_n\\F_n&F_{n+1}\end{pmatrix}.
 \label{eq:ym-av-fibonacci-powers}
\end{equation}
El vector de Perron es proporcional a \((1,\varphi)\). Como
\(F_{\rm av}\) es simétrica, la medida de Parry asigna a las dos hojas
pesos proporcionales a \((1,\varphi^2)\), y por tanto
\begin{equation}
 \frac{\mu_{\rm ar}}{\mu_{\rm vol}}=\varphi^{-2},
 \qquad
 \mu_{\rm vol}=\frac{\varphi^2}{1+\varphi^2}.
 \label{eq:ym-av-parry-weights}
\end{equation}
Aquí \(\mu_{\rm vol}\) es un carácter de frecuencia adimensional; no
denota un volumen de dimensión \(L^4\).

El lector regional de \(\pi\) se introduce una sola vez mediante
\begin{equation}
 D_\pi
 =\pi_{\rm HMT}|e_{\rm ar}\rangle\langle e_{\rm ar}|
  +|e_{\rm vol}\rangle\langle e_{\rm vol}|.
 \label{eq:ym-pi-regional-reader}
\end{equation}
La razón expresada a profundidad \(n\) es
\begin{equation}
 \Theta_n
 :=\frac{\langle e_{\rm ar},D_\pi F_{\rm av}^{n}e_{\rm ar}\rangle}
 {\langle e_{\rm vol},F_{\rm av}^{n}e_{\rm vol}\rangle}
 =\pi_{\rm HMT}\frac{F_{n-1}}{F_{n+1}}
 \longrightarrow
 r_{\rm av}:=\frac{\pi_{\rm HMT}}{\varphi_{\rm HMT}^{2}}.
 \label{eq:ym-target-free-av-limit}
\end{equation}
La matriz, la ruta y el cociente se fijan antes de consultar un salto de
masa o un parámetro de teoría gauge. En la profundidad nativa \(n=108\),
\begin{equation}
 F_{107}=10284720757613717413913,
 \qquad
 F_{109}=26925748508234281076009,
 \label{eq:ym-fibonacci-107-109}
\end{equation}
y
\begin{equation}
 \left|\Theta_{108}
 -\frac{\pi_{\rm HMT}}{\varphi_{\rm HMT}^{2}}\right|
 <2\cdot10^{-45}.
 \label{eq:ym-av-depth-108-error}
\end{equation}

\subsection{Transducción dimensional}

Sea \(\ell_0\) la unidad longitudinal del lector dimensional. La escala
inversa de longitud y el parámetro de una losa de espesor \(a_m\) son
\begin{equation}
 \kappa_{\rm av}
 :=\frac{\mu_{\rm vol}}{\ell_0}
 \left(\frac{\pi_{\rm HMT}}{\varphi_{\rm HMT}^{2}}-1\right)>0,
 \qquad
 q_m:=e^{-3a_m\kappa_{\rm av}}.
 \label{eq:ym-dimensional-gap-scale}
\end{equation}
Por tanto \([\kappa_{\rm av}]=L^{-1}\) y la brecha energética es
\begin{equation}
 \Delta_E=3\hbar c\,\kappa_{\rm av};
 \label{eq:ym-energy-gap-units}
\end{equation}
en unidades \(\hbar=c=1\), se escribe \(\Delta_E=3\kappa_{\rm av}\).
La memoria electrónica \(\Delta_4\) pertenece a otro dominio y a otro
lector; no interviene en \eqref{eq:ym-dimensional-gap-scale}.

En los semigrupos de las secciones siguientes, \(t,s,a_m\) son
longitudes euclídeas y los generadores \(D\) y \(H^{\rm ov}\) tienen
dimensión \(L^{-1}\). El Hamiltoniano energético es
\(H_E=\hbar_{\rm int}c_{\rm int}D\). Si
\(\tau=t/c_{\rm int}\), entonces
\[
 e^{-tD}=e^{-\tau H_E/\hbar_{\rm int}},\qquad
 \Delta=3\kappa_{\rm av},\quad
 \Delta_E=\hbar_{\rm int}c_{\rm int}\Delta,\quad
 m_{\rm gap}=\frac{\hbar_{\rm int}\Delta}{c_{\rm int}}.
\]
La misma brecha queda así expresada, respectivamente, como longitud
inversa, energía y masa en la coordenatización dimensional declarada.

Como \(a_m=9a_{m+1}\), la definición da inmediatamente
\begin{equation}
 q_{m+1}^{9}=q_m.
 \label{eq:ym-nonadic-gap-root}
\end{equation}
Esta identidad es la raíz nonádica que transportará el mismo semigrupo y la
misma escala espectral a través de todos los refinamientos.

\section[Proceso gauge euclídeo Yang--Mills en dimensión 4]{Construcción
proyectiva de un proceso gauge euclídeo Yang--Mills, reconstrucción de
Osterwalder--Schrader y brecha colectiva en dimensión \(4\)}


\noindent La especialización de la dinámica nonádica enriquecida,
construida en las secciones precedentes, es
\[
 (\widehat X_\infty,\Gamma_9)
 \xrightarrow{\ \mathcal R_{\rm YM}\ }
 (\mu_m,T_m,H_m,\Omega_m)
 \longrightarrow(\mu_\infty,\mathrm{OS},E(4))
 \longrightarrow\operatorname{gap}(H_{\rm phys})>0.
\]
El lector clover que expresa \(F^2\) reconoce la curvatura sobre este mismo
proceso proyectivo; no construye retrospectivamente \(\mu_m\). La teoría
gauge, la reconstrucción OS, el vacío simple y la brecha ya quedan fijados
por esta cadena. La identidad
\eqref{eq:amp-ym-identidad-accion-requerida} compara después una
parametrización opcional por \(g_R\); es un control no gobernante y no una
condición para identificar la teoría Yang--Mills construida.

\subsection*{Objeto gauge y dominio de refinamiento}

Sea \(G\) un grupo de Lie compacto, conectado y simple. Para cada
\(m\geq0\), se considera la red euclídea
\begin{equation}\label{eq:realizaciones-ym-red}
 \Lambda_m=a_m\mathbb Z^4,
 \qquad a_{m+1}=\frac{a_m}{9}.
\end{equation}
La álgebra cilíndrica gauge se genera por variables de enlace en \(G\),
operadores de entrelazamiento de representación y las coordenadas
enriquecidas de color,
orientación, holonomía, acarreo de fusión, ruta y memoria. Sea \(\mu_m\) la
medida invariante común y sea
\[
 \mathcal H_m=L^2(\mathcal A_m,\mu_m)^{\rm gauge}.
\]
El vector constante \(\Omega_m\) define
\begin{equation}\label{eq:realizaciones-ym-proyectores}
 P_{\Omega_m}=|\Omega_m\rangle\langle\Omega_m|,
 \qquad Q_m^{\rm phys}=I-P_{\Omega_m}.
\end{equation}
El segundo proyector es el complemento físico del vacío; no se identifica
con un residual empleado únicamente para organizar la escala de
renormalización.

El circuito de actualización es local:
\begin{equation}\label{eq:realizaciones-ym-circuito}
 K_m^{\rm loc}
 =\prod_{r=1}^{9\cdot81}
   \prod_{B\in\mathcal B_{m,r}}K_{m,B}.
\end{equation}
La partición \(\{\mathcal B_{m,r}\}_r\) enumera cada coordenada una sola vez; los factores
del circuito se identificarán abajo con el semigrupo de sus expectativas
condicionales, y no con nueve dinámicas independientes.
Los bloques de un mismo color poseen soportes disjuntos y el diámetro físico
de cada soporte está uniformemente acotado al variar \(m\). Esta localidad
forma parte del objeto, de modo que el refinamiento no convierte una
actualización elemental en una operación instantáneamente global.

Sea \(T_m\) la transferencia inducida por
\eqref{eq:realizaciones-ym-circuito}. Las inclusiones isométricas
\(J_m:\mathcal H_m\to\mathcal H_{m+1}\) conservan ancestro,
representación, orientación, acarreo y memoria. Para las cuatro reflexiones
euclídeas \(\Theta_{\nu,m}\), \(\nu=1,\ldots,4\), se cumple
\begin{equation}\label{eq:realizaciones-ym-entrelazamiento}
 (T_{m+1})^9J_m=J_mT_m,
 \qquad
 \Theta_{\nu,m+1}J_m=J_m\Theta_{\nu,m}
 \quad(\nu=1,\ldots,4).
\end{equation}
Las cuatro direcciones pertenecen, por tanto, al mismo refinamiento, la misma
medida y la misma dinámica.

\subsection*{Semigrupo común, corona relativa y cuatro formas de reflexión}

En la coordenatización triangular del estado completo, sean \(E_{m,c}\) las esperanzas
condicionales ortogonales y conmutantes asociadas a coordenadas
independientes. El generador positivo de solapamiento se construye primero en
un nivel semilla \(m_0\):
\begin{equation}\label{eq:realizaciones-ym-generador}
 \widehat H_{m_0}^{\rm ov}
 =3\kappa_{\rm av}
  \sum_{c\in\mathcal I_{m_0}^{\rm seed}}(I-E_{m_0,c}),
\end{equation}
y se prolonga mediante
\begin{equation}\label{eq:realizaciones-ym-generador-refinado}
 \begin{aligned}
 \widehat H_{m+1}^{\rm ov}
 &=\widehat\Gamma_m^*
   \bigl(\widehat H_m^{\rm ov}\otimes I+
         I\otimes\widehat H_{m+1}^{\rm det}\bigr)
   \widehat\Gamma_m,\\
 \widehat H_{m+1}^{\rm det}
 &=3\kappa_{\rm av}
   \sum_{\iota\in\mathcal I_{m+1/m}^{\rm det}}
   (I-E_{m+1,\iota}^{\rm det}),
 \end{aligned}
\end{equation}
donde
\begin{equation}\label{eq:realizaciones-ym-kappa}
 \kappa_{\rm av}
 =\frac{\mu_{\rm vol}}{\ell_0}
  \left(\frac{\pi_{\rm HMT}}{\varphi_{\rm HMT}^{\,2}}-1\right)>0.
\end{equation}
La frecuencia volumétrica \(\mu_{\rm vol}\) es adimensional; la dimensión
procede de \(\ell_0\).

El semigrupo, la transferencia de un paso físico y su compatibilidad de
escala se definen con tipos visibles:
\begin{equation}\label{eq:realizaciones-ym-semigrupo}
 \widehat K_{m,t}^{\rm ov}=e^{-t\widehat H_m^{\rm ov}},
 \qquad
 \widehat T_m=\widehat K_{m,a_m}^{\rm ov},
 \qquad a_{m+1}=a_m/9,
\end{equation}
Como la partición enumera cada coordenada una sola vez y las expectativas
condicionales de sus bloques conmutan, para cada bloque y para el circuito
completo se obtiene
\begin{equation}\label{eq:realizaciones-ym-circuito-semigrupo}
 K_{m,B}=e^{-3\kappa_{\rm av}a_m(I-E_{m,B})},
 \qquad
 K_m^{\rm loc}=e^{-a_m\widehat H_m^{\rm ov}}=\widehat T_m.
\end{equation}
\begin{equation}\label{eq:realizaciones-ym-semigrupo-entrelazado}
 \widehat H_{m+1}^{\rm ov}\widehat J_m
 =\widehat J_m\widehat H_m^{\rm ov},
 \qquad
 (\widehat T_{m+1})^9\widehat J_m
 =\widehat J_m\widehat T_m.
\end{equation}
Así, la novena potencia se compara mediante \(\widehat J_m\); no se igualan
operadores que actúan en Hilbert distintos.

La medida común de ocho caras procede de ese mismo kernel reversible:
\begin{equation}\label{eq:realizaciones-ym-medida-ocho}
 \widehat\Omega_m^{(8)}(dy_1\cdots dy_8)
 =\int\widehat\mu_m^{\rm ov}(dz)
   \prod_{r=1}^{8}
   \widehat K_{m,a_m/2}^{\rm ov}(z,dy_r).
\end{equation}
Para cada reflexión \(\nu=1,\ldots,4\) y todo observable cilíndrico \(F\)
en el semiespacio positivo,
\begin{equation}\label{eq:realizaciones-ym-os}
 \int\overline{F(\Theta_{\nu,m}y)}F(y)\,
 d\widehat\Omega_m^{(8)}
 =\|\widehat{\mathcal B}_{m,\nu}F\|^2\geq0,
\end{equation}
y la marginal de caras opuestas da la identidad de enlace
\begin{equation}\label{eq:realizaciones-ym-os-transferencia}
 \widehat T_{m,\nu}^{\rm OS}
 =\widehat K_{m,a_m}^{\rm ov}
 =e^{-a_m\widehat H_m^{\rm ov}}
 =\widehat T_m,
 \qquad \nu=1,\ldots,4.
\end{equation}
Las cuatro reflexiones expresan, por tanto, el mismo par
transferencia--medida.

\subsection*{Persistencia celular del terminal y brecha exacta}

Sean
\[
 K_9=C_*(Q_{9,3};\mathbb Z),\qquad
 K_{10}=C_*(Q_{10,3};\mathbb Z).
\]
La corona relativa tiene dimensiones
\[
 C_*(Q_{10,3},Q_{9,3})=(271,756,702,217),
 \qquad 271+702=756+217=973.
\]
La inclusión \(j:K_9\to K_{10}\), el colapso celular
\(r:K_{10}\to K_9\) y el homotopio integral \(H_{\rm cell}\) satisfacen
\begin{equation}\label{eq:realizaciones-ym-sdr-corona}
 rj=I,\qquad
 \partial H_{\rm cell}+H_{\rm cell}\partial=I-jr,\qquad
 rH_{\rm cell}=0,\quad H_{\rm cell}j=0,\quad H_{\rm cell}^2=0.
\end{equation}
Para todo mapa de cadenas
\(F:K_{10}\otimes V_{q=6}\to\mathcal T\), con
\(P=(jr)\otimes I_{q=6}\),
\begin{equation}\label{eq:realizaciones-ym-homotopia-terminal}
 F-FP
 =d_{\mathcal T}\bigl(F(H_{\rm cell}\otimes I)\bigr)
  +F(H_{\rm cell}\otimes I)(\partial\otimes I).
\end{equation}
La corona completa es así nulhomotópica y el terminal \(FP\) permanece
literalmente. Los \(271\) vértices no sustituyen las aristas, caras y cubos
de la corona; el factor celular tridimensional tampoco se identifica con el
factor gauge físico \(q=6\).

Para cada coordenada persistente \(c\), en la coordenatización local
\(q=(q_1,\ldots,q_6)\in\mathbb F_3^6\), se fija el testigo centrado
\begin{equation}\label{eq:realizaciones-ym-testigo-definicion}
 \bigcap_c\operatorname{Ran}E_{m,c}=\mathbb C\Omega_m,
 \qquad
 W_c(q)=\mathbf 1_{\{q_1+\cdots+q_6=0\}}-\frac13,
\end{equation}
donde la suma del indicador se toma en \(\mathbb F_3\). La descomposición de
Hoeffding de las expectativas conmutantes prueba
\begin{equation}\label{eq:realizaciones-ym-gap-finito}
 \ker\widehat H_m^{\rm ov}=\mathbb C\Omega_m,
 \qquad
 \widehat H_m^{\rm ov}\succeq
 3\kappa_{\rm av}(I-P_{\Omega_m}).
\end{equation}
La igualdad no se deduce sólo de que el núcleo sea trivial: cada sector no
constante de la descomposición recibe al menos un proyector
\(I-E_{m,c}\) y, por \eqref{eq:realizaciones-ym-generador}, su coste mínimo
es \(3\kappa_{\rm av}\). El testigo material \(W_c\) satisface
\begin{equation}\label{eq:realizaciones-ym-testigo-gap}
 \widehat H_m^{\rm ov}W_c=3\kappa_{\rm av}W_c,
 \qquad \|W_c\|^2=\frac29,
\end{equation}
de modo que el complemento es no vacío y la brecha es exactamente
\(3\kappa_{\rm av}\).

La escala se expresa por
\begin{equation}\label{eq:realizaciones-ym-qm}
 q_m=e^{-3\kappa_{\rm av}a_m},
 \qquad q_{m+1}^{\,9}=q_m,
 \qquad -a_m^{-1}\log q_m=3\kappa_{\rm av}.
\end{equation}
Aunque \(q_m\to1\) al refinar, el cociente espectral por unidad física
permanece constante. La contracción celular
\eqref{eq:realizaciones-ym-homotopia-terminal} conserva el terminal; la
brecha procede del generador positivo de solapamiento y de su descomposición de
Hoeffding, no de la aciclicidad por sí sola.

Las inclusiones de \eqref{eq:realizaciones-ym-semigrupo-entrelazado}
forman el límite inductivo
\begin{equation}\label{eq:realizaciones-ym-colimite}
 \mathcal H_\infty=\varinjlim(\mathcal H_m,\widehat J_m).
\end{equation}
Sobre la unión inductiva algebraica se define la forma
\begin{equation}\label{eq:realizaciones-ym-forma-limite}
 \mathfrak h_\infty[\widehat J_{m,\infty}\psi]
 :=\langle\psi,\widehat H_m^{\rm ov}\psi\rangle_{\mathcal H_m}.
\end{equation}
El entrelazamiento de los generadores hace que esta definición sea
independiente del representante. La forma es positiva y cerrable; su
clausura determina el generador autoadjunto \(\widehat H_\infty^{\rm ov}\).
La compatibilidad de las formas transporta
\begin{equation}\label{eq:realizaciones-ym-gap-forma-limite}
 \overline{\mathfrak h}_\infty[\psi]
 \geq3\kappa_{\rm av}
 \|(I-P_{\Omega_\infty})\psi\|^2,
\end{equation}
y el vector persistente
\(\widehat J_{m,\infty}W_c\) alcanza la igualdad. En consecuencia, la
brecha exacta del generador límite es \(3\kappa_{\rm av}\).

\subsection*{Hipótesis exactas y encadenamiento del teorema}

La representación gauge utiliza:
\begin{enumerate}
 \item el grupo compacto, conectado y simple y la red cuatro-dimensional
       \eqref{eq:realizaciones-ym-red};
 \item el circuito local \eqref{eq:realizaciones-ym-circuito}, con diámetro
       físico uniforme, y una medida invariante común;
 \item las inclusiones isométricas y las cuatro reflexiones entrelazadas por
       \eqref{eq:realizaciones-ym-entrelazamiento} y
       \eqref{eq:realizaciones-ym-semigrupo-entrelazado};
 \item las cuatro factorizaciones
       \eqref{eq:realizaciones-ym-os} y el enlace exacto
       \eqref{eq:realizaciones-ym-os-transferencia} sobre el mismo par
       transferencia--medida;
 \item el generador de expectativas conmutantes
       \eqref{eq:realizaciones-ym-generador}--
       \eqref{eq:realizaciones-ym-generador-refinado} y su descomposición de
       Hoeffding;
 \item el SDR integral de la corona y la conservación del terminal
       \eqref{eq:realizaciones-ym-sdr-corona}--
       \eqref{eq:realizaciones-ym-homotopia-terminal};
 \item el terminal \(P_{\Omega_m}\), el complemento físico
       \(Q_m^{\rm phys}\) y la ley de escala
       \eqref{eq:realizaciones-ym-qm};
 \item el límite inductivo y la clausura compatible de formas
       \eqref{eq:realizaciones-ym-colimite}--
       \eqref{eq:realizaciones-ym-gap-forma-limite};
 \item las coordenatizaciones internas de acción y velocidad
       \(\hbar_{\rm int}\), \(c_{\rm int}\), usadas sólo para expresar la
       dimensión de masa.
\end{enumerate}

\begin{theorem}[Teoría gauge local con vacío simple y brecha colectiva]
\label{thm:realizaciones-ym}
Bajo las premisas anteriores, el límite de la colección gauge es local y
positivo por reflexión en las cuatro direcciones euclídeas. Su generador
posee vacío simple y satisface, en el sector colectivo gauge-invariante,
\begin{equation}\label{eq:realizaciones-ym-gap-limite}
 \Delta_{\rm YM}^{\rm HMT}=3\kappa_{\rm av}>0,
 \qquad
 m_{\rm gap}^{\rm HMT}
 =\frac{\hbar_{\rm int}}{c_{\rm int}}
  \Delta_{\rm YM}^{\rm HMT}>0.
\end{equation}
La brecha pertenece al sector colectivo; no asigna una masa elemental al
gluón.
\end{theorem}

\begin{proof}
La localidad uniforme del circuito conserva la red de álgebras locales. Las
igualdades \eqref{eq:realizaciones-ym-semigrupo-entrelazado} identifican la
novena potencia después de aplicar la inclusión correcta. La construcción
\eqref{eq:realizaciones-ym-medida-ocho} usa ese mismo semigrupo en las ocho
caras y \eqref{eq:realizaciones-ym-os-transferencia} identifica sus cuatro
marginales OS con \(\widehat T_m\). Por tanto, las cuatro formas positivas
\eqref{eq:realizaciones-ym-os} son compatibles en el límite y conservan una
sola medida y una sola dinámica.

Las expectativas de \eqref{eq:realizaciones-ym-generador} son ortogonales y
conmutantes. Su descomposición de Hoeffding separa el modo constante de cada
sector no constante: el primero es \(\mathbb C\Omega_m\), mientras todo
sector del complemento recibe al menos un factor \(I-E_{m,c}\) con peso
\(3\kappa_{\rm av}\). Esto prueba
\eqref{eq:realizaciones-ym-gap-finito}. El testigo
\eqref{eq:realizaciones-ym-testigo-gap} alcanza la cota, de modo que no es
sólo un piso: es la primera energía no nula.

La homotopía \eqref{eq:realizaciones-ym-homotopia-terminal} contrae el
complemento celular y deja intacto \(FP\); por ello la corona no introduce un
segundo vacío ni borra el acarreo. La compatibilidad isométrica transporta el
vacío y su complemento al nivel siguiente sin mezclarlos con el residual de
escala. La ley
\eqref{eq:realizaciones-ym-qm} muestra que el cociente espectral por unidad
física es exactamente \(3\kappa_{\rm av}\) en todo nivel. Las formas
\eqref{eq:realizaciones-ym-forma-limite}--
\eqref{eq:realizaciones-ym-gap-forma-limite} transportan la cota y el testigo
al generador autoadjunto del límite, por lo que la brecha sigue siendo exacta.
La reconstrucción por positividad de
Osterwalder--Schrader identifica el Hamiltoniano con el generador de esta
misma transferencia, conserva el vacío simple y produce
\eqref{eq:realizaciones-ym-gap-limite}. La coordenatización dimensional final transforma
la brecha en la masa colectiva indicada.
\end{proof}

\subsection*{Compresión gauge y testigo físico de la brecha}

Sea
\[
 \widehat{\mathscr H}_m
 :=L^2(\widehat X_m,\widehat\mu_m^{\rm ov})
\]
el espacio de la realización material completa. Además de la acción gauge
usual sobre enlaces y marcos, cada collar \(c\) porta la acción finita de
incidencia
\begin{equation}\label{eq:amp-ym-accion-gauge-incidencia}
 q_c\longmapsto q_c+Bx_c,
 \qquad x_c\in\mathbb F_3^4,
\end{equation}

donde las seis filas de \(B\) están indexadas por
\((01,02,03,12,13,23)\). El producto de ambas acciones preserva
\(\widehat\mu_m^{\rm ov}\). Con \(\mathcal C_m\) el conjunto finito de
collares presentes en el nivel, se escribe
\[
 \mathcal G_m^{\rm full}
 :=G^{V(\Lambda_m)}
   \times\prod_{c\in\mathcal C_m}\mathbb F_3^4.
\]
Su promedio de Haar define la proyección
ortogonal
\begin{equation}\label{eq:amp-ym-proyeccion-fisica-haar}
 \mathbb P_m^{\rm phys}F
 :=\int_{\mathcal G_m^{\rm full}}F(g^{-1}\!\cdot\,)\,dg,
 \qquad
 \widehat{\mathscr H}_m^{\rm phys}
 :=\operatorname{Ran}\mathbb P_m^{\rm phys}.
\end{equation}
Ésta es una compresión sobre la subálgebra gauge del estado completo; la
marginal que olvida las coordenadas de incidencia es otro mapa y no se
confunde con \(\mathbb P_m^{\rm phys}\).

\begin{lemma}[Reducción del Hamiltoniano a la subálgebra física]
\label{lem:amp-ym-compresion-fisica}
La proyección \(\mathbb P_m^{\rm phys}\) reduce el generador y su
semigrupo. Las inclusiones de refinamiento reducen simultáneamente las
subálgebras físicas:
\begin{equation}\label{eq:amp-ym-compresion-entrelazada}
 \begin{gathered}
 \relax
 [\mathbb P_m^{\rm phys},\widehat H_m^{\rm ov}]=0,
 \qquad
 [\mathbb P_m^{\rm phys},e^{-t\widehat H_m^{\rm ov}}]=0,\\
 \mathbb P_{m+1}^{\rm phys}\widehat J_m
 =\widehat J_m\mathbb P_m^{\rm phys},\qquad
 H_m^{\rm phys}
 :=\widehat H_m^{\rm ov}\big|_{\widehat{\mathscr H}_m^{\rm phys}},\\
 H_{m+1}^{\rm phys}\widehat J_m
 =\widehat J_mH_m^{\rm phys}.
 \end{gathered}
\end{equation}
\end{lemma}

\begin{proof}
El generador físico sobre enlaces es gauge equivariante por la
biinvariancia de Haar. En la coordenatización producto, \(E_{q_c}\) es la esperanza
uniforme en \(\mathbb F_3^6\); por invariancia de la medida uniforme bajo
las traslaciones de \eqref{eq:amp-ym-accion-gauge-incidencia}, conmuta con
esa acción. Las expectativas de las subcolas de marcación actúan en
factores disjuntos. Por tanto cada sumando de
\[
 \widehat H_m^{\rm ov}
 =H_m^{\rm ov}\otimes I
 +3\kappa_{\rm av}\sum_c
 \bigl[(I-E_{q_c})+(I-E_{u_c^{\rm P}})\bigr]
\]
conmuta con el promedio \eqref{eq:amp-ym-proyeccion-fisica-haar}. El
cálculo funcional da la conmutación con el semigrupo. Finalmente,
\(\widehat J_m\) conserva los factores ancestrales e inserta la constante
en los factores nuevos; promediar antes o después de esa inclusión produce
la misma función. Esto prueba las cuatro identidades de
\eqref{eq:amp-ym-compresion-entrelazada}.
\end{proof}

\begin{lemma}[Testigo gauge no nulo y entrelazado]
\label{lem:amp-ym-testigo-fisico-q6}
Para todo collar \(c\), la función
\begin{equation}\label{eq:amp-ym-testigo-fisico-def}
 W_c(q)
 :=\mathbf 1_{\{q_1+\cdots+q_6=0\}}-\frac13
\end{equation}
pertenece a \(\widehat{\mathscr H}_m^{\rm phys}\), no es nula y verifica
\begin{equation}\label{eq:amp-ym-testigo-fisico-momentos}
 \mathbb E W_c=0,\qquad
 \|W_c\|^2=\frac29,\qquad
 \mathbb E(W_c^3)=\frac2{27}.
\end{equation}
Además,
\begin{equation}\label{eq:amp-ym-testigo-fisico-gap}
 H_m^{\rm phys}W_c=3\kappa_{\rm av}W_c,\qquad
 H_{m+1}^{\rm phys}\widehat J_mW_c
 =3\kappa_{\rm av}\widehat J_mW_c.
\end{equation}
En consecuencia, el valor \(3\kappa_{\rm av}\) pertenece al espectro de la
restricción física y no sólo al de una dilatación material.
\end{lemma}

\begin{proof}
La acción gauge sobre enlaces actúa en el primer factor y deja fijo
\(W_c\). Para la acción de incidencia, cada \(x_i\) aparece en tres de las
seis coordenadas, luego
\[
 \sum_{i<j}(Bx)_{ij}=3\sum_i x_i=0
 \quad\text{en }\mathbb F_3.
\]
La condición de \eqref{eq:amp-ym-testigo-fisico-def} es, por tanto,
invariante; también lo es bajo las permutaciones de las seis incidencias.
Así \(\mathbb P_m^{\rm phys}W_c=W_c\).

Bajo la medida uniforme, la suma de los seis trits es uniforme en
\(\mathbb F_3\). La función toma el valor \(2/3\) con probabilidad \(1/3\)
y \(-1/3\) con probabilidad \(2/3\), lo que da
\eqref{eq:amp-ym-testigo-fisico-momentos}. El primer sumando de
\(\widehat H_m^{\rm ov}\) anula \(W_c\); para \(c'\ne c\),
\((I-E_{q_{c'}})W_c=0\), mientras
\(E_{q_c}W_c=0\). Todas las expectativas de marcación anulan la
diferencia correspondiente porque \(W_c\) es constante en esas subcolas.
Queda exactamente
\(\widehat H_m^{\rm ov}W_c=3\kappa_{\rm av}W_c\). La primera identidad de
\eqref{eq:amp-ym-testigo-fisico-gap} sigue de la reducción; la segunda,
del entrelazamiento de \eqref{eq:amp-ym-compresion-entrelazada}.
\end{proof}

\subsection*{Semigrupo gauge, vacío y cota espectral exacta}

Sobre la subálgebra física se escribe
\[
 D_m^{\rm YM}:=H_m^{\rm phys},\qquad
 K_{m,t}:=\exp(-tD_m^{\rm YM}),\qquad
 P_{\Omega_m}:=|\Omega_m\rangle\langle\Omega_m|.
\]
La descomposición de Hoeffding del mismo generador y el cálculo funcional
dan, sin cambiar de medida ni de proceso,
\begin{equation}\label{eq:amp-ym-owner-semigrupo-vacio}
 K_{m,t}P_{\Omega_m}=P_{\Omega_m},
 \qquad
 \bigl\|K_{m,t}(I-P_{\Omega_m})\bigr\|
 \leq e^{-3\kappa_{\rm av}t}.
\end{equation}
El testigo de \eqref{eq:amp-ym-testigo-fisico-def} es no nulo y alcanza la
cota:
\begin{equation}\label{eq:amp-ym-owner-testigo-wc}
 D_m^{\rm YM}W_c=3\kappa_{\rm av}W_c,
 \qquad \|W_c\|^2=\frac29.
\end{equation}
Por tanto el vacío es simple y la brecha no sólo está acotada inferiormente,
sino que es exactamente
\begin{equation}\label{eq:amp-ym-owner-brecha-exacta}
 \Delta_{\rm YM}^{\rm HMT}=3\kappa_{\rm av}>0.
\end{equation}
El semigrupo, el proyector de vacío, la contracción del complemento y
\(W_c\) pertenecen a la misma colección proyectiva y constituyen la cadena
propietaria del cierre.

\subsection*{Reconstrucción euclídea, campos de Schwinger y lector de curvatura}

La medida proyectiva común del teorema determina inclusiones isométricas
\(\widehat J_m\) entre los espacios de nivel. Para cada una de las cuatro
reflexiones coordenadas \(\theta_\mu\), \(\mu=1,\ldots,4\), existe una
factorización
\begin{equation}\label{eq:amp-ym-cuatro-os}
 \int\overline{F(\theta_\mu y)}F(y)\,
 d\widehat\Omega_m(y)
 =\|\mathcal B_{m,\mu}F\|_2^2\geq0.
\end{equation}
Las cuatro formas proceden de la misma medida y del mismo semigrupo. Los
conectores OS
\[
 \mathcal J_{m,\mu}^{\rm OS}
 =\mathcal V_{m+1,\mu}^{-1}\widehat J_m\mathcal V_{m,\mu}
\]
entrelazan sus Hamiltonianos. En el colímite,
\begin{equation}\label{eq:amp-ym-hamiltoniano-comun}
 \mathcal V_{\infty,\mu}H_{{\rm OS},\infty}^{(\mu)}
 \mathcal V_{\infty,\mu}^{-1}
 =\widehat H_\infty,
 \qquad \mu=1,\ldots,4.
\end{equation}

La forma límite se diagonaliza por la descomposición de Hoeffding:
\begin{equation}\label{eq:amp-ym-forma-limite}
 \mathcal E_\infty(u)=3\kappa_{\rm av}
 \sum_{S\Subset\mathcal I_\infty}|S|\,\|u_S\|^2.
\end{equation}
Las sumas parciales proporcionan una sucesión de recuperación y la
semicontinuidad produce el límite inferior. La convergencia de Mosco de las
formas cerradas conserva el núcleo unidimensional y el umbral espectral. El
vector material \(W_c\) satisface
\begin{equation}\label{eq:amp-ym-testigo-gap}
 \widehat H_mW_c=3\kappa_{\rm av}W_c,
 \qquad \|W_c\|^2=\frac29,
 \qquad \mathbb E(W_c^3)=\frac2{27},
\end{equation}
de modo que la cota inferior se alcanza y la brecha es exactamente
\(3\kappa_{\rm av}\).

\begin{theorem}[Terminación euclídea y representación de curvatura]
\label{thm:amp-ym-terminacion}
La colección proyectiva del teorema \ref{thm:realizaciones-ym} determina una
acción fuertemente continua de \(E(4)\), una red local gauge no escalar,
funcionales de Schwinger compatibles y una representación de campo en
\(H^{-s}_{\rm loc}(\mathbb R^4)\) para \(s>2\). El producto superficial
ordenado de los enlaces define lectores clover y \(F^2\); en el dominio
suave, el lector de acción converge a
\begin{equation}\label{eq:amp-ym-f2}
 \mathcal S_{\rm YM}(A)
 =\frac1{4g_R^2}\int_{\mathbb R^4}\|F_A(x)\|^2\,d^4x.
\end{equation}
Estas representaciones conservan el vacío simple y la brecha exacta del mismo
Hamiltoniano común.
\end{theorem}

\begin{proof}
Los Gram cofinales de los lectores suavizados y la conmutatividad de los
cuadrados de rotación y refinamiento preservan los ideales nulos; su
colímite construye la representación fuerte de \(E(4)\) y la red local.
Las estimaciones martingala de los incrementos de enlace dan tightness en
\(H^{-s}_{\rm loc}\), \(s>2\), y las características mixtas identifican los
funcionales de Schwinger con el mismo proceso cilíndrico. El producto
orientado alrededor de cada plaqueta construye el clover. Su expansión en
el régimen suave y la convergencia fuerte del lector marginal dan
\eqref{eq:amp-ym-f2}. Ninguna de estas operaciones modifica la medida, el
semigrupo ni la forma límite usados para obtener
\eqref{eq:amp-ym-testigo-gap}; por ello el vacío y la brecha se conservan.
\end{proof}

\subsection*{Identificación exacta del límite gauge}

La identificación de la teoría no descansa en que cuatro reflexiones
generen por sí solas el grupo euclídeo.  Las reflexiones proporcionan la
positividad OS; la acción continua procede, separadamente, del grupoide de
pantallas y de la compatibilidad de refinamiento.  Para
\(g=(g_v)_{v\in\Lambda_m}\in G^{\Lambda_m}\), la acción sobre una variable
de enlace es
\begin{equation}\label{eq:amp-ym-accion-gauge-enlace}
 (g\cdot U)_e=g_{s(e)}U_e g_{t(e)}^{-1}.
\end{equation}
La invariancia de Haar, la desintegración triangular del kernel y la
cancelación de las aristas interiores prueban
\begin{equation}\label{eq:amp-ym-invariancia-gauge-medida}
 (g\cdot)_*\widehat\mu_m^{\rm ov}=\widehat\mu_m^{\rm ov},
 \qquad
 K_{m,t}^{\rm ov}(g\cdot U,g\cdot dV)
 =K_{m,t}^{\rm ov}(U,dV).
\end{equation}
Por diferenciación sobre el núcleo cilíndrico, todo generador vertical
\(X\) satisface la identidad de Ward
\begin{equation}\label{eq:amp-ym-ward}
 \int \nabla_XF\,d\widehat\mu_m^{\rm ov}=0.
\end{equation}

La curvatura se obtiene del mismo sistema de enlaces.  Si una plaqueta
gruesa \(p\) se subdivide en sus 81 subplaquetas \(p'\), entonces
\begin{equation}\label{eq:amp-ym-producto-superficial}
 \operatorname{Hol}^{(m)}_{p}
 =\prod_{p'\subset p}^{\longrightarrow}
 T_{p'}\operatorname{Hol}^{(m+1)}_{p'}T_{p'}^{-1}.
\end{equation}
Las contribuciones de toda arista interior aparecen con orientaciones
opuestas y se cancelan exactamente.  La expresión es gauge-covariante y
conmuta con las inclusiones \(\widehat J_m\).  En una coordenatización suave se define
\begin{equation}\label{eq:amp-ym-clover-definido}
 F_{{\rm cl},m}(x)
 =\frac1{4a_m^2}\sum_{p\ni x}
 \operatorname{Ad}_{T_p}\log\operatorname{Hol}_{p},
 \qquad
 F_{{\rm cl},m}=F_A+O(a_m^2).
\end{equation}
De aquí, para \(A\in C^4_{\rm loc}\) y
\(a_m^2/g_m^2\to0\), se obtiene por suma de Riemann
\begin{equation}\label{eq:amp-ym-f2-convergencia-calculada}
 \frac{a_m^4}{4g_m^2}
 \sum_{x\in\Lambda_m}\|F_{{\rm cl},m}(x)\|^2
 \longrightarrow
 \frac1{4g_R^2}\int_{\mathbb R^4}\|F_A(x)\|^2\,d^4x.
\end{equation}

\begin{theorem}[Proceso gauge euclídeo y lector de curvatura de
Yang--Mills]
\label{thm:amp-ym-identificacion-gauge}
El límite construido en el teorema
\ref{thm:realizaciones-ym} es un proceso gauge euclídeo en dimensión
\(4\) para el grupo compacto, conectado y simple fijado. Posee
covariancia gauge local, una red local no escalar, positividad por
reflexión, una representación fuertemente continua de \(E(4)\),
funcionales de Schwinger compatibles, curvatura gauge-covariante con límite
clásico \eqref{eq:amp-ym-f2-convergencia-calculada}, vacío simple y brecha
espectral
\(3\kappa_{\rm av}>0\).
\end{theorem}

\begin{proof}
Las identidades \eqref{eq:amp-ym-invariancia-gauge-medida}--
\eqref{eq:amp-ym-ward} construyen la simetría gauge y sus identidades de
Ward en todos los niveles; su naturalidad las transporta al límite.  La
localidad uniforme de los kernels produce la red isótona y local.  Las
cuatro formas \eqref{eq:amp-ym-cuatro-os}, junto con la acción del grupoide
de pantallas sobre el núcleo Weyl, producen respectivamente la positividad
OS y la acción fuerte de \(E(4)\).  La tightness en
\(H^{-s}_{\rm loc}\), \(s>2\), identifica una única ley límite y sus
funcionales de Schwinger.  Finalmente,
\eqref{eq:amp-ym-producto-superficial}--
\eqref{eq:amp-ym-f2-convergencia-calculada} identifican el lector de campo
y su límite clásico, mientras
\eqref{eq:amp-ym-forma-limite}--
\eqref{eq:amp-ym-testigo-gap} prueban vacío simple y brecha exacta para la
misma ley.

La medida queda determinada constructivamente por sus marginales
proyectivas, su kernel reversible y sus funcionales cilíndricos. El lector
clover no repondera esa medida: expresa una función del mismo proceso y su
límite clásico.
\end{proof}

\begin{proposition}[Control no gobernante de comparación Gibbs--acción]
\label{prop:amp-ym-puente-accion-requerido}
La ley proyectiva, su vacío simple y su brecha exacta ya están construidos
por \eqref{eq:amp-ym-owner-semigrupo-vacio}--
\eqref{eq:amp-ym-owner-brecha-exacta}. Como comparación exterior opcional de
tipo Gibbs--acción, una parametrización por \(g_R\) puede estudiar una colección
\[
 g_R\longmapsto
 \bigl(\widehat\mu_{m,g_R}^{\rm ov},H_{m,g_R}^{\rm phys}\bigr)
\]
y una relación entre la medida o su forma de Dirichlet y el funcional de
acción, por ejemplo
\begin{equation}\label{eq:amp-ym-identidad-accion-requerida}
 \frac{d\widehat\mu_{m,g_R}^{\rm ov}}{d\nu_m}(U)
 =Z_{m,g_R}^{-1}e^{-S_{m,g_R}^{F^2}(U)},
 \qquad
 \mathfrak h_{m,g_R}[F]
 =\langle F,H_{m,g_R}^{\rm phys}F\rangle,
\end{equation}
con compatibilidad de refinamiento y paso al límite. Aquí \(\nu_m\) puede
ser la medida producto de Haar de la red finita. Esta comprobación se rotula
\texttt{CONTROL\_NO\_GOBERNANTE}: no es premisa de la teoría gauge HMT ni
de su brecha.
\end{proposition}

\begin{proof}
En \eqref{eq:amp-ym-f2-convergencia-calculada}, variar \(g_R\) reescala el
lector de curvatura, mientras las fórmulas que construyen
\(\widehat\mu_m^{\rm ov}\), \(\widehat H_m^{\rm ov}\) y la brecha
\(3\kappa_{\rm av}\) no contienen \(g_R\). Son, por ello, construcciones
independientes. La compresión
\eqref{eq:amp-ym-compresion-entrelazada} y el testigo
\eqref{eq:amp-ym-testigo-fisico-gap} prueban la persistencia física de la
brecha del proceso ya construido; el control sólo compara después una
dependencia dinámica adicional en \(g_R\).
\end{proof}

La secuencia demostrada en esta sección es
\[
 \text{medida común}\to\text{cuatro OS}\to\text{Hamiltoniano común}
 \to E(4),\ \text{Schwinger},\ H^{-s}_{\rm loc}
 \to\text{clover y }F^2.
\]
Esta secuencia demuestra una teoría Yang--Mills cuatridimensional con vacío
simple y brecha positiva exacta \(3\kappa_{\rm av}\) en el dominio HMT
registrado. La identidad
\eqref{eq:amp-ym-identidad-accion-requerida} queda fuera de la cadena
rectora como comparación posterior independiente.

La realización dependiente de \(g\) se desarrolla en
\eqref{eq:pdf2-ym-gibbs-measure} y
\eqref{eq:pdf2-ym-g-dynamics}: la medida, las expectativas
condicionales y el generador se transportan conjuntamente,
y \eqref{eq:pdf2-ym-g-gap-before-limit} conserva la cota y su testigo.


\subsection*{Lectura de Wilson}

El observable de Wilson se obtiene después del generador. Para una coordenatización
arquimediana de enlace,
\begin{equation}\label{eq:realizaciones-ym-wilson}
 A_e=\sum_a\operatorname{ev}_{\mathbb R}(x_{e,a})T_a,
 \qquad
 U_e=e^{a_mA_e},
 \qquad
 W_\gamma=\operatorname{Tr}\!\prod_{e\in\gamma}U_e.
\end{equation}
La lectura conserva la representación, el operador de entrelazamiento y el
orden de la holonomía. Sirve como observable terminal de la teoría expresada; no define
retrospectivamente la medida, el semigrupo ni la brecha.

\subsection*{Necesidad de las hipótesis y obstrucciones exactas}

\begin{ablation}[Necesidad de la dinámica y la medida comunes]
\label{abl:realizaciones-ym}
El teorema \ref{thm:realizaciones-ym} no se conserva si se sustituyen las
nueve subfases por actualizaciones independientes, se emplean medidas o
núcleos diferentes en las cuatro reflexiones, se pierde la localidad física
uniforme, \(\widehat J_m\) deja de entrelazar transferencia y reflexión, se
borran acarreo o memoria, se elimina \(P_{\Omega_m}\), se identifica
\(Q_m^{\rm phys}\) con un residual de escala, se exige
\(q_m\leq q_*<1\) en unidades de retícula, se reemplaza la corona completa
por sus \(271\) vértices, se borra \(FP\), o se usa
\eqref{eq:realizaciones-ym-wilson} para seleccionar el generador.
\end{ablation}

\begin{proof}
Actualizaciones, núcleos o medidas distintos destruyen la factorización
simultánea \eqref{eq:realizaciones-ym-os}. Sin localidad o entrelazamiento,
la positividad de un corte no define una red local compatible. Borrar acarreo,
memoria o \(FP\) elimina información que
\eqref{eq:realizaciones-ym-homotopia-terminal} conserva. Usar sólo los
vértices cambia el complejo relativo y puede fabricar modos cero. Confundir
los dos complementos permite que el sector físico desaparezca por una
operación puramente escalar. La cota fija en unidades de retícula contradice
\eqref{eq:realizaciones-ym-qm}. Finalmente, Wilson sólo es una lectura del
estado construido y no demuestra por sí mismo ni
\eqref{eq:realizaciones-ym-os} ni
\eqref{eq:realizaciones-ym-gap-finito}.
\end{proof}

Los residuos independientes son: el fallo de
\(\widehat T_{m,\nu}^{\rm OS}=e^{-a_m\widehat H_m^{\rm ov}}\), el fallo del
entrelazamiento \eqref{eq:realizaciones-ym-semigrupo-entrelazado}, una forma
OS negativa, un vector de núcleo ortogonal a \(\Omega_m\), o un sector de
Hoeffding no constante con energía inferior a \(3\kappa_{\rm av}\).

\subsection*{Alcance de la construcción gauge}

La prueba combina la conexión conjunta, la conservación del terminal, la
doble lectura, el circuito local, la medida común, las cuatro factorizaciones
de reflexión, el refinamiento entrelazado y el generador positivo. En el
dominio gauge regular declarado, la
localidad, la teoría Yang--Mills cuatridimensional, el vacío simple y la
brecha colectiva son resultados exactos. La reconstrucción de
Osterwalder--Schrader y el observable de Wilson son
representaciones posteriores del mismo estado; ninguna de ellas actúa como
selector retrospectivo de la medida, del semigrupo, del vacío ni de la
brecha. La comparación opcional con una colección parametrizada por \(g_R\) se estudia
mediante \eqref{eq:amp-ym-identidad-accion-requerida}; esa comprobación no
gobierna ni reabre la medida, el semigrupo, el vacío o la brecha ya
demostrados.


\subsection{Incidencia proyectiva de calibre del constructor común y descenso al cociente}

La imagen gauge no es la marginal de enlaces. Su objeto completo por nivel es
\begin{equation}\label{eq:pdf2-g-gau-explicito}
 \mathfrak G_{{\rm gau},m}(d_{\rm gau})=
 \bigl(
 \begin{gathered}
  \mathfrak G_{\rm gau}^{\rm int};\\
  \Lambda_m,\widehat X_m,\widehat\mu_m^{\rm ov},
  \{\widehat\mu_{m,g}^{\rm YM},\nu_{m,g},K_{m,g}^{\rm phys}\}_{g>0},\\
  \mathcal G_m^{\rm full},\mathbb P_m^{\rm phys},
  \widehat J_m,\widehat\Gamma_m,
  \{E_{m,c}\}_c,\widehat H_m^{\rm ov},D_m^{\rm YM},\\
  \{\Theta_{\nu,m}\}_{\nu=1}^4,
  \{\mathfrak A_{+,\nu,m},\operatorname{ev}_{t}^{(\nu)},
      \mathcal V_{m,\nu}^{\rm OS}\}_{\nu=1}^4,\\
  \widehat\Omega_m^{(8)},P_{\Omega_m},W_c,
  \operatorname{Carr}_m,\operatorname{Mem}_m,
  \operatorname{Surv}_m,\tau_m,\mathcal P_m^{F^2}
 \end{gathered}
 \bigr).
\end{equation}
El grupo $\mathcal G_m^{\rm full}$ actúa simultáneamente en enlaces, marcos y
coordenadas de incidencia $q_c\in\mathbb F_3^6$. La proyección física es su promedio
de Haar. El generador reducido es
\begin{equation}\label{eq:pdf2-g-gau-generador-estructural}
 D_m^{\rm YM}
 =\left.3\kappa_{\rm av}\sum_c(I-E_{m,c})
  \right|_{\operatorname{Ran}\mathbb P_m^{\rm phys}},
\end{equation}
y $W_c$ pertenece a esa subálgebra. Las cuatro reflexiones actúan sobre una misma
medida, un mismo kernel y un mismo terminal.
El propietario exacto anterior construye la acción gauge de incidencia
\eqref{eq:amp-ym-accion-gauge-incidencia}, la proyección física de Haar
\eqref{eq:amp-ym-proyeccion-fisica-haar} y la compresión entrelazada
\eqref{eq:amp-ym-compresion-entrelazada}. De esa construcción proceden el
kernel completo equivariable $\widehat K_{m,t}^{\rm ov}$, su descenso al
cociente físico $K_{m,t}^{\rm phys}$, los cuatro operadores de reflexión y
la acción $U_m(g)$ de $E(4)$ en el núcleo cilíndrico suave. Este bloque
ensambla, sin introducirlos como datos, el generador, el vacío, el testigo
de curvatura y las formas OS:
\begin{equation}\label{eq:pdf2-g-gau-entrelazamiento-dato}
 D_{m+1}^{\rm YM}\widehat J_m=\widehat J_mD_m^{\rm YM},
 \quad
 \mathbb P_{m+1}^{\rm phys}\widehat J_m
  =\widehat J_m\mathbb P_m^{\rm phys},
 \quad
 \widehat J_mW_{m,c}=W_{m+1,c}.
\end{equation}
La primera y la tercera identidades son la representación, en la notación
común, de \eqref{eq:amp-ym-compresion-entrelazada} y
\eqref{eq:amp-ym-testigo-fisico-gap}; el semigrupo, el vacío simple y la
brecha exacta ya se obtienen en
\eqref{eq:amp-ym-owner-semigrupo-vacio}--
\eqref{eq:amp-ym-owner-brecha-exacta}. Aquí
$\nu_{m,g}:=(\pi_m^{\rm phys})_*\widehat\mu_{m,g}^{\rm YM}$ y
$K_{m,g}^{\rm phys}(t)=e^{-tD_{m,g}^{\rm YM}}$ en el cociente físico.
La completitud gauge incluye además la desintegración por cada colección de
hiperplanos euclídeos. Si $A_j\in\mathfrak A_{+,\nu,m}$ y
$t_1\le\cdots\le t_n$, la marginal de la misma medida física satisface
\begin{equation}\label{eq:pdf2-g-gau-transfer-euclidean-dato}
\begin{aligned}
 &\int\prod_{j=1}^n
   (\tau_{t_je_\nu}A_j)(x)\,d\widehat\mu_{m,g}^{\rm YM}(x)\\
 &\quad=
 \int\prod_{j=1}^nA_j(y_j)\,
  \nu_{m,g}(dy_1)
  \prod_{j=1}^{n-1}K_{m,g}^{\rm phys}
       (t_{j+1}-t_j;y_j,dy_{j+1})\\
 &\quad=
 \left\langle\mathbf1,
  M_{A_1}e^{-(t_2-t_1)D_{m,g}^{\rm YM}}M_{A_2}\cdots
  e^{-(t_n-t_{n-1})D_{m,g}^{\rm YM}}M_{A_n}\mathbf1
 \right\rangle.
\end{aligned}
\end{equation}
La reconstrucción anterior prueba la identidad en las cuatro direcciones
mediante \eqref{eq:amp-ym-cuatro-os} y construye un único Hamiltoniano
entrelazado en \eqref{eq:amp-ym-hamiltoniano-comun}. En consecuencia, la
traslación OS se expresa aquí como
\begin{equation}\label{eq:pdf2-g-gau-os-generator-dato}
 \mathcal V_{m,\nu}^{\rm OS}T_{m,\nu}^{\rm OS}(s)
 (\mathcal V_{m,\nu}^{\rm OS})^{-1}
 =e^{-sD_{m,g}^{\rm YM}}.
\end{equation}
Así, el generador de expectativas y el Hamiltoniano de traslaciones no se
identifican por nombre: son dos realizaciones ya construidas de la misma
desintegración, cuya forma límite y cuya representación de curvatura constan en
\eqref{eq:amp-ym-forma-limite} y \eqref{eq:amp-ym-f2}.

\section{Imagen gauge y composición hacia Yang--Mills}

\subsection{Estado gauge enriquecido y generadores de transferencia}

El objeto ya generado que recibe este capítulo es la estructura gauge completa
$\mathfrak G_{\rm gau}$; no constituye un dato externo ni una hipótesis terminal.
En cada nivel $m$ conserva la red
$\Lambda_m=a_m\mathbb Z^4$, $a_{m+1}=a_m/9$, la medida $\mu_m$, los
entrelazadores $J_m$, cuatro reflexiones $\Theta_{\nu,m}$, holonomía, color,
representación, acarreo de fusión, ruta, memoria, terminal y lector de curvatura.
Sus dos dinámicas se tipan con símbolos distintos antes de cualquier límite.

El semigrupo de una sola fibra de memoria es
\begin{equation}\label{eq:pdf2-ym-fibra-semigrupo}
 K^{\rm fib}_{m,t}
 =P_{\Omega_m}+e^{-3\kappa_{\rm av}t}(I-P_{\Omega_m})
 =e^{-tD_m^{\rm fib}},
 \qquad
 D_m^{\rm fib}:=3\kappa_{\rm av}(I-P_{\Omega_m}).
\end{equation}
Su espectro está contenido en $\{0,3\kappa_{\rm av}\}$. El generador de la
coordenatización producto completa es, en cambio,
\begin{equation}\label{eq:pdf2-ym-generador-producto}
 D_m^{\rm prod}:=3\kappa_{\rm av}
 \sum_{c\in\mathcal I_m}(I-E_{m,c}),
 \qquad
 K^{\rm prod}_{m,t}:=e^{-tD_m^{\rm prod}},
\end{equation}
donde las expectativas ortogonales $E_{m,c}$ conmutan. Si hay dos o más
coordenadas no constantes, $D_m^{\rm prod}$ posee niveles
$6\kappa_{\rm av},9\kappa_{\rm av},\ldots$; por ello
\eqref{eq:pdf2-ym-fibra-semigrupo} no puede declararse igual a
\eqref{eq:pdf2-ym-generador-producto}. El primero es la fibra simple; el segundo,
el producto de solapamiento.

\subsection{Teorema finito de Hoeffding, vacío y brecha alcanzada}

En la coordenatización producto construida en
\eqref{eq:pdf2-ym-owner-product-factorization}, los $E_{m,c}$ son expectativas
condicionales ortogonales y conmutantes, y la prueba de
\eqref{eq:pdf2-ym-owner-vacuum} establece
\begin{equation}\label{eq:pdf2-ym-interseccion-vacio}
 \bigcap_{c\in\mathcal I_m}\operatorname{Ran}E_{m,c}
 =\mathbb C\Omega_m.
\end{equation}
La descomposición de Hoeffding es la suma ortogonal
\[
 \mathcal H_m=\bigoplus_{S\subseteq\mathcal I_m}\mathcal H_{m,S},
 \qquad u=\sum_Su_S,
\]
donde $E_{m,c}u_S=0$ si $c\in S$ y $E_{m,c}u_S=u_S$ si $c\notin S$.
Entonces
\begin{equation}\label{eq:pdf2-ym-hoeffding}
 D_m^{\rm prod}u_S=3\kappa_{\rm av}|S|u_S,
 \qquad
 \langle u,D_m^{\rm prod}u\rangle
 =3\kappa_{\rm av}\sum_S|S|\|u_S\|^2.
\end{equation}
La componente $S=\varnothing$ es exactamente el vacío de
\eqref{eq:pdf2-ym-interseccion-vacio}. Por tanto
\begin{equation}\label{eq:pdf2-ym-cota-finita}
 D_m^{\rm prod}\succeq
 3\kappa_{\rm av}(I-P_{\Omega_m}),
 \qquad
 \|K^{\rm prod}_{m,t}(I-P_{\Omega_m})\|
 \le e^{-3\kappa_{\rm av}t}.
\end{equation}

En una coordenada de incidencia
$q=(q_1,\ldots,q_6)\in\mathbb F_3^6$, el testigo
\begin{equation}\label{eq:pdf2-ym-wc}
 W_c(q):=\mathbf1_{\{q_1+\cdots+q_6=0\}}-\frac13
\end{equation}
satisface
\begin{equation}\label{eq:pdf2-ym-wc-momentos}
 \mathbb EW_c=0,\qquad \|W_c\|^2=\frac29,\qquad
 D_m^{\rm prod}W_c=3\kappa_{\rm av}W_c.
\end{equation}
Así, en cada nivel finito, el umbral de
\eqref{eq:pdf2-ym-cota-finita} se alcanza. Este cálculo fija el espectro del
generador producto finito; la medida cuatridimensional continua de Yang--Mills se
construye materialmente en
\eqref{eq:pdf2-ym-gibbs-measure}--\eqref{eq:pdf2-ym-g-dynamics}.

\subsection{Compresión gauge física}

Sea $G$ compacto, conectado y simple. El grupo que actúa en el nivel $m$ es
\[
 \mathcal G_m^{\rm full}
 =G^{V(\Lambda_m)}\times
  \prod_{c\in\mathcal C_m}\mathbb F_3^4,
\]
donde el segundo factor actúa sobre las seis incidencias mediante una matriz
que debe quedar visible. Si las filas se indexan por
$(01,02,03,12,13,23)$ y las columnas por $0,1,2,3$, se fija
\begin{equation}\label{eq:pdf2-ym-incidence-matrix}
 B^+_{(ij),k}:=\delta_{ik}+\delta_{jk}\in\mathbb F_3,
 \qquad
 (B^+x)_{ij}=x_i+x_j.
\end{equation}
Esta es la incidencia no orientada del collar; no se confunde con el
coborde firmado de una arista. La acción finita es
\begin{equation}\label{eq:pdf2-ym-incidence-action}
 q_c\longmapsto q_c+B^+x_c,
 \qquad x_c\in\mathbb F_3^4.
\end{equation}
La proyección de Haar
\begin{equation}\label{eq:pdf2-ym-proyeccion-fisica}
 \mathbb P_m^{\rm phys}F
 :=\int_{\mathcal G_m^{\rm full}}F(g^{-1}\!\cdot\,)\,dg,
 \qquad
 \mathcal H_m^{\rm phys}:=\operatorname{Ran}\mathbb P_m^{\rm phys}
\end{equation}
es distinta de la marginal que olvida incidencia. La invariancia de las
expectativas $E_{m,c}$ da
\begin{equation}\label{eq:pdf2-ym-reduccion}
 [\mathbb P_m^{\rm phys},D_m^{\rm prod}]=0,
 \qquad
 [\mathbb P_m^{\rm phys},K_{m,t}^{\rm prod}]=0.
\end{equation}
Además, cada variable de $x_c\in\mathbb F_3^4$ aparece tres veces en la suma de
las seis incidencias, de modo que
\begin{equation}\label{eq:pdf2-ym-incidence-conservation}
 \sum_{i<j}(B^+x_c)_{ij}
 =3\sum_{i=0}^3x_{c,i}=0
 \quad\text{en }\mathbb F_3.
\end{equation}
Por ello $W_c$ es gauge
invariante y pertenece a $\mathcal H_m^{\rm phys}$. La brecha finita no es un
artefacto de una dilatación no física.

\subsection{Positividad por reflexión en cada nivel}

Para un kernel reversible común, defínase la medida de ocho caras
\begin{equation}\label{eq:pdf2-ym-ocho-caras}
 \Omega_m^{(8)}(dy_1\cdots dy_8)
 :=\int\mu_m(dz)\prod_{r=1}^{8}K^{\rm prod}_{m,a_m/2}(z,dy_r).
\end{equation}
Si $\Theta_{\nu,m}$ intercambia las dos mitades respecto de la dirección
$\nu$, la propiedad de Markov y la reversibilidad factorizan, para todo observable
cilíndrico $F$ soportado en el semiespacio positivo,
\begin{equation}\label{eq:pdf2-ym-os-finita}
 \int\overline{F(\Theta_{\nu,m}y)}F(y)\,d\Omega_m^{(8)}
 =\|\mathcal B_{m,\nu}F\|_2^2\ge0,
 \qquad \nu=1,\ldots,4.
\end{equation}
La marginal de caras opuestas recupera el mismo
$K^{\rm prod}_{m,a_m}$. De este modo las cuatro formas OS finitas son
consecuencia de una sola medida y un solo semigrupo, no de cuatro procesos
independientes.

\subsection{Estado gauge completo y coordenadas físicas de incidencia}

El propietario vigente no toma como espacio físico la marginal que conserva sólo
los enlaces. Su dominio es $\mathfrak G_{\rm gau}$, con enlace, marco,
representación, incidencia, color, ruta, orientación,
acarreo, memoria, supervivencia y terminal. La proyección a enlaces es un mapa de
olvido posterior. Juzgar el testigo físico después de aplicar ese olvido sustituye el
objeto y no transfiere la conclusión.

\subsubsection{Medida proyectiva y representación de enlaces}

Para una pantalla finita $\Gamma_m=(V_m,E_m)$ se parte del espacio levantado
\begin{equation}\label{eq:pdf2-ym-owner-lifted-space}
 \widetilde X_m:=G^{V_m}\times G^{E_m},\qquad
 d\widetilde\mu_m(h,z)
 =\prod_{v\in V_m}dh_v\prod_{e\in E_m}k_{\tau_{m,e}}(z_e)\,dz_e,
\end{equation}
y se expresa cada enlace mediante
\begin{equation}\label{eq:pdf2-ym-owner-link-publication}
 U_e=h_{s(e)}z_eh_{t(e)}^{-1}.
\end{equation}
Una arista gruesa $e$ se refina en la cadena ordenada
$e_1\cdots e_9$. Póngase
\(\boldsymbol\tau=(\tau_1,\ldots,\tau_9)\), con
\(\tau_j=\tau_{m+1,e_j}>0\) y
\(\tau=\tau_{m,e}=\sum_{j=1}^9\tau_j\). Sea
\(d\lambda_z^{(9)}\) la desintegración de Haar sobre la fibra del producto,
caracterizada por
\[
 \int_{G^9}F(z_1,\ldots,z_9)\,dz_1\cdots dz_9
 =\int_G\int_{z_1\cdots z_9=z}
   F\,d\lambda_z^{(9)}\,dz.
\]
El puente correcto para tiempos no necesariamente iguales es
\begin{equation}\label{eq:pdf2-ym-owner-heat-bridge}
 dB_{z,\boldsymbol\tau}^{(9)}(z_1,\ldots,z_9)
 =\frac{\prod_{j=1}^9k_{\tau_j}(z_j)}{k_\tau(z)}
   \,d\lambda_z^{(9)},
 \qquad z_1\cdots z_9=z,
\end{equation}
que es una probabilidad porque
\(k_{\tau_1}*\cdots*k_{\tau_9}=k_\tau\). En el refinamiento nonádico uniforme se
recupera \(\tau_j=\tau/9\), pero la fórmula no presupone esa igualdad.

Para hacer independiente la innovación sin borrar el producto grueso, se fija, para
cada $(e,z)$, un transporte boreliano
\begin{equation}\label{eq:pdf2-ym-owner-bridge-triangularization}
 \chi_{m,e,z}:\bigl((0,1)^{\mathbb N},du^{\otimes\mathbb N}\bigr)
 \longrightarrow
 \bigl(\{(z_j):z_1\cdots z_9=z\},B_{z,\boldsymbol\tau}^{(9)}\bigr)
\end{equation}
que es isomorfismo módulo conjuntos nulos. Existe porque ambas son realizaciones
estándar no atómicas; una elección queda fijada una vez y se transporta bajo
relabelados de aristas. La variable gruesa $z$ y la innovación $u^{\rm det}_{m+1,e}$
son entonces factores independientes en la coordenatización triangular. Los marcos nuevos se
integran con Haar. Ruta, acarreo, orientación y memoria se transportan por los
morfismos $\operatorname{Ext}_{\rm gau}$ del propio dato gauge.

La proyección gruesa multiplica los nueve incrementos y olvida exclusivamente las
innovaciones nuevas. Con esta definición se obtiene
\begin{equation}\label{eq:pdf2-ym-owner-projectivity}
 (\widehat\pi_{m+1,m})_*\widehat\mu_{m+1}^{\rm ov}
 =\widehat\mu_m^{\rm ov},\qquad
 \widehat J_mF=F\circ\widehat\pi_{m+1,m}.
\end{equation}
La misma arista aparece una sola vez en la medida; las seis caras incidentes son seis
lectores orientados de esa arista. Las rutas solapadas conservan el incremento
compartido y no se reemplazan por copias independientes.

\subsubsection{La fibra de incidencia es una coordenada física del estado completo}

La cola de incidencia incluida en $\mathfrak G_{\rm gau}$ se descompone sin perder
información:
\begin{equation}\label{eq:pdf2-ym-owner-seed-split}
 \Sigma:(\mathbb F_3^2)^{\mathbb N}\longrightarrow
 (\mathbb F_3^2)^{\mathbb N}\times\mathbb F_3^6,
 \qquad \Sigma_*\sigma=\sigma\otimes u_6.
\end{equation}
El primer factor continúa alimentando las coordenatizaciones de innovación y el segundo conserva
las seis incidencias $q_c=(q_{01},q_{02},q_{03},q_{12},q_{13},q_{23})$ del collar
$c$. En vez de dejar implícito qué coordenadas gobierna el generador, se fija una
coordenatización triangular recursiva
\begin{equation}\label{eq:pdf2-ym-owner-product-factorization}
 \widehat X_m
 =\widehat X_{m_0}^{\rm seed}
  \times\prod_{\ell=m_0+1}^{m}
       \mathcal U_{\ell/\ell-1}^{\rm det},
 \qquad
 \widehat\mu_m^{\rm ov}
 =\widehat\mu_{m_0}^{\rm seed}
  \otimes\bigotimes_{\ell=m_0+1}^{m}
       \upsilon_{\ell/\ell-1}^{\rm det}.
\end{equation}
El factor semilla enumera, una vez, los marcos y enlaces gruesos, las incidencias
$q_c$ y sus coordenatizaciones de marcación. Cada
$\mathcal U_{\ell/\ell-1}^{\rm det}$ enumera, una vez, los marcos nuevos, los puentes
\eqref{eq:pdf2-ym-owner-bridge-triangularization}, las incidencias nuevas y las
subcolas de marcación de ese refinamiento. La factorización es interna a
$\mathfrak G_{\rm gau}$ y no añade otro dato de partida.

Se registra el índice exhaustivo, disjunto y numerable
\begin{equation}\label{eq:pdf2-ym-owner-exhaustive-index}
 \mathcal I_m
 :=\mathcal I_{m_0}^{\rm frame}
 \sqcup\mathcal I_{m_0}^{\rm link}
 \sqcup\mathcal I_{m_0}^{q}
 \sqcup\mathcal I_{m_0}^{\rm mark}
 \sqcup\bigsqcup_{\ell=m_0+1}^{m}
 \bigl(
   \mathcal I_{\ell/\ell-1}^{\rm frame}
   \sqcup\mathcal I_{\ell/\ell-1}^{\rm bridge}
   \sqcup\mathcal I_{\ell/\ell-1}^{q}
   \sqcup\mathcal I_{\ell/\ell-1}^{\rm mark}
 \bigr).
\end{equation}
Cada factor independiente de \eqref{eq:pdf2-ym-owner-product-factorization} aparece
exactamente una vez; una subcola infinita se enumera por sus cilindros escalares. No
queda un factor genealógico fuera de $\mathcal I_m$.
En particular, las abreviaturas de la recurrencia significan literalmente
\begin{equation}\label{eq:pdf2-ym-owner-index-abbreviations}
 \begin{aligned}
 \mathcal I_{m_0}^{\rm seed}
 &:={\mathcal I}_{m_0}^{\rm frame}\sqcup
    {\mathcal I}_{m_0}^{\rm link}\sqcup
    {\mathcal I}_{m_0}^{q}\sqcup
    {\mathcal I}_{m_0}^{\rm mark},\\
 \mathcal I_{\ell/\ell-1}^{\rm det}
 &:={\mathcal I}_{\ell/\ell-1}^{\rm frame}\sqcup
    {\mathcal I}_{\ell/\ell-1}^{\rm bridge}\sqcup
    {\mathcal I}_{\ell/\ell-1}^{q}\sqcup
    {\mathcal I}_{\ell/\ell-1}^{\rm mark}.
 \end{aligned}
\end{equation}

La marginal de enlaces olvida $q_c$; la subálgebra física, en cambio, se obtiene
promediando la acción gauge completa
\begin{equation}\label{eq:pdf2-ym-full-gauge-group}
 \mathcal G_m^{\rm full}
 :=G^{V(\Lambda_m)}\times
 \prod_{c\in\mathcal C_m}\mathbb F_3^4,
 \qquad q_c\longmapsto q_c+B^+x_c.
\end{equation}
Por tanto la proyección física es
\begin{equation}\label{eq:pdf2-ym-full-physical-projection}
 \mathbb P_m^{\rm phys}F
 :=\int_{\mathcal G_m^{\rm full}}F(g^{-1}\!\cdot x)\,dg,
 \qquad
 \widehat{\mathscr H}_m^{\rm phys}
 :=\operatorname{Ran}\mathbb P_m^{\rm phys}.
\end{equation}
No coincide con el mapa de olvido
$\widehat X_m\to X_m^{\rm link}$. Esta distinción es el cortafuegos de tipo
que impide expulsar $W_c$ del sector físico.

\subsection{Generador gauge y circuito local de transporte}

La coordenatización triangular completa proporciona una unitaria
\begin{equation}\label{eq:pdf2-ym-owner-triangular-unitary}
 \widehat\Gamma_m:
 L^2(\widehat X_{m+1},\widehat\mu_{m+1}^{\rm ov})
 \longrightarrow
 L^2(\widehat X_m,\widehat\mu_m^{\rm ov})
 \widehat\otimes\mathcal K_{m+1}^{\rm det},
 \qquad
 \widehat\Gamma_m\widehat J_mF=F\otimes\mathbf1_{\rm det}.
\end{equation}
Para $\iota\in\mathcal I_m$, sea $E_{m,\iota}$ la esperanza que integra
exactamente el factor $\iota$ de
\eqref{eq:pdf2-ym-owner-product-factorization} y deja fijos todos los demás.
Son proyecciones Markov ortogonales, conservan la unidad y conmutan. Sobre el
núcleo denso $\mathscr C_m$ de cilindros que dependen de un número finito de
coordenadas se define la forma
\begin{equation}\label{eq:pdf2-ym-owner-full-form}
 \widehat{\mathcal E}_m[F]
 :=3\kappa_{\rm av}\sum_{\iota\in\mathcal I_m}
   \|(I-E_{m,\iota})F\|_2^2.
\end{equation}
La suma es finita para cada cilindro. Su clausura es una forma de Dirichlet y
determina el generador autoadjunto Markov
$\widehat H_m^{\rm ov}$. En un nivel semilla esto se escribe, en sentido de
formas,
\begin{equation}\label{eq:pdf2-ym-owner-base-generator}
 \widehat H_{m_0}^{\rm ov}
 =3\kappa_{\rm av}
  \sum_{\iota\in\mathcal I_{m_0}^{\rm seed}}(I-E_{m_0,\iota}),
\end{equation}
y cada refinamiento añade exactamente las expectativas de sus nuevas coordenadas:
\begin{equation}\label{eq:pdf2-ym-owner-refined-generator}
 \begin{aligned}
 \widehat H_{m+1}^{\rm ov}
 &=\widehat\Gamma_m^*
   (\widehat H_m^{\rm ov}\otimes I
    +I\otimes\widehat H_{m+1}^{\rm det})\widehat\Gamma_m,\\
 \widehat H_{m+1}^{\rm det}
 &=3\kappa_{\rm av}
   \sum_{\iota\in\mathcal I_{m+1/m}^{\rm det}}
   (I-E_{m+1,\iota}^{\rm det}).
 \end{aligned}
\end{equation}
La partición de colores no se deja nominal. Si una celda se indexa por
$n=(n_1,n_2,n_3,n_4)\in\mathbb Z^4$, se fija
\begin{equation}\label{eq:pdf2-ym-owner-729-coloring}
 \chi_m(n):=
 \bigl(n_1\bmod9,\ n_2\bmod9,\ n_3+3n_4\bmod9\bigr)
 \in(\mathbb Z/9\mathbb Z)^3,
 \qquad |(\mathbb Z/9\mathbb Z)^3|=9\cdot81.
\end{equation}
Si dos collares cerrados de radio celular uno se intersectan, su diferencia
$\delta$ satisface $|\delta_j|\le2$. La igualdad de colores fuerza
$\delta_1=\delta_2=0$ y
$\delta_3+3\delta_4\equiv0\pmod9$. Como el último entero pertenece a
$[-8,8]$, vale cero; las cotas $|\delta_3|,|\delta_4|\le2$ dan entonces
$\delta_3=\delta_4=0$. Por tanto, celdas distintas del mismo color tienen
collares disjuntos. Se define materialmente
\begin{equation}\label{eq:pdf2-ym-owner-729-blocks}
 \mathcal B_{m,r}:=\{\iota\in\mathcal I_m:
   \operatorname{cell}(\iota)=n, \chi_m(n)=r\},
 \qquad r\in(\mathbb Z/9\mathbb Z)^3.
\end{equation}
El transporte de una pantalla por $g\in E(4)$ transporta también esta
partición; no exige que una rotación preserve las coordenadas del rótulo.
Así, $\{\mathcal B_{m,r}\}_{r=1}^{9\cdot81}$ agrupa soportes espaciales
disjuntos;
el índice interno de cada bloque enumera sus factores de
\eqref{eq:pdf2-ym-owner-exhaustive-index}. Para un cilindro $F$, el producto sólo
contiene los bloques de los que depende $F$; la clausura define el producto fuerte
para toda la forma. El generador de un bloque conserva las multiplicidades de
Hoeffding; no se sustituye por un solo proyector que las borraría. Por
conmutatividad en cada coordenatización,
\begin{equation}\label{eq:pdf2-ym-owner-local-circuit}
\begin{aligned}
 H_{m,B}&:=3\kappa_{\rm av}\sum_{\iota\in B}(I-E_{m,\iota}),\\
 K_{m,B}&:=e^{-a_mH_{m,B}}
 =\prod_{\iota\in B}e^{-3\kappa_{\rm av}a_m(I-E_{m,\iota})},\\
 K_m^{\rm loc}&:=\mathop{\prod_{r,B}}^{\rm strong}K_{m,B}
 =e^{-a_m\widehat H_m^{\rm ov}}.
\end{aligned}
\end{equation}
El diámetro físico de cada bloque permanece uniformemente acotado. De aquí procede
la localidad: el generador no introduce una actualización instantánea sobre toda la
red.
La partición es compatible con Bianchi, no sólo con disjunción geométrica. Para una
celda cúbica $c$ se fija el producto orientado, con todos los factores transportados
al mismo vértice,
\begin{equation}\label{eq:pdf2-ym-owner-bianchi-block-product}
 \mathscr B_{m,c}(x):=
 \prod_{p\subset\partial c}^{\longrightarrow}
 T_p\operatorname{Hol}_p(x)T_p^{-1}.
\end{equation}
Cada arista interna de la frontera aparece dos veces y con orientaciones contrarias;
por tanto $\mathscr B_{m,c}=e$ punto a punto. Un bloque
$\mathcal B_{m,r}$ contiene el collar completo de toda celda que actualiza: si toca
una cara de $c$, actualiza simultáneamente el par inverso que comparte esa arista.
La cancelación anterior sigue siendo literal después de cada factor del circuito:
\begin{equation}\label{eq:pdf2-ym-owner-bianchi-block-invariance}
 K_{m,B}\mathbf1_{\{\mathscr B_{m,c}=e\}}
 =\mathbf1_{\{\mathscr B_{m,c}=e\}},
 \qquad
 K_m^{\rm loc}\mathbf1_{\{\mathscr B_{m,c}=e\}}
 =\mathbf1_{\{\mathscr B_{m,c}=e\}}.
\end{equation}
Así el calendario de $9\cdot81$ capas preserva Bianchi en cada paso y en el producto
fuerte, no sólo al final de una barrida.

La acción gauge permuta o traslada las coordenadas integradas por cada
$E_{m,\iota}$ y
preserva sus medidas. Por ello
\begin{equation}\label{eq:pdf2-ym-owner-physical-reduction}
 U_m(g)\widehat H_m^{\rm ov}=\widehat H_m^{\rm ov}U_m(g)
 \quad(g\in\mathcal G_m^{\rm full}),
 \qquad
 [\mathbb P_m^{\rm phys},\widehat H_m^{\rm ov}]=0,
 \qquad
 \mathbb P_{m+1}^{\rm phys}\widehat J_m
 =\widehat J_m\mathbb P_m^{\rm phys}.
\end{equation}
El semigrupo completo y su kernel Markov se denotan por
\begin{equation}\label{eq:pdf2-ym-owner-full-kernel}
 \widehat K_{m,t}^{\rm ov}:=e^{-t\widehat H_m^{\rm ov}},
 \qquad
\widehat K_{m,t}^{\rm ov}(x,dy)
 \quad(x,y\in\widehat X_m).
\end{equation}
La primera identidad de \eqref{eq:pdf2-ym-owner-physical-reduction} implica
\begin{equation}\label{eq:pdf2-ym-owner-full-kernel-equivariance}
 \widehat K_{m,t}^{\rm ov}(gx,gA)
 =\widehat K_{m,t}^{\rm ov}(x,A)
 \quad(g\in\mathcal G_m^{\rm full}).
\end{equation}
Para tipar la reducción sin usar un kernel restringido sobre el espacio equivocado,
sea $q_m:\widehat X_m\to Y_m:=\widehat X_m/\mathcal G_m^{\rm full}$ el mapa de
órbitas y $\nu_m=(q_m)_*\widehat\mu_m^{\rm ov}$. El pullback
\begin{equation}\label{eq:pdf2-ym-owner-quotient-unitary}
 \mathcal W_m:L^2(Y_m,\nu_m)\longrightarrow
 \widehat{\mathscr H}_m^{\rm phys},
 \qquad \mathcal W_m\phi=\phi\circ q_m
\end{equation}
es unitario. La equivariancia del kernel completo permite definir
\begin{equation}\label{eq:pdf2-ym-owner-physical-kernel}
 K_{m,t}^{\rm phys}(q_mx,A)
 :=\widehat K_{m,t}^{\rm ov}(x,q_m^{-1}A),
 \qquad A\subseteq Y_m,
\end{equation}
independientemente del representante $x$. El objeto físico queda entonces tipado por
\begin{equation}\label{eq:pdf2-ym-owner-D}
 D_m^{\rm YM}
 :=\left.\widehat H_m^{\rm ov}\right|_{\widehat{\mathscr H}_m^{\rm phys}},
 \qquad
 e^{-tD_m^{\rm YM}}
 =\mathcal W_mK_{m,t}^{\rm phys}\mathcal W_m^{-1}.
\end{equation}
Éste es el generador producto completo; no es el generador de una sola fibra de
\eqref{eq:pdf2-ym-fibra-semigrupo}.

\subsection{Acción de Yang--Mills, medida dependiente de $g$ y transporte exacto}

Hasta aquí se ha escrito la coordenatización cinemática de referencia. La denotaremos
$\widehat\mu_m^0:=\widehat\mu_m^{\rm ov}$ y
$\widehat H_m^0:=\widehat H_m^{\rm ov}$, con
$\widehat K_{m,t}^0:=e^{-t\widehat H_m^0}=\widehat K_{m,t}^{\rm ov}$.
Para convertir el lector estructural
$\mathcal P_m^{F^2}$ en una ley, y no en una comparación posterior, se descarga
primero su raíz positiva sobre la factorización triangular. La polarización de
$\mathcal P_m^{F^2}$ en un collar $b$ es el Gram positivo
\[
 Q_{m,b}(u,v):=\frac14\bigl(
 \mathcal P_{m,b}^{F^2}(u+v)-\mathcal P_{m,b}^{F^2}(u-v)
 \bigr).
\]
Se divide el espacio finito de lectores del collar por $\ker Q_{m,b}$, se completa
con $Q_{m,b}$ y se denota por $\mathcal E_{m,b}^{F}$ el Hilbert resultante. Si
$\zeta_{m,b}(x_b)$ es el vector de lectores centrados de la coordenada triangular
$x_b$, se fija materialmente
\begin{equation}\label{eq:pdf2-ym-action-root}
 \xi_{m,b}(x_b):=Q_{m,b}^{1/2}\zeta_{m,b}(x_b)
 \in\mathcal E_{m,b}^{F},
 \qquad
 p_{m,b}(x_b):=\|\xi_{m,b}(x_b)\|^2.
\end{equation}
El índice $b$ recorre los collares Gram de soporte mínimo. Un collar se asigna
al primer color de \eqref{eq:pdf2-ym-owner-729-coloring} que contiene su celda base;
así cada término aparece una vez y los collares de un mismo color son disjuntos.
Más precisamente, la diagonalización positiva agrupa los factores triangulares en
una partición disjunta
\begin{equation}\label{eq:pdf2-ym-action-factor-partition}
 \mathcal I_m^{F}
 =\bigsqcup_{r\in(\mathbb Z/9\mathbb Z)^3}
   \bigsqcup_{b\in\mathcal B_{m,r}^{F}}\mathcal I_{m,b},
 \qquad x_b:=(x_\iota)_{\iota\in\mathcal I_{m,b}},
 \qquad
 \mathcal B_{m,r}^{F}:=\{b:Q_{m,b}\ne0,\ \chi_m(b)=r\}.
\end{equation}
Ninguna coordenada triangular pertenece a dos $\mathcal I_{m,b}$; si un lector
geométrico toca dos collares, su modo de Gram se asigna al primero y su complemento
ortogonal al segundo. Ésta es la descomposición espectral finita de la forma positiva,
no una duplicación de la variable física.
La raíz de la suma ortogonal, calculada color por color, da la identidad, no una
aproximación,
\begin{equation}\label{eq:pdf2-ym-action-reader-sum}
 \mathcal P_m^{F^2}(x)
 =\sum_{r\in(\mathbb Z/9\mathbb Z)^3}
   \sum_{b\in\mathcal B_{m,r}^{F}}p_{m,b}(x_b),
 \qquad
 S_{m,g}(x):=\frac1{4g^2}\mathcal P_m^{F^2}(x),
 \quad g\in(0,\infty).
\end{equation}
Aquí $x_b:\widehat X_m\to X_{m,b}$ es la proyección de coordenadas del
collar, de modo que dominio, acción y codominio de cada sumando quedan fijados.
El lector es gauge invariante porque $Q_{m,b}$ usa el producto invariante de
$\mathfrak g$; es invariante por reflexión porque una reflexión sólo permuta los
seis pares orientados; y es natural por relabelado euclídeo.
La partición \eqref{eq:pdf2-ym-action-factor-partition} y la coordenatización producto dan
explícitamente
\begin{equation}\label{eq:pdf2-ym-action-factor-measures}
 \widehat\mu_m^0=\lambda_m^{\rm inert}\otimes
   \bigotimes_{r,b}\lambda_{m,b}^0,
 \qquad
 d\lambda_{m,b}^g
 :=(Z_{m,b}^g)^{-1}e^{-p_{m,b}/(4g^2)}d\lambda_{m,b}^0,
 \qquad
 \widehat\mu_{m,g}^{\rm YM}=\lambda_m^{\rm inert}\otimes
   \bigotimes_{r,b}\lambda_{m,b}^g.
\end{equation}
El factor inerte contiene exactamente las coordenadas sobre las que el lector es
cero; no se elimina de la teoría.

La separación grueso--detalle de
\eqref{eq:pdf2-ym-owner-product-factorization} es ortogonal para la misma raíz:
la polarización directa da
$Q_{m+1}=Q_m\oplus Q_{m+1/m}^{\rm det}$ y
$\zeta_{m+1}=\zeta_m\oplus\zeta_{m+1/m}^{\rm det}$.
Si $x_{m+1}=\widehat\Gamma_m^{-1}(x_m,u)$, entonces
\begin{equation}\label{eq:pdf2-ym-action-cocycle}
 \xi_{m+1}(x_{m+1})
 =\xi_m(x_m)\oplus\xi_{m+1/m}^{\rm det}(u),
 \qquad
 S_{m+1,g}\circ\widehat\Gamma_m^{-1}
 =S_{m,g}+S_{m+1/m,g}^{\rm det}.
\end{equation}
Esta es la forma precisa en que carry y memoria impiden contar dos veces una cara
ancestral: la parte gruesa no se vuelve a sumar en el detalle. En particular,
$Z_{m+1,g}=Z_{m,g}Z_{m+1/m,g}^{\rm det}$ para los normalizadores y se obtiene la
colección de probabilidades genuinamente dependiente del acoplamiento
\begin{equation}\label{eq:pdf2-ym-gibbs-measure}
 d\widehat\mu_{m,g}^{\rm YM}(x)
 :=Z_{m,g}^{-1}e^{-S_{m,g}(x)}d\widehat\mu_m^0(x),
 \qquad
 (\widehat\pi_{m+1,m})_*\widehat\mu_{m+1,g}^{\rm YM}
 =\widehat\mu_{m,g}^{\rm YM}.
\end{equation}
Para escribir correctamente Schwinger--Dyson respecto de Haar, y no fingir que la
medida de referencia es plana, sea $\lambda_m$ el producto de Haar, Lebesgue y
conteo uniforme de la coordenatización levantada, y sea
$R_m:=d\widehat\mu_m^0/d\lambda_m>0$. La acción total regulada es
\begin{equation}\label{eq:pdf2-ym-total-regulated-action}
 \mathscr S_{m,g}:=S_{m,g}-\log R_m,
 \qquad
 d\widehat\mu_{m,g}^{\rm YM}
 =\mathscr Z_{m,g}^{-1}e^{-\mathscr S_{m,g}}d\lambda_m.
\end{equation}
El término $-\log R_m$ contiene los kernels de calor y las coordenatizaciones de puente de la
estructura gauge de partida; es independiente de $g$. El término de Yang--Mills que
cambia el acoplamiento y cuyo límite clásico se calcula abajo es exactamente
$S_{m,g}=\mathcal P_m^{F^2}/(4g^2)$.
La consistencia de la segunda igualdad y la estandaridad de los espacios finitos
permiten aplicar Kolmogórov y definen la única medida cilíndrica
$\widehat\mu_{\infty,g}^{\rm YM}$ cuyas marginales son
$\widehat\mu_{m,g}^{\rm YM}$.
No se ha usado aquí ningún autovalor: primero se han fijado $g$, la acción y
la medida.

Para definir la dinámica sobre esa medida sin perder producto, gauge ni las
identidades superficiales, se construye el transporte, no se lo postula. En cada
factor no atómico $X_{m,b}$ se fija una coordenatización boreliana
$\rho_{m,b}:X_{m,b}/\mathcal G_{m,b}\to(0,1)$ y se desintegra primero sobre el
cociente gauge y después sobre la órbita de Haar. Sean $F_{m,b}^0$ y
$F_{m,b}^g$ las funciones de distribución continuas de $\rho_{m,b}$ para la
medida de referencia y la medida reponderada. El transporte triangular es
\begin{equation}\label{eq:pdf2-ym-gibbs-transport-factor}
 t_{m,b,g}:=(F_{m,b}^g)^{-1}\circ F_{m,b}^0,
 \qquad
 T_{m,b,g}(y,h):=(t_{m,b,g}(y),h),
 \qquad
 (T_{m,b,g})_*\lambda_{m,b}^0=\lambda_{m,b}^g.
\end{equation}
En un factor puramente atómico sobre el que $p_{m,b}=0$ se toma
$T_{m,b,g}=I$. En presencia de átomos y una coordenada continua, la coordenatización es el
orden lexicográfico átomo--intervalo y la misma fórmula se aplica al espacio
estándar no atómico total. Los transportes de los $9\cdot81$ colores se componen
en el orden fijado por \eqref{eq:pdf2-ym-owner-729-coloring}; dentro de un color
conmutan porque los collares son disjuntos. Se obtiene
\begin{equation}\label{eq:pdf2-ym-gibbs-transport-global}
 T_{m,g}:=\prod_{r\in(\mathbb Z/9\mathbb Z)^3}
           \prod_{b\in\mathcal B_{m,r}^{F}}T_{m,b,g},
 \quad
 (T_{m,g})_*\widehat\mu_m^0=\widehat\mu_{m,g}^{\rm YM},
 \quad
 \widehat\pi_{m+1,m}T_{m+1,g}=T_{m,g}\widehat\pi_{m+1,m}.
\end{equation}
La última igualdad se comprueba factor a factor usando
\eqref{eq:pdf2-ym-action-cocycle}: el transporte del factor grueso es el ya fijado y
los transportes nuevos actúan sólo en la innovación. Las coordenatizaciones de cociente se
eligen en una pantalla representante y se trasladan; por ello $T_{m,g}$ conmuta con
gauge y reflexión y, para $a\in E(4)$, satisface
$T_{a\alpha,g}=U_\alpha(a)T_{\alpha,g}U_\alpha(a)^{-1}$.

La composición
\begin{equation}\label{eq:pdf2-ym-full-probability-isomorphism}
 \mathcal T_{m,g}:L^\infty(\widehat X_m,\widehat\mu_{m,g}^{\rm YM})
 \longrightarrow L^\infty(\widehat X_m,\widehat\mu_m^0),
 \qquad \mathcal T_{m,g}F:=F\circ T_{m,g}
\end{equation}
es, por construcción, un isomorfismo de espacios de probabilidad y de
$*$-álgebras, y se prolonga unitariamente a $L^2$. En particular,
\begin{equation}\label{eq:pdf2-ym-full-algebra-identities}
 \mathcal T_{m,g}(FG)=\mathcal T_{m,g}F\,\mathcal T_{m,g}G,
 \quad
 \mathcal T_{m,g}(F^*)=(\mathcal T_{m,g}F)^*,
 \quad
 \int F\,d\widehat\mu_{m,g}^{\rm YM}
 =\int\mathcal T_{m,g}F\,d\widehat\mu_m^0.
\end{equation}
El dominio de cada $T_{m,b,g}$ es el propio espacio de holonomías que satisface
la ley de subdivisión y la restricción de Bianchi. Por tanto su imagen permanece
en ese espacio. Aplicando la multiplicatividad anterior generador a generador se
obtienen las identidades puntuales
\begin{equation}\label{eq:pdf2-ym-transport-surface-bianchi}
 \begin{aligned}
 \mathcal T_{m,g}(\operatorname{Hol}_p)
 &=\prod_{p'\subset p}^{\longrightarrow}
   \mathcal T_{m,g}(T_{p'}\operatorname{Hol}_{p'}T_{p'}^{-1}),\\
 \mathcal T_{m,g}(F_{{\rm cl},m})
 &=F_{{\rm cl},m}\circ T_{m,g},\\
 \prod_{p\subset\partial c}^{\longrightarrow}
 \mathcal T_{m,g}(T_p\operatorname{Hol}_pT_p^{-1})&=e,
 \qquad
 \operatorname{Alt}_{\lambda\mu\nu}
 D_\lambda(F_{{\rm cl},m}\circ T_{m,g})_{\mu\nu}=0.
 \end{aligned}
\end{equation}
La tercera línea es la cancelación arista por arista de la frontera orientada de
la frontera de $c$; la cuarta es su polarización infinitesimal. Así el transporte
preserva productos, estado, clover y Bianchi, no sólo el primer caos.

Finalmente se transportan expectativas, generador, kernel, proyección e inclusión:
\begin{equation}\label{eq:pdf2-ym-g-dynamics}
 \begin{aligned}
 E_{m,\iota}^{g}&:=\mathcal T_{m,g}^{-1}E_{m,\iota}\mathcal T_{m,g},&
 \widehat H_{m,g}^{\rm YM}&:=\mathcal T_{m,g}^{-1}\widehat H_m^0\mathcal T_{m,g},\\
 \widehat K_{m,t}^{g}&:=\mathcal T_{m,g}^{-1}\widehat K_{m,t}^{0}\mathcal T_{m,g},&
 \mathbb P_{m,g}^{\rm phys}&:=\mathcal T_{m,g}^{-1}
     \mathbb P_m^{\rm phys}\mathcal T_{m,g},\\
 \widehat J_{m,g}&:=\mathcal T_{m+1,g}^{-1}
     \widehat J_m\mathcal T_{m,g},&
 D_{m,g}^{\rm YM}&:=\left.\widehat H_{m,g}^{\rm YM}
     \right|_{\operatorname{Ran}\mathbb P_{m,g}^{\rm phys}}.
 \end{aligned}
\end{equation}
La primera línea no es sólo una conjugación abstracta. Usando
\eqref{eq:pdf2-ym-action-factor-measures}, su acción sobre un cilindro es la
actualización de baño térmico de la acción:
\begin{equation}\label{eq:pdf2-ym-gibbs-conditional-expectation}
 (E_{m,b}^{g}F)(x_{\ne b})
 =\frac{\displaystyle
   \int_{X_{m,b}}F(x_{\ne b},u)e^{-p_{m,b}(u)/(4g^2)}
       d\lambda_{m,b}^0(u)}{\displaystyle
   \int_{X_{m,b}}e^{-p_{m,b}(u)/(4g^2)}d\lambda_{m,b}^0(u)},
 \qquad
 \widehat H_{m,g}^{\rm YM}
 =3\kappa_{\rm av}\sum_b(I-E_{m,b}^{g}).
\end{equation}
Así tanto la medida como cada transición local contienen $g$ y la acción
$p_{m,b}/(4g^2)$ de forma explícita.
Son Markov porque $\mathcal T_{m,g}\mathbf1=\mathbf1$, locales porque cada
$T_{m,b,g}$ tiene el mismo collar que $p_{m,b}$, y dependen de $g$ por
\eqref{eq:pdf2-ym-gibbs-transport-factor}. La equivalencia es espectral y no
selecciona la medida a partir del espectro:
Para cerrar además las ecuaciones dinámicas, sea $X_{m,b}^a$ el campo
izquierdo-invariante en el factor de enlace $b$ asociado a
$e_a\in\mathfrak g$; en una coordenatización continua de innovación se usa la derivada
direccional ordinaria, y las coordenadas finitas se tratan por diferencia exacta.
Para todo cilindro suave $F$ el dominio común es
$\mathscr C_m^1:=\bigcap_{b,a}\operatorname{Dom}X_{m,b}^a$, tomando soporte
compacto en las coordenatizaciones abiertas de innovación; en los factores de grupo no hay
frontera. La invariancia de Haar y la regla del producto calculan
\begin{equation}\label{eq:pdf2-ym-schwinger-dyson-finite}
 \int X_{m,b}^aF\,d\widehat\mu_{m,g}^{\rm YM}
 =\int F\,X_{m,b}^a\mathscr S_{m,g}\,d\widehat\mu_{m,g}^{\rm YM},
 \qquad F\in\mathscr C_m^1.
\end{equation}
En una coordenada finita, para
$\Delta_\varepsilon F(q):=F(q+\varepsilon)-F(q)$, el cambio de variable
$q\mapsto q+\varepsilon$ da la versión exacta, sin derivada formal,
\begin{equation}\label{eq:pdf2-ym-schwinger-dyson-discrete}
 \int\Delta_\varepsilon F(q)\,d\widehat\mu_{m,g}^{\rm YM}
 =\int F(q)
  \left(e^{\mathscr S_{m,g}(q)-\mathscr S_{m,g}(q-\varepsilon)}-1\right)
  d\widehat\mu_{m,g}^{\rm YM}(q).
\end{equation}
Si $\delta_\varepsilon$ es la suma orientada de esos campos en las aristas que
inciden en un vértice, la invariancia gauge de la medida de referencia y de
$S_{m,g}$ da $\delta_\varepsilon\mathscr S_{m,g}=0$ y, por tanto, la identidad de Ward
\begin{equation}\label{eq:pdf2-ym-ward-finite}
 \int\delta_\varepsilon F\,d\widehat\mu_{m,g}^{\rm YM}=0.
\end{equation}
Para el factor gauge finito $\mathbb F_3^4$, la invariancia
$\mathscr S_{m,g}(q+\varepsilon)=\mathscr S_{m,g}(q)$ reduce
\eqref{eq:pdf2-ym-schwinger-dyson-discrete} a la misma identidad Ward con
$\delta_\varepsilon$ reemplazada por $\Delta_\varepsilon$.
Las dos expresiones son compatibles con $\widehat J_{m,g}$ porque
\eqref{eq:pdf2-ym-action-cocycle} separa las derivadas gruesas de las de detalle.
Todo cilindro del límite reside en algún nivel; así
\eqref{eq:pdf2-ym-schwinger-dyson-finite} y
\eqref{eq:pdf2-ym-ward-finite} pasan literalmente al estado cofinal. Ni la
ecuación Schwinger--Dyson ni Ward contienen la brecha como premisa.

\begin{equation}\label{eq:pdf2-ym-g-gap-before-limit}
 D_{m,g}^{\rm YM}\succeq3\kappa_{\rm av}
 (I-P_{\Omega_{m,g}}),
 \qquad
 W_{c,g}:=\mathcal T_{m,g}^{-1}W_c,
 \qquad
 D_{m,g}^{\rm YM}W_{c,g}=3\kappa_{\rm av}W_{c,g}.
\end{equation}
En lo que sigue se fija un $g\in(0,\infty)$ y se suprime ese subíndice; las
medidas, expectativas, inclusiones y generadores sin superíndice son los objetos
de \eqref{eq:pdf2-ym-g-dynamics}. Las holonomías geométricas permanecen como
funciones coordenadas canónicas sobre el espacio de configuraciones y, por ello, su
ley sí cambia con $g$; $\mathcal T_{m,g}$ las lleva a funciones compuestas de la
coordenatización de referencia. Sólo el testigo espectral se transporta explícitamente en
\eqref{eq:pdf2-ym-g-gap-before-limit}. Esta convención no borra $g$: su residencia explícita es
\eqref{eq:pdf2-ym-gibbs-measure}--\eqref{eq:pdf2-ym-g-dynamics}.

\subsection{Vacío simple y testigo que alcanza la brecha}

La intersección de todas las expectativas del estado completo es la constante:
\begin{equation}\label{eq:pdf2-ym-owner-vacuum}
 \bigcap_{\iota\in\mathcal I_m}\operatorname{Ran}E_{m,\iota}
 =\mathbb C\Omega_m.
\end{equation}
Esta identidad no se introduce como hipótesis. En la coordenatización producto
\eqref{eq:pdf2-ym-owner-product-factorization}, cada $E_{m,\iota}$ integra una
coordenada del índice exhaustivo
\eqref{eq:pdf2-ym-owner-exhaustive-index} y deja fijas las restantes. Si
$u$ pertenece a la intersección, sus expectativas condicionales respecto de todo
cilindro finito son constantes. La convergencia martingala de esos cilindros a la
$\sigma$-álgebra completa da que $u$ es constante casi en todas partes. La inclusión
inversa es inmediata porque toda expectativa conserva la unidad. Esto prueba
\eqref{eq:pdf2-ym-owner-vacuum} y, después de la proyección gauge, deja un vacío
físico unidimensional.
La descomposición de Hoeffding de
\eqref{eq:pdf2-ym-owner-base-generator}--
\eqref{eq:pdf2-ym-owner-refined-generator} da
\begin{equation}\label{eq:pdf2-ym-owner-hoeffding}
 \widehat{\mathcal E}_m(u)
 =3\kappa_{\rm av}\sum_{S\Subset\mathcal I_m}|S|\,\|u_S\|^2,
 \qquad
 D_m^{\rm YM}\succeq
 3\kappa_{\rm av}(I-P_{\Omega_m}).
\end{equation}
La función
\begin{equation}\label{eq:pdf2-ym-owner-wc}
 W_c(q):=\mathbf1_{\{q_1+\cdots+q_6=0\}}-\frac13
\end{equation}
es invariante bajo $G^{V(\Lambda_m)}$ porque no depende de enlaces ni marcos. Para
la acción de incidencia, cada $x_i$ aparece tres veces, luego
\begin{equation}\label{eq:pdf2-ym-owner-wc-invariance}
 \sum_{i<j}(B^+x)_{ij}=3\sum_i x_i=0
 \quad\text{en }\mathbb F_3.
\end{equation}
Así $\mathbb P_m^{\rm phys}W_c=W_c$. Bajo $u_6$, la suma de los seis trits es
uniforme; por cálculo directo,
\begin{equation}\label{eq:pdf2-ym-owner-wc-moments}
 \mathbb EW_c=0,\qquad
 \|W_c\|^2=\frac29,\qquad
 \mathbb E(W_c^3)=\frac2{27}.
\end{equation}
Sólo el proyector $I-E_{m,q_c}$ actúa no trivialmente sobre $W_c$. Por tanto
\begin{equation}\label{eq:pdf2-ym-owner-wc-eigenvalue}
 D_m^{\rm YM}W_c=3\kappa_{\rm av}W_c.
\end{equation}
La cota inferior de \eqref{eq:pdf2-ym-owner-hoeffding} no es sólo una cota: el
complemento físico contiene un vector no nulo que alcanza el umbral.

\subsection{Refinamiento nonádico y persistencia en el límite}

Las inclusiones físicas satisfacen
\begin{equation}\label{eq:pdf2-ym-owner-intertwining}
 D_{m+1}^{\rm YM}\widehat J_m
 =\widehat J_mD_m^{\rm YM},\qquad
 \bigl(e^{-a_{m+1}D_{m+1}^{\rm YM}}\bigr)^9\widehat J_m
 =\widehat J_me^{-a_mD_m^{\rm YM}},\qquad a_{m+1}=a_m/9.
\end{equation}
En el límite inductivo
\begin{equation}\label{eq:pdf2-ym-owner-inductive-limit}
 \widehat{\mathscr H}_\infty^{\rm phys}
 =\varinjlim(\widehat{\mathscr H}_m^{\rm phys},\widehat J_m)
\end{equation}
se define sobre la unión algebraica
\begin{equation}\label{eq:pdf2-ym-owner-limit-form}
 \mathcal E_\infty[\widehat J_{m,\infty}u]
 :=\langle u,D_m^{\rm YM}u\rangle.
\end{equation}
Su generador se denota
$D_{\infty,g}^{\rm YM}$; de acuerdo con la convención posterior a
\eqref{eq:pdf2-ym-g-gap-before-limit}, se escribe también
$D_\infty^{\rm YM}:=D_{\infty,g}^{\rm YM}$.
El primer entrelazamiento de \eqref{eq:pdf2-ym-owner-intertwining} hace que la
definición sea independiente del representante. Las truncaciones de Hoeffding son
un núcleo de forma y dan una sucesión de recuperación; Fatou en la suma no negativa
da la semicontinuidad inferior. La clausura posee, pues,
\begin{equation}\label{eq:pdf2-ym-owner-limit-gap}
 \overline{\mathcal E}_\infty[u]
 \ge3\kappa_{\rm av}\|(I-P_{\Omega_\infty})u\|^2.
\end{equation}
La inclusión conserva las coordenadas ancestrales, de modo que
$\widehat J_{m,\infty}W_c$ sigue siendo no nulo y satisface
\begin{equation}\label{eq:pdf2-ym-owner-limit-witness}
 D_\infty^{\rm YM}\widehat J_{m,\infty}W_c
 =3\kappa_{\rm av}\widehat J_{m,\infty}W_c.
\end{equation}
En consecuencia,
\begin{equation}\label{eq:pdf2-ym-owner-exact-gap}
 \inf\bigl(\operatorname{Spec}D_\infty^{\rm YM}\setminus\{0\}\bigr)
 =3\kappa_{\rm av}>0,
 \qquad \ker D_\infty^{\rm YM}=\mathbb C\Omega_\infty.
\end{equation}

\subsection{Cuatro reconstrucciones OS de la misma dinámica}

El kernel completo de \eqref{eq:pdf2-ym-owner-full-kernel} define una ley reversible
sobre $\widehat X_m$. Para la reconstrucción física se usa su descenso bien tipado
\eqref{eq:pdf2-ym-owner-physical-kernel}: sobre el cociente $Y_m$ se define
\begin{equation}\label{eq:pdf2-ym-owner-eight-face-law}
 \Omega_{m,{\rm phys}}^{(8)}(d\bar y_1\cdots d\bar y_8)
 =\int\nu_m(d\bar z)
  \prod_{r=1}^8K_{m,a_m/2}^{\rm phys}(\bar z,d\bar y_r).
\end{equation}
Para cada dirección $\nu=1,\ldots,4$, la reflexión intercambia las cuatro caras
positivas con las negativas. Condicionando respecto del estado central $z$ y usando
reversibilidad,
\begin{equation}\label{eq:pdf2-ym-owner-four-os}
 \int\overline{F(\Theta_{\nu,m}y)}F(y)\,
 d\Omega_{m,{\rm phys}}^{(8)}(y)
 =\|\mathcal B_{m,\nu}F\|_2^2\ge0.
\end{equation}
La marginal de la pareja opuesta normal a $\nu$ es
\begin{equation}\label{eq:pdf2-ym-owner-os-marginal}
 \nu_m(d\bar x)K_{m,a_m}^{\rm phys}(\bar x,d\bar y).
\end{equation}
Para identificar el Hamiltoniano de una dirección se toma el subespacio OS de
observables de una cara positiva simple,
$F(y)=F_\nu(\bar y_{\nu,+})$; los observables de varias caras conservan la
positividad anterior, pero no se confunden con este espacio de transferencia. En el
cociente de cara simple, el mapa
\begin{equation}\label{eq:pdf2-ym-owner-os-unitary}
 \mathcal V_{m,\nu}[F]
 :=e^{-a_mD_m^{\rm YM}/2}\mathcal W_mF
\end{equation}
es isométrico para la norma OS y tiene rango denso. Su extensión identifica el
cociente OS con $\widehat{\mathscr H}_m^{\rm phys}$. Póngase
$\overline D_m:=\mathcal W_m^{-1}D_m^{\rm YM}\mathcal W_m$ y
$\overline J_m:=\mathcal W_{m+1}^{-1}\widehat J_m\mathcal W_m$. El
entrelazamiento y $a_m=9a_{m+1}$ dan
\begin{equation}\label{eq:pdf2-ym-owner-os-range-inclusion}
 \overline J_me^{-a_m\overline D_m/2}
 =e^{-a_m\overline D_{m+1}/2}\overline J_m
 =\bigl(e^{-a_{m+1}\overline D_{m+1}/2}\bigr)^9\overline J_m.
\end{equation}
El último miembro pertenece al rango de
$e^{-a_{m+1}\overline D_{m+1}/2}$; por eso el inverso de $\mathcal V_{m+1,\nu}$
que aparece a continuación se aplica sólo en su rango y está bien definido. Los
conectores
\begin{equation}\label{eq:pdf2-ym-owner-os-connectors}
 \mathcal J_{m,\nu}^{\rm OS}
 :=\mathcal V_{m+1,\nu}^{-1}\widehat J_m\mathcal V_{m,\nu}
\end{equation}
son isometrías y entrelazan los cuatro Hamiltonianos reconstruidos:
\begin{equation}\label{eq:pdf2-ym-owner-common-os-hamiltonian}
 H_{{\rm OS},m+1}^{(\nu)}\mathcal J_{m,\nu}^{\rm OS}
 =\mathcal J_{m,\nu}^{\rm OS}H_{{\rm OS},m}^{(\nu)},
 \qquad
 \mathcal V_{\infty,\nu}H_{{\rm OS},\infty}^{(\nu)}
 \mathcal V_{\infty,\nu}^{-1}=D_\infty^{\rm YM}.
\end{equation}
Las cuatro reflexiones no fabrican cuatro medidas: son cuatro representaciones del mismo
estado, del mismo kernel y del mismo generador.

\subsection{Localidad, acción euclídea y representación de curvatura}

La categoría dirigida de pantallas incluye refinamientos, overlays y transportes
euclídeos. Si $a\in E(4)$ y $\alpha$ es una pantalla, el relabelado completo define
\begin{equation}\label{eq:pdf2-ym-owner-e4-finite-transport}
 U_\alpha(a):\widehat{\mathscr H}_\alpha^{\rm phys}
 \longrightarrow\widehat{\mathscr H}_{a\alpha}^{\rm phys},
 \qquad
 U_\beta(a)\widehat J_{\alpha\beta}
 =\widehat J_{a\alpha,a\beta}U_\alpha(a).
\end{equation}
El transporte preserva medida, generador, ideales nulos y soportes. La identidad de
naturalidad permite definir, sin elegir representante, una unitaria sobre el colímite
algebraico:
\begin{equation}\label{eq:pdf2-ym-owner-e4-colimit-action}
 U(a)\widehat J_{\alpha,\infty}F
 :=\widehat J_{a\alpha,\infty}U_\alpha(a)F.
\end{equation}
Los relabelados satisfacen $U_{a\alpha}(b)U_\alpha(a)=U_\alpha(ba)$ y preservan
norma; por ello \eqref{eq:pdf2-ym-owner-e4-colimit-action} se prolonga a una
representación unitaria de $E(4)$.

\subsubsection{Curvatura local del canal de incidencia}

Todo grupo compacto, conectado y simple posee un toro maximal no trivial. Se fija
$t_3\in G$ de orden exactamente tres. Para un collar $c$, sea
\begin{equation}\label{eq:pdf2-ym-owner-incidence-flux}
 \omega_c(q):=\sum_{i<j}q_{ij}\in\mathbb F_3.
\end{equation}
La identidad \eqref{eq:pdf2-ym-incidence-conservation} prueba que $\omega_c$ es
invariante bajo la acción finita. Si $h_c$ es el marco en el vértice base del collar,
se define la holonomía geométrica y, para el testigo espectral, su imagen transportada
en la ley a acoplamiento $g$:
\begin{equation}\label{eq:pdf2-ym-owner-incidence-holonomy}
 \operatorname{Hol}_{m,c}^{\rm inc,geom}:=h_ct_3^{\,\omega_c(q)}h_c^{-1}\in G,
 \qquad
 \operatorname{Hol}_{m,c,g}^{\rm inc,spec}
 :=\mathcal T_{m,g}^{-1}\operatorname{Hol}_{m,c}^{\rm inc,geom},
 \qquad
 \operatorname{Hol}_{m,c}^{\rm inc}:=\operatorname{Hol}_{m,c,g}^{\rm inc,spec}.
\end{equation}
Bajo $G^{V_m}$ se transforma por conjugación y bajo
$\prod_c\mathbb F_3^4$ permanece fija. Como $t_3$ tiene orden tres, el cálculo
funcional de esta variable de tres valores da la identidad literal
\begin{equation}\label{eq:pdf2-ym-owner-wc-curvature-functional}
 W_{c,g}:=\mathcal T_{m,g}^{-1}W_c^0
 =\mathbf1_{\{e\}}\!\left(\operatorname{Hol}_{m,c}^{\rm inc}\right)-\frac13.
\end{equation}
Se escribe $W_c:=W_{c,g}$ en el resto del capítulo.
Como $T_{m,b,g}$ actúa dentro del cociente del mismo collar, la holonomía espectral
es una función de su álgebra local de holonomías geométricas; no añade una variable
ni cambia el soporte. Las holonomías geométricas, usadas en $S_{m,g}$ y en el clover,
mantienen su distribución dependiente de $g$.
El proyector del miembro derecho es una función de clase local; por tanto $W_c$ no
es una coordenada auxiliar desacoplada, sino un observable de la holonomía de
incidencia de la misma restricción gauge.

Para un abierto acotado $O\subset\mathbb R^4$, sea
$\mathfrak A_{\rm YM}(O)$ el álgebra generada por funciones de clase acotadas de
\eqref{eq:pdf2-ym-owner-incidence-holonomy}, holonomías de plaqueta y lectores
clover con soporte en $O$. Se define el sector de curvatura física por
\begin{equation}\label{eq:pdf2-ym-owner-curvature-cyclic-sector}
 \mathscr H_{\rm YM}^{\rm curv}
 :=\overline{\bigcup_{O\Subset\mathbb R^4}
   \mathfrak A_{\rm YM}(O)\Omega_\infty}
 \subseteq\widehat{\mathscr H}_\infty^{\rm phys}.
\end{equation}
Por \eqref{eq:pdf2-ym-owner-wc-curvature-functional}, cada
$\widehat J_{m,\infty}W_c$ pertenece materialmente a este sector.
Las expectativas de coordenada envían cilindros de curvatura a cilindros de
curvatura; por tanto el semigrupo deja invariante
$\mathscr H_{\rm YM}^{\rm curv}$. Al ser autoadjunto, el sector es reductor. La
cota de \eqref{eq:pdf2-ym-owner-limit-gap} se restringe a él y el vector anterior
alcanza la igualdad:
\begin{equation}\label{eq:pdf2-ym-owner-curvature-sector-gap}
 D_{\rm curv,g}^{\rm YM}
 :=\left.D_{\infty,g}^{\rm YM}\right|_{\mathscr H_{\rm YM}^{\rm curv}},
 \qquad D_{\rm curv}^{\rm YM}:=D_{\rm curv,g}^{\rm YM},
 \qquad
 \inf\bigl(\operatorname{Spec}D_{\rm curv,g}^{\rm YM}\setminus\{0\}\bigr)
 =3\kappa_{\rm av}.
\end{equation}

\subsubsection{Sistema cofinal de representación y productos cruzados}

La inclusión $\widehat J_{\alpha\beta}$ conserva una coordenada ancestral; no debe
confundirse con el promedio de bloque que expresa un campo continuo. Para escribir
este segundo mapa sin alterar la genealogía, sea $\mathfrak P$ el conjunto dirigido
de pantallas celulares acotadas, cerrado bajo overlay, refinamiento nonádico y
transporte euclídeo. Para $\alpha\in\mathfrak P$ póngase
\begin{equation}\label{eq:pdf2-ym-owner-cell-step-space}
 V_\alpha:=\operatorname{span}\left\{
 e_{\alpha,c}:=|c|^{-1/2}\mathbf1_c:
 c\in\mathcal C_\alpha^{\rm cell}\right\}
 \subset L^2(\mathbb R^4),
 \qquad P_\alpha:L^2(\mathbb R^4)\to V_\alpha.
\end{equation}
Los overlays hacen que $P_\alpha h\to h$ para todo $h\in L^2$ a lo largo de
$\mathfrak P$. En el lado gauge se normaliza el testigo por
\begin{equation}\label{eq:pdf2-ym-owner-normalized-cell-witness}
 Z_{\alpha,c}:=\sqrt{\frac92}\,W_{\alpha,c},
 \qquad
 \langle Z_{\alpha,c},Z_{\alpha,d}\rangle=\delta_{cd},
 \qquad
 D_\alpha^{\rm YM}Z_{\alpha,c}=3\kappa_{\rm av}Z_{\alpha,c}.
\end{equation}
La segunda igualdad usa exactamente los momentos de
\eqref{eq:pdf2-ym-owner-wc-moments} y la factorización de coordenadas distintas.
Sea $\mathcal K_\alpha^{\rm cell}$ el espacio generado por esos vectores. Si
$\beta\succeq\alpha$, el entrelazador de representación, distinto de
$\widehat J_{\alpha\beta}$, se define generador a generador por
\begin{equation}\label{eq:pdf2-ym-owner-block-intertwiner}
 I_{\alpha\beta}^{\rm curv}Z_{\alpha,c}
 :=\sum_{\substack{d\in\mathcal C_\beta^{\rm cell}\\d\subset c}}
       \sqrt{\frac{|d|}{|c|}}\,Z_{\beta,d}.
\end{equation}
Como las subceldas son disjuntas y suman $|c|$, se obtiene, sin tomar límites,
\begin{equation}\label{eq:pdf2-ym-owner-block-intertwiner-identities}
 (I_{\alpha\beta}^{\rm curv})^*I_{\alpha\beta}^{\rm curv}=I,
 \quad
 I_{\beta\gamma}^{\rm curv}I_{\alpha\beta}^{\rm curv}
 =I_{\alpha\gamma}^{\rm curv},
 \quad
 D_\beta^{\rm YM}I_{\alpha\beta}^{\rm curv}
 =I_{\alpha\beta}^{\rm curv}D_\alpha^{\rm YM}.
\end{equation}
El relabelado celular conserva volúmenes y descendientes; para $a\in E(4)$ se tiene además
\begin{equation}\label{eq:pdf2-ym-owner-block-e4-naturality}
 U_\beta(a)I_{\alpha\beta}^{\rm curv}
 =I_{a\alpha,a\beta}^{\rm curv}U_\alpha(a).
\end{equation}

Para $h\in C_c^\infty(\mathbb R^4)$ defínase ahora el lector de bloque
\begin{equation}\label{eq:pdf2-ym-owner-smeared-witness}
 \Phi_\alpha(h):=
 \sum_{c\in\mathcal C_\alpha^{\rm cell}}
 \langle h,e_{\alpha,c}\rangle_{L^2}Z_{\alpha,c}.
\end{equation}
En una malla uniforme esta fórmula es
$\sqrt{9/2}\,a_\alpha^2\sum_c(P_\alpha h)(x_c)W_{\alpha,c}$; el promedio celular,
no una evaluación puntual, fija el producto cruzado. En efecto, si $\gamma$ refina
dos pantallas $\alpha$ y $\beta$, las definiciones anteriores dan la identidad
calculada
\begin{equation}\label{eq:pdf2-ym-owner-smeared-cross-gram}
 \left\langle
  I_{\alpha\gamma}^{\rm curv}\Phi_\alpha(h),
  I_{\beta\gamma}^{\rm curv}\Phi_\beta(k)
 \right\rangle
 =\langle P_\alpha h,P_\beta k\rangle_{L^2}.
\end{equation}
Por polarización, con $k=h$,
\begin{equation}\label{eq:pdf2-ym-owner-smeared-cauchy-derived}
 \left\|
  I_{\alpha\gamma}^{\rm curv}\Phi_\alpha(h)
  -I_{\beta\gamma}^{\rm curv}\Phi_\beta(h)
 \right\|_2^2
 =\|P_\alpha h-P_\beta h\|_{L^2}^2\longrightarrow0.
\end{equation}
Así la propiedad de Cauchy es consecuencia de la convergencia fuerte de
proyecciones, no una identidad adicional postulada. El colímite de
\eqref{eq:pdf2-ym-owner-block-intertwiner} contiene, por tanto,
\begin{equation}\label{eq:pdf2-ym-owner-smeared-gram}
 \Phi(h):=\lim_\alpha I_{\alpha,\infty}^{\rm curv}\Phi_\alpha(h),
 \qquad
 \langle\Phi(h),\Phi(k)\rangle=\langle h,k\rangle_{L^2},
 \qquad
 D_{\rm pub}^{\rm YM}\Phi(h)=3\kappa_{\rm av}\Phi(h).
\end{equation}

Esta representación no añade grados de libertad. En cada paso, la descomposición
$V_\beta=V_\alpha\oplus(V_\beta\ominus V_\alpha)$ tiene una base de Haar
nonádica. El orden exhaustivo de incidencias nuevas de
\eqref{eq:pdf2-ym-owner-exhaustive-index} empareja esa base, etiqueta por etiqueta,
con coordenadas $Z_{\beta,d}$ del factor de innovación, conservando el soporte de la
función de Haar. Las unitarias triangulares $\widehat\Gamma_\alpha$ pegan esos
emparejamientos y producen una isometría basada
\begin{equation}\label{eq:pdf2-ym-owner-publication-into-physical}
 \mathcal U_{\rm curv}:
 \varinjlim(\mathcal K_\alpha^{\rm cell},I_{\alpha\beta}^{\rm curv})
 \longrightarrow\mathscr H_{\rm YM}^{\rm curv},
 \quad
 D_{\rm curv}^{\rm YM}\mathcal U_{\rm curv}
 =\mathcal U_{\rm curv}D_{\rm pub}^{\rm YM},
 \quad U(a)\mathcal U_{\rm curv}=\mathcal U_{\rm curv}U_{L^2}(a).
\end{equation}
En adelante $\Phi(h)$ denota también su imagen por $\mathcal U_{\rm curv}$.
La compatibilidad de \eqref{eq:pdf2-ym-gibbs-transport-global} define
$\mathcal T_{\infty,g}$ en el límite. La representación física al acoplamiento fijado
no es una segunda isometría escogida sólo sobre el primer caos, sino la restricción
\begin{equation}\label{eq:pdf2-ym-full-publication-restriction}
 \mathcal U_{{\rm curv},g}
 :=\mathcal T_{\infty,g}^{-1}\mathcal U_{{\rm curv},0},
 \qquad
 \widetilde{\mathcal U}_g
 :=\mathcal T_{\infty,g}^{-1}:
 L^\infty(\widehat X_\infty,\widehat\mu_\infty^0)
 \longrightarrow
 L^\infty(\widehat X_\infty,\widehat\mu_{\infty,g}^{\rm YM}).
\end{equation}
Sobre la $*$-álgebra polinómica generada por holonomías, plaquetas, clover y
$\Phi(h)$ se verifican, generador a generador y después por linealidad,
\begin{equation}\label{eq:pdf2-ym-full-publication-products}
 \widetilde{\mathcal U}_g(AB)
 =\widetilde{\mathcal U}_g(A)\widetilde{\mathcal U}_g(B),
 \quad
 \widetilde{\mathcal U}_g(A^*)=\widetilde{\mathcal U}_g(A)^*,
 \quad
 \omega_g(\widetilde{\mathcal U}_gA)=\omega_0(A).
\end{equation}
Las igualdades de superficie, clover y Bianchi son preservadas por
\eqref{eq:pdf2-ym-transport-surface-bianchi}; por densidad monótona, la extensión
vale en toda el álgebra local acotada. Así
\eqref{eq:pdf2-ym-owner-smeared-cross-gram} es la restricción de un isomorfismo
probabilístico y algebraico completo, no un sustituto de él.

\subsubsection{Núcleo denso, continuidad fuerte y red local}

Sea $M_{\Phi(h)}$ el operador de multiplicación afiliado al álgebra local y
\begin{equation}\label{eq:pdf2-ym-owner-curvature-weyl}
 \mathsf W(h,t):=\exp\!\bigl(itM_{\Phi(h)}\bigr),
 \qquad h\in C_c^\infty(\mathbb R^4),\quad t\in\mathbb R.
\end{equation}
Se completa $\mathfrak A_{\rm YM}(O)$ con estos unitarios para
$\operatorname{supp}h\subset O$ y con las funciones de clase acotadas de las
holonomías de plaqueta. Denótese por $\mathfrak A_{\rm YM}(O)^{\rm sm}$ la
$*$-álgebra así generada con funciones de test suaves. El núcleo
\begin{equation}\label{eq:pdf2-ym-owner-smooth-core}
 \mathscr D_{\rm sm}:=\operatorname{span}\left\{
 A_1\cdots A_n\Omega_\infty:
 A_j\in\mathfrak A_{\rm YM}(O_j)^{\rm sm},\quad
 O_j\Subset\mathbb R^4\right\}
\end{equation}
es denso por la definición cíclica de
\eqref{eq:pdf2-ym-owner-curvature-cyclic-sector}. Las ecuaciones
\eqref{eq:pdf2-ym-owner-block-e4-naturality} y
\eqref{eq:pdf2-ym-owner-publication-into-physical} dan, para $a\in E(4)$,
\begin{equation}\label{eq:pdf2-ym-owner-field-e4-covariance}
 U(a)\Phi(h)=\Phi(a\!\cdot h),
 \qquad
 \|(U(a)-I)\Phi(h)\|=\|a\!\cdot h-h\|_{L^2}.
\end{equation}
Además, $|e^{ix}-e^{iy}|\le|x-y|$ y
\eqref{eq:pdf2-ym-owner-smeared-gram} implican
\begin{equation}\label{eq:pdf2-ym-owner-weyl-continuity}
 \|\bigl(\mathsf W(a\!\cdot h,t)-\mathsf W(h,t)\bigr)\Omega_\infty\|
 \le |t|\,\|a\!\cdot h-h\|_{L^2}.
\end{equation}
La continuidad se comprueba también para los generadores de holonomía del core;
\eqref{eq:pdf2-ym-owner-weyl-continuity} por sí sola no los incluye. Para una ruta poligonal orientada $\gamma$,
una plaqueta $p$ y $f\in C^1(G)^{\rm class}$ se ponen
\begin{equation}\label{eq:pdf2-ym-owner-holonomy-core-generators}
 A_{\gamma,f}:=f(\operatorname{Hol}_\gamma),
 \qquad A_{p,f}:=f(\operatorname{Hol}_p).
\end{equation}
En un overlay de $\gamma$ y $a\gamma$, con $a\in E(4)$, la ley de puente
\eqref{eq:pdf2-ym-owner-heat-bridge} cancela los incrementos comunes. Para cada
incremento restante $z$ se usa
$|f(xz)-f(x)|\le\|\nabla f\|_\infty d_G(z,e)$ y la cota de segundo momento del
kernel de calor
\(
 \int_Gd_G(z,e)^2k_\tau(z)\,dz\le C_G\tau
\).
Para la ley reponderada, $0<e^{-p/(4g^2)}\le1$ y la continuidad de $p$ en el
factor compacto dan
\begin{equation}\label{eq:pdf2-ym-owner-gibbs-heat-moment}
 \frac{\int_Gd_G(z,e)^2e^{-p(z)/(4g^2)}k_\tau(z)\,dz}
      {\int_Ge^{-p(z)/(4g^2)}k_\tau(z)\,dz}
 \le C_{G,g}\tau,
\end{equation}
uniformemente en las pantallas de un compacto y para $0<\tau\le\tau_0$.
Sumando los tiempos de los incrementos de la diferencia simétrica se obtiene la
estimación calculada
\begin{equation}\label{eq:pdf2-ym-owner-holonomy-e4-estimate}
 \begin{aligned}
 \|U(a)A_{\gamma,f}\Omega_\infty-A_{\gamma,f}\Omega_\infty\|_2^2
 &\le C_{G,g}\|\nabla f\|_\infty^2
     \tau(\gamma\triangle a\gamma),\\
 \|U(a)A_{p,f}\Omega_\infty-A_{p,f}\Omega_\infty\|_2^2
 &\le C_{G,g}\|\nabla f\|_\infty^2
     \tau(\partial p\triangle\partial(ap)).
 \end{aligned}
\end{equation}
Aquí la diferencia simétrica no se mide como dos curvas disjuntas: los conectores
de superficie la llenan por la cinta orientada mínima
$\Sigma(\gamma,a\gamma)$ y
\begin{equation}\label{eq:pdf2-ym-owner-bridge-time-ribbon}
 \tau(\gamma\triangle a\gamma)
 :=\sum_{p\subset\Sigma(\gamma,a\gamma)}\tau_p,
 \qquad
 \tau(\gamma\triangle a\gamma)
 \le C_\gamma d_{E(4)}(a,e),
\end{equation}
con la misma definición para las fronteras de plaqueta. La desigualdad sale de
sumar los tiempos aditivos de \eqref{eq:pdf2-ym-owner-heat-bridge} sobre una cinta
de anchura $d_{E(4)}(a,e)$ y longitud acotada por la de $\gamma$.
La medida de tiempo de puente del miembro derecho tiende así a cero cuando $a\to e$;
esto es exactamente la continuidad de los conectores de ruta de la imagen gauge.
Cada lector clover acotado $\chi(F_{\rm cl})$, $\chi\in C_b^1$, es función de una
suma finita de plaquetas transportadas; la misma cota, multiplicada por
$\|\chi'\|_\infty$ y por el número de términos, prueba también su continuidad.

Para funciones de clase meramente acotadas se hace visible la regularización que
entra en el core. Si $d b$ es Haar izquierdo en $E(4)$ y
$\eta\in C_c^\infty(E(4))$, se define el integral de Bochner
\begin{equation}\label{eq:pdf2-ym-owner-orbit-smearing}
 A[\eta]:=\int_{E(4)}\eta(b)U(b)AU(b)^*\,db,
 \qquad A\in\{A_{\gamma,f},A_{p,f},\chi(F_{{\rm cl}})\},
 \quad \chi\in C_b^1.
\end{equation}
Con $(L_a\eta)(b):=\eta(a^{-1}b)$, el cambio de variable $b\mapsto ab$ da
\begin{equation}\label{eq:pdf2-ym-owner-orbit-smearing-continuity}
 U(a)A[\eta]U(a)^*=A[L_a\eta],
 \qquad
 \|(A[L_a\eta]-A[\eta])\Omega_\infty\|
 \le\|A\|\,\|L_a\eta-\eta\|_{L^1(E(4))}\longrightarrow0.
\end{equation}
Se entiende desde ahora que los generadores de holonomía, plaqueta y clover de
$\mathfrak A_{\rm YM}(O)^{\rm sm}$ son los de
\eqref{eq:pdf2-ym-owner-holonomy-core-generators} con $f\in C^1$, o sus
regularizaciones \eqref{eq:pdf2-ym-owner-orbit-smearing}, y que el soporte barrido
permanece compacto dentro de $O$. Las aproximaciones de identidad en $E(4)$,
junto con \eqref{eq:pdf2-ym-owner-holonomy-e4-estimate}, recuperan los generadores
no regularizados en $L^2$; por ello no se reduce el sector cíclico.

Un telescopio de productos que usa
\eqref{eq:pdf2-ym-owner-weyl-continuity} y
\eqref{eq:pdf2-ym-owner-orbit-smearing-continuity} prueba la continuidad de todos
los generadores de $\mathscr D_{\rm sm}$; por
densidad y unitariedad,
\begin{equation}\label{eq:pdf2-ym-owner-strong-e4}
 \lim_{a\to e}\|(U(a)-I)\Psi\|=0
 \qquad(\Psi\in\mathscr H_{\rm YM}^{\rm curv}).
\end{equation}
Los soportes disjuntos usan factores disjuntos y conmutan. Queda así la red
\begin{equation}\label{eq:pdf2-ym-owner-local-net}
\begin{aligned}
 O_1\subseteq O_2&\Rightarrow
 \mathfrak A_{\rm YM}(O_1)\subseteq\mathfrak A_{\rm YM}(O_2),\\
 [\mathfrak A_{\rm YM}(O_1),\mathfrak A_{\rm YM}(O_2)]&=0
 \quad (O_1\cap O_2=\varnothing),\\
 U(a)\mathfrak A_{\rm YM}(O)U(a)^*&=\mathfrak A_{\rm YM}(aO).
\end{aligned}
\end{equation}
La identidad de Gram para $h\ne0$ y
\eqref{eq:pdf2-ym-owner-wc-curvature-functional} prueban que la red no es escalar.

\subsubsection{Axiomas de Osterwalder--Schrader y reconstrucción continua}

Para una dirección $\nu$, sea $\mathfrak A_{+,\nu}$ la clausura de la unión de
las álgebras anteriores con soporte compacto en $\{x_\nu>0\}$. El estado
estático, la ley de pilas y el operador de transferencia no se definen por
separado. En una pantalla $\alpha$, sea
$\widehat\mu_{\alpha,g}^{\rm YM}$ la marginal de
\eqref{eq:pdf2-ym-gibbs-measure}, sea
$\nu_{\alpha,g}:=(\pi_\alpha^{\rm phys})_*
\widehat\mu_{\alpha,g}^{\rm YM}$ su ley en el cociente y sea
$K_{\alpha,g}(t)=e^{-tD_{\alpha,g}^{\rm YM}}$ su kernel.
Para tiempos ordenados $t_1\le\cdots\le t_n$, la marginal euclídea de la
misma medida es
\begin{equation}\label{eq:pdf2-ym-owner-euclidean-time-marginal}
 \mathbb P_{\alpha,g}^{E,(\nu;t_1,\ldots,t_n)}
 :=(\operatorname{ev}_{t_1}^{(\nu)},\ldots,
    \operatorname{ev}_{t_n}^{(\nu)})_*
    \widehat\mu_{\alpha,g}^{\rm YM}.
\end{equation}
La identidad de desintegración que forma parte de
\eqref{eq:pdf2-g-gau-transfer-euclidean-dato} la calcula, no la reemplaza por
una cadena auxiliar:
\begin{equation}\label{eq:pdf2-ym-owner-euclidean-markov-identification}
 d\mathbb P_{\alpha,g}^{E,(\nu;t_1,\ldots,t_n)}
 =\nu_{\alpha,g}(dy_1)
   \prod_{j=1}^{n-1}K_{\alpha,g}
      (t_{j+1}-t_j;y_j,dy_{j+1}).
\end{equation}
Si $M_A$ denota multiplicación por un lector acotado $A$, sus correlaciones
son por tanto
\begin{equation}\label{eq:pdf2-ym-owner-schwinger-transfer}
 \begin{aligned}
 S_{\alpha,n}^{(\nu,g)}(A_1,t_1;\ldots;A_n,t_n)
 &:=\int\prod_{j=1}^n(\tau_{t_je_\nu}A_j)(x)\,
       d\widehat\mu_{\alpha,g}^{\rm YM}(x)\\
 &=\langle\mathbf1,
 M_{A_1}K_{\alpha,g}(t_2-t_1)M_{A_2}\cdots
 K_{\alpha,g}(t_n-t_{n-1})M_{A_n}\mathbf1
 \rangle_{L^2(\nu_{\alpha,g})}.
 \end{aligned}
\end{equation}
La ley de la pila estacionaria en esos mismos tiempos es esa marginal
euclídea, escrita mediante su desintegración:
\begin{equation}\label{eq:pdf2-ym-owner-markov-stack-law}
 d\mathbb P_{\alpha,g}^{(\nu;t_1,\ldots,t_n)}
 :=d\mathbb P_{\alpha,g}^{E,(\nu;t_1,\ldots,t_n)}
 =\nu_{\alpha,g}(dy_1)
   \prod_{j=1}^{n-1}K_{\alpha,g}(t_{j+1}-t_j;y_j,dy_{j+1}).
\end{equation}
Integrar sucesivamente $y_n,y_{n-1},\ldots,y_1$ da la identidad exacta
\begin{equation}\label{eq:pdf2-ym-owner-stack-transfer-identity}
 S_{\alpha,n}^{(\nu,g)}(A_1,t_1;\ldots;A_n,t_n)
 =\int\prod_{j=1}^nA_j(y_j)\,
   d\mathbb P_{\alpha,g}^{(\nu;t_1,\ldots,t_n)}(y).
\end{equation}
Si todos los tiempos coinciden, los kernels son la identidad y
\begin{equation}\label{eq:pdf2-ym-owner-static-is-stack}
 S_{\alpha,n}^{(\nu,g)}(A_1,t;\ldots;A_n,t)
 =\int A_1\cdots A_n\,d\nu_{\alpha,g}.
\end{equation}
Así el estado Schwinger estático es la marginal de tiempo igual de la misma
medida euclídea, y su desintegración es la ley Markov; no son dos funcionales
que luego se identifican por nombre. La
compatibilidad de $\widehat J_{\alpha,g}$ con medida y kernel hace compatible
\eqref{eq:pdf2-ym-owner-stack-transfer-identity} con todo overlay. Su límite define
\begin{equation}\label{eq:pdf2-ym-owner-schwinger-functionals}
 S_n^{\rm YM}(A_1,t_1;\ldots;A_n,t_n)
 :=\lim_{\alpha}S_{\alpha,n}^{(\nu,g)}
 (A_{1,\alpha},t_1;\ldots;A_{n,\alpha},t_n).
\end{equation}

Las inserciones de $\Phi$ se obtienen derivando en $t=0$ los unitarios
\eqref{eq:pdf2-ym-owner-curvature-weyl}. Puesto que los $Z_{\alpha,c}$ transportados
son centrados, independientes y uniformemente acotados, el lema de Hoeffding da,
nivel a nivel,
\begin{equation}\label{eq:pdf2-ym-owner-subgaussian-bound}
 \left\langle\Omega_\alpha,
 e^{t\Phi_\alpha(h)}\Omega_\alpha\right\rangle
 \le \exp\!\left(\frac9{16}t^2\|P_\alpha h\|_{L^2}^2\right).
\end{equation}
La cota pasa al límite por
\eqref{eq:pdf2-ym-owner-smeared-cauchy-derived}; todas las derivadas son continuas en
la topología de Schwartz y definen distribuciones temperadas.

La forma de ocho caras se prolonga primero a toda pila finita. Si $F$ depende de
las $N$ capas positivas de una pantalla $\alpha$, la reflexión de
\eqref{eq:pdf2-ym-owner-markov-stack-law} y la condición en la capa central dan
\begin{equation}\label{eq:pdf2-ym-owner-finite-stack-os}
 \begin{aligned}
 &\int \overline{F(\Theta_\nu y_{-N},\ldots,
                         \Theta_\nu y_{-1})}
          F(y_1,\ldots,y_N)\,d\Omega_{\alpha,N}^{(\nu,g)}(y)\\
 &\qquad=
 \int\nu_{\alpha,g}(dz)\left|
  \int K_{\alpha,g}(a_\alpha/2;z,dy_1)
       \prod_{j=1}^{N-1}K_{\alpha,g}(a_\alpha;y_j,dy_{j+1})
       F(y_1,\ldots,y_N)
 \right|^2\ge0.
 \end{aligned}
\end{equation}
Todo cilindro de semiespacio pertenece a una pila finita; por ello
\eqref{eq:pdf2-ym-owner-finite-stack-os} materializa la forma OS antes de la
clausura. En el cociente OS, si $F$ tiene inserciones en tiempos positivos, sea
$(\tau_sF)(A_1,t_1;\ldots;A_n,t_n):=
F(A_1,t_1+s;\ldots;A_n,t_n+s)$. Se define el operador de traslado por
\begin{equation}\label{eq:pdf2-ym-owner-os-transfer-operator}
 T_{\rm OS}^{(\nu)}(s)[F]:=[\tau_sF],
 \qquad s\ge0.
\end{equation}
Para cilindros $F,G$ de semiespacio, la igualdad euclídea--Markov
\eqref{eq:pdf2-ym-owner-euclidean-markov-identification} y la propiedad de
semigrupo dan generador a generador
\begin{equation}\label{eq:pdf2-ym-owner-os-markov-intertwining}
 \begin{aligned}
 \langle[F],T_{\rm OS}^{(\nu)}(s)[G]\rangle_{\rm OS}
 &=S_g^{\rm YM}((\Theta_\nu F)^*\tau_sG)\\
 &=\left\langle
    \mathcal V_{\infty,\nu}[F],
    e^{-sD_{\rm curv,g}^{\rm YM}}
    \mathcal V_{\infty,\nu}[G]
   \right\rangle .
 \end{aligned}
\end{equation}
La identidad muestra a la vez que las clases nulas permanecen nulas y que,
por densidad de los cilindros,
\begin{equation}\label{eq:pdf2-ym-owner-os-generator-identification}
 \mathcal V_{\infty,\nu}T_{\rm OS}^{(\nu)}(s)
 \mathcal V_{\infty,\nu}^{-1}
 =e^{-sD_{\rm curv,g}^{\rm YM}}.
\end{equation}
Ésta es la flecha explícita medida euclídea $\to$ desintegración $\to$
transferencia OS $\to$ Hamiltoniano.

\begin{proposition}[Paquete Osterwalder--Schrader continuo]
\label{prop:pdf2-ym-owner-continuum-os-package}
Los funcionales \eqref{eq:pdf2-ym-owner-schwinger-functionals} satisfacen:
\begin{enumerate}
 \item normalización, hermiticidad y simetría bosónica;
 \item covariancia euclídea y regularidad temperada;
 \item positividad por reflexión en cada una de las cuatro direcciones,
 \begin{equation}\label{eq:pdf2-ym-owner-continuum-reflection-positive}
  \sum_{i,j}\overline{a_i}a_j\,
  S^{\rm YM}\!\left((\Theta_\nu F_i)^*F_j\right)\ge0,
  \qquad F_i\in\mathfrak A_{+,\nu};
 \end{equation}
 \item localidad y cluster exponencial: si $A,B$ son locales centrados y una
 traslación de longitud $r$ separa sus soportes, entonces
 \begin{equation}\label{eq:pdf2-ym-owner-continuum-cluster}
  |S^{\rm YM}(A\,\tau_rB)|
  \le e^{-3\kappa_{\rm av}r}\,
       \|A\Omega_\infty\|\,\|B\Omega_\infty\|.
 \end{equation}
\end{enumerate}
La reconstrucción OS de cualquiera de las cuatro semirredes produce el mismo
espacio cíclico, el mismo vacío y el Hamiltoniano
$D_{\rm curv,g}^{\rm YM}$. La continuación OS produce campos de Wightman locales y una
representación de Poincaré cuyo espectro satisface
\begin{equation}\label{eq:pdf2-ym-owner-wightman-spectrum}
 \operatorname{Spec}(P^0,\mathbf P)\subseteq\overline V_+,
 \qquad P^0=D_{\rm curv,g}^{\rm YM}\ge0.
\end{equation}
\end{proposition}

\begin{proof}
Normalización y hermiticidad proceden del estado; simetría y localidad, de la
conmutatividad euclídea a soportes disjuntos. La covariancia y continuidad son
\eqref{eq:pdf2-ym-owner-field-e4-covariance}--
\eqref{eq:pdf2-ym-owner-strong-e4}; la regularidad es
\eqref{eq:pdf2-ym-owner-subgaussian-bound}. Para
$F_i$ cilíndricos, todos viven en un overlay y una pila finitos;
\eqref{eq:pdf2-ym-owner-finite-stack-os} y la polarización prueban
\eqref{eq:pdf2-ym-owner-continuum-reflection-positive}. La aproximación en norma OS
y Cauchy--Schwarz la extienden a $\mathfrak A_{+,\nu}$.

Tras una rotación euclídea, la separación espacial se convierte en tiempo positivo.
La identidad pila--transferencia
\eqref{eq:pdf2-ym-owner-stack-transfer-identity} y el operador
\eqref{eq:pdf2-ym-owner-os-transfer-operator} escriben la correlación centrada como
\[
 \langle A^*\Omega_\infty,
 e^{-rD_{\rm curv,g}^{\rm YM}}B\Omega_\infty\rangle.
\]
La cota espectral \eqref{eq:pdf2-ym-owner-curvature-sector-gap} da
\eqref{eq:pdf2-ym-owner-continuum-cluster}. Finalmente, los conectores
\eqref{eq:pdf2-ym-owner-os-connectors} forman un sistema dirigido; su unión es densa
por \eqref{eq:pdf2-ym-owner-smooth-core}. La identidad
\eqref{eq:pdf2-ym-owner-common-os-hamiltonian} identifica las cuatro clausuras y sus
generadores con $D_{\rm curv,g}^{\rm YM}$. Las hipótesis ya verificadas de regularidad,
covariancia, reflexión y cluster son las entradas de la continuación OS; su teorema de
reconstrucción da localidad de Wightman y
\eqref{eq:pdf2-ym-owner-wightman-spectrum}.
\end{proof}

\subsubsection{Reconstrucción cuatridimensional de la teoría de Yang--Mills}

La identificación no se apoya sólo en evaluar un clover sobre una conexión dada.
Sea $\mathcal R_G^{\rm loc}$ la $*$-álgebra de lectores formada por todas las
funciones de clase de las holonomías de incidencia y plaqueta, sus transportes a un
vértice base y las polarizaciones del lector $\mathcal P^{F^2}$. Sea
\begin{equation}\label{eq:pdf2-ym-owner-curvature-fibre}
 \mathcal E_G:=\Lambda^2(\mathbb R^4)^*\otimes\mathfrak g
\end{equation}
con el producto invariante. En cada celda, la polarización de
$\mathcal P^{F^2}$ da un Gram positivo $G_{\alpha,c}^{F}$ sobre
$\mathcal E_G$. Se divide por su kernel y se aplica su raíz positiva; el lector
normalizado resultante se denota $Y_{\alpha,c}(v)$ y verifica
\begin{equation}\label{eq:pdf2-ym-owner-curvature-cell-gram}
 \langle Y_{\alpha,c}(v),Y_{\alpha,d}(w)\rangle
 =\delta_{cd}\langle v,w\rangle_{\mathcal E_G}.
\end{equation}
Si $d\subset c$, sea $T_{d\leftarrow c}$ el transporte de marco ya presente en la
holonomía superficial. El entrelazador de curvatura completa es
\begin{equation}\label{eq:pdf2-ym-owner-curvature-intertwiner}
 I_{\alpha\beta}^{F}Y_{\alpha,c}(v)
 :=\sum_{d\subset c}\sqrt{\frac{|d|}{|c|}}\,
 Y_{\beta,d}\!\left(\operatorname{Ad}_{T_{d\leftarrow c}}v\right).
\end{equation}
La invariancia del producto de $\mathfrak g$, la disjunción de subceldas y la
composición de transportes prueban
\begin{equation}\label{eq:pdf2-ym-owner-curvature-intertwiner-identities}
 (I_{\alpha\beta}^{F})^*I_{\alpha\beta}^{F}=I,
 \qquad
 I_{\beta\gamma}^{F}I_{\alpha\beta}^{F}=I_{\alpha\gamma}^{F}.
\end{equation}
Para $h\in C_c^\infty(\mathbb R^4;\mathcal E_G)$ se pone
\begin{equation}\label{eq:pdf2-ym-owner-f-alpha}
 \mathbf F_\alpha(h):=
 \sum_{c\in\mathcal C_\alpha^{\rm cell}}
 Y_{\alpha,c}\!\left(
   |c|^{-1/2}\int_c h(x)\,d^4x\right).
\end{equation}
En un overlay común, el mismo cálculo de
\eqref{eq:pdf2-ym-owner-smeared-cross-gram} da ahora
\begin{equation}\label{eq:pdf2-ym-owner-curvature-cross-gram}
 \left\langle I_{\alpha\gamma}^{F}\mathbf F_\alpha(h),
 I_{\beta\gamma}^{F}\mathbf F_\beta(k)\right\rangle
 =\langle P_\alpha h,P_\beta k\rangle_{L^2(\mathcal E_G)}.
\end{equation}
Por tanto $\mathbf F_\alpha(h)$ es Cauchy por convergencia fuerte de las proyecciones,
y no por una hipótesis nueva. El mismo emparejamiento de Haar e innovaciones de
\eqref{eq:pdf2-ym-owner-publication-into-physical}, ahora aplicado a la base
ortonormalizada de $G_{\alpha,c}^{F}$, realiza el colímite mediante una isometría
local y gauge-covariante en
$L^2(\widehat X_\infty,\widehat\mu_{\infty,g}^{\rm YM})$. Su límite define la
distribución aleatoria
gauge-covariante
\begin{equation}\label{eq:pdf2-ym-owner-distributional-curvature}
\begin{aligned}
 \mathbf F_{\mu\nu}&\in
 \mathcal S'(\mathbb R^4;\mathfrak g)\ \widehat\otimes\
 L^2(\widehat X_\infty,\widehat\mu_{\infty,g}^{\rm YM}),\\
 U(a)\mathbf F(h)U(a)^*&=\mathbf F(a\!\cdot h),\\
 u\mathbf F(h)u^{-1}&=\mathbf F(\operatorname{Ad}_u h),
 \qquad a\in E(4),\ u\in G.
\end{aligned}
\end{equation}
Su canal de clase espectral de orden tres es precisamente $\Phi$. La relación
superficial finita
\begin{equation}\label{eq:pdf2-ym-owner-surface-product}
 \operatorname{Hol}^{(\alpha)}_p
 =\prod_{p'\subset p}^{\longrightarrow}
  T_{p'}\operatorname{Hol}^{(\beta)}_{p'}T_{p'}^{-1}
\end{equation}
pasa a la clausura cofinal porque todos sus productos cruzados se calculan en un
overlay. En particular, el lector clover aleatorio, y no sólo su evaluación en una
conexión suave, satisface la identidad de pequeños lazos
\begin{equation}\label{eq:pdf2-ym-owner-holonomy-curvature-relation}
 \mathbf F(h)=L^2\!\mathord{-}\!\lim_{\alpha}
 a_\alpha^4\sum_x
 \left\langle(P_\alpha h)(x),F_{{\rm cl},\alpha}(x)\right\rangle,
 \qquad
 \delta_{\rm gauge}S_n^{\rm YM}=0,
 \qquad
 \operatorname{Alt}_{\lambda\mu\nu}D_\lambda\mathbf F_{\mu\nu}=0.
\end{equation}
La primera igualdad es \eqref{eq:pdf2-ym-owner-curvature-cross-gram} escrita en el
representante clover; la segunda es la identidad de Ward
\eqref{eq:pdf2-ym-ward-finite} pasada al límite, y la
tercera es la cancelación de la frontera de una frontera en el producto superficial.
Ninguna usa el término de masa.

Para evitar definir el blanco por la conclusión que se quiere obtener, sea
$\mathbf{YM}_4(G,g)$ la clase de quíntuplas locales
$(\mathfrak A,\mu_g,K_g,\operatorname{Hol},\mathbf F)$ que verifican las siguientes
ecuaciones comprobables: (i) la densidad de Gibbs es
\eqref{eq:pdf2-ym-gibbs-measure}; (ii) $K_g$ es Markov, reversible y satisface
\eqref{eq:pdf2-ym-schwinger-dyson-finite}--\eqref{eq:pdf2-ym-ward-finite};
(iii) holonomía, clover y curvatura satisfacen
\eqref{eq:pdf2-ym-transport-surface-bianchi} y
\eqref{eq:pdf2-ym-owner-holonomy-curvature-relation}; y (iv) sus funcionales son la
misma ley de pilas de \eqref{eq:pdf2-ym-owner-stack-transfer-identity} y cumplen el
paquete OS\@. El objeto construido es, componente a componente,
\begin{equation}\label{eq:pdf2-ym-owner-qym-object}
\begin{aligned}
 \mathfrak Q_{\rm YM}^{(4)}(G,g):={}&
 \bigl(\{\mathfrak A_{\rm YM}(O)\}_O,
 S_g^{\rm YM},U,\Theta,\Omega_{\infty,g},\\
 &\widehat\mu_{\infty,g}^{\rm YM},
 \{S_{m,g},\mathscr S_{m,g}\}_m,K_g,
 \operatorname{Hol},\mathbf F,\\
 &\mathscr H_{\rm OS},D_{\rm curv,g}^{\rm YM}\bigr)
 \in\mathbf{YM}_4(G,g).
\end{aligned}
\end{equation}
El mapa de representación no omite ninguna de esas componentes:
\begin{equation}\label{eq:pdf2-ym-owner-typed-publication-map}
 \begin{aligned}
 \mathcal P_{{\rm YM},g}:\mathfrak G_{\rm gau}(G)&\longrightarrow
 \mathbf{YM}_4(G,g),\\
 (\widehat\mu^0,\mathbb P^{\rm phys},D^0,\Theta,W,
   \operatorname{Hol},\mathcal P^{F^2})
 &\longmapsto
 (\widehat\mu_g^{\rm YM},S_g,\mathscr S_g,K_g,\mathbb P_g^{\rm phys},
  S_g^{\rm YM},\mathfrak A_{\rm YM},D_{\rm curv,g}^{\rm YM},
  \mathbf F,\operatorname{Hol}).
 \end{aligned}
\end{equation}
Los $S_n^{\rm YM}$ son así exactamente las funciones de Schwinger de la medida
dependiente de $g$, la ley Markov y el operador de transferencia ya construidos.

Queda por comprobar que $S_{m,g}$ es la acción de Yang--Mills y no sólo una
función con ese nombre. Sea
$A\in C_c^3(\mathbb R^4;T^*\mathbb R^4\otimes\mathfrak g)$ una conexión suave en
una coordenatización donde el logaritmo de los lazos pequeños es único. Para la plaqueta
$p_{\mu\nu}(x)$, la expansión ordenada de Magnus da
\begin{equation}\label{eq:pdf2-ym-owner-plaquette-magnus}
 \log\operatorname{Hol}_{p_{\mu\nu}(x)}(A)
 =a_m^2F_{A,\mu\nu}(x)
  +\frac{a_m^3}{2}(D_\mu+D_\nu)F_{A,\mu\nu}(x)
  +R_{m,\mu\nu}(x),
\end{equation}
con
$\|R_{m,\mu\nu}(x)\|\le C_G a_m^4
(\|D^2F_A\|_{\infty,p}+\|F_A\|_{\infty,p}^2)$.
El promedio de las cuatro orientaciones cancela el término impar y define
\begin{equation}\label{eq:pdf2-ym-owner-clover}
 F_{{\rm cl},m}(x)
 :=\frac1{4a_m^2}\sum_{p\ni x}
  \operatorname{Ad}_{T_p}\log\operatorname{Hol}_p,
 \qquad
 \|F_{{\rm cl},m}(x)-F_A(x)\|
 \le C_G a_m^2(\|D^2F_A\|_\infty+\|F_A\|_\infty^2).
\end{equation}
La definición de la raíz
\eqref{eq:pdf2-ym-action-root} es precisamente la polarización de esos seis
componentes. Sustituyendo \eqref{eq:pdf2-ym-owner-plaquette-magnus} en cada
generador de esa polarización se calcula
\begin{equation}\label{eq:pdf2-ym-owner-action-root-smooth-chart}
 \xi_{m,b}(A)
 =a_m^2\operatorname{vec}_{\mu<\nu}F_{A,\mu\nu}(x_b)
  +r_{m,b}(A),
 \qquad
 \|r_{m,b}(A)\|\le C_{G,A}a_m^4.
\end{equation}
Los términos cruzados impares se cancelan por las cuatro orientaciones del clover;
los pares están incluidos una sola vez por
\eqref{eq:pdf2-ym-action-factor-partition}. Por suma de Riemann y
Cauchy--Schwarz,
\begin{equation}\label{eq:pdf2-ym-owner-action-error}
 \left|S_{m,g}(A)-\frac1{4g^2}
  a_m^4\sum_x\|F_{{\rm cl},m}(x)\|^2\right|
 \le C_{G,A,g}a_m^2,
\end{equation}
y, usando la cota de \eqref{eq:pdf2-ym-owner-clover},
\begin{equation}\label{eq:pdf2-ym-owner-f2}
 S_{m,g}(A)\longrightarrow
 S_{{\rm YM},g}(A):=\frac1{4g^2}
 \int_{\mathbb R^4}\|F_A(x)\|^2\,d^4x.
\end{equation}
Las ecuaciones de Schwinger--Dyson y Ward ya se calcularon sobre la medida que usa
esta acción. La brecha, en cambio, se deduce después por la conjugación unitaria de
\eqref{eq:pdf2-ym-g-dynamics} y el testigo
\eqref{eq:pdf2-ym-g-gap-before-limit}; no fue usada para escoger $g$, $S_{m,g}$ ni
$\widehat\mu_{m,g}^{\rm YM}$. Éste es el enlace no circular entre acción, ley y gap.

\begin{theorem}[Cierre Yang--Mills para todo grupo de Lie compacto, conectado y simple]
\label{thm:pdf2-ym-cierre-global}
La estructura completa $\mathfrak G_{\rm gau}$ construye una colección proyectiva
local para todo grupo de Lie compacto, conectado y simple. El mapa
\eqref{eq:pdf2-ym-owner-typed-publication-map} expresa una teoría cuántica de
Yang--Mills no trivial sobre $\mathbb R^4$: sus funcionales satisfacen el paquete OS
continuo y su continuación satisface los axiomas locales de Wightman por la
proposición~\ref{prop:pdf2-ym-owner-continuum-os-package}; sus
holonomías y su curvatura
distribucional satisfacen
\eqref{eq:pdf2-ym-owner-holonomy-curvature-relation}, y las cuatro reconstrucciones
coinducen el mismo Hamiltoniano $D_{\infty,g}^{\rm YM}$, con vacío simple. Su ley es
$\widehat\mu_{\infty,g}^{\rm YM}$, construida por la colección compatible de acciones
locales $\{S_{m,g},\mathscr S_{m,g}\}_m$, y
satisface las ecuaciones Schwinger--Dyson y Ward
\eqref{eq:pdf2-ym-schwinger-dyson-finite}--\eqref{eq:pdf2-ym-ward-finite}. El sector
cíclico de curvatura
\eqref{eq:pdf2-ym-owner-curvature-cyclic-sector} es reductor, contiene el testigo
\eqref{eq:pdf2-ym-owner-wc-curvature-functional} y posee
\begin{equation}\label{eq:pdf2-ym-final-gap}
 \boxed{\Delta_{\rm YM}^{\rm HMT}=3\kappa_{\rm av}>0},
 \qquad
 m_{\rm gap}^{\rm HMT}
 =\frac{\hbar_{\rm int}}{c_{\rm int}}\Delta_{\rm YM}^{\rm HMT}>0.
\end{equation}
La red local, la acción fuerte de $E(4)$, las distribuciones de Schwinger, la
holonomía y $\mathbf F$ proceden del mismo estado completo. El lector exacto
$\mathcal P^{F^2}$ define la densidad de la ley; el clover es su coordenatización suave y
\eqref{eq:pdf2-ym-owner-f2} identifica su acción clásica. La brecha pertenece al sector colectivo
gauge-invariante de curvatura y no asigna una masa elemental al gluón.
\end{theorem}

\begin{proof}
La proyectividad cinemática es \eqref{eq:pdf2-ym-owner-projectivity}; la raíz
\eqref{eq:pdf2-ym-action-root}, el cociclo
\eqref{eq:pdf2-ym-action-cocycle} y la densidad
\eqref{eq:pdf2-ym-gibbs-measure} construyen primero la acción y la medida a
acoplamiento $g$. El transporte factor a factor
\eqref{eq:pdf2-ym-gibbs-transport-factor}--
\eqref{eq:pdf2-ym-full-algebra-identities} prueba proyectividad, preservación del
estado, productos y $*$; 
\eqref{eq:pdf2-ym-transport-surface-bianchi} descarga clover y Bianchi. El índice exhaustivo
\eqref{eq:pdf2-ym-owner-exhaustive-index}, la forma
\eqref{eq:pdf2-ym-owner-full-form} y el circuito
\eqref{eq:pdf2-ym-owner-local-circuit} construyen el generador común. La reducción
física y su kernel cociente son
\eqref{eq:pdf2-ym-owner-physical-reduction}--
\eqref{eq:pdf2-ym-owner-D}. La conjugación
\eqref{eq:pdf2-ym-g-dynamics} los transporta a la medida dependiente de $g$ sin
alterar su espectro. La descomposición de Hoeffding prueba el vacío y la cota, y
\eqref{eq:pdf2-ym-owner-wc-invariance}--
\eqref{eq:pdf2-ym-owner-wc-eigenvalue} exhiben un vector físico que alcanza el
umbral. El entrelazamiento nonádico transporta forma, vacío y testigo al límite, donde
\eqref{eq:pdf2-ym-owner-limit-gap}--
\eqref{eq:pdf2-ym-owner-exact-gap} fijan la brecha exacta. Las cuatro
factorizaciones \eqref{eq:pdf2-ym-owner-four-os}, la inclusión de rangos
\eqref{eq:pdf2-ym-owner-os-range-inclusion} y los conectores
\eqref{eq:pdf2-ym-owner-os-connectors} identifican sus reconstrucciones con ese
mismo generador. La holonomía
\eqref{eq:pdf2-ym-owner-incidence-holonomy} coloca materialmente $W_c$ en el sector
de curvatura, donde \eqref{eq:pdf2-ym-owner-curvature-sector-gap} prueba la brecha
exacta. El sistema \eqref{eq:pdf2-ym-owner-block-intertwiner} controla los productos
cruzados mediante \eqref{eq:pdf2-ym-owner-smeared-cross-gram};
\eqref{eq:pdf2-ym-full-publication-restriction}--
\eqref{eq:pdf2-ym-full-publication-products} elevan ese primer caos al isomorfismo
de toda el álgebra. La convergencia fuerte de $P_\alpha$ produce $\Phi$; las cotas
\eqref{eq:pdf2-ym-owner-holonomy-e4-estimate} y
\eqref{eq:pdf2-ym-owner-orbit-smearing-continuity} prueban además la continuidad
$E(4)$ de holonomías, plaquetas y clover. La desintegración euclídea
\eqref{eq:pdf2-ym-owner-euclidean-markov-identification}, la identidad exacta
\eqref{eq:pdf2-ym-owner-stack-transfer-identity} y el entrelazamiento
\eqref{eq:pdf2-ym-owner-os-markov-intertwining} enlazan la misma medida,
estado Schwinger, pila Markov, transferencia OS y generador del gap. La
proposición~\ref{prop:pdf2-ym-owner-continuum-os-package} prolonga las
cuatro formas finitas a las
semirredes completas, prueba regularidad y cluster y reconstruye el Hamiltoniano. Por
último, \eqref{eq:pdf2-ym-owner-distributional-curvature}--
\eqref{eq:pdf2-ym-owner-typed-publication-map} identifican esos mismos funcionales
como los Schwinger de la teoría gauge local; las integraciones por partes
\eqref{eq:pdf2-ym-schwinger-dyson-finite}--\eqref{eq:pdf2-ym-ward-finite} prueban
Schwinger--Dyson y Ward. Finalmente,
\eqref{eq:pdf2-ym-owner-plaquette-magnus}--\eqref{eq:pdf2-ym-owner-f2} identifican
la acción clásica antes de aplicar la equivalencia espectral. La cadena no usa el
gap para escoger la medida.
\end{proof}

\begin{falsifier}[Colapso de la dinámica producto]
Si se reemplaza $D_m^{\rm prod}$ por $D_m^{\rm fib}$, todos los sectores no
constantes reciben el mismo autovalor. Un vector de Hoeffding con $|S|=2$ debería
tener energía $6\kappa_{\rm av}$ por \eqref{eq:pdf2-ym-hoeffding}, pero el
semigrupo simple le asignaría $3\kappa_{\rm av}$. La identificación de ambos
generadores es falsa.
\end{falsifier}

\begin{falsifier}[Confusión entre inclusión genealógica y promedio de bloque]
Si se sustituye $I_{\alpha\beta}^{\rm curv}$ por
$\widehat J_{\alpha\beta}$ en
\eqref{eq:pdf2-ym-owner-smeared-cross-gram}, una coordenada gruesa persiste con
coeficiente uno mientras el coeficiente fino cambia por el cociente de escalas. Para
$a_{m+1}=a_m/9$, el producto cruzado ancestral lleva un factor $1/81$ y no converge
a la norma diagonal. La identidad
\eqref{eq:pdf2-ym-owner-smeared-cauchy-derived} falla; por eso los dos mapas se
mantienen separados.
\end{falsifier}

\begin{falsifier}[Ablación de la fibra de incidencia]
Si se sustituye $\widehat X_m$ por su marginal de enlaces, se borra $q_c$ y con ella
el observable físico $W_c$. El mapa de olvido no refleja la subálgebra gauge completa;
por tanto no puede utilizarse para negar el autovalor del objeto fuente. Objeto
sustituido; conclusión no transferible.
\end{falsifier}

\begin{falsifier}[Sustitución del lector exacto por su coordenatización suave]
Si se usa directamente el término asintótico de
\eqref{eq:pdf2-ym-owner-f2} como densidad en una pantalla finita, se borra el resto
controlado de \eqref{eq:pdf2-ym-owner-action-error} y se pierde el cociclo exacto
\eqref{eq:pdf2-ym-action-cocycle}. La densidad finita gobernante es
$Z_{m,g}^{-1}e^{-\mathcal P_m^{F^2}/(4g^2)}$; el clover y el límite suave la
identifican con la acción clásica, pero no la sustituyen.
\end{falsifier}

\paragraph{Articulación de las demostraciones.}
La separación entre fibra simple y generador producto, la medida común, la compresión
gauge, las cuatro condiciones de Osterwalder--Schrader, la forma límite y el testigo físico son
resultados exactos previamente establecidos en el dominio gauge
regular HMT. La desintegración con tiempos no uniformes, el índice
exhaustivo de expectativas, el kernel cociente, la matriz $B^+$, la acción explícita
en el colímite, el sistema cofinal de bloque, el paquete OS continuo, la curvatura
distribucional y la holonomía de incidencia son
formalizaciones explícitas: determinan los dominios y demuestran los enlaces antes
enunciados, manteniendo la brecha como resultado. Su composición completa
establece el teorema precedente.
La raíz positiva del lector $F^2$, su densidad Gibbs, el transporte local compatible,
las ecuaciones Schwinger--Dyson/Ward, el isomorfismo algebraico completo, las cotas
de continuidad de holonomías y la identidad pila--transferencia son
formalizaciones explícitas de las cuatro correspondencias. La identificación
entre el observable clover y la acción clásica se obtiene como consecuencia
de la expansión de Magnus.

\section{Generadores tríticos, acciones cuadráticas y transporte de energía}
\label{sec:hamiltonianos-curvatura}

La geometría dinámica conserva la distinción entre el estado que genera una
realización y las coordenadas con las que ésta expresa su evolución. APP
construye y conserva las hojas aditiva y multiplicativa; el TRIT orienta el
modo operativo; el TPK compone selección, transporte y actualización sobre
el estado enriquecido mediante
\[
 U_t=\operatorname{Upd}_t\circ\operatorname{Tra}_t\circ
 \operatorname{Sel}_t.
\]
La holonomía nonádica prolonga este estado a través de
\(w_6\to w_{12}\to w_{18}\to w_{24}\to w_{30}\to R_{36}\to G_9\).
El retorno de fase conserva el avance de memoria. Sus representaciones
\(\pi_{\rm HMT},\varphi_{\rm HMT},e_{\rm HMT},\alpha_{\rm HMT}\)
pertenecen a la misma generación correlacionada, con los lectores y el
orden constructivo interno ya establecidos. La estructura discreta conjunta
del continuo sostiene después las representaciones geométricas y físicas.
En este punto de la construcción, un Hamiltoniano es el funcional que
realiza un transporte ya tipado sobre una fibra simpléctica; su dominio,
su orientación y su parámetro temporal forman parte de la afirmación.

El artículo sobre la extrema y media razón holográfica recupera tres
operadores concretos. La multiplicación por \(4\) módulo \(9\) actúa sobre
las órbitas ternarias de unidades; en su módulo de aumento tiene matriz
\[
 C=\begin{pmatrix}0&-1\\1&-1\end{pmatrix},\qquad C^2+C+I=0.
\]
El transporte parabólico y la incidencia areal--volumétrica están
representados, en sus respectivos planos, por
\[
 N=\begin{pmatrix}0&1\\0&0\end{pmatrix},\qquad
 F_{\rm av}=\begin{pmatrix}0&1\\1&1\end{pmatrix}.
\]
La incidencia procede del calendario de hojas; sus potencias cuentan
composiciones de esos modos. Los generadores reales expresados son
\begin{equation}
 J_+=\frac{2C+I}{\sqrt3},\qquad J_0=N,\qquad
 J_-=\frac{2F_{\rm av}^{2}-3I}{\sqrt5},\qquad
 J_\tau^2=-\tau I.
 \label{eq:hc-generadores}
\end{equation}
Aquí \(\tau=+1,0,-1\) discrimina el régimen elíptico, parabólico e
hiperbólico. Las fibras tienen procedencias y bases propias. Las fórmulas
\(e^{\pi J_+}=-I\), \(e^{J_0}=I+J_0\) y
\(e^{2\log\varphi J_-}=F_{\rm av}^{2}\) evalúan coordenadas HMT
previamente generadas en esos operadores~\cite{hmtX}.

\subsection{Forma alternante y Hamiltoniano del generador}

En una fibra real orientada de dimensión dos se fija
\[
 \omega=dQ\wedge dP,\qquad
 \Omega=\begin{pmatrix}0&1\\-1&0\end{pmatrix},\qquad z=(Q,P)^T.
\]
Se adopta \(\iota_{X_H}\omega=dH\); por tanto
\(\dot Q=\partial_PH\), \(\dot P=-\partial_QH\).
Esta convención fija también el signo de la acción orientada.

\begin{proposition}[Realización hamiltoniana de un generador de traza nula]
\label{prop:hc-cuadratica}
Para una matriz real \(J\) de tamaño dos y traza nula, la matriz
\(G_J=-\Omega J\) es simétrica y el Hamiltoniano
\begin{equation}
 H_J(z)=\frac12z^TG_Jz
 \label{eq:hc-hj}
\end{equation}
produce exactamente \(\dot z=Jz\). Si \(J^2=-\tau I\),
entonces \(\det G_J=\tau\). Para un isomorfismo simpléctico
\(B:V\to V'\), el transporte
\[
 J'=BJB^{-1},\qquad H'(z')=H_J(B^{-1}z')
\]
satisface \(G'=B^{-T}G_JB^{-1}\) y entrelaza ambos flujos.
\end{proposition}
\begin{proof}
Escribiendo \(J=\bigl(\begin{smallmatrix}a&b\\c&-a\end{smallmatrix}\bigr)\),
se obtiene \(G_J=\bigl(\begin{smallmatrix}-c&a\\a&b\end{smallmatrix}\bigr)\).
Así \(\Omega\nabla H_J=\Omega(-\Omega J)z=Jz\).
Cayley--Hamilton da \(J^2=-\det(J)I\), y
\(\det(-\Omega)=1\). Finalmente, \(B^T\Omega B=\Omega\)
permite sustituir \(z=B^{-1}z'\) tanto en la forma como en la ecuación.
\end{proof}

Aplicada a \eqref{eq:hc-generadores}, la proposición da
\begin{equation}
 H_{J_+}=-\frac{Q^2-QP+P^2}{\sqrt3},\qquad
 H_{J_0}=\frac{P^2}{2},\qquad
 H_{J_-}=\frac{-Q^2-QP+P^2}{\sqrt5}.
 \label{eq:hc-hamiltonianos-nativos}
\end{equation}
El signo elíptico registra la orientación de \(C\) en la base de aumento.
El flujo inverso \(-J_+\) posee la forma positiva opuesta. La elipse
de acción del corpus utiliza precisamente una orientación de fase
decreciente para su Hamiltoniano positivo. Conservar este signo hace
compatible el operador cíclico con la acción \(P\,dQ\).

La conjugación también proporciona un criterio exacto de comparación.
Si \(BJ_\tau=J_{\tau'}B\) y \(B\) es invertible, elevar al cuadrado
da \((\tau-\tau')B=0\), luego \(\tau=\tau'\). Un enlace puede
conservar \(\omega\) entre fibras de regímenes distintos y conservar
sus memorias respectivas; la conjugación de los generadores exige además
igualdad del régimen. Esta diferencia permite estudiar una transición
trítica sin sustituirla por un cambio de nombre del mismo flujo.

\subsection{Deformación de la elipse y transformación de Legendre}

La coordenada \(\eta=\operatorname{artanh}(C^*/A)\), en el dominio
\(|C^*/A|<1\), define el marco simpléctico
\(M_\eta=\operatorname{diag}(e^{\eta/2},e^{-\eta/2})\).
La deformación de las dos direcciones es recíproca y conserva su producto.
En la coordenatización normalizada \((q,p)=M_\eta^{-1}(Q,P)\), consideremos
\[
 L_\tau=\begin{pmatrix}0&1\\-\tau&0\end{pmatrix},\qquad
 I_{\tau,\eta}=\frac12\bigl(\tau e^{-\eta}Q^2+e^\eta P^2\bigr).
\]
Estos representantes clasifican flujos cuadráticos orientados. Su empleo
en una fibra nativa se hace mediante un marco simpléctico declarado,
conforme a la proposición~\ref{prop:hc-cuadratica}. El representante
elíptico positivo es el usado por la elipse de acción; para él, \(I\)
tiene dimensión de acción.

Los cambios de marco pueden escribirse explícitamente. Si
\(b_+=\sqrt{\sqrt3/2}\) y \(b_-=\sqrt{\sqrt5/2}\), las matrices
\[
 B_+=\begin{pmatrix}b_+&b_+/\sqrt3\\0&2b_+/\sqrt3\end{pmatrix},
 \qquad
 B_-=\begin{pmatrix}b_-&-b_-/\sqrt5\\0&2b_-/\sqrt5\end{pmatrix}
\]
tienen determinante uno y satisfacen
\( (-J_+)B_+=B_+L_{+1}\) y
\(J_-B_-=B_-L_{-1}\). Para el régimen parabólico se toma
\(B_0=I\). La multiplicación de las dos columnas verifica las
igualdades, y el determinante uno prueba la conservación de la forma
alternante.

\begin{theorem}[Transporte variable y lagrangiano regular]
\label{thm:hc-legendre}
Sean \(\eta(t)\) de clase \(C^1\) y \(\omega(t)>0\). El Hamiltoniano
\begin{equation}
 H_{\tau,\rm tr}
 =\omega I_{\tau,\eta}+\frac{\dot\eta}{2}QP
 \label{eq:hc-htr}
\end{equation}
conserva \(I_{\tau,\eta(t)}\) a lo largo de sus soluciones. Su
transformación de Legendre respecto de \(P\) es invertible y da
\begin{equation}
 \mathcal L_{\tau,\rm tr}(Q,\dot Q,t)
 =\frac{(\dot Q-\dot\eta Q/2)^2}{2\omega e^\eta}
 -\frac{\omega\tau e^{-\eta}}2Q^2.
 \label{eq:hc-ltr}
\end{equation}
\end{theorem}
\begin{proof}
Las ecuaciones son
\[
 \dot Q=\omega e^\eta P+\frac{\dot\eta}2Q,
 \qquad
 \dot P=-\omega\tau e^{-\eta}Q-\frac{\dot\eta}2P.
\]
Derivar \(I_{\tau,\eta(t)}\) y sustituirlas cancela por separado
los términos \(\omega QP\) y los términos de variación de \(\eta\).
La derivada segunda \(\partial_P^2H=\omega e^\eta>0\)
permite resolver
\(P=(\dot Q-\dot\eta Q/2)/(\omega e^\eta)\).
Sustituir en \(P\dot Q-H\) prueba \eqref{eq:hc-ltr}.
Además,
\[
 P\dot Q-H_{\tau,\rm tr}
 =p\dot q-\frac\omega2(\tau q^2+p^2),
\]
lo que prueba que el término mixto es la conexión variacional del marco
móvil. La conservación procede del transporte completo y no de mantener
constantes los semiejes.
\end{proof}

El resultado elíptico y su acción están desarrollados en el corpus
variacional~\cite{hmtX}; la misma prueba exhibe sus dos realizaciones
cuadráticas conjugadas. La nulidad del determinante de la forma parabólica
convive con una transformación de Legendre regular: ésta depende de la
derivada segunda en el momento, no del determinante de la forma completa.
En cambio, para el levantamiento cotangente lineal
\(H_A(q,p)=p^TAq\), se tiene \(\partial_p^2H_A=0\).
Su acción de primer orden conserva \(\dot q=Aq\) mediante un
multiplicador; es un objeto variacional distinto de
\eqref{eq:hc-ltr}.

\subsection{Retorno de memoria y carga variacional conservada}

En el registro de memoria del calendario, los contadores \((g,s)\)
reciben el incremento \((4,72)\) en el retorno completo. La relación
\(72=18\cdot4\) fija
\[
 v=(1,18)^T,\qquad I_{\rm mem}=s-18g,
 \qquad B=\begin{pmatrix}1&0\\18&1\end{pmatrix}.
\]
El estabilizador nativo es
\begin{equation}
 N_v=BNB^{-1}=
 \begin{pmatrix}-18&1\\-324&18\end{pmatrix},\qquad
 N_vm=vI_{\rm mem}(m),\qquad N_v^2=0.
 \label{eq:hc-nv}
\end{equation}
En el plano realificado \(m=(g,s)^T\), dotado de \(dg\wedge ds\),
el Hamiltoniano \(H_{\rm mem}=I_{\rm mem}^2/2\) produce
\(dm/du=N_vm\). El parámetro \(u\) pertenece a esta realización
adimensional del registro. La coordenatización
\(q=g,\ p=s-18g\) es simpléctica y transforma la uno-forma de acción
según
\begin{equation}
 s\,dg=p\,dq+d(9q^2).
 \label{eq:hc-memoria-borde}
\end{equation}
Por tanto, la eliminación de \(s\) en la acción de primer orden da
\begin{equation}
 \mathcal L_{\rm mem}
 =\frac12\left(\frac{dg}{du}\right)^2
 +18g\frac{dg}{du}
 =\frac12\left(\frac{dg}{du}\right)^2
 +\frac{d}{du}(9g^2).
 \label{eq:hc-lmem}
\end{equation}
La ecuación variacional es \(d^2g/du^2=0\), y el momento
conservado en la coordenatización simpléctica es \(p=I_{\rm mem}\).
El retorno afín completo procede del Hamiltoniano lineal
\(H_{\rm ret}=4I_{\rm mem}\): su aplicación al tiempo \(u=1\)
es \((g,s)\mapsto(g+4,s+72)\). Ambos Hamiltonianos conmutan
porque son funciones de la misma carga, pero realizan operaciones
distintas: estabilización parabólica y traslación del registro.
Estas fórmulas recuperan el mecanismo variacional de la memoria, incluido
su término de frontera~\cite{hmtX}.

\subsection{Curvatura métrica y energía geodésica}

La realización métrica de los regímenes utiliza una coordenatización radial y una
escala inversa de longitud \(\rho>0\):
\begin{equation}
 g_{\tau,\rho}=dr^2+S_{\tau,\rho}(r)^2d\theta^2,
 \qquad
 S_{\tau,\rho}(r)=
 \begin{cases}
 \sin(\rho r)/\rho,&\tau=+1,\\
 r,&\tau=0,\\
 \sinh(\rho r)/\rho,&\tau=-1.
 \end{cases}
 \label{eq:hc-metrique}
\end{equation}
La curvatura gaussiana es \(K_G=-S''/S=\tau\rho^2\).
La coordenatización se considera en su interior regular, \(S>0\); para la
realización esférica, \(0<r<\pi/\rho\). Los polos se tratan con
coordenatizaciones regulares diferentes. El parámetro \(\rho\) y la realización
dimensional de la métrica mantienen su declaración propia.

Para una masa realizada \(m>0\), el lagrangiano geodésico y su
Hamiltoniano son
\begin{equation}
 \mathcal L_{\rm geo}=\frac m2(\dot r^2+S^2\dot\theta^2),\qquad
 H_{\rm geo}=\frac1{2m}\left(p_r^2+\frac{p_\theta^2}{S^2}\right).
 \label{eq:hc-geodesique}
\end{equation}
En efecto, \(p_r=m\dot r\), \(p_\theta=mS^2\dot\theta\);
el hessiano de velocidades es \(m\operatorname{diag}(1,S^2)\)
y es positivo e invertible. La ecuación radial es
\(\ddot r-SS'\dot\theta^2=0\), mientras
\(d(mS^2\dot\theta)/dt=0\) conserva el momento angular.
La positividad de esta energía vale para las tres curvaturas. La forma
indefinida del generador hiperbólico de \eqref{eq:hc-hamiltonianos-nativos}
y la energía geodésica positiva de una métrica de curvatura negativa
son funciones sobre espacios distintos. Esta comparación fija exactamente
qué significado de curvatura interviene en cada afirmación.

\subsection{Respuesta constitutiva y realización modal de Maxwell}

Las secciones de acción y las coordenadas angulares preceden a la
respuesta constitutiva \cite{hmtIntegral}. En la orientación real fijada \(x>y>0\),
con los proyectores de las dos hojas,
\[
 T=q_+P_++q_-P_-,\quad q_\pm=e^{-x\mp y},\quad
 \mathcal V=(I-T^{90})(I-T^{120})^{-1},
 \qquad 0<q_+<q_-<1,
\]
el cálculo funcional produce
\[
 r_\pm=\frac{1-q_\pm^{90}}{1-q_\pm^{120}},\qquad
 \widehat\varepsilon=r_+^2,\quad \widehat\mu=r_-^2,
 \quad \widehat c=(r_+r_-)^{-1},\quad \widehat Z=r_-/r_+.
\]
Las coordenatizaciones dimensionales correspondientes conservan exactamente
\(\varepsilon\mu=c^{-2}\) y \(\mu/\varepsilon=Z^2\).
La permitividad, la permeabilidad y la velocidad son así representaciones
del mismo operador, con los sectores de incidencia fijados por el TPK;
la dinámica siguiente las emplea como resultados ya construidos.

Para exhibir el transporte hacia un Hamiltoniano de campo, se fija una
realización espacial orientada, condiciones de frontera que permitan
integración por partes, y un modo transversal real \(f\) normalizado por
\[
 \nabla\cdot f=0,\qquad \int|f|^2=1,\qquad
 \nabla\times\nabla\times f=k^2f,\qquad k>0.
\]
En calibre transversal y sin fuentes, \(A=af\) y el momento
\(\Pi=pf\) definen una inmersión del plano modal en el espacio de
campos. Para constantes constitutivas homogéneas, el Hamiltoniano y la
forma simpléctica de Maxwell se restringen a
\begin{equation}
 H_{\rm EM}(a,p)=\frac{p^2}{2\varepsilon}
 +\frac{k^2a^2}{2\mu},\qquad \omega_{\rm EM}=da\wedge dp.
 \label{eq:hc-maxwell-mode}
\end{equation}
La integración por partes da
\(\int|\nabla\times f|^2=k^2\), que prueba esta restricción.
Con \(\omega_k=ck\), el cambio
\begin{equation}
 Q=\sqrt{\varepsilon\omega_k}\,a,
 \qquad P=\frac{p}{\sqrt{\varepsilon\omega_k}}
 \label{eq:hc-maxwell-symplectic}
\end{equation}
preserva \(da\wedge dp=dQ\wedge dP\) y da
\(H_{\rm EM}=\omega_k(Q^2+P^2)/2\). Su inversa y sus ecuaciones
\(\dot Q=\omega_kP,\ \dot P=-\omega_kQ\)
constituyen el entrelazamiento modal con el régimen elíptico positivo.
La restricción es invariante porque el operador de Maxwell conserva el
modo propio transversal. La transformación de Legendre produce
\(\mathcal L_{\rm EM}=\varepsilon\dot a^2/2-k^2a^2/(2\mu)\).

Dos modos transversales de igual frecuencia, combinados en cuadratura,
realizan las dos orientaciones circulares del corpus. El campo satisface
\(\mathbf B=c^{-1}\widehat{\mathbf k}\times\mathbf E\) y
\(\omega=ck\). La relación con el oscilador se demuestra aquí mediante
la inmersión de campos, la normalización modal y el transporte de la forma
simpléctica, manteniendo explícitas las condiciones de frontera.
Un cambio de \(k\), de la geometría de la cavidad o de la preparación
modifica el modo realizado; mantiene intacta la genealogía de
\(\varepsilon,\mu,c,Z\).

En la coordenatización lorentziana del mismo sector, el acoplamiento a una sección
material se escribe con \(F=dA\) y una carga \(q_e\) ya tipada.
Para esta realización variacional se toma \(\nabla\) igual a la
conexión de Levi--Civita de \(g\). Sobre
\(T^*\mathcal M\), el Hamiltoniano covariante
\[
 H_A=\frac1{2m}g^{\mu\nu}(p_\mu-q_eA_\mu)(p_\nu-q_eA_\nu)
\]
es la transformación de Legendre de
\(\mathcal L_A=\frac m2g_{\mu\nu}\dot x^\mu\dot x^\nu
+q_eA_\mu\dot x^\mu\). El hessiano \(mg\) es invertible.
La ecuación de Euler--Lagrange da
\(m\nabla_u u^\mu=q_eF^\mu{}_{\nu}u^\nu\).
La antisimetría de \(F\) conserva \(g(u,u)\). Se trata de la
realización covariante con parámetro afín, mientras
\eqref{eq:hc-maxwell-mode} utiliza el tiempo de una descomposición
espacial transversal: ambos parámetros y ambos dominios quedan definidos.

\subsection{Canales graduados, representación de Clifford y operador de masa}

El atlas material produce primero canales completos de ruta, hoja, sabor,
color y representación. En una profundidad finita, su realización
hermítica es \(\mathcal H_K=\bigoplus_{s\in S_K}V_s\), con base
ortonormal de canales y graduación de paridad transportada. El álgebra
exterior del sector impar genera las CAR: la inserción de un modo
\(a_j^*\) y su retirada adjunta \(a_j\) satisfacen
\(\{a_i,a_j^*\}=\delta_{ij}I\) y \(\{a_i,a_j\}=0\).
La prueba cuenta los signos de intercambio de modos completos; las masas
se evalúan después~\cite{hmtVI}.

En el sector exterior de un solo modo, la base \((|0\rangle,|1\rangle)\)
identifica el operador de retirada con
\(a=N\). El mapa \(e_1\mapsto|0\rangle,\ e_2\mapsto|1\rangle\)
es una isometría del plano complejo normalizado y entrelaza ambos
nilpotentes. Esta identificación emplea el producto interno y la
graduación declarados; conserva el distinto significado de sus bases.
Con el \(i\) interno ya realizado, se obtienen
\begin{equation}
 \sigma_1=N+N^*,\qquad
 \sigma_2=-i(N-N^*),\qquad
 \sigma_3=NN^*-N^*N.
 \label{eq:hc-pauli}
\end{equation}
La multiplicación directa da
\(\sigma_j^*=\sigma_j\) y
\(\sigma_i\sigma_j=\delta_{ij}I+i\epsilon_{ijk}\sigma_k\).
En particular, la pareja nilpotente y su adjunto producen una
representación de Clifford; el adjunto es el de la fibra hermítica y
se transporta junto con ella.

El operador de masa sobre una colección finita de canales \(\mathcal F\)
está construido por el registro, el carácter y sus refinamientos:
\begin{equation}
 M_\beta=S M_0 S,\qquad
 S=R_{\rm act}^{B_M/2},\quad
 R_{\rm act}=\hbar_{\rm pre}/\hbar_{\rm ret}>0,
 \qquad M_0>0.
 \label{eq:hc-mass}
\end{equation}
Aquí \(B_M\) es el operador diagonal del exponente de ruta y \(M_0\)
contiene el carácter refinado relativo al origen electrónico.
La positividad sigue de \(\langle v,M_\beta v\rangle
=\langle Sv,M_0Sv\rangle>0\) para \(v\ne0\).
El producto de hojas conserva el sector nulo por su realización propia;
la exponencial del carácter material permanece positiva.

En esta realización se escribe \(\hbar:=\hbar_{\rm ret}\),
la sección de acción ya construida, y \(c:=c_{\rm vac}\) en la
coordenatización dimensional declarada. Sobre \(\mathbb C^4\otimes\mathcal F\), definamos
\[
 \alpha_j^{D}=\sigma_1\otimes\sigma_j,\qquad
 \beta^D=\sigma_3\otimes I_2.
\]
El superíndice distingue estas matrices de la constante de estructura fina.
La coordenatización espacial \(1+3\) asigna tres componentes de momento
\(\mathbf p\). La aplicación de multiplicación
\begin{equation}
 H_D(\mathbf p)
 =c\sum_{j=1}^3\alpha_j^Dp_j\otimes I_{\mathcal F}
 +c^2\beta^D\otimes M_\beta
 \label{eq:hc-dirac}
\end{equation}
es la realización libre de Dirac de ese operador material.

\begin{theorem}[Factorización espinorial de la dispersión material]
\label{thm:hc-dirac-square}
La matriz \eqref{eq:hc-dirac} es autoadjunta y satisface
\begin{equation}
 H_D(\mathbf p)^2
 =I_4\otimes\bigl(c^2|\mathbf p|^2I_{\mathcal F}
 +c^4M_\beta^2\bigr).
 \label{eq:hc-dirac-square}
\end{equation}
En un autoespacio \(M_\beta v=m v\), sus energías son
\(\pm E_m(\mathbf p)\), donde
\(E_m=\sqrt{c^2|\mathbf p|^2+m^2c^4}\).
\end{theorem}
\begin{proof}
Las relaciones de \eqref{eq:hc-pauli} dan
\(\{\alpha_i^D,\alpha_j^D\}=2\delta_{ij}I_4\),
\(\{\alpha_i^D,\beta^D\}=0\) y \((\beta^D)^2=I_4\).
El factor de sabor conmuta con las matrices espinoriales.
Los términos cruzados se anulan y los cuadrados dan
\eqref{eq:hc-dirac-square}. La traza nula y las relaciones de Clifford
producen ambas ramas, con multiplicidad espinorial dos cuando \(E_m>0\).
\end{proof}

En representación de momentos, el dominio natural es
\(\{\psi\in L^2(\mathbb R^3;\mathbb C^4\otimes\mathcal F):
H_D\psi\in L^2\}\). Al ser una matriz hermítica mensurable actuando
por multiplicación, el operador es autoadjunto en ese dominio. La
representación espacial procede de la transformación unitaria de Fourier
posterior a la construcción de los canales. El teorema fija la realización
libre y su dependencia exacta de la masa HMT; la interacción conserva sus
operadores de conexión y su dominio propio.

En la rama positiva \(m>0\), la transformación de Legendre ordinaria
del símbolo \(E_m(\mathbf p)\) es
\[
 \mathbf v=\frac{c^2\mathbf p}{E_m},\qquad
 \mathcal L_m=-mc^2\sqrt{1-|\mathbf v|^2/c^2},
 \qquad |\mathbf v|<c.
\]
Resolver \(\mathbf p=m\mathbf v/\sqrt{1-|\mathbf v|^2/c^2}\)
y sustituir en \(\mathbf p\cdot\mathbf v-E_m\) demuestra la
igualdad. La acción espinorial, en cambio, es de primer orden en el campo:
\(\mathcal L_D=\frac{i\hbar}{2}(\psi^*\dot\psi-\dot\psi^*\psi)
-\psi^*H_D\psi\). Su variación da
\(i\hbar\dot\psi=H_D\psi\); su hessiano en velocidades del campo
es singular. Ambas afirmaciones son compatibles porque actúan sobre
espacios variacionales diferentes.

Para una conexión abeliana estática en una coordenatización espacial plana, con
\(\Pi_j=-i\hbar\partial_j-q_eA_j\), masa constante y potencial
escalar nulo, se conserva sobre funciones lisas de soporte compacto
\[
 [\Pi_i,\Pi_j]=iq_e\hbar\epsilon_{ijk}B_k,
\]
y la misma factorización produce
\begin{equation}
 H_D(A)^2=c^2\boldsymbol\Pi^2+c^4M_\beta^2
 -q_e\hbar c^2\boldsymbol\Sigma\cdot\mathbf B,
 \qquad \Sigma_j=I_2\otimes\sigma_j.
 \label{eq:hc-dirac-curvature}
\end{equation}
Los factores identidad de sabor quedan implícitos. El término de curvatura
es el producto de los dos antisimétricos: el conmutador de derivadas
covariantes y la multiplicación de Pauli. Su coeficiente queda fijado
al elegir la representación de carga y la coordenatización de acción; no se
selecciona mediante un valor espectral objetivo.

Si \(V\) es la transformación unitaria de sabor ya construida, entonces
\[
 (I_4\otimes V)H_D(M_\beta)(I_4\otimes V^*)
 =H_D(VM_\beta V^*).
\]
La mezcla unitaria conserva el espectro completo. La congruencia positiva
de \eqref{eq:hc-mass} transporta otra estructura y puede modificar los
autovalores en una métrica fija. Por ejemplo,
\(M_0=\operatorname{diag}(1,2)\), \(S=\operatorname{diag}(2,1)\)
da \(SM_0S=\operatorname{diag}(4,2)\). Si se transporta también el
producto interno a \(G'=S^*S\), el problema generalizado
\(SM_0Sv=\lambda G'v\) recupera los autovalores de \(M_0\)
mediante \(w=Sv\). De este modo, cambio de base de sabor, deformación
material y transporte de métrica tienen consecuencias espectrales precisas.

\subsection{Curvatura del potencial escalar en la realización electrodébil}

El sector electrodébil utiliza el determinante del mismo operador angular,
\(D_A=\det T=e^{-2A_{\rm rad}}\), junto con las firmas materiales ya
construidas. En la realización angular con refinamientos comunes de
\(W\) y \(Z\), sus firmas \((94,-1)\) y \((95,-1)\) dan
\(m_W^\angle/m_Z^\angle=e^{-A_{\rm rad}}\).
En la convención racionalizada de árbol del corpus,
\[
 g^2=\frac{4\pi_{\rm HMT}\alpha_{\rm HMT}}{1-D_A},\qquad
 \lambda_H=\frac{\pi_{\rm HMT}\alpha_{\rm HMT}}{2(1-D_A)}
 \exp\!\left(6A_{\rm rad}+\frac{10}3C^*_{\rm rad}\right)>0.
\]
La segunda fórmula resulta de la firma escalar \((97,9)\) y de
\(m_H^2/m_W^2=8\lambda_H/g^2\). La normalización del potencial es
\(V(\mathsf H)=-\mu_H^2\mathsf H^*\mathsf H
+\lambda_H(\mathsf H^*\mathsf H)^2\). Estas composiciones conservan
su escala y su convención de árbol; los refinamientos que se cancelan
en una razón quedan declarados en la propia realización.

En una coordenatización radial real \(\mathsf H^*\mathsf H=h^2/2\), el potencial
es \(V(h)=-\mu_H^2h^2/2+\lambda_Hh^4/4\). Para
\(\mu_H^2>0\), sus puntos críticos son \(0\) y
\(\pm v\), con \(v^2=\mu_H^2/\lambda_H\), y
\[
 V''(0)=-\mu_H^2,\qquad V''(\pm v)=2\mu_H^2.
\]
En unidades naturales de esta coordenatización, el modo homogéneo normalizado posee
\(H_h=p_h^2/2+V(h)\). Su linealización en un punto crítico \(h_0\)
tiene generador
\[
 A_{h_0}=\begin{pmatrix}0&1\\-V''(h_0)&0\end{pmatrix},
 \qquad A_{h_0}^2=-V''(h_0)I.
\]
Por tanto, el origen es hiperbólico y los mínimos son elípticos; la
frontera \(V''=0\) tiene el tipo parabólico. Esta afirmación se deduce
de la derivada segunda del potencial y de una normalización modal
explícita. La energía material positiva del mínimo y la dirección
inestable de otro fondo pertenecen así a la misma función con distintos
puntos de expansión. El modo linealizado es un transporte particular de
los tres tipos cuadráticos; sus frecuencias y unidades proceden del
potencial, no de identificar su hessiano con la curvatura gaussiana de
\eqref{eq:hc-metrique}.

\subsection{Modo de curvatura de Yang--Mills y oscilación elíptica}

La realización de Yang--Mills construida en las secciones precedentes
posee un generador positivo \(D^{\rm YM}\) de dimensión inversa de
longitud. La escala
\[
 \kappa_{\rm av}=\frac{\mu_{\rm vol}}{\ell_0}
 \left(\frac{\pi_{\rm HMT}}{\varphi_{\rm HMT}^2}-1\right)>0,
 \qquad \delta=3\kappa_{\rm av}
\]
desciende de la incidencia de hojas y del lector regional, con la unidad
longitudinal \(\ell_0\) declarada. El vacío es simple y el testigo de
curvatura en la coordenatización producto es
\[
 W(q)=\mathbf1_{\{q_1+\cdots+q_6=0\}}-\frac13,
 \quad \mathbb EW=0,\quad \|W\|^2=\frac29,
 \quad D^{\rm YM}W=\delta W.
\]
El transporte unitario hacia la medida de interacción lleva \(W\) a
\(W_g\) y conserva esas identidades. Las inclusiones de refinamiento
\(J_m\) entrelazan los generadores y conservan el testigo, de modo que
éste permanece en el límite físico. Aquí se utiliza el resultado completo
de las pruebas de vacío, brecha uniforme y reconstrucción, no un
autovalor aislado de una matriz finita~\cite{HMT-YM}.

Sea \(w=W_g/\sqrt{2/9}\) el vector unitario en la fibra física y
\(\hbar=\hbar_{\rm ret}>0\). El Hamiltoniano energético es
\(H_E=\hbar cD^{\rm YM}\). Sobre el espacio de Hilbert realificado
se adopta la forma \(\omega_{\mathcal H}(u,v)=2\hbar
\operatorname{Im}\langle u,v\rangle\), con producto interno antilineal
en el primer argumento. Definamos
\begin{equation}
 \mathfrak F(Q,P)=\frac{Q+iP}{\sqrt{2\hbar}}\,w.
 \label{eq:hc-ym-embedding}
\end{equation}

\begin{theorem}[Inmersión simpléctica del modo físico de curvatura]
\label{thm:hc-ym-oscillator}
La aplicación \(\mathfrak F:\mathbb R^2\to\mathcal H^{\rm phys}\)
es una inmersión real con imagen invariante por la evolución física y
satisface
\begin{equation}
 \mathfrak F^*\omega_{\mathcal H}=dQ\wedge dP,
 \qquad
 \langle\mathfrak F,H_E\mathfrak F\rangle
 =\frac{c\delta}{2}(Q^2+P^2).
 \label{eq:hc-ym-energy-pullback}
\end{equation}
La ecuación de Schrödinger restringida entrelaza exactamente el flujo
elíptico \(\dot Q=c\delta P,\ \dot P=-c\delta Q\).
El mapa es compatible con las inclusiones de refinamiento.
\end{theorem}
\begin{proof}
Para dos variaciones \((a,b),(a',b')\), el producto interno de sus
imágenes tiene parte imaginaria \((ab'-ba')/(2\hbar)\).
Esto prueba la primera identidad. La ecuación propia y \(\|w\|=1\)
dan la segunda. Puesto que
\(\dot\psi=-icD^{\rm YM}\psi\), el coeficiente \(Q+iP\)
satisface \(\dot Q+i\dot P=-ic\delta(Q+iP)\).
Finalmente, \(J_mw_m=w_{m+1}\) implica
\(J_m\mathfrak F_m=\mathfrak F_{m+1}\). Estas identidades pasan al
límite por las isometrías ya construidas.
\end{proof}

El plano completo describe amplitudes de un modo; su círculo
\(Q^2+P^2=2\hbar\) representa los vectores normalizados de ese
autoespacio complejo. Sobre tal círculo la fase global evoluciona y el
rayo cuántico permanece fijo. Superposiciones de autoespacios distintos
conservan, además, sus fases relativas. La inmersión mantiene esta
diferencia entre espacio de amplitudes y espacio proyectivo.

La energía asociada al salto espectral y el salto de masa son
\(\Delta_E=\hbar c\delta\) y \(m_{\rm gap}=\hbar\delta/c\).
La frecuencia modal \(c\delta\) tiene dimensión de tiempo inverso.
Las longitudes euclídeas de \(e^{-sD^{\rm YM}}\) y el tiempo físico
de \(e^{-itH_E/\hbar}\) se relacionan mediante la coordenatización y la
reconstrucción, conservando sus distintos grupos y semigrupos.
La composición \(\mathfrak F\circ M_\eta^{-1}\) transporta esta
misma energía a la elipse de acción; si el marco varía, el término de
conexión de \eqref{eq:hc-htr} completa su evolución.

El enlace obtenido es una realización exacta sobre el modo invariante de
curvatura y sus prolongaciones. La positividad del Hamiltoniano cuántico
genera una rotación en su realificación simpléctica; la denominación
hiperbólica de una métrica espacial pertenece al tensor métrico, y el
transporte unipotente del registro pertenece a su memoria. En dimensión
finita, un grupo \(I+tN\) unitario para todo \(t\in\mathbb R\), con
\(N^2=0\), fuerza \(N=0\): derivar la unitariedad da \(N^*=-N\),
y entonces \(N^*N=-N^2=0\). El flujo parabólico clásico conserva su
forma simpléctica y posee su cuantización en espacios apropiados; esta
identidad sólo delimita la identificación con una matriz unitaria finita.

Así queda articulado un conjunto de transportes concretos. La memoria
produce una carga y una acción parabólica; el marco de la elipse produce
la conexión variacional y su lagrangiano; las respuestas constitutivas
determinan los coeficientes de Maxwell; la graduación de canales construye
Clifford y permite factorizar el operador material; la reconstrucción de
Yang--Mills aporta una inmersión modal simpléctica que conserva exactamente
su brecha y su refinamiento. Cada identificación dispone de una aplicación,
una forma conservada y una ecuación entrelazada.


\section{Discusión: reconstrucción, invariantes y dualidad}
\label{sec:discussion}
El resultado estructural prioritario es la recuperabilidad de la forma.
La carga uniforme y las dos memorias nonádicas reconstruyen el registro
mediante una inversa explícita; el cono positivo describe exactamente
el dominio donde esa reconstrucción induce menores estrictamente positivos.
La cota de estabilidad cuantifica la persistencia de ese dominio. Por
tanto, los pesos geométricos admiten una procedencia y una recuperación
operatorias, además de su evaluación numérica.

La composición desarrollada tiene una secuencia explícita. El registro
terminal genera doce pesos positivos; sus sumas parciales dan trece puntos
ordenados; una red planar realiza sus menores; la curva de momentos
construye polígonos y politopos; sus triangulaciones proporcionan cobertura
y formas con residuos orientados. El refinamiento de una celda fija conserva
la suma de sus formas. La introducción de nuevos vértices en la curva de
momentos produce, por su parte, otra geometría y exige comparar dominios.

La dirección algebraica del registro distingue dos operaciones. El
reescalado positivo de columnas conserva las razones proyectivas pertinentes.
Su acción normalizada sobre los intervalos genera en cambio una deformación
efectiva, certificada por una derivada exacta no nula. Esta deformación admite
un potencial convexo cuya transformada de Legendre se conserva al subdividir
con dirección heredada. La invariancia tiene así un mecanismo comprobable:
cada sucesor conserva el generador de su predecesor y la suma de sus pesos recupera
el peso anterior.

El refinamiento energético proporciona una segunda conservación exacta.
La reducción de Schur recupera el laplaciano efectivo y la prolongación
minimizante separa cada configuración en su parte observada y un detalle
reconstructible. Las energías de los detalles se suman ortogonalmente.
Con mallas decrecientes y energía uniformemente acotada, esa sucesión
converge en la realización de Sobolev declarada. La invariancia de la
forma canónica, la identidad del potencial convexo y la conservación de
energía son tres resultados distintos de compatibilidad, demostrados
para los refinamientos especificados.

La prolongación a los doce planos reticulares conserva una simetría
ternaria y el volumen. El resultado decisivo es la reciprocidad
\(\Lambda_s^*=\Lambda_{-s}\), que induce la inversión
\(\Theta_s(t)=t^{-12}\Theta_{-s}(t^{-1})\). La realización partida de
signatura \((24,24)\) conserva integralidad, paridad y autodualidad para
todo el flujo. En la proyección euclídea, la variación explícita de una
norma distingue deformación y automorfismo de la red inicial. Esta
distinción determina el contenido exacto de la conexión con la dualidad
de momentos y devanados.

La realización material conserva sus propios lectores. La representación de
acción y el carácter de masas reciben el registro por composiciones ya
construidas; la congruencia preserva positividad e inercia, mientras la
métrica determina la interpretación de sus autovalores. Estas diferencias
precisan el significado de multimorfología: varias formas matemáticas
proceden de una genealogía común y se relacionan por operaciones verificables.
La identificación completa de amplituedros de rango superior requerirá
las pruebas de cobertura y residuos para esos mapas concretos. El resultado
de rango uno demostrado aquí fija un caso completo, con coordenadas
racionales, triangulación y controles reproducibles.

La orientación temporal del TRIT integra el retorno con memoria, la
actualización de umbral y las prolongaciones conjugadas. Sus generadores
elíptico, nilpotente e hiperbólico proporcionan leyes distintas de
evolución; la composición del TPK conserva la genealogía que las
concatena. En particular, la actividad del régimen cero queda expresada
por un transporte unipotente, y la doble proyección conserva simultáneamente
la convergencia de las coordenadas y la incidencia excepcional. La tesis
geométrica del artículo reside en esta articulación: la estructura de
partida genera registros recuperables y éstos sostienen realizaciones
positivas cuyas transformaciones, fronteras y memorias se pueden calcular
y comparar con precisión.

El entrelazamiento \(B_FM=MB_V\) añade un transporte exacto desde el
bloque modular hasta el cociente incidencial. En él, las multiplicidades
uno y tres se deducen de la matriz de Gram, y la forma firmada declarada
se expresa como \(J_L=(7I_Q-\overline B_F)/2\). Las matrices de
Vandermonde y de Hankel aportan, por otra parte, realizaciones explícitas
\(Y^{(k)}\) para \(2\le k\le9\), con positividad y rango probados
en la partición inicial. La prolongación nonádica extiende después
esa construcción a todo par finito admisible de rango y dimensión
externa, mediante los sistemas de nodos y momentos desarrollados aquí.
Estos resultados amplían el sistema de realizaciones sin identificar
la dimensión de un cociente, el rango de una matriz y el número de
generadores de una superálgebra.

En la realización gauge, la escala areal--volumétrica precede al
Hamiltoniano. El generador producto conserva todos los sectores de la
descomposición ortogonal, mientras su restricción física contiene un
testigo que alcanza el primer nivel positivo. Los entrelazadores de
refinamiento conservan la coercividad y ese testigo. El desarrollo
continuo reúne además la medida común, la curvatura, la continuidad
euclídea, las formas de reflexión y el Hamiltoniano reconstruido. Así,
el argumento de brecha queda situado junto a la construcción que le da
dominio y significado, y la positividad de Plücker mantiene su realización
geométrica propia. Su articulación exige esos transportes explícitos;
la palabra positividad designa propiedades tipadas que el estado común
permite comparar sin sustituir unas por otras.

\clearpage
\part{Comprobación exacta y enumeración de las emisiones}
\section{Comprobación exacta de las construcciones finitas}
\label{sec:exact-reproduction}

La reproducción distingue la demostración general y la evaluación de
una especialización. Las proposiciones precedentes utilizan identidades
algebraicas, positividad e hipótesis de convergencia explícitas. Sus
controles computacionales se realizan con enteros y fracciones racionales;
el desarrollo reticular conserva, además, su representación exacta en
la extensión algebraica declarada. Los datos del registro son salidas
expresadas por la construcción anterior. Las ejecuciones de esta
sección comprueban su transporte geométrico y operatorio, y conservan
por separado la procedencia de su generación.

El primer control recupera las separaciones a partir de los menores,
comprueba la relación de Plücker en las triangulaciones y compara
cambios de coordenatización con deformaciones efectivas. Un cambio invertible
de las filas multiplica todos los menores por el mismo determinante;
los cocientes marcados permanecen iguales. Un reescalado independiente
de las columnas conserva ciertas razones dobles y altera, en general,
las separaciones reconstruidas. Evaluar ambas operaciones sobre los
mismos datos proporciona un control negativo de la equivalencia
que se afirma.

El segundo control subdivide las longitudes por nueve, conserva las
marcas heredadas y verifica las identidades de promedios y de Gram.
Los determinantes se calculan por eliminación exacta y se contrastan
con sus productos de Vandermonde. La positividad en cualquier
dimensión finita procede de la fórmula demostrada; los casos finitos
ejecutados comprueban la implementación de esa fórmula, sus índices
y sus normalizaciones. Las medidas de conteo y las medidas
cronológicas se mantienen separadas para que una igualdad de
momentos no dependa de cambiar silenciosamente la normalización.

El tercer control reproduce los ocho primeros momentos mediante
el operador de Jacobi de orden cuatro. La matriz simétrica se
representa de forma racional en la base de polinomios mónicos;
su métrica es \(\operatorname{diag}(h_0,\ldots,h_3)\). Se
comprueba la autoadjunción en esa métrica, la recurrencia de tres
términos y el acuerdo entre el resolvente matricial y su fracción
continua en \(z=2\). El resultado mantiene explícitos el vector
constante observado y la carga \(Q\), que intervienen en el
teorema de fidelidad espectral marcada.

El cuarto grupo verifica los cambios simplécticos y sus
Hamiltonianos. La derivada respecto del momento recupera la
velocidad; la transformación de Legendre se comprueba como
identidad polinómica o de Laurent. La inercia de una forma
cuadrática y la positividad de un Hamiltoniano cuántico
pertenecen a sus espacios respectivos. La representación
de Clifford se comprueba mediante sus anticonmutadores,
y la composición con el operador de masa se prueba por
el cuadrado del operador de Dirac. Estas igualdades
identifican cada término y sus productos cruzados.

Finalmente se conservan las comprobaciones anteriores de
triangulación, forma canónica, reducción de Schur, congruencia
de masas y construcción reticular. En Yang--Mills, las
proposiciones sobre medida común, refinamiento y reconstrucción
se exponen en el cuerpo con sus premisas y falsadores. Un
control finito de matrices no sustituye el paso al límite
ni la reconstrucción de campos: sus dominios se registran
separadamente en la entrega reproducible.

El paquete de fuentes enlaza cada programa con el enunciado
que evalúa y conserva su salida íntegra. Las condiciones se
comprueban mediante bifurcaciones explícitas, de modo que
permanecen activas al ejecutar el intérprete en modo
optimizado. Se reproducen asimismo en modo aislado, sin
cargar paquetes de usuario. La correspondencia entre
construcción, prueba y ejecución puede inspeccionarse
paso a paso sin convertir la salida de un programa en
una afirmación de alcance superior a la que comprueba.

\import{foundation/}{libro/censo_firmas_emision.tex}
\clearpage
\section{Conclusiones}
\label{sec:conclusions-final}
\begingroup
\fontsize{10.2}{12.8}\selectfont
\setlength{\parskip}{3pt plus .5pt}

La tesis geométrica queda establecida mediante reconstrucción y transporte:
una representación positiva de la estructura discreta del continuo conserva,
con sus marcas, el registro del que procede y los lectores que éste
determina. La carga total y los cocientes de menores recuperan cada
coordenada; las dos memorias nonádicas proporcionan otra reconstrucción
operatoria. El valor y la forma se articulan así por inversas explícitas,
después de la generación APP--TRIT--TPK y de sus prolongaciones compatibles.
La doble proyección mantiene unidas la realización aritmética y la
incidencia excepcional dentro de esa genealogía.

El refinamiento supera la limitación de una matriz con trece columnas.
Para cualquier rango y dimensión externa finitos, un nivel nonádico
suficiente produce los menores positivos y la imagen amplituhedral
correspondiente. Los nodos heredados conservan sus banderas; los
promedios celulares conservan la medida y registran el detalle
ortogonal. En rango uno, la cobertura simplicial y la recursión de
residuos se extienden a toda dimensión finita. En rangos mayores,
el resultado determina colecciones positivas compatibles con rango y
límite probados; su amplitud dimensional queda separada de la
clasificación global de todas las celdas amplituhedrales.

La conservación admite tres expresiones complementarias. La forma
canónica se suma bajo subdivisión de un soporte fijo; la energía
efectiva se conserva mediante reducción de Schur; y la medida
cronológica determina momentos cuyo incremento se descompone en
detalles recuperables. Estos momentos construyen operadores de
Jacobi cuyas compresiones finitas tienen espectro simple y residuos positivos. Se obtiene así
un recorrido efectivo desde las separaciones aritméticas hasta una
representación espectral, con fracción continua determinada por la
misma medida. Las mutaciones que cambian únicamente la coordenatización
conservan los lectores del punto marcado.

La realización física gana precisión mediante sus operadores. Los
Hamiltonianos de los tres regímenes tríticos poseen formas
simplécticas y transformaciones de Legendre explícitas. La coordenatización
elíptica variable incorpora su término de conexión; la memoria
parabólica conserva su invariante; la apertura hiperbólica tiene
su evolución propia. Los lectores de acción y respuesta
constitutiva se transportan por la inversa grassmanniana. Con
las representaciones regionales y las marcas enteras declaradas,
la respuesta constitutiva etiquetada recupera a su vez el registro.
La condición de igualdad de trazas gobierna el pegado de formas
con valores en la recta de acción.

La deformación excepcional conserva el volumen y la simetría
ternaria, y convierte la inversión del parámetro en dualidad
reticular. Su reciprocidad theta y su realización de signatura
partida enlazan la variación geométrica con la dualidad de la
red. En Yang--Mills, la escala areal--volumétrica, la medida común,
la coercividad uniforme y el testigo persistente sitúan la brecha
junto al Hamiltoniano que la realiza. El desarrollo conserva
el espacio físico y los entrelazadores necesarios para transportar
esa afirmación a través del refinamiento.

El resultado conjunto es una arquitectura de formas con memoria:
la dimensión puede aumentar, la representación puede cambiar y
la información pertinente permanece recuperable mediante mapas
determinados. Éste es el contenido diferencial del artículo.
La geometría positiva, la dualidad reticular y la realización
espectral quedan enlazadas por construcciones que especifican
qué generan, qué conservan y cómo se reconstruyen desde un mismo
origen discreto.
La representación de incidencia conserva también la componente uniforme
del registro: su inversa reconstruye las doce coordenadas y permite
recuperar la dirección geométrica después del transporte de memoria.
En la realización sectorial, el área total y la ocupación ponderada
determinan conjuntamente el reparto héxada--óctada. La purificación
especificada convierte ese reparto en una estructura de entrelazamiento,
con identidad finita entre entropía, áreas sectoriales y entropía relativa.
La conservación de información queda así precisada para cada aplicación:
registro completo, dirección, área y composición poseen inversas o
fibras explícitamente distintas.
La escala de área se determina mediante una longitud generada por la composición
lineal del retorno registrado y la norma local de Hadamard. La
reciprocidad radial determina el coeficiente gravitatorio en la
realización reducida declarada y reconoce después esa longitud como
longitud de Planck. Las mutaciones del mismo punto marcado y los
refinamientos heredados conservan el operador de área, sus
multiplicidades y la identidad modular. Las formas canónicas con
valores en ese operador extienden la conservación a los residuos,
mientras el radio de cada horizonte mantiene su factor geométrico
propio. Esta composición hace explícita la dependencia entre
longitud, incidencia y área sin alterar los estados sectoriales.
\par\endgroup
\clearpage


\begin{thebibliography}{99}
\bibitem{hmtI} O. Haidara Fall y R. Ramos Balsa,
\emph{Generación coinductiva de \(\pi,\varphi,e\) a profundidad arbitraria,
derivación de la constante de estructura fina y estructura algebraica del
electrón}. Serie HMT--MD, artículo I, edición consultada de septiembre de 2026.
\bibitem{hmtX} O. Haidara Fall y R. Ramos Balsa,
\emph{Extrema y media razón holográfica}. Desarrollo del acoplamiento
aritmético--geométrico, registro dodecafásico y dirección dimensional,
edición consultada de septiembre de 2026.
\bibitem{hmtII} O. Haidara Fall y R. Ramos Balsa,
\emph{Origen estructural del parámetro de Barbero--Immirzi y de los
ángulos de mezcla CKM; derivación de las constantes de Planck, Boltzmann
y gravitación universal}. Serie HMT--MD, artículo II, edición de
septiembre de 2026. Relación entre acción, área, información y entropía
de entrelazamiento.
\bibitem{hmtVI} O. Haidara Fall y R. Ramos Balsa,
\emph{Partículas, persistencia y espectro de masas}. Serie HMT--MD, artículo VI,
edición consultada de septiembre de 2026.
\bibitem{HMT-IV} O. Haidara Fall y R. Ramos Balsa,
\emph{Moonshine, dualidad T y teoría M desde la estructura discreta del
continuo}. Serie HMT--MD, artículo IV, edición consultada de septiembre de 2026.
\bibitem{HMT-VIII} O. Haidara Fall y R. Ramos Balsa,
\emph{Aritmética genealógica de los números primos y estructura espectral
de la función zeta}. Serie HMT--MD, artículo VIII, edición consultada de septiembre de 2026.
\bibitem{HMT-IX} O. Haidara Fall y R. Ramos Balsa,
\emph{Resolución de la hipótesis del continuo en la clausura genealógica
holográfica}. Serie HMT--MD, artículo IX, edición consultada de septiembre de 2026.
\bibitem{Elkies} N. D. Elkies,
\emph{Rational lattices and their theta functions}. Arizona Winter School,
2009, teorema 2.
\href{https://people.math.harvard.edu/~elkies/aws09.pdf}{Notas del autor}.
\bibitem{ConwaySloane} J. H. Conway y N. J. A. Sloane,
\emph{Sphere Packings, Lattices and Groups}, tercera edición, Springer, 1999.
\href{https://doi.org/10.1007/978-1-4757-6568-7}{DOI: 10.1007/978-1-4757-6568-7}.
\bibitem{postnikov} A. Postnikov,
\emph{Total positivity, Grassmannians, and networks}, 2006,
\href{https://arxiv.org/abs/math/0609764}{arXiv:math/0609764}.
\bibitem{scott} J. S. Scott,
\emph{Grassmannians and Cluster Algebras}, 2003,
\href{https://arxiv.org/abs/math/0311148}{arXiv:math/0311148}.
\bibitem{amplituhedron} N. Arkani-Hamed y J. Trnka,
\emph{The Amplituhedron}, 2013,
\href{https://arxiv.org/abs/1312.2007}{arXiv:1312.2007}.
\bibitem{positive} N. Arkani-Hamed, Y. Bai y T. Lam,
\emph{Positive Geometries and Canonical Forms}, 2017,
\href{https://arxiv.org/abs/1703.04541}{arXiv:1703.04541}.
\bibitem{HMT-YM} O. Haidara Fall y R. Ramos Balsa,
\emph{El cierre holográfico del infinito}, edición integral de 2.249 páginas,
capítulo 99, «Imagen gauge y composición hacia Yang--Mills», pp. 1742--1762;
y \emph{Yang--Mills: medida común, reconstrucción y brecha espectral},
monografía especializada, desarrollo de la escala areal--volumétrica.
\bibitem{hmtIntegral} O. Haidara Fall y R. Ramos Balsa,
\emph{El cierre holográfico del infinito}, edición integral de 2.249 páginas,
septiembre de 2026. Desarrollos de los regímenes cuadráticos, el vacío
electromagnético, la realización electrodébil y las correspondencias
espirales de Mecánica Dimensional.

\phantomsection
\addcontentsline{toc}{section}{Referencias}
\bibitem{feynman1948} R. P. Feynman, ``Space-Time Approach to Non-Relativistic Quantum Mechanics'', \emph{Reviews of Modern Physics} \textbf{20}, 367--387 (1948). \href{https://doi.org/10.1103/RevModPhys.20.367}{doi:10.1103/RevModPhys.20.367}.
\bibitem{feynman1949} R. P. Feynman, ``The Theory of Positrons'', \emph{Physical Review} \textbf{76}, 749--759 (1949). \href{https://doi.org/10.1103/PhysRev.76.749}{doi:10.1103/PhysRev.76.749}.
\bibitem{feynman1965} R. P. Feynman, ``The Development of the Space-Time View of Quantum Electrodynamics'', conferencia Nobel, 11 de diciembre de 1965. \href{https://www.nobelprize.org/prizes/physics/1965/feynman/lecture/}{Texto de la conferencia}.
\bibitem{mustovicary} B. Musto y J. Vicary, ``Quantum Latin Squares and Unitary Error Bases'', \emph{Quantum Information and Computation} \textbf{16}, 1318--1332 (2016). \href{https://doi.org/10.26421/QIC16.15-16-4}{doi:10.26421/QIC16.15-16-4}.
\bibitem{delascuevas2020} G. De las Cuevas, T. Drescher y T. Netzer, ``Quantum Magic Squares: Dilations and Their Limitations'', \emph{Journal of Mathematical Physics} \textbf{61}, 111704 (2020). \href{https://doi.org/10.1063/5.0022344}{doi:10.1063/5.0022344}.
\bibitem{delascuevas2023} G. De las Cuevas, T. Netzer e I. Valentiner-Branth, ``Magic Squares: Latin, Semiclassical, and Quantum'', \emph{Journal of Mathematical Physics} \textbf{64}, 022201 (2023). \href{https://doi.org/10.1063/5.0127393}{doi:10.1063/5.0127393}.
\bibitem{dirac1933} P. A. M. Dirac, ``The Lagrangian in Quantum Mechanics'', \emph{Physikalische Zeitschrift der Sowjetunion} \textbf{3}, 64--72 (1933). \href{https://www.informationphilosopher.com/solutions/scientists/dirac/Lagrangian_1933.pdf}{Reproducción del artículo original}.
\bibitem{pauli1946} W. Pauli, ``Exclusion Principle and Quantum Mechanics'', conferencia Nobel, 13 de diciembre de 1946, en \emph{Nobel Lectures, Physics 1942--1962}, Elsevier, 1964, pp. 27--43, especialmente p. 42. \href{https://www.nobelprize.org/uploads/2018/06/pauli-lecture.pdf}{Texto de la conferencia}.
\bibitem{codata2018} E. Tiesinga, P. J. Mohr, D. B. Newell y B. N. Taylor, ``CODATA Recommended Values of the Fundamental Physical Constants: 2018'', \emph{Reviews of Modern Physics} \textbf{93}, 025010 (2021). \href{https://doi.org/10.1103/RevModPhys.93.025010}{doi:10.1103/RevModPhys.93.025010}. \href{https://physics.nist.gov/cuu/Constants/ArchiveASCII/allascii_2018.txt}{Tablas archivadas de 2018}.
\bibitem{codata2022} P. J. Mohr, D. B. Newell, B. N. Taylor y E. Tiesinga, ``CODATA Recommended Values of the Fundamental Physical Constants: 2022'', \emph{Reviews of Modern Physics} \textbf{97}, 025002 (2025). \href{https://doi.org/10.1103/RevModPhys.97.025002}{doi:10.1103/RevModPhys.97.025002}. \href{https://physics.nist.gov/cuu/Constants/Table/allascii.txt}{Tabla de constantes}.
\bibitem{flm} I. B. Frenkel, J. Lepowsky y A. Meurman, \emph{Vertex Operator Algebras and the Monster}, Pure and Applied Mathematics, vol. 134, Academic Press, 1988.
\bibitem{borcherds} R. E. Borcherds, ``Monstrous Moonshine and Monstrous Lie Superalgebras'', \emph{Inventiones Mathematicae} \textbf{109}, 405--444 (1992). \href{https://doi.org/10.1007/BF01232032}{doi:10.1007/BF01232032}.
\bibitem{conwaysloane} J. H. Conway y N. J. A. Sloane, \emph{Sphere Packings, Lattices and Groups}, 3.\textsuperscript{a} ed., Grundlehren der mathematischen Wissenschaften 290, Springer, 1999. \href{https://doi.org/10.1007/978-1-4757-6568-7}{doi:10.1007/978-1-4757-6568-7}.
\bibitem{bipm2018} Conférence générale des poids et mesures, \emph{Resolution 1 of the 26th CGPM: On the revision of the International System of Units}, 2018. \href{https://www.bipm.org/en/committees/cg/cgpm/26-2018/resolution-1}{Resolución oficial}. \href{https://doi.org/10.59161/CGPM2018RES1E}{doi:10.59161/CGPM2018RES1E}.
\bibitem{dirac1931} P. A. M. Dirac, ``Quantised singularities in the electromagnetic field'', \emph{Proceedings of the Royal Society of London. Series A} \textbf{133}, 60--72 (1931). \href{https://doi.org/10.1098/rspa.1931.0130}{doi:10.1098/rspa.1931.0130}.
\bibitem{lep2006} ALEPH, DELPHI, L3, OPAL y SLD Collaborations; LEP Electroweak Working Group; SLD Electroweak Group y SLD Heavy Flavour Group, ``Precision electroweak measurements on the Z resonance'', \emph{Physics Reports} \textbf{427}, 257--454 (2006). \href{https://doi.org/10.1016/j.physrep.2005.12.006}{doi:10.1016/j.physrep.2005.12.006}.
\bibitem{chenlamshimakura2016} H.-Y. Chen, C. H. Lam y H. Shimakura, ``\(\mathbb Z_3\)-orbifold construction of the Moonshine vertex operator algebra and some maximal \(3\)-local subgroups of the Monster'', prepublicación, arXiv:1606.05961 (2016), versión 2 (2017). \href{https://arxiv.org/abs/1606.05961}{arXiv:1606.05961}.
\bibitem{abelamyamada2017} T. Abe, C. H. Lam y H. Yamada, ``A remark on \(\mathbb Z_p\)-orbifold constructions of the Moonshine vertex operator algebra'', prepublicación, arXiv:1705.09022 (2017), versión 4. \href{https://arxiv.org/abs/1705.09022v4}{arXiv:1705.09022}.
\bibitem{kmconwaynorton1979} J. H. Conway y S. P. Norton, ``Monstrous Moonshine'', \emph{Bulletin of the London Mathematical Society} \textbf{11}, 308--339 (1979). \href{https://doi.org/10.1112/blms/11.3.308}{doi:10.1112/blms/11.3.308}.
\bibitem{bennett2003} C. H. Bennett, ``Notes on Landauer's principle, reversible computation, and Maxwell's Demon'', \emph{Studies in History and Philosophy of Modern Physics} \textbf{34}, 501--510 (2003). \href{https://doi.org/10.1016/S1355-2198(03)00039-X}{doi:10.1016/S1355-2198(03)00039-X}.

\end{thebibliography}
\end{document}
